\documentclass[11pt,oneside]{article}

\usepackage[a4paper,margin=27mm,headheight=15pt]{geometry}
\usepackage[T1]{fontenc}
\usepackage[utf8]{inputenc}
\IfFileExists{libertinus.sty}{\usepackage{libertinus}}{\usepackage{lmodern}}
\usepackage{microtype}
\usepackage{amsmath,amssymb,amsthm,mathtools,mathrsfs}
% Canonical mathematical notation for active FabiusFunction LaTeX documents.
%
% This file is intentionally a plain TeX input rather than a LaTeX package.
% Every standalone document loads it by an explicit path relative to that
% document's directory; included fragments inherit the definitions from their
% root document.  Keep this file independent of document layout and metadata.
%
% Contract:
%   * one command has one mathematical meaning throughout the active corpus;
%   * names encode semantic roles rather than a document-local choice of letter;
%   * visually similar but differently normalized objects have different print;
%   * every command below is defined normatively in the notation catalogue.
%
% The active corpus excludes docs/archive/.  Historical sources may retain
% their original notation only inside an explicitly labelled quotation.

\makeatletter
\@ifundefined{mathbb}{%
  \PackageError{fabius-notation}{The amssymb package must be loaded first}%
    {Load amsmath and amssymb before inputting fabius-notation.tex.}%
}{}
\@ifundefined{operatorname}{%
  \PackageError{fabius-notation}{The amsmath package must be loaded first}%
    {Load amsmath before inputting fabius-notation.tex.}%
}{}
\makeatother

% Version stamp used by the catalogue and by static audits.
\newcommand{\FabiusNotationVersion}{2026-08-31}

% Number systems.  NaturalNumbers is explicitly the nonnegative convention.
\newcommand{\NaturalNumbers}{\mathbb{N}_{0}}
\newcommand{\PositiveIntegers}{\mathbb{N}_{>0}}
\newcommand{\IntegerNumbers}{\mathbb{Z}}
\newcommand{\RationalNumbers}{\mathbb{Q}}
\newcommand{\RealNumbers}{\mathbb{R}}
\newcommand{\ComplexNumbers}{\mathbb{C}}
\newcommand{\ExtendedRealNumbers}{\overline{\mathbb{R}}}

% The three principal Fabius--Rvachev functions.  Subscripts are deliberate:
% a bare F and a calligraphic F have several unrelated uses in the corpus.
\newcommand{\FabiusBounded}{F_{\mathrm{bd}}}
\newcommand{\FabiusGlobal}{F_{\mathrm{gl}}}
\newcommand{\FabiusQuantile}{Q_{F}}
\newcommand{\FabiusClampedQuantile}{Q_{F,\mathrm{cl}}}
\newcommand{\RvachevUp}{\operatorname{up}}

% Constants, elementary operators, and delimiters.
\newcommand{\EulerE}{\mathrm{e}}
\newcommand{\ImaginaryUnit}{\mathrm{i}}
\newcommand{\Differential}{\mathop{}\!\mathrm{d}}
\newcommand{\AbsoluteValue}[1]{\left\lvert #1\right\rvert}
\newcommand{\Norm}[1]{\left\lVert #1\right\rVert}
\newcommand{\Floor}[1]{\left\lfloor #1\right\rfloor}
\newcommand{\Ceiling}[1]{\left\lceil #1\right\rceil}
\newcommand{\FractionalPart}[1]{\left\{#1\right\}}
\newcommand{\PositivePart}[1]{\left(#1\right)_{+}}
\newcommand{\IndicatorOf}[1]{\mathbf{1}_{#1}}
\newcommand{\SetOf}[1]{\left\{#1\right\}}
\newcommand{\InnerProduct}[1]{\left\langle #1\right\rangle}
\newcommand{\CardinalityOf}[1]{\left\lvert #1\right\rvert}
\newcommand{\DefinitionEquals}{\mathrel{\coloneqq}}
\newcommand{\Sign}{\operatorname{sgn}}
\newcommand{\IdentityOperator}{\operatorname{id}}
\newcommand{\SpanOperator}{\operatorname{span}}
\newcommand{\TraceOperator}{\operatorname{tr}}
\newcommand{\RankOperator}{\operatorname{rank}}
\newcommand{\DiagonalOperator}{\operatorname{diag}}
\newcommand{\ResidueOperator}{\operatorname*{Res}}
\newcommand{\LeastCommonMultiple}{\operatorname{lcm}}
\newcommand{\OrderOperator}{\operatorname{ord}}
\newcommand{\PrincipalLogarithm}{\operatorname{Log}}
\newcommand{\FallingFactorial}[2]{\left(#1\right)^{\underline{#2}}}
\newcommand{\RisingFactorial}[2]{\left(#1\right)^{\overline{#2}}}

% Support, topology, convolution, and metric notation.
\newcommand{\SupportOperator}{\operatorname{supp}}
\newcommand{\SupportOf}[1]{\SupportOperator\!\left(#1\right)}
\newcommand{\TopologicalSupportOf}[1]{\operatorname{tsupp}\!\left(#1\right)}
\newcommand{\Convolution}{\mathbin{*}}
\newcommand{\MetricDistance}{\operatorname{dist}}
\newcommand{\TotalVariationDistance}{d_{\mathrm{TV}}}

% Probability.  Equality and convergence in law are not metric distance.
\newcommand{\Probability}{\mathbb{P}}
\newcommand{\Expectation}{\mathbb{E}}
\newcommand{\Variance}{\operatorname{Var}}
\newcommand{\Covariance}{\operatorname{Cov}}
\newcommand{\LawOperator}{\operatorname{Law}}
\newcommand{\LawOf}[1]{\LawOperator\!\left(#1\right)}
\newcommand{\EqualInLaw}{\mathrel{\overset{\mathrm{law}}{=}}}
\newcommand{\ConvergesInLaw}{\mathrel{\xrightarrow{\mathrm{law}}}}

% Fourier and integral transforms.  The printed labels expose normalization.
% FourierTwoPi uses exp(-2*pi*i*x*xi); FourierAngular uses exp(-i*x*omega).
\newcommand{\FourierTwoPi}[1]{\widehat{#1}_{\,2\pi}}
\newcommand{\FourierAngular}[1]{\widehat{#1}_{\,\mathrm{ang}}}
\newcommand{\LaplaceTransformSymbol}{\mathcal{L}}
\newcommand{\MellinTransformSymbol}{\mathcal{M}}
\newcommand{\LaplaceTransformOf}[1]{\LaplaceTransformSymbol\!\left[#1\right]}
\newcommand{\MellinTransformOf}[1]{\MellinTransformSymbol\!\left[#1\right]}
\newcommand{\MomentGeneratingFunction}[1]{M_{#1}}
\newcommand{\CharacteristicFunction}[1]{\varphi_{#1}}

% The two standard sinc functions must never share one symbol.
\newcommand{\SincRad}{\operatorname{sinc}_{\mathrm{rad}}}
\newcommand{\SincPi}{\operatorname{sinc}_{\pi}}

% Binary arithmetic and Thue--Morse notation.
\newcommand{\BinaryDigitSum}{\operatorname{wt}_{2}}
\newcommand{\ThueMorseSignSymbol}{\varepsilon}
\newcommand{\ThueMorseSign}[1]{\varepsilon_{#1}}
\newcommand{\PadicValuation}[1]{\nu_{#1}}
\newcommand{\TwoAdicValuation}{\nu_{2}}

% q-series notation.  Every base is an explicit argument.
\newcommand{\QPochhammer}[3]{\left(#1;#2\right)_{#3}}
\newcommand{\GaussianBinomial}[3]{%
  \genfrac{[}{]}{0pt}{}{#1}{#2}_{#3}%
}
\newcommand{\QInteger}[2]{\left[#1\right]_{#2}}

% Named special functions.
\newcommand{\Polylogarithm}{\operatorname{Li}}
\newcommand{\LambertW}{W}
\newcommand{\GammaFunction}{\Gamma}
\newcommand{\RiemannZeta}{\zeta}

% Asymptotics.  BigO and LittleO are operator tokens for existing formula
% grammar; prefer the At forms so that the limiting regime is printed.
\newcommand{\BigO}{\mathcal{O}}
\newcommand{\LittleO}{o}
\newcommand{\BigOAt}[2]{\mathcal{O}_{#2}\!\left(#1\right)}
\newcommand{\LittleOAt}[2]{o_{#2}\!\left(#1\right)}

\usepackage{bm}
\usepackage{booktabs,longtable,array,tabularx}
\usepackage{graphicx}
\usepackage{pgfplots}
\usepackage{xcolor}
\usepackage{enumitem}
\usepackage{fancyhdr}
\usepackage{titlesec}
\usepackage{listings}
\usepackage{xurl}
\usepackage{ragged2e,needspace}
\usepackage{hyperref}
\usepackage[nameinlink,noabbrev]{cleveref}

\definecolor{linkblue}{RGB}{25,75,135}
\definecolor{shadegray}{RGB}{246,247,249}
\definecolor{rulegray}{RGB}{195,201,208}
\pgfplotsset{compat=1.18}
\hypersetup{
  colorlinks=true,
  linkcolor=linkblue,
  citecolor=linkblue,
  urlcolor=linkblue,
  pdfauthor={Vladimir Reshetnikov; ProveIt notation catalogue},
  pdftitle={Mathematical Notation Catalogue for Analysis/FabiusFunction},
  pdfsubject={Normative symbols, conventions, aliases, and formalization crosswalk for the active Fabius-function corpus},
  pdfkeywords={Fabius function, Rvachev up-function, notation, Fourier transform, q-series, Lambert W, Thue-Morse}
}

\pagestyle{fancy}
\fancyhf{}
\fancyhead[L]{\small FabiusFunction notation catalogue}
\fancyhead[R]{\small\nouppercase{\leftmark}}
\fancyfoot[C]{\thepage}
\renewcommand{\headrulewidth}{0.4pt}

\titleformat{\section}{\Large\bfseries}{\thesection}{0.7em}{}
\titleformat{\subsection}{\large\bfseries}{\thesubsection}{0.7em}{}
\titleformat{\subsubsection}{\normalsize\bfseries}{\thesubsubsection}{0.7em}{}
\titleformat{\paragraph}[block]{\normalsize\bfseries}{}{0pt}{}
\titlespacing*{\paragraph}{0pt}{1.0\baselineskip}{0.25\baselineskip}
\setcounter{secnumdepth}{3}
\setcounter{tocdepth}{2}
\setlist[itemize]{leftmargin=2em,itemsep=0.25em,topsep=0.4em}
\setlist[enumerate]{leftmargin=2.3em,itemsep=0.25em,topsep=0.4em}
\widowpenalty=10000
\clubpenalty=10000
\displaywidowpenalty=10000

% Long command names and file paths occasionally leave no natural break point.
\setlength{\emergencystretch}{2em}

\newtheorem{theorem}{Theorem}[section]
\newtheorem{proposition}[theorem]{Proposition}
\newtheorem{lemma}[theorem]{Lemma}
\newtheorem{corollary}[theorem]{Corollary}
\newtheorem{conjecture}[theorem]{Conjecture}
\theoremstyle{definition}
\newtheorem{definition}[theorem]{Definition}
\newtheorem{algorithm}[theorem]{Algorithm}
\newtheorem{example}[theorem]{Example}
\theoremstyle{remark}
\newtheorem{remark}[theorem]{Remark}
\newtheorem{warning}[theorem]{Warning}

\newcolumntype{L}[1]{>{\RaggedRight\arraybackslash}p{#1}}
\newcolumntype{S}[1]{>{\RaggedRight\arraybackslash\scriptsize}p{#1}}

% Repository-facing helpers are not mathematical notation and therefore do not
% belong in the shared notation input.
\DeclareRobustCommand{\shortpath}[1]{\mbox{\path{#1}}}
\DeclareRobustCommand{\modulepath}[1]{\shortpath{#1}}
\DeclareRobustCommand{\repo}{\modulepath{Analysis/FabiusFunction}}
\newcommand{\lean}[1]{%
  \ifmmode\text{\nolinkurl{#1}}\else\mbox{\nolinkurl{#1}}\fi}
\newcommand{\leanline}[1]{%
  \par\nopagebreak[3]\smallskip
  \noindent\hspace*{1.5em}{\footnotesize\lean{#1}}\par\nopagebreak[2]}
\DeclareRobustCommand{\commandname}[1]{\texttt{\textbackslash#1}}
% Machine-readable token for the exhaustive document-local accent register.
\DeclareRobustCommand{\localaccentname}[1]{\commandname{#1}}

\newenvironment{proofidea}{\par\smallskip\noindent\textit{Proof architecture.}\ }{\hfill$\square$\par\smallskip}
\newenvironment{boxedremark}{%
  \par\smallskip\noindent\begin{center}\begin{minipage}{0.94\textwidth}%
  \color{black}\hrule\vspace{0.6em}\small
}{%
  \vspace{0.6em}\hrule\end{minipage}\end{center}\smallskip
}

\lstdefinestyle{pseudo}{
  basicstyle=\ttfamily\scriptsize,
  backgroundcolor=\color{shadegray},
  frame=single,
  rulecolor=\color{rulegray},
  columns=fullflexible,
  breaklines=true,
  breakatwhitespace=true,
  keepspaces=true,
  showstringspaces=false,
  xleftmargin=0.7em,
  xrightmargin=0.7em,
  aboveskip=0.8em,
  belowskip=0.8em
}

% Catalogue-only presentation helpers.  They add no mathematical notation.
\newenvironment{catalogueentry}[2]{%
  \Needspace{6\baselineskip}%
  \par\addvspace{1.05\baselineskip}%
  \phantomsection\hypertarget{notation:#1}{}%
  \noindent{\large\bfseries\RaggedRight #2\par}\smallskip
}{%
  \par\nopagebreak[2]\vspace{0.2\baselineskip}%
  \noindent\color{rulegray}\rule{\linewidth}{0.3pt}\color{black}%
}
\newcommand{\entryfield}[2]{\Needspace{3\baselineskip}\par\noindent\textit{#1.}\enspace #2}
\newcommand{\notelink}[2]{\hyperlink{notation:#1}{#2}}
\AddToHook{cmd/paragraph/before}{\Needspace{6\baselineskip}}

\title{\bfseries Mathematical Notation Catalogue for
\mbox{\texttt{Analysis/FabiusFunction}}\\[0.35em]
\large Canonical symbols, exact definitions, local namespaces, and historical aliases}
\author{Encyclopedic companion to Vladimir Reshetnikov's \texttt{ProveIt} formalization}
\date{Notation contract \FabiusNotationVersion\\
\small Active corpus at the validation snapshot: 174 recursively inventoried
\texttt{.tex} sources}

\begin{document}
\maketitle

\begin{abstract}
This catalogue is the normative dictionary for mathematical notation in the active
\repo\ documentation corpus.  It distinguishes the bounded Fabius function, the signed
global extension, and Rvachev's up-function; fixes number-system, probability, transform,
binary, $q$-series, asymptotic, fractional-calculus, approximation, and operator
conventions; and records every recurring collision discovered by a recursive audit of the
current 174-source active \texttt{.tex} corpus.  Each entry states its type, domain,
codomain, indices, parameters, normalization, exceptional values, and formalization
status.  No phrase such as ``with the usual convention'' supplies missing data: when two
usual conventions exist, both are defined and their conversion is printed.  Historical
papers and frontier reports remain mathematically interpretable through explicit alias and
local-namespace crosswalks rather than through overloaded bare letters.
\end{abstract}

\begin{boxedremark}
\textbf{Source/PDF version notice.}
This TeX source contains the post-replay notation-catalogue union and its incoming
presentation polish.  The retained PDF predates that union, so source/PDF parity is not
claimed; a fresh uninterrupted three-pass Libertinus rebuild is pending.
\end{boxedremark}

\begin{boxedremark}
\textbf{Normative scope.}
The command definitions in \modulepath{docs/fabius-notation.tex}, version
\FabiusNotationVersion, are normative for active project-authored LaTeX.  This catalogue
defines their mathematical meaning.  The directory \modulepath{docs/archive/} is not an active
migration target and is used only as provenance.  A historical source under
\modulepath{docs/papers/} may display its original notation only inside an explicit historical
scope.  A source-faithful whole-paper transcription opens that scope with the exact
preamble marker \texttt{HISTORICAL-NOTATION-SCOPE}; a shorter quotation labels the
historical convention at the quotation itself.  The marker never exempts surrounding
project prose, which uses the canonical symbols.  Included tables and generated fragments
inherit the namespace of their standalone parent document.
\end{boxedremark}

\tableofcontents
\newpage

% Keep each heading with enough opening material to identify the new scope,
% without forcing a nearly empty predecessor page merely because a major
% division follows.  The running mark still changes at the actual heading.
\AddToHook{cmd/section/before}{\Needspace{12\baselineskip}}
\AddToHook{cmd/subsection/before}{\Needspace{9\baselineskip}}
\AddToHook{cmd/subsubsection/before}{\Needspace{9\baselineskip}}

\section{Scope contract and how to use this catalogue}

\subsection{What counts as notation}

The inventory includes a glyph or compound expression whenever its interpretation requires
one or more mathematical choices: a domain, codomain, index set, endpoint rule, transform
normalization, branch, sign, scaling, convergence mode, or proof-status boundary.  Thus it
includes not only named functions but also interval brackets, empty sums, $0^0$, arrows of
convergence, subscripts on transforms, and superscripts denoting derivatives or powers.
Typography that changes no mathematical meaning---for example the insertion of braces in
\commandname{mathbb}\texttt{\{R\}}---is not a separate notation.

The catalogue has two levels.
\begin{enumerate}
\item A \emph{global semantic symbol} has one meaning throughout the active corpus and,
      where available, one command in \modulepath{fabius-notation.tex}.
\item A \emph{local symbol} is scoped to a definition, theorem, table, or explicitly named
      document part.  A local symbol must be declared where its scope begins and must not
      be cited outside that scope without qualification.
\end{enumerate}
The local level is essential: a corpus this broad legitimately uses $A$, $D$, $I$, $P$,
$Q$, and $T$ for many unrelated temporary objects.  Unification means making those scopes
and meanings explicit, not forcing every matrix and polynomial to share a single name.

\subsection{The no-implicit-parts checklist}

Every new or revised notation declaration answers all applicable questions below.
\begin{enumerate}
\item What mathematical object is denoted: scalar, set, map, measure, sequence, matrix,
      operator, equivalence class, or proposition?
\item What are its domain and codomain?  If it is indexed, what is the complete index set?
\item What is the exact defining formula, including all summation and product bounds?
\item Which parameters are fixed, which vary, and what inequalities or nonvanishing
      conditions do they satisfy?
\item Which normalization, sign, Fourier frequency, logarithm, complex-power branch, and
      measure convention is used?
\item What happens at endpoints, at an empty index set, at zero, and at a removable
      singularity?
\item Does equality mean pointwise equality, equality almost everywhere, equality of
      measures, equality in distribution, or asymptotic equivalence?
\item Is the assertion formally proved, a standard definition, derived only in a frontier
      report, conjectural, refuted, numerical, or historical?
\end{enumerate}

\subsection{Status vocabulary}
\label{sec:status-register}

\begin{longtable}{@{}p{0.19\textwidth}p{0.75\textwidth}@{}}
\toprule
Status & Exact meaning in this catalogue\\
\midrule
\endhead
\textbf{Definition} & A definition or typography contract; no theorem is asserted.\\
\textbf{Lean-proved} & The stated mathematical content has a named proved Lean declaration
in the module identified by the crosswalk.\\
\textbf{Classical} & A standard mathematical convention used by the prose; its presence in
the catalogue is not a claim that the project has formalized the general theorem.\\
\textbf{Frontier-derived} & Derived in a research-frontier document but not certified here as
having an exact current Lean counterpart.\\
\textbf{Conjectural} & Proposed and explicitly unproved.\\
\textbf{Refuted} & A historical formula or normalization for which the project records a
counterexample or negative theorem.\\
\textbf{Historical alias} & A symbol needed to read a source artifact; it is not canonical
for new project prose.\\
\bottomrule
\end{longtable}

\section{Canonical quick reference}

The first column is the literal shared command name.  The second column is its rendering;
the third fixes its semantics.  Commands taking arguments show metavariables in braces.

\begingroup\small
\renewcommand{\commandname}[1]{\texttt{\textbackslash#1}\allowbreak}
\begin{longtable}{@{}S{0.35\textwidth}L{0.14\textwidth}L{0.45\textwidth}@{}}
\toprule
Shared command & Rendering & Normative meaning\\
\midrule
\endfirsthead
\toprule
Shared command & Rendering & Normative meaning\\
\midrule
\endhead
\commandname{FabiusNotationVersion} & \FabiusNotationVersion & Date-valued version of the
shared notation contract; it is metadata, not a mathematical variable.\\
\commandname{NaturalNumbers} & $\NaturalNumbers$ & Nonnegative integers
$\{0,1,2,\ldots\}$.\\
\commandname{PositiveIntegers} & $\PositiveIntegers$ & Strictly positive integers
$\{1,2,3,\ldots\}$.\\
\commandname{IntegerNumbers} & $\IntegerNumbers$ & Integers.\\
\commandname{RationalNumbers} & $\RationalNumbers$ & Rational numbers.\\
\commandname{RealNumbers} & $\RealNumbers$ & Real numbers.\\
\commandname{ComplexNumbers} & $\ComplexNumbers$ & Complex numbers.\\
\commandname{ExtendedRealNumbers} & $\ExtendedRealNumbers$ & Extended reals
$\RealNumbers\cup\{-\infty,+\infty\}$.\\
\commandname{FabiusBounded} & $\FabiusBounded$ & Bounded, CDF-normalized Fabius
function.\\
\commandname{FabiusGlobal} & $\FabiusGlobal$ & Signed global Fabius extension.\\
\commandname{FabiusQuantile} & $\FabiusQuantile$ & Ordinary inverse of
$\FabiusBounded|_{[0,1]}$.\\
\commandname{FabiusClampedQuantile} & $\FabiusClampedQuantile$ & Totalized clamped
inverse on all real inputs.\\
\commandname{RvachevUp} & $\RvachevUp$ & Rvachev's even compactly supported
up-function.\\
\commandname{EulerE} & $\EulerE$ & Base of the natural logarithm.\\
\commandname{ImaginaryUnit} & $\ImaginaryUnit$ & The complex number satisfying
$\ImaginaryUnit^2=-1$.\\
\commandname{Differential} & $\Differential x$ & Upright differential in an integral or
differential form.\\
\commandname{AbsoluteValue}\texttt{\{x\}} & $\AbsoluteValue{x}$ & Absolute value or,
when a field is explicitly complex, complex modulus.\\
\commandname{Norm}\texttt{\{x\}} & $\Norm{x}$ & The norm supplied by the explicitly
named normed space.\\
\commandname{Floor}\texttt{\{x\}} & $\Floor{x}$ & Greatest integer not exceeding $x$.\\
\commandname{Ceiling}\texttt{\{x\}} & $\Ceiling{x}$ & Least integer not below $x$.\\
\commandname{FractionalPart}\texttt{\{x\}} & $\FractionalPart{x}$ & $x-\Floor{x}$ for
$x\in\RealNumbers$.\\
\commandname{PositivePart}\texttt{\{x\}} & $\PositivePart{x}$ & $\max\{x,0\}$ for
an ordered scalar $x$.\\
\commandname{IndicatorOf}\texttt{\{A\}} & $\IndicatorOf{A}$ & Indicator of a set $A$,
equal to one on $A$ and zero off $A$.\\
\commandname{SetOf}\texttt{\{x:P(x)\}} & $\SetOf{x:P(x)}$ & Delimiters for a set
specified by membership data or a predicate; the ambient type is stated locally.\\
\commandname{InnerProduct}\texttt{\{x,y\}} & $\InnerProduct{x,y}$ & Inner-product
delimiters; the scalar field, space, and linear-slot convention are stated locally.\\
\commandname{CardinalityOf}\texttt{\{A\}} & $\CardinalityOf{A}$ & Finite cardinality
or cardinal number according to the declared type; never measure or absolute value.\\
\commandname{DefinitionEquals} & $a\DefinitionEquals b$ & Definitional assignment:
the left-hand symbol is introduced to mean exactly the right-hand object.\\
\commandname{Sign} & $\Sign$ & Real sign function with $\Sign(0)=0$.\\
\commandname{IdentityOperator} & $\IdentityOperator$ & Identity map on the domain stated
at the point of use.\\
\commandname{SpanOperator} & $\SpanOperator$ & Linear span over the stated scalar field.\\
\commandname{TraceOperator} & $\TraceOperator$ & Matrix or trace-class operator trace.\\
\commandname{RankOperator} & $\RankOperator$ & Dimension of an image, over the stated
field.\\
\commandname{DiagonalOperator} & $\DiagonalOperator$ & Diagonal matrix/operator constructor.\\
\commandname{ResidueOperator} & $\ResidueOperator$ & Complex residue at a stated point.\\
\commandname{LeastCommonMultiple} & $\LeastCommonMultiple$ & Nonnegative least common
multiple; finite-family and empty-family rules must be stated locally.\\
\commandname{OrderOperator} & $\OrderOperator$ & An order invariant whose species
(vanishing, pole, multiplicative, group-element, or recurrence order) is printed.\\
\commandname{PrincipalLogarithm} & $\PrincipalLogarithm$ & Principal complex logarithm
with argument in $(-\pi,\pi]$ and domain $\ComplexNumbers\setminus(-\infty,0]$.\\
\commandname{FallingFactorial}\texttt{\{x\}\{n\}} & $\FallingFactorial{x}{n}$ &
Falling factorial $\prod_{j=0}^{n-1}(x-j)$, equal to one when $n=0$.\\
\commandname{RisingFactorial}\texttt{\{x\}\{n\}} & $\RisingFactorial{x}{n}$ &
Rising factorial $\prod_{j=0}^{n-1}(x+j)$, equal to one when $n=0$.\\
\commandname{SignedStirlingFirstKind}\texttt{\{n\}\{k\}} &
\(\SignedStirlingFirstKind{n}{k}\) & Signed coefficient of \(x^k\) in the
falling factorial \(\FallingFactorial{x}{n}\).\\
\commandname{UnsignedStirlingFirstKind}\texttt{\{n\}\{k\}} &
\(\UnsignedStirlingFirstKind{n}{k}\) & Number of permutations of \(n\)
labels having \(k\) cycles.\\
\commandname{StirlingSecondKind}\texttt{\{n\}\{k\}} &
\(\StirlingSecondKind{n}{k}\) & Number of partitions of \(n\) labels into
\(k\) nonempty unordered blocks.\\
\commandname{LahNumber}\texttt{\{n\}\{k\}} & \(\LahNumber{n}{k}\) &
Number of partitions of \(n\) labels into \(k\) nonempty linearly ordered
lists.\\
\commandname{TypeAEulerianNumber}\texttt{\{n\}\{k\}} &
\(\TypeAEulerianNumber{n}{k}\) & Number of permutations of \(n\) labels with
\(k\) ordinary descents.\\
\commandname{TypeBEulerianNumber}\texttt{\{n\}\{k\}} &
\(\TypeBEulerianNumber{n}{k}\) & Number of signed permutations of size \(n\)
with \(k\) type-B descents, using the sentinel \(\pi_0=0\).\\
\commandname{SecondOrderEulerianNumber}\texttt{\{n\}\{k\}} &
\(\SecondOrderEulerianNumber{n}{k}\) & Number of order-\(n\) Stirling
permutations with \(k\) ordinary descents.\\
\commandname{TypeAEulerianPolynomial}\texttt{\{n\}} &
\(\TypeAEulerianPolynomial{n}\) & Type-A Eulerian row polynomial; its
evaluation argument follows the command.\\
\commandname{TypeBEulerianPolynomial}\texttt{\{n\}} &
\(\TypeBEulerianPolynomial{n}\) & Type-B Eulerian row polynomial; its
evaluation argument follows the command.\\
\commandname{SecondOrderEulerianPolynomial}\texttt{\{n\}} &
\(\SecondOrderEulerianPolynomial{n}\) & Second-order Eulerian row polynomial;
its evaluation argument follows the command.\\
\commandname{BellNumber}\texttt{\{n\}} & \(\BellNumber{n}\) & Number of all
set partitions of \(n\) labelled elements.\\
\commandname{ComplementaryBellNumber}\texttt{\{n\}} &
\(\ComplementaryBellNumber{n}\) & Signed block-parity sum
\(\sum_k(-1)^k\StirlingSecondKind{n}{k}\).\\
\commandname{TouchardPolynomial}\texttt{\{n\}} &
\(\TouchardPolynomial{n}\) & Polynomial
\(\sum_k\StirlingSecondKind{n}{k}x^k\); its evaluation argument follows.\\
\commandname{ExponentialPartialBellPolynomial}\texttt{\{n\}\{k\}} &
\(\ExponentialPartialBellPolynomial{n}{k}\) & Partial Bell polynomial in the
exponential-series normalization.\\
\commandname{ExponentialCompleteBellPolynomial}\texttt{\{n\}} &
\(\ExponentialCompleteBellPolynomial{n}\) & Sum over all partial exponential
Bell polynomials of total weight \(n\).\\
\commandname{OrdinaryPartialBellPolynomial}\texttt{\{n\}\{k\}} &
\(\OrdinaryPartialBellPolynomial{n}{k}\) & Partial Bell polynomial in the
ordinary-series normalization.\\
\commandname{CoefficientExtraction}\texttt{\{z\}\{n\}} &
\(\CoefficientExtraction{z}{n}\) & Operator extracting the coefficient of
\(z^n\) from the following series operand.\\
\commandname{CycleCountOperator} & \(\CycleCountOperator\) & Number of cycles
of the following permutation operand.\\
\commandname{SetPartitionOperator} & \(\SetPartitionOperator\) & Set of set
partitions of the following finite-set operand.\\
\commandname{BitwiseAnd} & \(a\BitwiseAnd b\) & Bitwise AND on declared
nonnegative integers.\\
\commandname{ExponentialGeneratingFunctionOf}\texttt{\{A\}} &
\(\ExponentialGeneratingFunctionOf{A}\) & Exponential generating function of
the named coefficient family.\\
\commandname{OrdinaryGeneratingFunctionOf}\texttt{\{A\}} &
\(\OrdinaryGeneratingFunctionOf{A}\) & Ordinary generating function of the
named coefficient family.\\
\commandname{NormalizedReversionCoefficient}\texttt{\{f\}\{j\}} &
\(\NormalizedReversionCoefficient{f}{j}\) & Normalized exponential-series
input \(f_{j+1}/((j+1)f_1)\) used in compositional-reversion formulas.\\
\commandname{OrdinaryGeneratingCoefficient}\texttt{\{x\}\{j\}} &
\(\OrdinaryGeneratingCoefficient{x}{j}\) & Explicit family-and-index ledger
coefficient in an ordinary generating series.\\
\commandname{PartitionLatticeMinimum} & \(\PartitionLatticeMinimum\) &
Discrete singleton partition, the minimum of a declared partition lattice.\\
\commandname{PartitionLatticeMaximum} & \(\PartitionLatticeMaximum\) &
One-block partition, the maximum of a declared partition lattice.\\
\commandname{TraceBellArgument}\texttt{\{k\}} &
\(\TraceBellArgument{k}\) & Scaled trace input
\(-(k-1)!\TraceOperator(A^k)\) for determinant/Bell identities.\\
\commandname{FiniteField} & \(\FiniteField_q\) & Finite field of the declared
prime-power order \(q\).\\
\commandname{CatalanNumber}\texttt{\{n\}} & \(\CatalanNumber{n}\) & Catalan
number \(\binom{2n}{n}/(n+1)\), visibly separated from complementary Bell
numbers.\\
\commandname{GregoryCoefficient}\texttt{\{n\}} &
\(\GregoryCoefficient{n}\) & Coefficient of \(z^n\) in
\(z/\log(1+z)\).\\
\commandname{GenocchiNumber}\texttt{\{n\}} & \(\GenocchiNumber{n}\) &
Genocchi number in the exponential generating convention specified below.\\
\commandname{GenocchiPolynomial}\texttt{\{n\}} &
\(\GenocchiPolynomial{n}\) & Genocchi polynomial; its evaluation argument
follows the command.\\
\commandname{BarnesGFunction} & \(\BarnesGFunction\) & Barnes \(G\)-function
normalized by \(\BarnesGFunction(1)=1\) and
\(\BarnesGFunction(z+1)=\GammaFunction(z)\BarnesGFunction(z)\).\\
\commandname{BellHessenbergMatrixH}\texttt{\{n\}} &
\(\BellHessenbergMatrixH{n}\) & First explicit Hessenberg determinant model
for the complete exponential Bell polynomial.\\
\commandname{BellHessenbergMatrixK}\texttt{\{n\}} &
\(\BellHessenbergMatrixK{n}\) & Second explicit Hessenberg determinant model
for the complete exponential Bell polynomial.\\
\commandname{AssociahedronByDimension}\texttt{\{d\}} &
\(\AssociahedronByDimension{d}\) & Stasheff associahedron indexed by its
actual dimension \(d\), never by an inherited ambiguous source index.\\
\commandname{LagrangeShiftCoordinate} & \(\LagrangeShiftCoordinate\) & Local
translated coordinate \(z-b\) in a declared Lagrange-inversion scope.\\
\commandname{SupportOperator} & $\SupportOperator$ & Pointwise-support operator token;
its operand and ambient space must follow.\\
\commandname{SupportOf}\texttt{\{f\}} & $\SupportOf{f}$ & Pointwise support
$\{x:f(x)\ne0\}$.\\
\commandname{CoefficientSupportOf}\texttt{\{F\}} &
\(\CoefficientSupportOf{F}\) & Set of exponent indices or monomials carrying
nonzero coefficients; the ambient index set is stated locally.\\
\commandname{TopologicalSupportOf}\texttt{\{f\}} & $\TopologicalSupportOf{f}$ & Closure
of the pointwise support.\\
\commandname{Convolution} & $f\Convolution g$ & Convolution, with the underlying group and
measure stated locally.\\
\commandname{MetricDistance} & $\MetricDistance(x,y)$ & Distance of the named metric space.\\
\commandname{TotalVariationOf}\texttt{\{a\}} & $\TotalVariationOf{a}$ & Total
variation of one typed function, signed measure, or other locally declared object.\\
\commandname{TotalVariationDistance} & $\TotalVariationDistance(\mu,\nu)$ & Total
variation distance with the probability convention $\sup_A|\mu(A)-\nu(A)|$.\\
\commandname{Probability} & $\Probability$ & Probability measure.\\
\commandname{Expectation} & $\Expectation$ & Expectation with respect to the stated law.\\
\commandname{Variance} & $\Variance$ & Variance $\Expectation[(X-\Expectation X)^2]$.\\
\commandname{Covariance} & $\Covariance$ & Covariance of explicitly stated random
variables.\\
\commandname{LawOperator} & $\LawOperator$ & Pushforward-law operator token; it is not a
metric, density, or samplewise map.\\
\commandname{LawOf}\texttt{\{X\}} & $\LawOf{X}$ & Pushforward probability law of $X$.\\
\commandname{EqualInLaw} & $X\EqualInLaw Y$ & Equality of probability laws, not
pointwise equality.\\
\commandname{ConvergesInLaw} & $X_n\ConvergesInLaw X$ & Weak convergence of laws.\\
\commandname{DyadicSigmaField} & $\DyadicSigmaField_n$ & Level-$n$ dyadic filtration
$\sigma$-field; the generating random coordinates and base probability space are printed.\\
\commandname{DecayOptimizationObjective}\texttt{\{n\}} &
$\DecayOptimizationObjective{n}$ & Exact level-$n$ nonlinear objective in the
Fourier-decay saddle problem, with variables and feasible set declared locally.\\
\commandname{LinearizedDecayObjective}\texttt{\{n\}} &
$\LinearizedDecayObjective{n}$ & The explicitly linearized level-$n$ decay objective;
it is not identified with the nonlinear objective without an equality theorem.\\
\commandname{FourierTwoPiOperator} & $\FourierTwoPiOperator$ & Forward Fourier operator
with kernel $\EulerE^{-2\pi\ImaginaryUnit x\xi}$.\\
\commandname{FourierAngularOperator} & $\FourierAngularOperator$ & Forward Fourier
operator with kernel $\EulerE^{-\ImaginaryUnit x\omega}$.\\
\commandname{InverseFourierTwoPiOperator} & $\InverseFourierTwoPiOperator$ & Inverse of
$\FourierTwoPiOperator$ on a named inversion domain.\\
\commandname{InverseFourierAngularOperator} & $\InverseFourierAngularOperator$ & Inverse
of $\FourierAngularOperator$, including the $1/(2\pi)$ measure factor.\\
\commandname{FourierTwoPiTransformOf}\texttt{\{f\}} &
$\FourierTwoPiTransformOf{f}$ & Operator-application form of the $2\pi$-frequency
Fourier transform.\\
\commandname{FourierAngularTransformOf}\texttt{\{f\}} &
$\FourierAngularTransformOf{f}$ & Operator-application form of the angular-frequency
Fourier transform.\\
\commandname{FourierTwoPi}\texttt{\{f\}} & $\FourierTwoPi{f}$ & Cycles-per-unit Fourier
transform with kernel $\EulerE^{-2\pi\ImaginaryUnit x\xi}$.\\
\commandname{FourierAngular}\texttt{\{f\}} & $\FourierAngular{f}$ & Angular-frequency
Fourier transform with kernel $\EulerE^{-\ImaginaryUnit x\omega}$.\\
\commandname{FourierSeriesCoefficient}\texttt{\{h\}\{k\}} &
$\FourierSeriesCoefficient{h}{k}$ & The $k$th coefficient of a one-periodic function,
with kernel $\EulerE^{-2\pi\ImaginaryUnit kx}$ on $[0,1]$.\\
\commandname{FourierSeriesCoefficientAtPeriod}\texttt{\{h\}\{k\}\{T\}} &
$\FourierSeriesCoefficientAtPeriod{h}{k}{T}$ & The normalized $k$th coefficient of
an explicitly $T$-periodic function; the period is part of the symbol.\\
\commandname{LaplaceTransformSymbol} & $\LaplaceTransformSymbol$ & One-sided Laplace
operator token unless a bilateral subscript is printed.\\
\commandname{MellinTransformSymbol} & $\MellinTransformSymbol$ & Mellin-transform
operator token on a stated strip of convergence.\\
\commandname{LaplaceTransformOf}\texttt{\{f\}} & $\LaplaceTransformOf{f}$ & Laplace
transform; one-sided or bilateral domain must be printed.\\
\commandname{MellinTransformOf}\texttt{\{f\}} & $\MellinTransformOf{f}$ & Mellin
transform $\int_0^\infty x^{s-1}f(x)\,\Differential x$ on its stated strip.\\
\commandname{MomentGeneratingFunction}\texttt{\{X\}} & $\MomentGeneratingFunction{X}$ &
$z\mapsto\Expectation[\EulerE^{zX}]$ on its stated domain.\\
\commandname{CharacteristicFunction}\texttt{\{X\}} & $\CharacteristicFunction{X}$ &
$t\mapsto\Expectation[\EulerE^{\ImaginaryUnit tX}]$; a declared probability measure
$\mu$ may replace $X$, meaning its positive-sign Fourier--Stieltjes transform.\\
\commandname{SincRad} & $\SincRad$ & Radian sinc, $\sin t/t$, totalized to one at zero.\\
\commandname{SincPi} & $\SincPi$ & Normalized sinc, $\sin(\pi x)/(\pi x)$, totalized
to one at zero.\\
\commandname{BinaryDigitSum} & $\BinaryDigitSum$ & Number of one-bits in a nonnegative
integer.\\
\commandname{ThueMorseSignSymbol} & $\ThueMorseSignSymbol$ & The symbol reserved for the
Thue--Morse sign sequence; use the argument-taking command in formulas.\\
\commandname{ThueMorseSign}\texttt{\{n\}} & $\ThueMorseSign{n}$ &
$(-1)^{\BinaryDigitSum(n)}$.\\
\commandname{PadicValuation}\texttt{\{p\}} & $\PadicValuation{p}$ & $p$-adic valuation;
the argument and zero convention are stated locally.\\
\commandname{TwoAdicValuation} & $\TwoAdicValuation$ & Two-adic valuation.\\
\commandname{QPochhammer}\texttt{\{a\}\{q\}\{n\}} & $\QPochhammer{a}{q}{n}$ & Finite
or infinite $q$-shifted factorial with explicit base and length.\\
\commandname{GaussianBinomial}\texttt{\{n\}\{k\}\{q\}} &
$\GaussianBinomial{n}{k}{q}$ & Gaussian binomial coefficient with explicit base.\\
\commandname{QInteger}\texttt{\{n\}\{q\}} & $\QInteger{n}{q}$ & $q$-integer with
explicit base.\\
\commandname{Polylogarithm} & $\Polylogarithm$ & Polylogarithm with order supplied as a
subscript.\\
\commandname{LambertW} & $\LambertW$ & Lambert $W$; a branch subscript is mandatory when
more than one branch is possible.\\
\commandname{WrightOmega} & $\WrightOmega$ & Wright omega, the inverse of
\(y\mapsto y+\log y\) on the declared real or complex branch.\\
\commandname{LambertPolynomial}\texttt{\{n\}} &
\(\LambertPolynomial{n}\) & Canonical zero-based Lambert polynomial;
\(\LambertPolynomial{0}(u)=-u\), and the evaluation variable follows outside
the one-argument command.\\
\commandname{ScaledLambertPolynomial}\texttt{\{n\}} &
\(\ScaledLambertPolynomial{n}\) & Factorially scaled Lambert polynomial
\(n!\LambertPolynomial{n}\) for \(n\ge1\); the evaluation variable follows the
command.\\
\commandname{MultiplicativeInverseOf}\texttt{\{F\}} &
\(\MultiplicativeInverseOf{F}\) & Inverse of a unit \(F\) under multiplication.\\
\commandname{CompositionalInverseOf}\texttt{\{F\}} &
\(\CompositionalInverseOf{F}\) & Inverse of an admissible map \(F\) under
composition; distinct from a reciprocal and a time-\((-1)\) iterate.\\
\commandname{PolynomialLogPowerSeriesRing}\texttt{\{K\}} &
\(\PolynomialLogPowerSeriesRing{K}\) & Polynomial-log power-series ring
\(K[L][[t]]\), with \(t=X^{-1}\) and \(L=\log X\).\\
\commandname{PolynomialLogLaurentRing}\texttt{\{K\}} &
\(\PolynomialLogLaurentRing{K}\) & Polynomial-log Laurent ring
\(K[L]((t))\), generally not a field.\\
\commandname{RationalLogLaurentField}\texttt{\{K\}} &
\(\RationalLogLaurentField{K}\) & Laurent-series field \(K(L)((t))\)
with exact rational dependence on \(L\).\\
\commandname{IteratedLaurentLogField}\texttt{\{K\}} &
\(\IteratedLaurentLogField{K}\) & Iterated Laurent field
\(K((L^{-1}))((t))\), with the completion order shown.\\
\commandname{NearIdentityPolynomialLogGroup}\texttt{\{K\}} &
\(\NearIdentityPolynomialLogGroup{K}\) & Near-identity composition group
\(\SetOf{X+A:A\in K[L][[t]]}\), subject to admissibility.\\
\commandname{GammaFunction} & $\GammaFunction$ & Euler Gamma function.\\
\commandname{RiemannZeta} & $\RiemannZeta$ & Riemann zeta function or its meromorphic
continuation, as stated.\\
\commandname{HurwitzZeta}\texttt{\{s\}\{a\}} & \(\HurwitzZeta{s}{a}\) &
Hurwitz zeta function with order \(s\) and shift \(a\) both explicit.\\
\commandname{BigO} (operator token; 0 TeX arguments) & $\BigO{g}$ & Big-$O$ applied to
the following mathematical operand $g$, in a limit declared in the same sentence.\\
\commandname{LittleO} (operator token; 0 TeX arguments) & $\LittleO{g}$ & Little-$o$
applied to the following mathematical operand $g$, in a limit declared in the same
sentence.\\
\commandname{BigOAt}\texttt{\{g\}\{a\}} & $\BigOAt{g}{a}$ & Big-$O$ with limiting
regime $a$ printed.\\
\commandname{LittleOAt}\texttt{\{g\}\{a\}} & $\LittleOAt{g}{a}$ & Little-$o$ with
limiting regime $a$ printed.\\
\commandname{FormalBigO} (operator token; 0 TeX arguments) &
\(\FormalBigO\) & Formal valuation-tail operator; it is not an analytic
estimate and, in a Laurent ring, not an ideal-quotient congruence.\\
\commandname{FormalBigOAt}\texttt{\{t\string^N\}\{t\}} &
\(\FormalBigOAt{t^N}{t}\) & Formal tail whose omitted difference has
\(t\)-valuation at least \(N\).\\
\bottomrule
\end{longtable}
\endgroup

\section{Foundational mathematical grammar}

\begin{catalogueentry}{contract-version}{Shared notation contract version}
\entryfield{Command}{\commandname{FabiusNotationVersion} expands to the ISO date
\FabiusNotationVersion.}
\entryfield{Meaning}{It identifies the semantic API described by this catalogue.}
\entryfield{Contract file}{The corresponding source is
\modulepath{docs/fabius-notation.tex}.  The version is not a publication date, source
revision, theorem date, numeric parameter, or claim that every historical source was
authored on that date.}
\entryfield{Update rule}{Any addition, removal, arity change, or mathematical change to
a shared command updates this date and the quick reference, alphabetical index, rendered
symbol index, retired-alias mapping, and relevant detailed entry in the same change.}
\end{catalogueentry}

\begin{catalogueentry}{elementary-constants}{The exponential base and imaginary unit}
\entryfield{Exponential base}{$\EulerE\in\RealNumbers$ is the unique positive number
with $\log\EulerE=1$.  The expression $\EulerE^z$ is the complex exponential for
$z\in\ComplexNumbers$, not a local error variable raised to a power.}
\entryfield{Imaginary unit}{$\ImaginaryUnit\in\ComplexNumbers$ is the root of
$z^2=-1$ with positive imaginary part.  Thus $\ImaginaryUnit^2=-1$ and
$\overline{\ImaginaryUnit}=-\ImaginaryUnit$.  The other root is
$-\ImaginaryUnit$.}
\entryfield{Typography}{Both constants are upright.  An italic $e$ or $i$ remains a
local variable or index and acquires no constant meaning from context alone.}
\end{catalogueentry}

\begin{catalogueentry}{number-systems}{Number systems and inclusions}
\entryfield{Canonical symbols}{
$\NaturalNumbers=\{0,1,2,\ldots\}$,
$\PositiveIntegers=\{1,2,3,\ldots\}$,
$\IntegerNumbers$, $\RationalNumbers$, $\RealNumbers$, $\ComplexNumbers$, and
$\ExtendedRealNumbers=\RealNumbers\cup\{-\infty,+\infty\}$.}
\entryfield{Inclusions}{
$\NaturalNumbers\subset\IntegerNumbers\subset\RationalNumbers\subset
\RealNumbers\subset\ComplexNumbers$.  These are the standard injective embeddings; a
formula involving coercion into a larger system is interpreted only after the target has
been named.}
\entryfield{Endpoint convention}{Neither $+\infty$ nor $-\infty$ is a real number.  They
belong only to $\ExtendedRealNumbers$.  Arithmetic involving them follows extended-order
conventions only when explicitly stated.}
\entryfield{Units of a ring or field}{For a ring \(K\) with identity,
\(K^\times=\{a\in K:\text{there exists }b\in K\text{ with }ab=ba=1\}\) is its
multiplicative group of units.  If \(K\) is a field, then
\(K^\times=K\setminus\{0\}\).  The superscript \(\times\) means neither Cartesian
product nor a nonzero natural power in this context.}
\entryfield{Retired ambiguity}{A bare $\mathbb N$ is forbidden in normative prose because
authors disagree about whether it contains zero.}
\end{catalogueentry}

\begin{catalogueentry}{metavariables}{Default metavariables and mandatory overrides}
\entryfield{Discrete indices}{Unless an entry says otherwise,
$j,k,m,n,r,N\in\NaturalNumbers$.  A statement requiring positivity writes
$n\in\PositiveIntegers$ or $n\ge1$.  The letter $p$ denotes a prime only when the same
sentence says ``prime''.}
\entryfield{Real variables}{The letters $x,y,t,u,s,\xi,\omega$ have no global type.  Each
definition declares whether they lie in $\RealNumbers$, $\ComplexNumbers$, an interval, or
a measure space.}
\entryfield{Parameters}{The letters $a,b,q,\alpha,\beta,\lambda$ are never assigned a
default domain.  Their admissible set and all inequalities, such as $\AbsoluteValue{q}<1$
or $\alpha>0$, are part of the declaration.}
\entryfield{Functions and measures}{Lower-case $f,g,h$ denote local functions only after a
domain and codomain are supplied.  Greek $\mu,\nu$ denote local measures only after a
measurable space is supplied.  No letter acquires mathematical meaning merely from its
font.}
\end{catalogueentry}

\begin{catalogueentry}{definition-equality}{Definition, equality, equivalence, and
identification}
\entryfield{Definition}{The symbol $a\coloneqq b$ introduces $a$ as an abbreviation or new
object whose value is exactly $b$.  It is not an asymptotic statement.}
\entryfield{Equality}{The symbol $=$ means equality in the ambient type: equality of
scalars, functions, sets, measures, matrices, or equivalence classes as determined by the
declared type.  Function equality is extensional; equality almost everywhere is written
separately.}
\entryfield{Almost-everywhere equality}{For a measure $\mu$, $f=g$ $\mu$-a.e. means
$\mu(\{x:f(x)\ne g(x)\})=0$.  It does not license substitution into a pointwise endpoint
formula without a continuous-representative theorem.}
\entryfield{Law equality}{For random variables $X$ and $Y$, $X\EqualInLaw Y$ abbreviates
$\LawOf{X}=\LawOf{Y}$ and does not imply $X(\omega)=Y(\omega)$ on a shared sample space.}
\entryfield{Isomorphism}{Symbols $\cong$ and $\simeq$ require the category---sets, groups,
rings, topological spaces, normed spaces, or probability spaces---to be stated locally.}
\end{catalogueentry}

\begin{catalogueentry}{sets-intervals}{Sets, set operations, and intervals}
\entryfield{Set grammar}{For a universe $U$, $x\in A$ means membership in $A\subseteq U$;
$A\subseteq B$ permits equality and $A\subsetneq B$ does not.  Union, intersection,
complement, and set difference are $A\cup B$, $A\cap B$, $U\setminus A$, and $A\setminus
B$.}
\entryfield{Intervals}{For $a,b\in\RealNumbers$ with $a\le b$,
$[a,b]=\{x:a\le x\le b\}$, $(a,b)=\{x:a<x<b\}$,
$[a,b)=\{x:a\le x<b\}$, and $(a,b]=\{x:a<x\le b\}$.}
\entryfield{Empty reversed intervals}{If $b<a$, each ordered interval above is empty.  An
oriented integral $\int_a^b$ is not an integral over that empty set; it changes sign under
endpoint reversal.}
\entryfield{Lean aliases}{The Lean set constructors are \lean{Set.Icc}, \lean{Set.Ioo},
\lean{Set.Ico}, and \lean{Set.Ioc}, in the same order as the four printed intervals.}
\end{catalogueentry}

\begin{catalogueentry}{maps}{Maps, restrictions, composition, images, and preimages}
\entryfield{Typed map}{The declaration $f:A\to B$ includes domain $A$ and codomain $B$.
The restriction to $S\subseteq A$ is $f|_S:S\to B$.}
\entryfield{Composition}{For $f:A\to B$ and $g:B\to C$,
$(g\circ f)(x)=g(f(x))$.  Composition is evaluated right to left.}
\entryfield{Identity map}{For a declared set or type $X$,
$\IdentityOperator_X:X\to X$ is the map $x\mapsto x$.  The zero-argument command
\commandname{IdentityOperator} supplies only the operator token; the domain subscript
$X$ is source notation and must be printed whenever the surrounding typed equation does
not determine it uniquely.  Thus $f\circ\IdentityOperator_X=f$ and
$\IdentityOperator_Y\circ f=f$ for every $f:X\to Y$.}
\entryfield{Image and preimage}{For $S\subseteq A$ and $T\subseteq B$,
$f(S)=\{f(x):x\in S\}$ and $f^{-1}(T)=\{x\in A:f(x)\in T\}$.  The latter does not
require $f$ to be invertible.}
\entryfield{Inverse map}{The symbol $f^{-1}:B\to A$ denotes an inverse function only after
$f$ has been declared bijective.  A superscript $-1$ on a set image is always the preimage,
not reciprocal arithmetic.}
\end{catalogueentry}

\begin{catalogueentry}{semantic-delimiters}{Set, inner-product, cardinality, and
definition delimiters}
\entryfield{Set constructor}{The command \commandname{SetOf}\texttt{\{body\}} supplies
only extensible braces.  The body supplies either an explicit finite list or a typed
predicate, for example
$\SetOf{n\in\NaturalNumbers:n<N}$.  A colon and a vertical bar have the same
set-builder meaning only after the ambient type has been printed.  The command does not
deduplicate a list, choose a universe, or turn a class into a set.}
\entryfield{Inner product}{The command
\commandname{InnerProduct}\texttt{\{x,y\}} supplies angle delimiters.  In this catalogue
$\InnerProduct{x,y}_{H}$ is linear in $x$ and conjugate-linear in $y$ on the complex
Hilbert space $H$; a source adopting the opposite convention must say so at the first
use and conjugate all adjoint formulas consistently.  A duality pairing is instead
labelled $\langle\phi,x\rangle_{X^\ast,X}$.}
\entryfield{Cardinality}{For a finite set $A$,
$\CardinalityOf{A}\in\NaturalNumbers$ is its number of elements and
$\CardinalityOf{\varnothing}=0$.  For an infinite set it may denote a cardinal number
only when the ambient set theory is explicit.  Length, Lebesgue measure, probability,
determinant, absolute value, and norm are never inferred from the same vertical bars.}
\entryfield{Definitional equality}{$a\DefinitionEquals b$ introduces the symbol $a$
with type and value supplied by $b$.  Reversing the direction uses
$b\mathrel{\eqqcolon}a$; a theorem already requiring proof uses ordinary equality.
Definitional equality in prose is not a claim of Lean kernel reduction unless an exact
Lean definition is cited.}
\entryfield{Empty and singleton cases}{$\SetOf{}=\varnothing$ when the body is an
explicit empty list, whereas $\SetOf{x}$ is a singleton.  A displayed set-builder with
an unsatisfiable predicate is empty by its predicate, not by typography.}
\end{catalogueentry}

\begin{catalogueentry}{families}{Sequences, families, ranges, and finite indexing}
\entryfield{Sequence}{A sequence $(a_n)_{n\in\NaturalNumbers}$ is a map
$a:\NaturalNumbers\to A$ into a declared codomain $A$.  A finite row
$(a_{n,k})_{0\le k\le n}$ is defined only for the printed triangular range.}
\entryfield{Half-open range}{The project and Lean frequently use $j<n$ to mean
$j\in\{0,1,\ldots,n-1\}$.  Thus $\sum_{j<n}$ and $\prod_{j<n}$ are half-open at $n$.}
\entryfield{Inclusive range}{The notation $0\le j\le n$ includes both endpoints and has
$n+1$ indices.  It is never silently identified with $j<n$.}
\entryfield{Indexed family}{For a possibly infinite set $I$, $(a_i)_{i\in I}$ is a family
indexed by $I$.  Summability over $I$ is unconditional unless an order is separately
specified.}
\end{catalogueentry}

\begin{catalogueentry}{sums-products}{Sums, products, and neutral-element conventions}
\entryfield{Finite operations}{For a finite set $S$, $\sum_{i\in S}a_i$ and
$\prod_{i\in S}a_i$ are taken in the stated additive and multiplicative structures.}
\entryfield{Empty cases}{The empty sum is $0$ and the empty product is $1$.  Consequently
$\sum_{j<n}$ is zero and $\prod_{j<n}$ is one when $n=0$.}
\entryfield{Infinite series}{The notation $\sum_{n=0}^{\infty}a_n$ includes an ordering by
$\NaturalNumbers$.  For a family over an arbitrary type, the project uses unconditional
summability; rearrangement requires absolute or unconditional convergence.}
\entryfield{Infinite products}{The product $\prod_{n=0}^{\infty}a_n$ is the limit of the
ordered partial products $\prod_{n<N}a_n$ as $N\to\infty$.  Local uniform convergence is
stated when the factors depend on a complex variable.}
\end{catalogueentry}

\begin{catalogueentry}{powers-factorials}{Powers, factorials, and binomial coefficients}
\entryfield{Natural powers}{For an element $x$ of a monoid and
$n\in\NaturalNumbers$, $x^0=1$ and $x^{n+1}=x^n x$.  In particular, the natural-power
convention is $0^0=1$.}
\entryfield{Integer powers}{For $n\in\IntegerNumbers$, $x^n$ requires an invertible
nonzero $x$ when $n<0$; the ambient field may totalize $0^{-1}$ algebraically, but an
analytic formula never relies on that totalization without saying so.}
\entryfield{Complex powers}{For $z^\alpha$ with noninteger $\alpha$, the chosen branch of
$\log z$ and its cut are mandatory.  The principal convention is
$z^\alpha=\exp(\alpha\PrincipalLogarithm z)$ with
$\arg z\in(-\pi,\pi]$.}
\entryfield{Factorial}{For $n\in\NaturalNumbers$, $0!=1$ and
$n!=\prod_{1\le j\le n}j$.  The odd double factorial satisfies
$(-1)!!=1$ and $(2n+1)!!=\prod_{j=0}^{n}(2j+1)$.}
\entryfield{Binomial coefficient}{For $n,k\in\NaturalNumbers$,
$\binom nk=n!/(k!(n-k)!)$ when $k\le n$ and $\binom nk=0$ when $k>n$.}
\entryfield{Falling factorial}{For a scalar $x$ and $n\in\NaturalNumbers$,
$\FallingFactorial{x}{n}\DefinitionEquals\prod_{j=0}^{n-1}(x-j)$; the empty product
gives $\FallingFactorial{x}{0}=1$.  Thus
$\FallingFactorial{x}{n}=\GammaFunction(x+1)/\GammaFunction(x-n+1)$ only where the
quotient is defined or after a separately justified continuation.}
\entryfield{Rising factorial}{For a scalar $x$ and $n\in\NaturalNumbers$,
$\RisingFactorial{x}{n}\DefinitionEquals\prod_{j=0}^{n-1}(x+j)$ and
$\RisingFactorial{x}{0}=1$.  It equals the ordinary Pochhammer symbol $(x)_n$ in this
catalogue; the base-$q$ analogue is always \commandname{QPochhammer} and never shares
this two-argument syntax.}
\end{catalogueentry}

\begin{catalogueentry}{logarithm-order}{The principal logarithm and order operators}
\entryfield{Principal logarithm}{The command \commandname{PrincipalLogarithm} denotes
the analytic branch
\[
 \PrincipalLogarithm z=\log|z|+\ImaginaryUnit\operatorname{Arg}z,
 \qquad -\pi<\operatorname{Arg}z\le\pi,
\]
on $\ComplexNumbers\setminus(-\infty,0]$.  In particular it is undefined at zero and
has distinct boundary values on the two sides of the negative real axis.  The real
symbol $\log x$ for $x>0$ is its restriction and carries no branch ambiguity.}
\entryfield{Complex powers}{The principal power is
$z^\alpha=\exp(\alpha\PrincipalLogarithm z)$ on the same cut domain.  Algebraic laws
such as $(zw)^\alpha=z^\alpha w^\alpha$ are not used across a branch crossing without
an explicit argument calculation.}
\entryfield{Vanishing order}{For a nonzero holomorphic $f$ near $z_0$,
$\OrderOperator_{z_0}(f)=m\in\NaturalNumbers$ when
$f(z)=(z-z_0)^m g(z)$ with $g(z_0)\ne0$.  For a pole, a signed divisor order is written
$\OrderOperator_{z_0}^{\mathrm{div}}(f)<0$ so that it cannot be mistaken for the
nonnegative zero order.}
\entryfield{Initial \(q\)-order}{For a nonzero formal Laurent series
\(f(q)=\sum_{n\ge n_0}c_nq^n\in R((q))\), where \(R\) is an integral domain and
\(c_{n_0}\ne0\),
\[
 \OrderOperator_{q=0}(f)\DefinitionEquals n_0.
\]
For \(f=0\) the order is \(+\infty\) when the extended convention is explicitly in
force; otherwise it is left undefined.  This subscripted operator records the expansion
variable and base point, and is not the \(2\)-adic valuation
\(\TwoAdicValuation(\,\cdot\,)\) or a group-element order.}
\entryfield{Multiplicative order}{For a group element $a$, put
$S_G(a)=\SetOf{n\in\PositiveIntegers:a^n=e_G}$.  When this set is nonempty,
\[
 \OrderOperator_G(a)=\min S_G(a).
\]
When it is empty, $\OrderOperator_G(a)=\infty$.  The group subscript is mandatory
whenever an analytic order also occurs.}
\entryfield{Other orders}{Differential-equation order, recurrence order, matrix order,
and asymptotic truncation order are prose-qualified quantities.  The shared operator
provides typography, not permission to leave the species or base point implicit.}
\end{catalogueentry}

\begin{catalogueentry}{scalar-size}{Absolute value, norms, and distance}
\entryfield{Absolute value}{For $x\in\RealNumbers$, $\AbsoluteValue{x}=\max\{x,-x\}$.
For $z\in\ComplexNumbers$, $\AbsoluteValue{z}=\sqrt{z\overline z}$.}
\entryfield{Norm}{The expression $\Norm{x}$ is meaningful only after a normed space has
been named.  For a function, $\Norm{f}_{\infty}$ means an essential supremum or a pointwise
supremum exactly as the surrounding declaration states.}
\entryfield{Distance}{In a metric space $(X,d)$,
$\MetricDistance(x,y)=d(x,y)$.  Probability metrics such as total variation and
Wasserstein distance are given their own decorated symbols and never use a bare
$\MetricDistance$ without the metric name.}
\end{catalogueentry}

\begin{catalogueentry}{rounding-positive-part}{Floor, ceiling, fractional part, sign, and
positive part}
\entryfield{Definitions}{For $x\in\RealNumbers$,
$\Floor{x}=\max\{n\in\IntegerNumbers:n\le x\}$,
$\Ceiling{x}=\min\{n\in\IntegerNumbers:x\le n\}$,
$\FractionalPart{x}=x-\Floor{x}\in[0,1)$,
$\PositivePart{x}=\max\{x,0\}$, and
$\Sign(x)=-1,0,1$ according as $x<0$, $x=0$, or $x>0$.}
\entryfield{Nearest integer}{A nearest-integer or nearest-dyadic operation must separately
state its tie-breaking rule; neither floor nor ceiling implies one.}
\entryfield{Use in powers}{For $\alpha>0$, $\PositivePart{x}^{\alpha}$ is zero when
$x\le0$.  For $\alpha=0$, the natural-power value at $x\le0$ is one only if a formula
explicitly uses natural exponentiation; real powers require a separate endpoint rule.}
\end{catalogueentry}

\begin{catalogueentry}{indicator}{Indicator and half-weighted interval indicator}
\entryfield{Ordinary indicator}{For a set $A$ in a universe $U$,
$\IndicatorOf{A}:U\to\{0,1\}$ is one on $A$ and zero on $U\setminus A$.}
\entryfield{Half-weighted indicator}{When adjacent closed intervals meet, the project uses
a separately named half-endpoint indicator taking value $1/2$ at each boundary and $1$ in
the interior.  It is not denoted by the ordinary $\IndicatorOf{A}$.}
\entryfield{Measure-theoretic caveat}{Changing endpoint values does not change a Lebesgue
integral, but it does change pointwise finite-spline identities and exact values at grid
points.}
\end{catalogueentry}

\begin{catalogueentry}{support}{Pointwise support, topological support, and essential
support}
\entryfield{Pointwise support}{For a function $f:X\to A$ with zero in $A$,
$\SupportOf{f}=\SetOf{x\in X:f(x)\ne0}$.  The bare command
\commandname{SupportOperator} is the same operator token before an explicitly supplied
operand; \commandname{SupportOf} supplies the parentheses and is preferred in prose.}
\entryfield{Topological support}{For topological $X$,
$\TopologicalSupportOf{f}=\overline{\SupportOf{f}}$.  Thus a boundary point can belong to
the topological support even when $f$ and all of its derivatives vanish there.}
\entryfield{Measure support}{For a Borel measure $\mu$, $\SupportOperator\mu$ is the
complement of the largest open $\mu$-null set.  This is not obtained by evaluating a
measure as a pointwise function.}
\entryfield{Essential support}{Essential support depends on a background measure and ignores
null changes; it must be labeled ``essential''.}
\end{catalogueentry}

\begin{catalogueentry}{calculus}{Derivatives, iterated derivatives, and integrals}
\entryfield{Derivative}{For $f:\RealNumbers\to E$ into a real normed space,
$f'(x)$ denotes the Fr\'echet derivative evaluated on $1$ when it exists.  Higher
derivatives are $f^{(n)}$, with $f^{(0)}=f$.}
\entryfield{Partial derivatives}{For a multivariable function, $\partial_j f$ differentiates
in the $j$th coordinate.  A gradient, Jacobian, Hessian, or directional derivative is named
explicitly and includes its inner-product or coordinate convention.}
\entryfield{Differential}{In $\int_a^b f(x)\,\Differential x$, the upright
$\Differential x$ identifies the variable of integration and is not a free multiplicative
variable.}
\entryfield{Oriented interval integral}{For integrable $f$,
$\int_b^a f(x)\,\Differential x=-\int_a^b f(x)\,\Differential x$ and
$\int_a^a f(x)\,\Differential x=0$.}
\entryfield{Whole-space integral}{An integral $\int_{\RealNumbers}$ is with respect to
Lebesgue measure unless another measure is printed.  An expectation is written
$\Expectation$, not as an unlabeled integral.}
\end{catalogueentry}

\begin{catalogueentry}{convolution}{Convolution of functions, measures, and sequences}
\entryfield{Functions}{For integrable functions $f,g:\RealNumbers\to\ComplexNumbers$,
$(f\Convolution g)(x)=\int_{\RealNumbers}f(x-y)g(y)\,\Differential y$.}
\entryfield{Measures}{For finite Borel measures $\mu,\nu$ on $\RealNumbers$,
$(\mu\Convolution\nu)(A)=\int\!\int\IndicatorOf{A}(x+y)\,\mu(\Differential x)
\nu(\Differential y)$.  It is the law of $X+Y$ for independent $X,Y$ with laws
$\mu,\nu$.}
\entryfield{Sequences}{For sequences indexed by $\NaturalNumbers$, the Cauchy convolution
is $(a\Convolution b)_n=\sum_{k=0}^{n}a_kb_{n-k}$.  It is distinct from bilateral
convolution over $\IntegerNumbers$.}
\entryfield{Power}{The notation $f^{\Convolution n}$ means $n$-fold convolution, with
$f^{\Convolution0}=\delta_0$ for measures and the appropriate convolution unit for the
ambient algebra.}
\end{catalogueentry}

\section{The semantic Fabius--Rvachev objects}

\begin{catalogueentry}{fabius-bounded}{The bounded Fabius function
$\FabiusBounded$}
\entryfield{Command}{\commandname{FabiusBounded}.}
\entryfield{Type}{A smooth function $\FabiusBounded:\RealNumbers\to[0,1]\subset
\RealNumbers$.}
\entryfield{Normalization}{
$\FabiusBounded(x)=0$ for $x\le0$,
$\FabiusBounded(x)=1$ for $x\ge1$,
$\FabiusBounded(1-x)=1-\FabiusBounded(x)$ for every real $x$, and
$\FabiusBounded'(x)=2\FabiusBounded(2x)$ on the left half of the unit interval.}
\entryfield{Shape}{Its restriction $\FabiusBounded|_{[0,1]}$ is continuous and strictly
increasing from zero to one.  Its constant tails are part of the object, not an informal
extension.}
\entryfield{Lean crosswalk}{The real-valued object is \lean{fabiusReal}; the abstract
specification is \lean{IsFabius}; construction and uniqueness are exposed through
\lean{fabius} and \lean{fabius_spec}.  Status: \textbf{Lean-proved}.}
\entryfield{Forbidden aliases}{The bare capital letter, the former calligraphic glyph,
and the former decorated-capital glyph are not canonical because the active historical
corpus also uses them for transforms, families, and the signed global extension.  This
catalogue names those aliases textually rather than re-emitting their retired control
sequences.}
\end{catalogueentry}

\begin{catalogueentry}{rvachev-up}{Rvachev's up-function $\RvachevUp$}
\entryfield{Command}{\commandname{RvachevUp}.}
\entryfield{Type and definition}{The function $\RvachevUp:\RealNumbers\to\RealNumbers$ is
defined by
\[
 \RvachevUp(x)=\FabiusBounded(1-\AbsoluteValue{x}).
\]
This single formula uses the constant tails of $\FabiusBounded$ and therefore gives
$\RvachevUp(x)=0$ whenever $\AbsoluteValue{x}\ge1$.}
\entryfield{Normalization}{It is even, nonnegative, smooth, has integral one, satisfies
$\RvachevUp(0)=1$, and has topological support $[-1,1]$.  Every derivative vanishes at
$x=\pm1$.}
\entryfield{Differential law}{For every $x\in\RealNumbers$,
$\RvachevUp'(x)=2(\RvachevUp(2x+1)-\RvachevUp(2x-1))$.}
\entryfield{Lean crosswalk}{\lean{rvachevUp} in \modulepath{Basic.lean}, with differential,
support, Fourier, and shape theorems in their dedicated modules.  Status:
\textbf{Lean-proved}.}
\entryfield{Historical aliases}{The source papers print an upright word ``up'', use
legacy commands named \commandname{up} or \commandname{Up}, or use the local glyph
$\varphi$.  These spellings are historical evidence, not commands to be re-emitted in
project-authored formulas; normative prose uses only $\RvachevUp$.}
\end{catalogueentry}

\begin{catalogueentry}{fabius-global}{The signed global extension $\FabiusGlobal$}
\entryfield{Command}{\commandname{FabiusGlobal}.}
\entryfield{Type}{A smooth function $\FabiusGlobal:\RealNumbers\to\RealNumbers$.  It is not
a cumulative distribution function and is not constrained to $[0,1]$.}
\entryfield{Canonical translate formula}{For $x\in\RealNumbers$,
\[
 \FabiusGlobal(x)=\sum_{n=0}^{\infty}
   \ThueMorseSign{n}\,\RvachevUp(x-2n-1).
\]
At each fixed $x$, compact support leaves only finitely many nonzero summands; hence this
formula is pointwise locally finite even when written as an infinite sum.}
\entryfield{Identifications}{For $x\in[0,1]$,
$\FabiusGlobal(x)=\FabiusBounded(x)$.  Globally,
$\FabiusGlobal'(x)=2\FabiusGlobal(2x)$.  Repeated differentiation gives
$\FabiusGlobal^{(m)}(x)=2^{\binom{m+1}{2}}\FabiusGlobal(2^m x)$ for
$m\in\NaturalNumbers$.}
\entryfield{Lean declaration}{\lean{extendedFabius}.}
\entryfield{Lean module and status}{Translate-cell and calculus results are in
\modulepath{GlobalExtension.lean}.  Status: \textbf{Lean-proved}.}
\entryfield{Retired aliases}{The former calligraphic glyph, the former decorated-capital
glyph, and the plain capital glyph of the 2017 arithmetic paper are historical aliases.
They must be translated semantically, never by a blind global replacement.}
\end{catalogueentry}

\begin{catalogueentry}{fabius-quantile}{The ordinary and clamped Fabius inverses}
\entryfield{Commands}{\commandname{FabiusQuantile} and
\commandname{FabiusClampedQuantile}.}
\entryfield{Ordinary inverse}{Because
$\FabiusBounded|_{[0,1]}:[0,1]\to[0,1]$ is a strictly increasing bijection,
\[
 \FabiusQuantile:[0,1]\to[0,1],\qquad
 \FabiusQuantile(y)=\bigl(\FabiusBounded|_{[0,1]}\bigr)^{-1}(y).
\]
Thus $\FabiusBounded(\FabiusQuantile(y))=y$ for $y\in[0,1]$ and
$\FabiusQuantile(\FabiusBounded(x))=x$ for $x\in[0,1]$.}
\entryfield{Clamped inverse}{For $y\in\RealNumbers$, define
$\operatorname{clamp}(y)=\min\{1,\max\{0,y\}\}$ and
\[
 \FabiusClampedQuantile(y)=\FabiusQuantile(\operatorname{clamp}(y)).
\]
Consequently it is zero on $(-\infty,0]$ and one on $[1,\infty)$.}
\entryfield{Lean crosswalk}{\lean{fabiusInv} and \lean{fabiusOrderIso} in
\modulepath{FabiusInverse.lean}.  Status: \textbf{Lean-proved}.}
\entryfield{Retired aliases}{Bare $I$, $G$, $J$, the formal inverse expression, and a
legacy inverse macro occur in older prose.  They are local or historical only.}
\end{catalogueentry}

\begin{catalogueentry}{is-fabius}{The abstract predicate \lean{IsFabius}}
\entryfield{Kind}{A proposition about a candidate bounded function, not a fourth Fabius
function.}
\entryfield{Role}{It packages the exact tails, reflection, smoothness/differentiability, and
left-half differential law required by the formal development.  A theorem quantified over
an arbitrary $f$ with \lean{IsFabius}~\(f\) is more general than a theorem stated only for
$\FabiusBounded$.}
\entryfield{Authority}{The Lean declaration's type is authoritative for its precise fields,
coercions, and domains.  A prose summary is not permission to omit a hypothesis.}
\entryfield{Status}{\textbf{Lean-proved}: existence supplies the canonical witness and
uniqueness identifies every witness with it.}
\end{catalogueentry}

\begin{catalogueentry}{fabius-relations}{Canonical relations among the three functions}
\entryfield{Fold}{For every $x\in\RealNumbers$,
$\RvachevUp(x)=\FabiusBounded(1-\AbsoluteValue{x})$.}
\entryfield{Density relation}{For $x\in\RealNumbers$,
$\FabiusBounded'(x)=2\RvachevUp(2x-1)$.  Both sides vanish outside the unit interval.}
\entryfield{Unit-interval translate}{For $x\in[0,1]$,
$\FabiusBounded(x)=\RvachevUp(x-1)$.  This is not a global identity.}
\entryfield{Global agreement}{For $x\in[0,1]$,
$\FabiusGlobal(x)=\FabiusBounded(x)$.  Outside that interval the left side is signed and
the right side has constant tails.}
\entryfield{Audit rule}{Every use of a differential, translation, or scaling identity states
whether its subject is $\FabiusBounded$, $\RvachevUp$, or $\FabiusGlobal$ and prints the
domain on which it holds.}
\end{catalogueentry}

\begin{catalogueentry}{fabius-dyadic-cells}{Dyadic points, reduced levels, and derivative
cells}
\entryfield{Dyadic point}{A dyadic rational is $x=m/2^n$ with
$m\in\IntegerNumbers$ and $n\in\NaturalNumbers$.  A unit-interval dyadic additionally has
$0\le m\le2^n$.}
\entryfield{Reduced positive level}{For an interior dyadic, ``reduced at level $n$'' means
$0<m<2^n$ and $m$ is odd.  The oddness prevents the same rational from being represented
at a smaller positive denominator level.}
\entryfield{Cell}{A level-$n$ dyadic cell is an interval between consecutive multiples of
$2^{-n}$; open, closed, or half-open endpoints are printed because derivative and spline
formulas can distinguish them.}
\entryfield{Derivative scaling}{At a reduced interior dyadic $m/2^n$, derivatives above a
finite order vanish according to the exact global-scaling theorem.  The catalogue does not
infer local analyticity from the terminating Taylor jet.}
\entryfield{Representation independence}{An exact evaluator accepting $(m,n)$ must prove
invariance under $(m,n)\mapsto(2m,n+1)$ rather than assuming a reduced input.}
\end{catalogueentry}

\section{Probability, measures, and distributional notation}

\begin{catalogueentry}{probability-space}{Probability spaces and random variables}
\entryfield{Probability space}{$(\Omega,\mathscr A,\Probability)$ denotes a probability
space: $\Omega$ is the sample set, $\mathscr A$ is a $\sigma$-algebra on $\Omega$, and
$\Probability:\mathscr A\to[0,1]$ is countably additive with
$\Probability(\Omega)=1$.}
\entryfield{Random variable}{A real random variable is an
$\mathscr A/\mathcal B(\RealNumbers)$-measurable map $X:\Omega\to\RealNumbers$.
A random vector $X:\Omega\to\RealNumbers^d$ is measurable for the Borel
$\sigma$-algebra of $\RealNumbers^d$.}
\entryfield{Events}{Braces inside probability are inverse images:
$\Probability\{X\le x\}$ abbreviates
$\Probability(X^{-1}((-\infty,x]))$.  They do not imply that $X$ has a density.}
\entryfield{Almost surely}{$X=Y$ almost surely, written $X=Y$ a.s., means
$\Probability\{\omega:X(\omega)=Y(\omega)\}=1$.  This is weaker than pointwise equality
and different from \notelink{law}{equality in law}.}
\entryfield{Independence}{Random variables $(X_j)_{j\in J}$ are independent when
\[
 \Probability\!\left(\bigcap_{j\in K}\{X_j\in B_j\}\right)
 =\prod_{j\in K}\Probability\{X_j\in B_j\}
\]
for every finite $K\subseteq J$ and every family of Borel sets $B_j$.
\emph{i.i.d.} means independent and identically distributed.}
\end{catalogueentry}

\begin{catalogueentry}{law}{Law, pushforward, equality in law, and weak convergence}
\entryfield{Law}{For a measurable $X:\Omega\to S$,
\[
 \LawOf{X}\DefinitionEquals\LawOperator(X)\DefinitionEquals X_{\#}\Probability,\qquad
 \LawOf{X}(B)=\Probability(X^{-1}(B)),
\]
where $S$ is a measurable space, $B$ is measurable, and $X_{\#}$ denotes pushforward.
Thus \commandname{LawOperator} is the operator token and \commandname{LawOf} is its
argument-delimited form; both take a random variable and return a probability measure.}
\entryfield{Equality in law}{$X\EqualInLaw Y$ means
$\LawOf{X}=\LawOf{Y}$.  It neither asserts that $X$ and $Y$ share a sample space nor
that they are equal almost surely.}
\entryfield{Convergence in law}{$X_n\ConvergesInLaw X$ means
\[
 \int_S h\,\Differential\LawOf{X_n}\longrightarrow
 \int_S h\,\Differential\LawOf{X}
\]
for every bounded continuous $h:S\to\RealNumbers$, with $S$ a metric space.  For
real-valued variables this is equivalent to convergence of distribution functions at
every continuity point of the limiting distribution function.}
\entryfield{Measure notation}{$\mu_n\Rightarrow\mu$ is allowed for weak convergence of
measures; it must not be reused for implication in the same displayed formula.}
\entryfield{Lean crosswalk}{Probability measures and pushforwards use Mathlib's
\lean{ProbabilityMeasure}, \lean{Measure.map}, and distribution APIs.  A prose
probabilistic representation is labelled \textbf{prose-derived} unless an exact Lean
theorem is cited.}
\end{catalogueentry}

\begin{catalogueentry}{dyadic-filtration}{The canonical dyadic filtration}
\entryfield{Command}{\commandname{DyadicSigmaField}.  The printed base glyph is
$\DyadicSigmaField$; it is never a generic event, matrix, Appell family, or algebra when
the command is used.}
\entryfield{Probability space}{Fix a probability space
$(\Omega,\mathscr A,\Probability)$ and an explicitly ordered coordinate family
$(U_j)_{j\in\PositiveIntegers}$.  For $n\in\NaturalNumbers$ define
\[
 \DyadicSigmaField_n\DefinitionEquals
 \sigma(U_1,\ldots,U_n),
 \qquad
 \DyadicSigmaField_0\DefinitionEquals\{\varnothing,\Omega\}.
\]
If a document instead generates from binary digits, Rademacher variables, or dyadic
partition cells, it replaces the displayed generators while retaining the level-zero
and monotonicity clauses.}
\entryfield{Filtration law}{$\DyadicSigmaField_n\subseteq\DyadicSigmaField_{n+1}
\subseteq\mathscr A$ for every $n$.  Conditional expectation
$\Expectation[X\mid\DyadicSigmaField_n]$ is defined only for an integrable $X$ (or the
appropriate nonnegative extended variant) and is an equivalence class modulo almost
sure equality.}
\entryfield{Limit field}{$\DyadicSigmaField_\infty\DefinitionEquals
\sigma(\bigcup_{n\in\NaturalNumbers}\DyadicSigmaField_n)$.  Equality with the ambient
$\mathscr A$ is an additional generation hypothesis, not part of the notation.}
\entryfield{Collision policy}{An unqualified calligraphic $A$ remains local.  The shared
command is reserved for this filtration, and every subscript is a level index in
$\NaturalNumbers\cup\{\infty\}$ rather than a matrix dimension or event label.}
\end{catalogueentry}

\begin{catalogueentry}{cdf-density}{Distribution functions, densities, and atoms}
\entryfield{Distribution function}{For a real random variable $X$,
\[
 F_X(x)\DefinitionEquals\Probability\{X\le x\},\qquad x\in\RealNumbers.
\]
It is right-continuous, nondecreasing, and has limits $0$ and $1$ at negative and
positive infinity.  The subscript is compulsory unless the distribution is uniquely
fixed in the surrounding entry.}
\entryfield{Density}{A Borel function $f_X:\RealNumbers\to[0,\infty)$ is a density when
$\LawOf{X}(B)=\int_B f_X(x)\,\Differential x$ for every Borel set $B$.  Then
$F_X(x)=\int_{-\infty}^{x}f_X(t)\,\Differential t$.  A distribution need not possess a
density.}
\entryfield{Atom}{The mass at $a$ is
$\Probability\{X=a\}=F_X(a)-F_X(a^-)$, where
$F_X(a^-)=\lim_{x\uparrow a}F_X(x)$.  A Dirac measure at $a$ is
$\delta_a(B)=\IndicatorOf{B}(a)$.}
\entryfield{Signed measures}{A signed measure $\nu$ is not a probability law.  Its total
variation measure is $|\nu|$; the scalar $\nu(\RealNumbers)$ is its total mass and may
be negative.}
\entryfield{Endpoints}{A uniform law on $[a,b]$, $a<b$, has density
$(b-a)^{-1}\IndicatorOf{[a,b]}(x)$.  Changing the values at the two endpoints changes
neither the associated absolutely continuous measure nor any Lebesgue integral.}
\end{catalogueentry}

\begin{catalogueentry}{expectation}{Expectation, moments, variance, and covariance}
\entryfield{Expectation}{For an integrable real random variable,
\[
 \Expectation[X]\DefinitionEquals\int_{\Omega}X(\omega)\,\Differential\Probability(\omega)
 =\int_{\RealNumbers}x\,\Differential\LawOf{X}(x).
\]
For a nonnegative $X$, the extended integral is permitted and may equal $+\infty$.}
\entryfield{Complex expectation}{For an integrable $X:\Omega\to\ComplexNumbers$,
meaning $\Expectation[\AbsoluteValue{X}]<\infty$,
\[
 \Expectation[X]
 \DefinitionEquals
 \Expectation[\Re X]+\ImaginaryUnit\Expectation[\Im X]
 =\int_\Omega X\,\Differential\Probability,
\]
where the last integral is the complex Bochner integral.  This is the interpretation used
when a characteristic-function integrand is complex-valued.}
\entryfield{Raw moments}{$m_n(X)\DefinitionEquals\Expectation[X^n]$ for
$n\in\NaturalNumbers$ whenever the expectation exists; $m_0(X)=1$, including on the
event $X=0$, by the polynomial convention $x^0=1$ for every $x$.}
\entryfield{Central moments}{$\mu_n(X)\DefinitionEquals
\Expectation[(X-\Expectation X)^n]$, assuming the needed integrability.  Raw $m_n$ and
central $\mu_n$ must not be interchanged.}
\entryfield{Variance and covariance}{For square-integrable real variables,
\[
 \Variance(X)=\Expectation[(X-\Expectation X)^2],\qquad
 \Covariance(X,Y)=\Expectation[(X-\Expectation X)(Y-\Expectation Y)].
\]
Thus $\Variance(X)=\Covariance(X,X)\ge0$.}
\entryfield{Cumulants}{When $\log\MomentGeneratingFunction{X}(t)$ exists near $0$,
$\kappa_n(X)$ is the coefficient defined by
$\log\MomentGeneratingFunction{X}(t)=\sum_{n\ge1}\kappa_n(X)t^n/n!$.
The logarithm is the real logarithm on the positive MGF, not an unstated complex branch.}
\end{catalogueentry}

\begin{catalogueentry}{fabius-random-series}{Canonical probabilistic model of the Fabius
distribution}
\entryfield{Input variables}{Let $(U_k)_{k\in\PositiveIntegers}$ be independent with
$U_k\sim\operatorname{Uniform}[0,1]$.  The index begins at $1$.}
\entryfield{Random series}{Define
\[
 X\DefinitionEquals\sum_{k=1}^{\infty}2^{-k}U_k.
\]
The series converges absolutely for every sample point because
$0\le U_k\le1$ and $\sum_{k\ge1}2^{-k}=1$; hence $X\in[0,1]$ almost surely.}
\entryfield{Distribution}{$\FabiusBounded(x)=\Probability\{X\le x\}$ for
$x\in\RealNumbers$.  The density is
$f_X(x)=2\RvachevUp(2x-1)$, so the distribution identity reproduces
$\FabiusBounded'(x)=2\RvachevUp(2x-1)$.}
\entryfield{Centered variable}{For $Y=2X-1$ one has $Y\in[-1,1]$ almost surely and
$Y$ has density $\RvachevUp$.  Equivalently,
$\RvachevUp(t)=\tfrac12 f_X((t+1)/2)$.}
\entryfield{Finite approximants}{$X_N=\sum_{k=1}^{N}2^{-k}U_k$ has support
$[0,1-2^{-N}]$ and $X_N\to X$ almost surely, in every $L^p$ for
$1\le p<\infty$, and in law.  These distinct modes are printed when a bound depends on
one of them.}
\entryfield{Status}{The random-series representation is a semantic bridge used in the
corpus.  Each downstream theorem retains its own proof-status label; the representation
alone does not formalize a probabilistic argument in Lean.}
\end{catalogueentry}

\begin{catalogueentry}{geometric-uniform-family}{Geometrically weighted uniform series}
\entryfield{Parameters}{Unless a complex analytic continuation is explicitly declared,
$q\in(0,1)$.  Let $U_0,U_1,\ldots$ be independent and uniformly distributed on
$[0,1]$.}
\entryfield{Normalized series}{The probability-normalized variable is
\[
 X_q\DefinitionEquals(1-q)\sum_{n=0}^{\infty}q^nU_n.
\]
The weights sum to $1$, so $X_q\in[0,1]$ almost surely.  At $q=\tfrac12$ its law agrees
with the canonical Fabius law after the index shift $U_n\leftrightarrow U_{n+1}$.}
\entryfield{Unnormalized series}{Some frontier documents instead study
$S_q=\sum_{n\ge0}q^nU_n=X_q/(1-q)$ with support $[0,(1-q)^{-1}]$.
The letter $S_q$ and its support must be printed whenever this normalization is used.}
\entryfield{Characteristic function}{Independence gives
\[
 \CharacteristicFunction{X_q}(t)
 =\prod_{n=0}^{\infty}
   \frac{\EulerE^{\ImaginaryUnit t(1-q)q^n}-1}
        {\ImaginaryUnit t(1-q)q^n},
\]
\nopagebreak[4]
with every factor assigned its continuous value $1$ at $t=0$.}
\entryfield{Complex \(q\)}{For $q\in\ComplexNumbers$, the same product may define an
analytic function for $|q|<1$, but it is no longer the characteristic function of a real
probability variable unless additional reality hypotheses are given.}
\end{catalogueentry}

\begin{catalogueentry}{probability-metrics}{Total variation, Wasserstein, and
Kolmogorov distances}
\entryfield{One-object total variation}{The command
\commandname{TotalVariationOf}\texttt{\{a\}} denotes the total variation of one object
$a$.  Its type must be stated locally.  For a real function on $[u,v]$,
\[
 \TotalVariationOf{f}
 \DefinitionEquals
 \sup_{u=x_0<\cdots<x_n=v}
 \sum_{j=1}^{n}\AbsoluteValue{f(x_j)-f(x_{j-1})}.
\]
For a finite signed measure $\sigma$, it denotes $|\sigma|(X)$, equivalently the
total-variation norm under the convention stated in the source.  These are one-object
functionals, not distances between two probability measures.}
\entryfield{Total variation}{For probability measures $\mu,\nu$ on the same measurable
space,
\[
 \TotalVariationDistance(\mu,\nu)
 \DefinitionEquals\sup_{A}|\mu(A)-\nu(A)|
 =\frac12\int|f-g|\,\Differential\lambda
\]
when both have densities $f,g$ with respect to $\lambda$.  Some literature omits the
factor $1/2$; that alternative is never silently substituted.}
\entryfield{Kolmogorov distance}{For real laws with CDFs $F_X,F_Y$,
$d_{\mathrm K}(F_X,F_Y)=\sup_{x\in\RealNumbers}|F_X(x)-F_Y(x)|$.}
\entryfield{Wasserstein distance}{For $p\in[1,\infty)$ and laws with finite $p$th moments,
\[
 W_p(\mu,\nu)=
 \left(\inf_{\pi\in\Pi(\mu,\nu)}
 \int |x-y|^p\,\Differential\pi(x,y)\right)^{1/p},
\]
where $\Pi(\mu,\nu)$ is the set of couplings.  The product coupling makes
$\Pi(\mu,\nu)$ nonempty for probability measures on standard Borel spaces.}
\entryfield{Type rule}{A metric symbol takes two laws or measures, not two random
variables, unless the law operation is explicitly stated in the local scope.}
\end{catalogueentry}

\section{Fourier, Laplace, Mellin, and generating transforms}

\subsection{The two Fourier conventions}

\begin{catalogueentry}{fourier-two-pi}{The \(2\pi\)-frequency Fourier transform}
\entryfield{Commands}{\commandname{FourierTwoPiOperator} is the operator token,
\commandname{FourierTwoPiTransformOf}\texttt{\{f\}} is its bracketed application, and
\commandname{FourierTwoPi}\texttt{\{f\}} is its normalization-labelled hat.  These are
three presentations of one forward transform, not three transforms.}
\entryfield{Forward transform}{For $f\in L^1(\RealNumbers)$,
\[
 \FourierTwoPi{f}(\xi)
 \DefinitionEquals\FourierTwoPiTransformOf{f}(\xi)
 \DefinitionEquals\FourierTwoPiOperator[f](\xi)
 \DefinitionEquals\int_{\RealNumbers}
 f(x)\EulerE^{-2\pi\ImaginaryUnit x\xi}\,\Differential x,
 \qquad \xi\in\RealNumbers.
\]
The transform variable is $\xi$ and carries cycles per unit.}
\entryfield{Inverse}{Under a hypothesis guaranteeing inversion,
\[
 f(x)=\int_{\RealNumbers}
 \FourierTwoPi{f}(\xi)\EulerE^{2\pi\ImaginaryUnit x\xi}\,\Differential\xi.
\]
There is no scalar prefactor in either direction.  The operator implementing this
inverse is \commandname{InverseFourierTwoPiOperator}.  It is written only on a named
inversion domain; it does not assert pointwise inversion for arbitrary $L^1$ data.  For
$L^2$, both formulas are interpreted through Plancherel extension.}
\entryfield{Rules}{With justified derivatives and integrals,
\[
 \FourierTwoPi{f'}(\xi)=2\pi\ImaginaryUnit\xi\,\FourierTwoPi f(\xi),
 \qquad
 \FourierTwoPi{f\Convolution g}=\FourierTwoPi f\,\FourierTwoPi g.
\]
Translation by $a$ contributes $\EulerE^{-2\pi\ImaginaryUnit a\xi}$; dilation obeys
$\FourierTwoPi{f(a\,\cdot)}(\xi)=|a|^{-1}\FourierTwoPi f(\xi/a)$ for real $a\ne0$.}
\entryfield{Measures}{For a finite Borel measure $\mu$,
$\FourierTwoPi{\mu}(\xi)=\int\EulerE^{-2\pi\ImaginaryUnit x\xi}\,\Differential\mu(x)$.
This has no density assumption.}
\end{catalogueentry}

\begin{catalogueentry}{fourier-angular}{The angular-frequency Fourier transform}
\entryfield{Operator token}{\commandname{FourierAngularOperator} denotes the transform
operator without supplying its operand.}
\entryfield{Bracketed application}{\commandname{FourierAngularTransformOf}\texttt{\{f\}}
denotes the transform applied to $f$ and prints the operand in brackets.}
\entryfield{Normalization-labelled hat}{\commandname{FourierAngular}\texttt{\{f\}}
denotes the same transform of $f$ with the angular-frequency normalization printed on the
hat.}
\entryfield{Forward transform}{For $f\in L^1(\RealNumbers)$,
\[
 \FourierAngular{f}(\omega)
 \DefinitionEquals\FourierAngularTransformOf{f}(\omega)
 \DefinitionEquals\FourierAngularOperator[f](\omega)
 \DefinitionEquals\int_{\RealNumbers}
 f(x)\EulerE^{-\ImaginaryUnit\omega x}\,\Differential x,
 \qquad\omega\in\RealNumbers.
\]}
\entryfield{Measures}{For a finite Borel measure $\mu$,
\[
 \FourierAngular{\mu}(\omega)
 \DefinitionEquals\int_{\RealNumbers}
 \EulerE^{-\ImaginaryUnit\omega x}\,\Differential\mu(x),
 \qquad\omega\in\RealNumbers.
\]
This is a Fourier--Stieltjes transform and requires no density.}
\entryfield{Inverse}{Under an inversion hypothesis,
\[
 f(x)=\frac1{2\pi}\int_{\RealNumbers}
 \FourierAngular{f}(\omega)\EulerE^{\ImaginaryUnit\omega x}\,\Differential\omega.
\]
The forward direction has no prefactor and the inverse has $1/(2\pi)$.}
\entryfield{Inverse command}{\commandname{InverseFourierAngularOperator} denotes precisely
this inverse on a stated inversion domain.}
\entryfield{Rules}{$\FourierAngular{f'}(\omega)
=\ImaginaryUnit\omega\FourierAngular f(\omega)$ and
$\FourierAngular{f\Convolution g}
=\FourierAngular f\,\FourierAngular g$.}
\entryfield{Conversion}{The two catalogued transforms satisfy
\[
 \FourierTwoPi f(\xi)=\FourierAngular f(2\pi\xi),\qquad
 \FourierAngular f(\omega)=\FourierTwoPi f(\omega/(2\pi)).
\]
The same identities hold, with $f$ replaced throughout by a finite Borel measure
$\mu$, for the two Fourier--Stieltjes transforms just defined.  A formula may change
variables between the conventions only by displaying this rescaling.}
\end{catalogueentry}

\Needspace{13\baselineskip}%
\begin{catalogueentry}{characteristic-function}{Characteristic functions and Fourier
signs}
\entryfield{Definition}{For a real random variable $X$ with probability law $\mu$,
\[
 \CharacteristicFunction{X}(t)
 \DefinitionEquals\Expectation[\EulerE^{\ImaginaryUnit tX}]
 =\CharacteristicFunction{\mu}(t)
 \DefinitionEquals\int_{\RealNumbers}\EulerE^{\ImaginaryUnit tx}\,\Differential\mu(x),
 \qquad t\in\RealNumbers.
\]
Thus the command accepts either a declared random variable or a declared probability
measure.  The argument's type fixes the interpretation, and
$\CharacteristicFunction{X}=\CharacteristicFunction{\LawOf X}$ pointwise.  It is not
used for an arbitrary signed or complex measure without a local declaration.  The
probabilistic sign is positive.}
\entryfield{Fourier conversion}{Because both catalogued Fourier transforms use a negative
sign,
\[
 \CharacteristicFunction{X}(t)
 =\FourierAngular{\LawOf X}(-t)
 =\FourierTwoPi{\LawOf X}\!\left(-\frac{t}{2\pi}\right).
\]}
\entryfield{Moments}{If $\Expectation|X|^n<\infty$, then
$\CharacteristicFunction{X}^{(n)}(0)=\ImaginaryUnit^n\Expectation[X^n]$.
Analyticity at zero requires stronger hypotheses than the existence of all moments.}
\entryfield{Independent sums}{For independent $X,Y$,
$\CharacteristicFunction{X+Y}=\CharacteristicFunction X\,
\CharacteristicFunction Y$.  The equality is pointwise in $t$.}
\end{catalogueentry}

\begin{catalogueentry}{sinc}{The two sinc normalizations}
\entryfield{Radian sinc}{For $z\in\ComplexNumbers$,
\[
 \SincRad(z)\DefinitionEquals
 \begin{cases}\sin z/z,&z\ne0,\\1,&z=0.\end{cases}
\]
The value at zero fills the removable singularity, and the result is entire.}
\entryfield{Normalized sinc}{For $z\in\ComplexNumbers$,
\[
 \SincPi(z)\DefinitionEquals
 \begin{cases}\sin(\pi z)/(\pi z),&z\ne0,\\1,&z=0.\end{cases}
\]}
\entryfield{Conversion}{$\SincPi(z)=\SincRad(\pi z)$ and
$\SincRad(z)=\SincPi(z/\pi)$.}
\entryfield{Zeros}{$\SincRad$ vanishes at $\pi k$, while $\SincPi$ vanishes at $k$, for
$k\in\IntegerNumbers\setminus\{0\}$.  This distinction is essential in zero-set and
sampling arguments.}
\entryfield{Retired alias}{A bare sinc operator is forbidden in normative prose.
Historical occurrences must be annotated with one of the two definitions.}
\end{catalogueentry}

\begin{catalogueentry}{up-fourier-product}{Fourier products for Rvachev's up-function}
\entryfield{\(2\pi\)-frequency form}{With $\int\RvachevUp=1$ and the
\notelink{fourier-two-pi}{$2\pi$ convention},
\[
 \FourierTwoPi{\RvachevUp}(\xi)
 =\prod_{n=0}^{\infty}\SincRad\!\left(\frac{\pi\xi}{2^n}\right)
 =\prod_{n=0}^{\infty}\SincPi\!\left(\frac{\xi}{2^n}\right).
\]}
\entryfield{Angular form}{Equivalently,
\[
 \FourierAngular{\RvachevUp}(\omega)
 =\prod_{j=1}^{\infty}\SincRad\!\left(\frac{\omega}{2^j}\right).
\]
The angular product begins at $j=1$; the $2\pi$ product begins at $n=0$.}
\entryfield{At zero}{Every factor is assigned value $1$ at zero, so both products equal
$1$ at the origin.}
\entryfield{Convergence}{The products converge locally uniformly on
$\ComplexNumbers$, hence define entire functions.  A logarithm of the product is used
only on a zero-free domain with an explicit branch.}
\entryfield{Zeros}{Multiplicities are computed by counting all pairs $(n,k)$ for which
$\xi/2^n=k\in\IntegerNumbers\setminus\{0\}$, not by treating the displayed factors as
having disjoint zero sets.}
\end{catalogueentry}

\begin{catalogueentry}{periodic-fourier}{Fourier series on the unit circle}
\entryfield{Command}{
\commandname{FourierSeriesCoefficient}\texttt{\{h\}\{k\}}.  The printed
subscript ``FS'' distinguishes a discrete Fourier-series coefficient from both continuous
line-transform normalizations.  For a period other than one, use
\commandname{FourierSeriesCoefficientAtPeriod}\texttt{\{h\}\{k\}\{T\}}; its third
argument is mandatory and appears in the printed symbol.}
\entryfield{Coefficient}{For a $1$-periodic integrable function $h$,
\[
 \FourierSeriesCoefficient{h}{k}\DefinitionEquals\int_0^1h(x)
 \EulerE^{-2\pi\ImaginaryUnit kx}\,\Differential x,
 \qquad k\in\IntegerNumbers.
\]}
\entryfield{Series}{The formal Fourier series is
$\sum_{k\in\IntegerNumbers}
\FourierSeriesCoefficient{h}{k}\EulerE^{2\pi\ImaginaryUnit kx}$.
Its mode of convergence---pointwise, $L^2$, uniform, distributional, or a named
summability method---is part of every theorem.}
\entryfield{Separation from line transform}{Square brackets and an integer index mark a
Fourier-series coefficient.  Parentheses and a real frequency mark
$\FourierTwoPi h(\xi)$ or $\FourierAngular h(\omega)$ on the line.}
\entryfield{Period \(T\)}{For a $T$-periodic function, $T>0$,
$\FourierSeriesCoefficientAtPeriod{h}{k}{T}
=T^{-1}\int_0^T h(x)\EulerE^{-2\pi\ImaginaryUnit kx/T}\,\Differential x$.
The period, the factor $T^{-1}$, and the kernel scale are never implicit.  A finite
measure or an unnormalized Fourier--Stieltjes coefficient is a different local object
and does not use either function-coefficient command.}
\end{catalogueentry}

\begin{catalogueentry}{laplace-transform}{One-sided and bilateral Laplace transforms}
\entryfield{Commands}{\commandname{LaplaceTransformSymbol} is the operator token and
\commandname{LaplaceTransformOf}\texttt{\{f\}} supplies a bracketed operand.  An
unsubscripted token is one-sided by this catalogue's convention.}
\entryfield{One-sided}{For $f:[0,\infty)\to\ComplexNumbers$,
\[
 \LaplaceTransformOf f(s)
 \DefinitionEquals\int_0^\infty \EulerE^{-sx}f(x)\,\Differential x,
 \qquad s\in\ComplexNumbers
\]
where the integral converges.}
\entryfield{Bilateral}{For $f:\RealNumbers\to\ComplexNumbers$,
\[
 \LaplaceTransformSymbol_{\mathrm{bi}}[f](s)
 \DefinitionEquals\int_{-\infty}^{\infty}
 \EulerE^{-sx}f(x)\,\Differential x.
\]
The subscript \(\mathrm{bi}\) is compulsory; it prevents a support assumption from being
smuggled into an identity.}
\entryfield{Abscissa}{The abscissa of absolute convergence is
$\sigma_a=\inf\{\sigma\in\RealNumbers:
\int_0^\infty \EulerE^{-\sigma x}|f(x)|\,\Differential x<\infty\}$,
with infimum in $\ExtendedRealNumbers$.}
\entryfield{Probability relation}{For a nonnegative random variable $X$,
$\Expectation[\EulerE^{-sX}]$ is the Laplace--Stieltjes transform of its law.
$\MomentGeneratingFunction X(t)=\Expectation[\EulerE^{tX}]$ is obtained at $s=-t$
only where both sides are finite.}
\end{catalogueentry}

\begin{catalogueentry}{mgf}{Moment-generating and exponential-generating functions}
\entryfield{Moment-generating function}{For a real random variable $X$,
\[
 \MomentGeneratingFunction X(t)\DefinitionEquals
 \Expectation[\EulerE^{tX}]
\]
for real or complex $t$ at which the expectation is absolutely finite.  It is not
declared on all of $\ComplexNumbers$ merely because every formal moment exists.}
\entryfield{Moment expansion}{If the MGF is analytic at zero,
$\MomentGeneratingFunction X(t)=\sum_{n=0}^{\infty}
\Expectation[X^n]t^n/n!$ locally.}
\entryfield{Exponential-generating function}{For an abstract sequence $(a_n)$,
$A_{\mathrm{egf}}(z)=\sum_{n\ge0}a_nz^n/n!$.  It is an MGF only after a random variable
with $a_n=\Expectation[X^n]$ has been specified.}
\entryfield{Ordinary-generating function}{$A_{\mathrm{ogf}}(z)=\sum_{n\ge0}a_nz^n$.
The subscripts \(\mathrm{egf}\) and \(\mathrm{ogf}\) are retained wherever both occur.}
\end{catalogueentry}

\begin{catalogueentry}{mellin-transform}{Mellin transform}
\entryfield{Commands}{\commandname{MellinTransformSymbol} is the operator token and
\commandname{MellinTransformOf}\texttt{\{f\}} supplies the operand.  Neither command
chooses a convergence strip or analytic continuation.}
\entryfield{Definition}{For $f:(0,\infty)\to\ComplexNumbers$,
\[
 \MellinTransformOf f(s)\DefinitionEquals
 \int_0^\infty f(x)x^{s-1}\,\Differential x,
 \qquad s\in\ComplexNumbers,
\]
on the maximal strip of absolute convergence.}
\entryfield{Complex power}{For $x>0$, $x^{s-1}$ means
$\exp((s-1)\log x)$ with the real logarithm $\log x$; no complex branch is involved.}
\entryfield{Scaling}{For real $a>0$,
$\MellinTransformOf{f(a\,\cdot)}(s)=a^{-s}\MellinTransformOf f(s)$.}
\entryfield{Continuation}{A meromorphic continuation of the integral-defined transform is
identified by the same symbol only after its initial strip and uniqueness of continuation
have been stated.  Poles, residues, and contour shifts refer to that continuation.}
\end{catalogueentry}

\begin{catalogueentry}{cauchy-stieltjes}{Cauchy and Stieltjes transforms}
\entryfield{Cauchy transform}{For a finite Borel measure $\mu$ on $\RealNumbers$,
\[
 G_\mu(z)\DefinitionEquals\int_{\RealNumbers}\frac{1}{z-x}\,
 \Differential\mu(x),\qquad z\in\ComplexNumbers\setminus\SupportOf\mu.
\]}
\entryfield{Alternative sign}{Some sources define
$\widetilde G_\mu(z)=\int(x-z)^{-1}\Differential\mu(x)=-G_\mu(z)$.
The catalogue uses the $z-x$ denominator; a quoted alternative carries a tilde and the
minus-sign conversion.}
\entryfield{Stieltjes transform}{For a measure on $[0,\infty)$,
$S_\mu(s)=\int_0^\infty(s+x)^{-1}\Differential\mu(x)$ for
$s\in\ComplexNumbers\setminus(-\infty,0]$ when finite.  Denominator signs must be
accounted for before relating it to $G_\mu$.}
\entryfield{Boundary values}{Symbols $G_\mu(x\pm\ImaginaryUnit0)$ mean non-tangential
limits from the upper or lower half-plane when those limits exist; they are not
evaluation at a complex number with an infinitesimal component.}
\end{catalogueentry}

\begin{catalogueentry}{transform-distributions}{Transforms of distributions and
generalized functions}
\entryfield{Tempered distributions}{For $T\in\mathscr S'(\RealNumbers)$, its
$2\pi$-Fourier transform is defined by duality
$\langle\FourierTwoPi T,\phi\rangle
=\langle T,\FourierTwoPi\phi\rangle$ after fixing the test-function transform above.
The pairing is bilinear, not a Hermitian inner product.}
\entryfield{Dirac identities}{$\FourierTwoPi{\delta_a}(\xi)
=\EulerE^{-2\pi\ImaginaryUnit a\xi}$ and
$\FourierTwoPi 1=\delta_0$ in $\mathscr S'$.}
\entryfield{Derivatives}{$\FourierTwoPi{D^mT}
=(2\pi\ImaginaryUnit\xi)^m\FourierTwoPi T$, where $D$ is distributional derivative.}
\entryfield{Audit rule}{A distributional identity is never presented as an absolutely
convergent integral identity.  The ambient test-function space and duality convention
appear at the first use.}
\end{catalogueentry}

\begin{catalogueentry}{fourier-decay-objectives}{Nonlinear and linearized Fourier-decay
objectives}
\entryfield{Commands}{For $n\in\NaturalNumbers$ the exact optimization functional is
$\DecayOptimizationObjective{n}$ and its first-order or otherwise explicitly specified
linearization is $\LinearizedDecayObjective{n}$.  Their semantic names prevent either
from being mistaken for a Fabius function, a Fourier operator, or a filtration.}
\entryfield{Required declaration}{At the first use a document supplies a variable space
$V_n$, a feasible set $\mathcal C_n\subseteq V_n$, a codomain (normally
$\RealNumbers\cup\{+\infty\}$), and formulas
\[
 \DecayOptimizationObjective{n}:\mathcal C_n\to\ExtendedRealNumbers,
 \qquad
 \LinearizedDecayObjective{n}:T_{v_0}\mathcal C_n\to\ExtendedRealNumbers.
\]
The base point $v_0$, constraints, frequency normalization, logarithm base, weights, and
all $n$-dependence are part of those formulas.  The shared commands intentionally do not
smuggle in one draft's choice.}
\entryfield{Linearization contract}{A label ``linearized'' identifies a new objective.
An expansion such as
$\DecayOptimizationObjective{n}(v_0+h)=
\DecayOptimizationObjective{n}(v_0)+
\LinearizedDecayObjective{n}(h)+r_n(h)$ states the norm and limit in which
$r_n(h)=\LittleOAt{\Norm{h}}{h\to0}$.  Dropping $r_n$ produces an approximation, not an
identity.}
\entryfield{Optimization values}{$\inf_{v\in\mathcal C_n}
\DecayOptimizationObjective{n}(v)$ is an extended-real value.  A symbol for a minimizer
is introduced only after existence; uniqueness is an additional theorem.  The optimum
of the linearized problem is not substituted for the nonlinear optimum without a bound.}
\entryfield{Status}{The commands are \textbf{definitions of notation}.  Particular
decay bounds, saddle locations, convexity assertions, and asymptotic optima inherit the
proof-status of the document that states them.}
\end{catalogueentry}

\section{Binary arithmetic, Thue--Morse data, and exact combinatorics}

\begin{catalogueentry}{binary-weight}{Binary digit sum and bits}
\entryfield{Binary expansion}{Every $n\in\NaturalNumbers$ has the unique finite expansion
$n=\sum_{j=0}^{\infty}b_j(n)2^j$ with $b_j(n)\in\{0,1\}$ and all but finitely many
$b_j(n)$ zero.  Leading zeros are allowed but contribute nothing.}
\entryfield{Digit sum}{The canonical binary weight is
\[
 \BinaryDigitSum(n)\DefinitionEquals\sum_{j=0}^{\infty}b_j(n),
 \qquad n\in\NaturalNumbers.
\]
Thus $\BinaryDigitSum(0)=0$.}
\entryfield{Bit extraction}{$b_j(n)=\Floor{n/2^j}-2\Floor{n/2^{j+1}}$ for
$j\in\NaturalNumbers$.  Negative integers require an explicitly selected finite-word or
$2$-adic representation and are outside this global definition.}
\entryfield{Retired aliases}{The former digit-sum, population-count, and bare-weight
spellings may occur in quotations.  Normative prose uses \commandname{BinaryDigitSum};
an unsubscripted weight is especially unsafe because it can denote an integration weight.}
\end{catalogueentry}

\begin{catalogueentry}{thue-morse}{Thue--Morse bit and sign sequences}
\entryfield{Bit sequence}{$t_n\DefinitionEquals\BinaryDigitSum(n)\bmod2\in\{0,1\}$ for
$n\in\NaturalNumbers$.}
\entryfield{Sign sequence}{The canonical sign is
\[
 \ThueMorseSign n\DefinitionEquals(-1)^{\BinaryDigitSum(n)}
 =(-1)^{t_n}\in\{-1,1\}.
\]
Consequently $\ThueMorseSign0=1$.}
\entryfield{Command layer}{\commandname{ThueMorseSignSymbol} is a zero-argument reserved
glyph token rendering $\ThueMorseSignSymbol$; it names the sign sequence inside operator
labels, DFT displays, and prose where no index is being evaluated.
\commandname{ThueMorseSign}\texttt{\{n\}} takes exactly one argument
$n\in\NaturalNumbers$ and renders the indexed value $\ThueMorseSign n$.  The former is
not a constant value and the latter must not receive a second raw subscript.}
\entryfield{Recurrences}{$t_{2n}=t_n$, $t_{2n+1}=1-t_n$ and
$\ThueMorseSign{2n}=\ThueMorseSign n$,
$\ThueMorseSign{2n+1}=-\ThueMorseSign n$.}
\entryfield{Generating product}{For every $N\in\NaturalNumbers$,
\[
 \sum_{n=0}^{2^N-1}\ThueMorseSign n\,z^n
 =\prod_{j=0}^{N-1}(1-z^{2^j}).
\]
At $N=0$ both sides are $1$: the sum contains the sole term $n=0$ and the product is
empty.}
\entryfield{Lean crosswalk}{The formal sign sequence uses declarations in
\modulepath{ThueMorse.lean} and related prefix modules.  Exact theorem names are recorded in
the crosswalk appendix; prose identities are not upgraded beyond their verified status.}
\end{catalogueentry}

\begin{catalogueentry}{padic-valuation}{\(p\)-adic and \(2\)-adic valuations}
\entryfield{Finite valuation}{For a prime $p$ and nonzero
$n\in\IntegerNumbers$,
\[
 \PadicValuation p(n)\DefinitionEquals
 \max\{k\in\NaturalNumbers:p^k\mid n\}.
\]
The value depends only on $|n|$.}
\entryfield{At zero}{In analytic formulas using extended values,
$\PadicValuation p(0)=+\infty\in\ExtendedRealNumbers$.  In Lean or finite-natural APIs,
the valuation of zero may instead be represented by a separate option or by a
library-specific totalization; each theorem states which codomain it uses.}
\entryfield{Two-adic abbreviation}{$\TwoAdicValuation(n)=\PadicValuation2(n)$.
For $n\ne0$, write $n=2^{\TwoAdicValuation(n)}n_{\mathrm{odd}}$ with
$n_{\mathrm{odd}}$ odd.}
\entryfield{Prime-power Pascal row}{For a prime \(p\), \(m\in\NaturalNumbers\), and
\(j\in\PositiveIntegers\) with \(j\le p^m\), the finite-natural valuation satisfies
\[
 \PadicValuation p\!\binom{p^m}{j}+\PadicValuation p(j)=m,
 \qquad
 \PadicValuation p\!\binom{p^m}{j}=m-\PadicValuation p(j).
\]
The right endpoint \(j=p^m\) and the boundary \(m=0\) are included.  The strict
dyadic-comb specialization has \(0<j<2^m\).  These formulas do not cover the distinct
endpoint-flat weights \(\binom{2^m-2}{j-1}\).}
\entryfield{Lean crosswalk}{In \modulepath{PrimePowerBinomialValuation.lean}, the
arbitrary-prime identities are the declarations
\leanline{primePowerChoose_padicValNat_add}
\leanline{primePowerChoose_padicValNat}
The strict dyadic specialization is
\leanline{twoPowChoose_padicValNat}
Lean's
\lean{padicValNat} is totalized in \(\NaturalNumbers\), so its exact zero behavior is the
declaration's behavior rather than the extended-value convention above.  Status:
\textbf{Lean-proved}.}
\entryfield{Order notation}{For a zero or pole of a meromorphic function at $z_0$,
$\OrderOperator_{z_0}(f)$ is a different valuation-like object; positive values denote
zeros and negative values poles in this catalogue.}
\end{catalogueentry}

\begin{catalogueentry}{dyadic-prefix}{Dyadic blocks and prefix sums}
\entryfield{Block}{The level-$N$ dyadic block of indices is
$B_N=\{0,1,\ldots,2^N-1\}$ for $N\in\NaturalNumbers$.}
\entryfield{Signed prefix sum}{For a function $P:\NaturalNumbers\to\ComplexNumbers$,
\[
 \Sigma_N(P)\DefinitionEquals
 \sum_{n=0}^{2^N-1}\ThueMorseSign n\,P(n).
\]}
\entryfield{Prouhet cancellation}{If $P$ is a polynomial of degree $<N$, then
$\Sigma_N(P)=0$.  At degree $N$ the surviving value depends on the leading coefficient
and the exact sign convention; it must not be quoted without that coefficient.}
\entryfield{Arbitrary prefix}{$S_M(P)=\sum_{n=0}^{M-1}\ThueMorseSign nP(n)$ for
$M\in\NaturalNumbers$.  The half-open upper bound ensures $S_0(P)=0$.}
\entryfield{Shifted block}{A sum over
$a2^N\le n<(a+1)2^N$ prints both $a\in\NaturalNumbers$ and $N\in\NaturalNumbers$;
its sign is not inferred from the unshifted identity without applying the binary carry
law.}
\end{catalogueentry}

\begin{catalogueentry}{finite-dft}{Finite discrete Fourier transforms}
\entryfield{Forward DFT}{For a vector $a=(a_0,\ldots,a_{N-1})\in\ComplexNumbers^N$,
$N\in\PositiveIntegers$,
\[
 \widehat a_k\DefinitionEquals
 \sum_{n=0}^{N-1}a_n\EulerE^{-2\pi\ImaginaryUnit kn/N},
 \qquad k=0,\ldots,N-1.
\]}
\entryfield{Inverse}{$a_n=N^{-1}\sum_{k=0}^{N-1}
\widehat a_k\EulerE^{2\pi\ImaginaryUnit kn/N}$.  The normalization is forward-unscaled,
inverse-$N^{-1}$.}
\entryfield{Root-of-unity form}{With
$\zeta_N=\EulerE^{2\pi\ImaginaryUnit/N}$,
$\widehat a_k=\sum_n a_n\zeta_N^{-kn}$.  The symbol $\zeta_N$ is a primitive root
chosen with positive angle.}
\entryfield{Thue--Morse specialization}{For $N=2^m$ and
$a_n=\ThueMorseSign n$, the product formula evaluates the DFT at
$z=\zeta_N^{-k}$.  Multiplicities of zeros require counting vanishing factors, not only
identifying a vanishing product.}
\Needspace{13\baselineskip}%
\entryfield{Document-local command}{In the Thue--Morse atlas,
\commandname{ThueMorseDFTCoefficient}\texttt{\{m\}\{k\}} has the two arguments
$m\in\NaturalNumbers$ and $k\in\IntegerNumbers$ in that order and means
\[
 \operatorname{DFT}_{2^m}[\ThueMorseSignSymbol][k]
 \DefinitionEquals
 \sum_{n=0}^{2^m-1}\ThueMorseSign n\,
 \EulerE^{-2\pi\ImaginaryUnit kn/2^m}.
\]
\nopagebreak[4]
The value depends only on $k$ modulo $2^m$; the canonical representative range is
$0\le k<2^m$.  The forward transform is unscaled, the exponential sign is negative,
and inversion therefore carries the factor $2^{-m}$ and the positive exponential.  This
operator is a finite DFT, not either line-Fourier transform and not a non-Fourier accent.}
\end{catalogueentry}

\begin{catalogueentry}{zero-runs}{Zero sets, multiplicities, and runs}
\entryfield{Order of a zero}{For analytic $f$ not identically zero near $z_0$,
$\OrderOperator_{z_0}(f)=m\in\NaturalNumbers$ means
$f(z)=(z-z_0)^mg(z)$ with $g(z_0)\ne0$.}
\entryfield{Run in a sequence}{For a sequence $(a_n)$, a zero run of length $\ell$ from
$r$ means $a_r=\cdots=a_{r+\ell-1}=0$ and, when the adjacent indices exist,
$a_{r-1}\ne0$ and $a_{r+\ell}\ne0$.}
\entryfield{Boundary runs}{At $r=0$, only the right maximality condition is required.
For a finite sequence ending at $M$, a terminal run has no right-neighbor condition.}
\entryfield{Separation}{A multiplicity of a zero of a generating polynomial is not a run
of zero coefficients.  The corpus contains both phenomena and keeps
$\OrderOperator_{z_0}$ separate from $\ell_{\mathrm{run}}$.}
\end{catalogueentry}

\begin{catalogueentry}{compositions-partitions}{Compositions, partitions, and multi-indices}
\entryfield{Weak composition}{A weak composition of $n$ into $k$ parts is
$\alpha=(\alpha_1,\ldots,\alpha_k)\in\NaturalNumbers^k$ with
$|\alpha|=\sum_{j=1}^k\alpha_j=n$.}
\entryfield{Composition}{A composition has $\alpha_j\in\PositiveIntegers$.
For $n=0$, the unique composition with zero parts is the empty tuple; there is no
positive-part composition into $k>0$ parts.}
\entryfield{Partition}{A partition
$\lambda\vdash n$ is a finite nonincreasing sequence of positive integers summing to
$n$.  Its length is $\ell(\lambda)$ and its multiplicity of part $j$ is $m_j(\lambda)$.}
\entryfield{Multi-index}{For $\alpha\in\NaturalNumbers^d$,
$|\alpha|=\sum_j\alpha_j$, $\alpha!=\prod_j\alpha_j!$,
$x^\alpha=\prod_jx_j^{\alpha_j}$, and
$\partial^\alpha=\prod_j\partial_{x_j}^{\alpha_j}$.  Products over $d=0$ are empty and
equal $1$.}
\entryfield{Factorials}{The shared
$\FallingFactorial{x}{n}=x(x-1)\cdots(x-n+1)$ and
$\RisingFactorial{x}{n}=x(x+1)\cdots(x+n-1)$ both equal $1$ at $n=0$.}
\end{catalogueentry}

\section{Combinatorial coefficient calculus and inversion}

\begin{catalogueentry}{stirling-and-lah-families}
  {Signed and unsigned Stirling numbers, second-kind numbers, and Lah numbers}
\entryfield{Index domain and zero extension}{Unless a formula says otherwise,
\(n,k\in\NaturalNumbers\).  Each triangular family is extended by zero when
\(k<0\), \(n<0\), or \(k>n\).  Every empty degree-zero basis expansion has the
single value at \((n,k)=(0,0)\) equal to \(1\).  Thus the boundary conventions
are part of the notation, not omitted side conditions.}
\entryfield{Signed first kind}{The signed Stirling number of the first kind is
defined by
\[
 \FallingFactorial{x}{n}
 =x(x-1)\cdots(x-n+1)
 =\sum_{k=0}^{n}\SignedStirlingFirstKind{n}{k}x^k.
\]
Consequently
\[
 \SignedStirlingFirstKind{0}{0}=1,
 \qquad
 \SignedStirlingFirstKind{n}{0}=0\quad(n>0),
\]
and
\[
 \SignedStirlingFirstKind{n+1}{k}
 =\SignedStirlingFirstKind{n}{k-1}
  -n\SignedStirlingFirstKind{n}{k}.
\]
The sign is \((-1)^{n-k}\) whenever the corresponding unsigned count is
nonzero.}
\entryfield{Unsigned first kind}{The unsigned cycle number is
\[
 \UnsignedStirlingFirstKind{n}{k}
 \DefinitionEquals
 (-1)^{n-k}\SignedStirlingFirstKind{n}{k}.
\]
It counts permutations of an \(n\)-element set having exactly \(k\) disjoint
cycles.  Its recurrence is
\[
 \UnsignedStirlingFirstKind{n+1}{k}
 =\UnsignedStirlingFirstKind{n}{k-1}
  +n\UnsignedStirlingFirstKind{n}{k}.
\]
The bracket is not a coefficient-extraction bracket and is not a Gaussian
binomial coefficient.}
\entryfield{Second kind}{The Stirling number
\(\StirlingSecondKind{n}{k}\) counts partitions of an \(n\)-element labelled
set into exactly \(k\) nonempty, unordered blocks.  It is characterized by
\[
 x^n=\sum_{k=0}^{n}\StirlingSecondKind{n}{k}
                 \FallingFactorial{x}{k}
\]
and by
\[
 \StirlingSecondKind{n+1}{k}
 =k\StirlingSecondKind{n}{k}
  +\StirlingSecondKind{n}{k-1}.
\]
The boundary values are
\(\StirlingSecondKind{0}{0}=1\),
\(\StirlingSecondKind{n}{0}=0\) for \(n>0\), and
\(\StirlingSecondKind{n}{n}=1\).}
\entryfield{Inverse basis changes}{The lower-triangular matrices of signed
first-kind and second-kind numbers are inverses:
\[
 \sum_{j=k}^{n}\SignedStirlingFirstKind{n}{j}
                   \StirlingSecondKind{j}{k}
 =\delta_{n,k}
 =\sum_{j=k}^{n}\StirlingSecondKind{n}{j}
                   \SignedStirlingFirstKind{j}{k}.
\]
This is matrix inversion over the integer-indexed triangular arrays; it is not
functional or multiplicative inversion.}
\entryfield{Unsigned Lah number}{For \(1\le k\le n\),
\[
 \LahNumber{n}{k}
 \DefinitionEquals
 \binom{n-1}{k-1}\frac{n!}{k!}.
\]
It counts partitions of an \(n\)-element labelled set into \(k\) nonempty
linearly ordered lists, with the collection of lists itself unordered.  Put
\(\LahNumber{0}{0}=1\) and set it to zero outside \(0\le k\le n\).}
\entryfield{Lah basis changes}{With the preceding boundary convention,
\[
 \RisingFactorial{x}{n}
 =\sum_{k=0}^{n}\LahNumber{n}{k}\FallingFactorial{x}{k},
\]
\[
 \FallingFactorial{x}{n}
 =\sum_{k=0}^{n}(-1)^{n-k}\LahNumber{n}{k}
                    \RisingFactorial{x}{k}.
\]
The label \(\mathrm{Lah}\) prevents the family from colliding with a Lambert
polynomial, logarithm, Laplace transform, or other \(L\)-family.}
\entryfield{Shared commands}{The normative source commands are
\commandname{SignedStirlingFirstKind},
\commandname{UnsignedStirlingFirstKind},
\commandname{StirlingSecondKind}, and \commandname{LahNumber}.}
\end{catalogueentry}

\begin{catalogueentry}{bell-and-touchard-families}
  {Bell numbers, Touchard polynomials, and the two Bell-polynomial normalizations}
\entryfield{Bell number}{The Bell number is
\[
 \BellNumber{n}
 \DefinitionEquals
 \sum_{k=0}^{n}\StirlingSecondKind{n}{k},
 \qquad \BellNumber{0}=1.
\]
It counts all set partitions of an \(n\)-element labelled set.  The superscript
\(\mathrm{Bell}\) is printed because an unlabelled \(B_n\) is also widely used
for Bernoulli numbers and polynomials.}
\entryfield{Touchard polynomial}{For \(n\in\NaturalNumbers\),
\[
 \TouchardPolynomial{n}(x)
 \DefinitionEquals
 \sum_{k=0}^{n}\StirlingSecondKind{n}{k}x^k.
\]
Thus
\(\TouchardPolynomial{n}(1)=\BellNumber{n}\) and
\(\TouchardPolynomial{0}(x)=1\).  Its exponential generating function is
\[
 \sum_{n\ge0}\TouchardPolynomial{n}(x)\frac{z^n}{n!}
 =\exp\!\bigl(x(\EulerE^z-1)\bigr).
\]
The tag \(\mathrm{Tou}\) distinguishes it from every other \(T_n\)-family.}
\entryfield{Complementary Bell number}{The signed block-parity sum is
\[
 \ComplementaryBellNumber{n}
 \DefinitionEquals
 \sum_{k=0}^{n}(-1)^k\StirlingSecondKind{n}{k}
 =\TouchardPolynomial{n}(-1).
\]
Its EGF and recurrence are
\[
 \sum_{n\ge0}\ComplementaryBellNumber{n}\frac{z^n}{n!}
 =\exp(1-\EulerE^z),
\]
\[
 \ComplementaryBellNumber{n+1}
 =-\sum_{j=0}^{n}\binom{n}{j}\ComplementaryBellNumber{j}.
\]
The source aliases \(C_n\), \(\mathsf{U}_n\), and a hatted or tilded Bell number
are retired because \(C_n\) also denoted a Catalan number and accents have
unrelated meanings elsewhere.}
\entryfield{Exponential partial Bell polynomial}{For
\(0\le k\le n\),
\[
 \ExponentialPartialBellPolynomial{n}{k}(x_1,x_2,\ldots)
 \DefinitionEquals
 \sum_{\substack{m_1,m_2,\ldots\ge0\\
                   \sum_{j\ge1}m_j=k\\
                   \sum_{j\ge1}jm_j=n}}
 \frac{n!}{\prod_{j\ge1}m_j!\,(j!)^{m_j}}
 \prod_{j\ge1}x_j^{m_j}.
\]
Only \(x_1,\ldots,x_{n-k+1}\) can occur.  The empty-index boundary is
\(\ExponentialPartialBellPolynomial{0}{0}=1\); it is zero for impossible
indices.  Equivalently,
\[
 \frac1{k!}\left(\sum_{j\ge1}x_j\frac{t^j}{j!}\right)^k
 =\sum_{n\ge k}
   \ExponentialPartialBellPolynomial{n}{k}(x_1,x_2,\ldots)
   \frac{t^n}{n!}.
\]}
\entryfield{Exponential complete Bell polynomial}{The complete polynomial is
\[
 \ExponentialCompleteBellPolynomial{n}(x_1,\ldots,x_n)
 \DefinitionEquals
 \sum_{k=0}^{n}
 \ExponentialPartialBellPolynomial{n}{k}(x_1,\ldots,x_{n-k+1}),
\]
with \(\ExponentialCompleteBellPolynomial{0}=1\).  Its generating identity is
\[
 \exp\!\left(\sum_{j\ge1}x_j\frac{t^j}{j!}\right)
 =\sum_{n\ge0}\ExponentialCompleteBellPolynomial{n}(x_1,\ldots,x_n)
   \frac{t^n}{n!}.
\]}
\entryfield{Ordinary partial Bell polynomial}{For an ordinary series
\(X(t)=\sum_{j\ge1}x_jt^j\), define
\[
 \OrdinaryPartialBellPolynomial{n}{k}(x_1,x_2,\ldots)
 \DefinitionEquals
 \sum_{\substack{m_1,m_2,\ldots\ge0\\
                   \sum m_j=k\\
                   \sum jm_j=n}}
 \frac{k!}{\prod_{j\ge1}m_j!}\prod_{j\ge1}x_j^{m_j}.
\]
Then
\[
 X(t)^k
 =\sum_{n\ge k}
   \OrdinaryPartialBellPolynomial{n}{k}(x_1,x_2,\ldots)t^n.
\]
The printed \(\mathrm{ord}\) is essential: this polynomial differs from the
exponential normalization by factorials.}
\entryfield{Conversion of normalizations}{For all valid \(n,k\),
\[
 \OrdinaryPartialBellPolynomial{n}{k}(x_1,x_2,\ldots)
 =\frac{k!}{n!}
 \ExponentialPartialBellPolynomial{n}{k}
   (1!x_1,2!x_2,3!x_3,\ldots).
\]
This equality explains the factorial conversion; it does not identify the two
polynomial families.}
\entryfield{Fa\`a di Bruno specialization}{For sufficiently differentiable
functions \(f,g\) and \(n\ge1\),
\[
 \frac{\Differential^n}{\Differential x^n}f(g(x))
 =\sum_{k=1}^{n}f^{(k)}(g(x))
  \ExponentialPartialBellPolynomial{n}{k}
  \bigl(g'(x),g''(x),\ldots,g^{(n-k+1)}(x)\bigr).
\]
The \(n=0\) case is simply \(f(g(x))\).  Derivatives therefore use the
exponential Bell normalization, whereas ordinary series composition uses the
ordinary normalization.}
\end{catalogueentry}

\begin{catalogueentry}{eulerian-families}
  {Type-A, type-B, and second-order Eulerian numbers and row polynomials}
\entryfield{Type-A number}{For \(n\ge1\),
\(\TypeAEulerianNumber{n}{k}\) counts permutations of
\(\{1,\ldots,n\}\) with exactly \(k\) ordinary descents, where a descent is an
index \(i\in\{1,\ldots,n-1\}\) satisfying \(\pi_i>\pi_{i+1}\).  It is zero
outside \(0\le k\le n-1\), with
\(\TypeAEulerianNumber{0}{0}=1\), and obeys
\[
 \TypeAEulerianNumber{n}{k}
 =(k+1)\TypeAEulerianNumber{n-1}{k}
 +(n-k)\TypeAEulerianNumber{n-1}{k-1}.
\]}
\entryfield{Type-A row polynomial}{Set
\[
 \TypeAEulerianPolynomial{0}(t)=1,
 \qquad
 \TypeAEulerianPolynomial{n}(t)
 =\sum_{k=0}^{n-1}\TypeAEulerianNumber{n}{k}t^k
 \quad(n\ge1).
\]
The type tag prevents this row polynomial from colliding with an arbitrary
\(A_n(t)\).}
\entryfield{Type-B descent convention}{A signed permutation of size \(n\) is
a word \(\pi_1\cdots\pi_n\) whose absolute values form a permutation of
\(\{1,\ldots,n\}\).  Put \(\pi_0=0\).  A type-B descent is an
\(i\in\{0,\ldots,n-1\}\) for which \(\pi_i>\pi_{i+1}\).  The number with \(k\)
such descents is \(\TypeBEulerianNumber{n}{k}\), not the type-A angle-bracket
number.}
\entryfield{Type-B recurrence and polynomial}{With
\(\TypeBEulerianNumber{0}{0}=1\) and zero extension,
\[
 \TypeBEulerianNumber{n}{k}
 =(2k+1)\TypeBEulerianNumber{n-1}{k}
 +(2n-2k+1)\TypeBEulerianNumber{n-1}{k-1}.
\]
The row polynomial is
\[
 \TypeBEulerianPolynomial{n}(t)
 \DefinitionEquals
 \sum_{k=0}^{n}\TypeBEulerianNumber{n}{k}t^k,
\]
and satisfies
\[
 \sum_{m\ge0}(2m+1)^nt^m
 =\frac{\TypeBEulerianPolynomial{n}(t)}{(1-t)^{n+1}}.
\]}
\entryfield{Second-order number}{A Stirling permutation of order \(n\) is a
permutation of the multiset
\(\{1,1,2,2,\ldots,n,n\}\) in which every entry between the two copies of
\(i\) is larger than \(i\).  The number having \(k\) ordinary descents is
\(\SecondOrderEulerianNumber{n}{k}\).  For \(n\ge1\),
\[
 \SecondOrderEulerianNumber{n}{k}
 =(2n-k-1)\SecondOrderEulerianNumber{n-1}{k-1}
 +(k+1)\SecondOrderEulerianNumber{n-1}{k}.
\]
The degree-zero convention is
\(\SecondOrderEulerianNumber{0}{0}=1\).}
\entryfield{Second-order row polynomial}{Define
\[
 \SecondOrderEulerianPolynomial{0}(t)=1,
 \qquad
 \SecondOrderEulerianPolynomial{n}(t)
 =\sum_{k=0}^{n-1}\SecondOrderEulerianNumber{n}{k}t^k
 \quad(n\ge1).
\]
Then
\[
 \SecondOrderEulerianPolynomial{n+1}(t)
 =(2nt+1)\SecondOrderEulerianPolynomial{n}(t)
 -t(t-1)\frac{\Differential}{\Differential t}
   \SecondOrderEulerianPolynomial{n}(t),
\]
and
\(\SecondOrderEulerianPolynomial{n}(1)=(2n-1)!!\) for \(n\ge1\).}
\entryfield{Collision boundary}{Single angle brackets mean type A; double
angle brackets mean second order; the explicitly tagged
\(\mathsf{E}^{\mathrm{Eul},B}\)
means type B.  None of these symbols is an inner product.}
\end{catalogueentry}

\begin{catalogueentry}{combinatorial-series-normalizations}
  {Coefficient extraction, EGF/OGF labels, reversion coefficients, and trace inputs}
\entryfield{Coefficient extraction}{If
\(F(z)=\sum_{m\in I}a_mz^m\), where \(I=\NaturalNumbers\) for an ordinary
formal series or \(I\subseteq\IntegerNumbers\) for a Laurent series, then
\[
 \CoefficientExtraction{z}{n}F(z)\DefinitionEquals a_n.
\]
The first argument is the series variable and the second is its exponent; the
series operand follows the bracket.  For a multivariate series,
\([z_1^{n_1}\cdots z_d^{n_d}]F\) extracts the indicated multi-index
coefficient.}
\entryfield{Exponential generating function}{For a sequence
\(A=(a_n)_{n\ge0}\),
\[
 \ExponentialGeneratingFunctionOf{A}(z)
 \DefinitionEquals
 \sum_{n\ge0}a_n\frac{z^n}{n!}.
\]
The superscript \(\mathrm{egf}\) is part of the printed symbol and is not a
Fourier hat.}
\entryfield{Ordinary generating function}{For the same sequence,
\[
 \OrdinaryGeneratingFunctionOf{A}(z)
 \DefinitionEquals
 \sum_{n\ge0}a_nz^n.
\]
The EGF and OGF are different transforms even when they are evaluated at the
same variable.}
\entryfield{Ordinary coefficient ledger}{If
\[
 X(z)=\sum_{j\ge1}x_j\frac{z^j}{j!}
 \quad\text{and}\quad
 \OrdinaryGeneratingFunctionOf{X}(z)
 =\sum_{j\ge1}\OrdinaryGeneratingCoefficient{x}{j}z^j,
\]
then
\[
 \OrdinaryGeneratingCoefficient{x}{j}=\frac{x_j}{j!}.
\]
The family argument \(x\) and the index \(j\) are both explicit; this is not a
reversion coefficient.}
\entryfield{Normalized reversion coefficient}{Suppose
\[
 f(w)=\sum_{j\ge1}f_j\frac{w^j}{j!},
 \qquad f_1\ne0,
 \qquad
 g(z)=f^{\langle-1\rangle}(z)
     =\sum_{n\ge1}g_n\frac{z^n}{n!},
\]
where \(f(g(z))=z=g(f(z))\).  For \(j\ge1\), the reports use
\[
 \NormalizedReversionCoefficient{f}{j}
 \DefinitionEquals
 \frac{f_{j+1}}{(j+1)f_1}
\]
as an input to Bell-polynomial inverse formulas.  Its superscript
\(\mathrm{rev,norm}\) records both reversion and normalization.  It is distinct
from an ordinary-series coefficient and from a factorially rescaled transfer
coefficient.}
\entryfield{Scaled trace input for Bell polynomials}{For a square matrix \(A\)
and \(k\ge1\),
\[
 \TraceBellArgument{k}
 \DefinitionEquals
 -(k-1)!\TraceOperator(A^k).
\]
This is the \(k\)-th scaled trace argument used in determinant/Bell identities.
In particular \(\TraceBellArgument{2}\) is not the binary digit sum
\(\BinaryDigitSum(2)\) and not a signed Stirling number.}
\entryfield{Finite field}{The shared symbol \(\FiniteField_q\) denotes the
finite field with \(q\) elements, where \(q=p^r\) is a prime power.  A bare
\(\FiniteField\) denotes an unspecified finite base field only when its order
is immaterial or declared in the same scope.}
\entryfield{Small operators}{\(\CycleCountOperator(\sigma)\) is the number of
cycles of a permutation \(\sigma\);
\(\SetPartitionOperator(S)\) is the set of set partitions of \(S\); and
\(a\BitwiseAnd b\) is bitwise AND of the nonnegative binary expansions of
integers \(a,b\).  These operators are not multiplication, intersection, or a
generic trace.}
\end{catalogueentry}

\begin{catalogueentry}{partition-lattice-bounds}
  {Partition-lattice endpoints and M\"obius inversion}
\entryfield{Ambient poset}{For a finite set \(S\), let \(\Pi(S)\) be the set of
partitions of \(S\), ordered by refinement:
\(\pi\le\sigma\) means every block of \(\pi\) is contained in a block of
\(\sigma\).  When \(S=\{1,\ldots,n\}\), write \(\Pi_n\).}
\entryfield{Minimum and maximum}{The minimum
\(\PartitionLatticeMinimum\) is the discrete partition into singleton blocks;
the maximum \(\PartitionLatticeMaximum\) is the one-block partition.  Their
\(\Pi\) subscripts are semantic and distinguish them from numeric \(0,1\),
Boolean values, or endpoints of an unrelated poset.}
\entryfield{M\"obius function}{For any locally finite poset, the M\"obius
function is determined by
\[
 \mu(x,x)=1,
 \qquad
 \sum_{x\le z\le y}\mu(x,z)=0\quad(x<y).
\]
For the partition lattice,
\[
 \mu_{\Pi_n}(\pi,\PartitionLatticeMaximum)
 =(-1)^{\CardinalityOf{\pi}-1}
  (\CardinalityOf{\pi}-1)!,
\]
where \(\CardinalityOf{\pi}\) is the number of blocks, and
\[
 \mu_{\Pi_n}(\PartitionLatticeMinimum,\PartitionLatticeMaximum)
 =(-1)^{n-1}(n-1)!\quad(n\ge1).
\]}
\entryfield{Moment--cumulant use}{For random variables indexed by \(S\), a
moment factors over blocks of \(\pi\), and M\"obius inversion on \(\Pi(S)\)
recovers the corresponding joint cumulant.  The lattice endpoints are
combinatorial objects; no numerical substitution for their glyphs is valid.}
\end{catalogueentry}

\begin{catalogueentry}{combinatorial-letter-collisions}
  {Catalan, complementary Bell, Gregory, Genocchi, Barnes, Hessenberg, associahedral,
   and shifted-coordinate symbols}
\entryfield{Catalan number}{For \(n\in\NaturalNumbers\),
\[
 \CatalanNumber{n}
 \DefinitionEquals
 \frac1{n+1}\binom{2n}{n}.
\]
It counts, among many equinumerous families, full plane binary trees with \(n\)
internal vertices and triangulations of a convex \((n+2)\)-gon.  It is not the
complementary Bell number, even though one source used \(C_n\) for both.}
\entryfield{Gregory coefficient}{The Gregory coefficients are defined by the
ordinary Taylor expansion
\[
 \frac{z}{\log(1+z)}
 =1+\sum_{n\ge1}\GregoryCoefficient{n}z^n
 \qquad(|z|<1\text{ analytically}).
\]
The equality also defines them formally over a characteristic-zero coefficient
field.}
\entryfield{Genocchi polynomial and number}{The Genocchi polynomials satisfy
\[
 \frac{2t\EulerE^{xt}}{\EulerE^t+1}
 =\sum_{n\ge0}\GenocchiPolynomial{n}(x)\frac{t^n}{n!},
\]
and
\[
 \GenocchiNumber{n}\DefinitionEquals\GenocchiPolynomial{n}(0).
\]
They are unrelated to the Gregory coefficients despite the legacy shared
letter \(G_n\).}
\entryfield{Barnes G-function}{The Barnes function is denoted
\(\BarnesGFunction(z)\) and normalized by
\[
 \BarnesGFunction(1)=1,
 \qquad
 \BarnesGFunction(z+1)=\GammaFunction(z)\BarnesGFunction(z).
\]
For \(N\in\PositiveIntegers\),
\(\BarnesGFunction(N+1)=\prod_{k=1}^{N-1}k!\).  It is a function, not either
\(G_n\) coefficient family.}
\entryfield{Bell Hessenberg matrices}{For \(1\le i,j\le n\), define
\[
 (\BellHessenbergMatrixH{n})_{ij}
 =\begin{cases}
   \binom{n-i}{j-i}x_{j-i+1},&j\ge i,\\
   -1,&j=i-1,\\
   0,&j<i-1,
  \end{cases}
\]
and
\[
 (\BellHessenbergMatrixK{n})_{ij}
 =\begin{cases}
   x_{j-i+1}/(j-i)!,&j\ge i,\\
   -(i-1),&j=i-1,\\
   0,&j<i-1.
  \end{cases}
\]
Then
\[
 \det\BellHessenbergMatrixH{n}
 =\det\BellHessenbergMatrixK{n}
 =\ExponentialCompleteBellPolynomial{n}(x_1,\ldots,x_n).
\]
The labels \(H\) and \(K\) describe two determinant models for the same Bell
polynomial; neither matrix is an associahedron.}
\entryfield{Associahedron indexed by dimension}{The symbol
\(\AssociahedronByDimension{d}\) denotes the \(d\)-dimensional Stasheff
associahedron.  Its vertices correspond to triangulations of a convex
\((d+3)\)-gon, equivalently to full parenthesizations of \(d+2\) factors.
The lossless legacy-to-canonical crosswalk is:
\(K_n\) of declared dimension \(n-1\) becomes
\(\AssociahedronByDimension{n-1}\);
\(K_n\) of declared dimension \(n-2\) becomes
\(\AssociahedronByDimension{n-2}\);
the corresponding shifted \(K_{n+1}\) becomes
\(\AssociahedronByDimension{n-1}\); and
\(\mathcal{K}_N\) for triangulations of an \(N\)-gon becomes
\(\AssociahedronByDimension{N-3}\).
The unified formulae-and-inversion report is internally mixed: one section
declares \(K_n\) to have dimension \(n-1\), while a later section declares
\(K_{n+1}\) to have dimension \(n-1\), which is the incompatible rule
\(\dim K_m=m-2\).
No source index may be copied without first applying its declared dimension
formula.}
\entryfield{Lagrange shift coordinate}{In a shifted Lagrange-inversion
calculation, the local coordinate
\[
 \LagrangeShiftCoordinate\DefinitionEquals z-b
\]
is a translated variable.  It is not the Riemann zeta function
\(\RiemannZeta(s)\), a Hurwitz zeta value \(\HurwitzZeta{s}{a}\), a root of
unity, an incidence-algebra kernel, or a generic complex coordinate outside
that declared scope.}
\entryfield{Zeta-family boundary}{Five roles must remain separate:
\(\RiemannZeta(s)\) is the one-argument Riemann zeta function;
\(\HurwitzZeta{s}{a}\) is the two-argument Hurwitz zeta function;
\(\LagrangeShiftCoordinate=z-b\) is a local translated coordinate;
\(\zeta_P(x,y)=1\) for \(x\le y\) is the incidence-algebra zeta function of a
locally finite poset \(P\), whose convolution inverse is the M\"obius
function; and a root-of-unity variable satisfies \(\zeta^N=1\) for a declared
positive integer \(N\).  Boolean-subset zeta transforms are the specialization
of the poset construction to the subset lattice.  Their shared historical
glyph does not identify the objects.}
\entryfield{Governance rule}{A raw \(C_n\), \(G_n\), \(K_n\), \(L(n,k)\),
\(T_n(x)\), or local \(\zeta\) is insufficient wherever two listed meanings
coexist.  The semantic commands in this entry are normative and their printed
tags are intentionally visible.}
\end{catalogueentry}

\section{\texorpdfstring{\(q\)-series}{q-series}, products, and arithmetic functions}

\begin{catalogueentry}{q-parameter}{The \(q\) parameter, formal series, and analytic
regimes}
\entryfield{Polynomial and formal regimes}{The letter \(q\) may be an indeterminate.
Then \(R[q]\) is the polynomial ring and
\[
 R[[q]]
 =\left\{\sum_{n=0}^{\infty}a_nq^n:a_n\in R\right\}
\]
is the ring of formal power series over a commutative ring \(R\).  The infinity sign in a
formal series records an unbounded coefficient sequence, not analytic convergence.}
\entryfield{Equality and Cauchy product}{Formal equality is coefficientwise:
\(\sum_{n\ge0}a_nq^n=\sum_{n\ge0}b_nq^n\) exactly when \(a_n=b_n\) for every
\(n\).  Addition is coefficientwise and multiplication is
\[
 \left(\sum_{n\ge0}a_nq^n\right)
 \left(\sum_{n\ge0}b_nq^n\right)
 =\sum_{n\ge0}\left(\sum_{j=0}^{n}a_jb_{n-j}\right)q^n.
\]
The inner sum is finite, so no convergence hypothesis is hidden in this definition.}
\entryfield{Formal Laurent series}{The ring \(R((q))\) consists of
\(\sum_{n\ge n_0}a_nq^n\) with \(n_0\in\IntegerNumbers\) and only finitely many
negative powers.  It is distinct from a bilateral analytic Laurent series with infinitely
many terms in both directions.}
\entryfield{Terminating rational regime}{A terminating identity may be interpreted in the
rational-function field \(K(q)\), or after clearing explicitly named denominators in
\(K[q]\).  Termination removes infinite-series convergence questions but does not remove
poles caused by a zero denominator.}
\entryfield{Analytic regime}{For infinite products and unilateral or bilateral series, an
explicit common regime is \(q\in\ComplexNumbers\) with \(0<|q|<1\); products permitting
\(q=0\) say so separately.  Probability models normally strengthen this to \(q\in(0,1)\).
Every convergence disk or annulus also includes the remaining parameters.}
\entryfield{Modular parameter and nome}{For theta and eta functions,
\(\tau\in\mathbb H=\{z\in\ComplexNumbers:\operatorname{Im}z>0\}\) and the nome is
\[
 q_\tau=\EulerE^{\pi\ImaginaryUnit\tau}.
\]
Thus \(|q_\tau|<1\).  Some sources use
\(\EulerE^{2\pi\ImaginaryUnit\tau}=q_\tau^2\); a formula must identify which nome is in
force.  Under \(\tau\mapsto-1/\tau\), the transformed nome is
\(\widetilde q_\tau=\exp(-\pi\ImaginaryUnit/\tau)\).}
\entryfield{Roots of unity}{At \(q=\zeta\) with \(\zeta^N=1\), denominators in rational
\(q\)-expressions can vanish.  Evaluation means polynomial evaluation after cancellation
only when a polynomial representative has been proved; otherwise the singular value is
excluded.}
\entryfield{Base-role rule}{If a document also uses \(q\) for a denominator, quantile,
frequency, conjugate exponent, or finite-field cardinality, the series base is
\(q_{\mathrm{base}}\) and the field cardinality is \(Q\).  Semantic subscripts take
priority over traditional brevity.}
\end{catalogueentry}

\begin{catalogueentry}{q-pochhammer}{Finite, infinite, integer-index, and complex-order
\(q\)-Pochhammer symbols}
\entryfield{Finite definition}{For \(a,q\) in a commutative ring and
\(n\in\NaturalNumbers\),
\[
 \QPochhammer a q n\DefinitionEquals
 \prod_{j=0}^{n-1}(1-aq^j).
\]
The product is empty at \(n=0\), so \(\QPochhammer a q0=1\) for every \(a,q\), including
values that make later factors vanish.}
\entryfield{Several parameters}{For a finite ordered list
\(\boldsymbol a=(a_1,\ldots,a_r)\),
\[
 (a_1,\ldots,a_r;q)_n
 \DefinitionEquals\prod_{i=1}^{r}\QPochhammer{a_i}qn.
\]
For \(r=0\) this is the empty product \(1\).  The punctuation separates a parameter list
from the base; it does not make the \(a_i\) into additional bases.}
\entryfield{Infinite definition}{For \(a,q\in\ComplexNumbers\) and \(|q|<1\),
\[
 \QPochhammer a q\infty\DefinitionEquals
 \prod_{j=0}^{\infty}(1-aq^j),
\]
with analytic meaning as a locally uniformly convergent limit of finite products.  The
literal index \(\infty\) denotes
an infinite product, not an extended-natural argument to the finite definition.}
\entryfield{Positive and negative integer indices}{For \(m\in\PositiveIntegers\),
\[
 (a;q)_{-m}
 \DefinitionEquals\frac{1}{(aq^{-m};q)_m}
 =\frac{(-q/a)^m q^{\binom m2}}{(q/a;q)_m},
\]
where all displayed denominators are nonzero.  Together with the finite definition this
gives \((a;q)_{m+n}=(a;q)_m(aq^m;q)_n\) for integer indices whenever both sides are
defined.}
\entryfield{Complex order}{Fix \(q\in\ComplexNumbers\setminus\{0\}\), \(|q|<1\), and a
logarithm value \(L_q\) satisfying \(\exp(L_q)=q\).  Put
\(q^\nu=\exp(\nu L_q)\).  For
\(\nu\in\ComplexNumbers\),
\[
 (a;q)_\nu
 \DefinitionEquals\frac{(a;q)_\infty}{(aq^\nu;q)_\infty}
\]
where the denominator is nonzero.  Changing the branch can change \(q^\nu\), hence the
value.  For \(\nu\in\IntegerNumbers\) the definition agrees with the integer-index
convention on their common domain.}
\entryfield{The base \(q=0\)}{For finite indices,
\((a;0)_0=1\) and \((a;0)_n=1-a\) for \(n\ge1\), using \(0^0=1\) in the first factor and
\(0^j=0\) for \(j>0\).  Negative and complex orders at \(q=0\) are not inferred from
formulas containing \(q^{-m}\) or a logarithm of \(q\).}
\entryfield{Factorization}{For \(m,n\in\NaturalNumbers\),
\[
 \QPochhammer a q{m+n}
 =\QPochhammer a q m\,\QPochhammer{aq^m}q n.
\]
This is a finite product identity over a commutative ring and needs no division.}
\entryfield{Retired signatures}{Former two-argument macros hid the base and are retired.
The shared \commandname{QPochhammer} always receives \((a,q,n)\) in that order.
Monograph-local extensions render the same glyph but declare integer or complex order at
the governing definition.}
\end{catalogueentry}

\begin{catalogueentry}{q-integer}{\(q\)-integers and \(q\)-factorials}
\entryfield{Definition}{For $n\in\NaturalNumbers$,
\[
 \QInteger n q\DefinitionEquals
 \begin{cases}(1-q^n)/(1-q),&q\ne1,\\n,&q=1.\end{cases}
\]
Equivalently $\QInteger n q=1+q+\cdots+q^{n-1}$ for $n>0$, and
$\QInteger0q=0$.}
\entryfield{\(q\)-factorial}{$[n]_q!\DefinitionEquals\prod_{j=1}^{n}\QInteger j q$,
with $[0]_q!=1$.  The exclamation mark belongs to the notation and does not apply an
ordinary factorial to the polynomial $[n]_q$.}
\entryfield{Limit}{$\lim_{q\to1}\QInteger n q=n$ and
$\lim_{q\to1}[n]_q!=n!$ for fixed $n$.}
\entryfield{Alternative symmetric integer}{The symmetric quantity
$(q^n-q^{-n})/(q-q^{-1})$ is written $[n]_{q,\mathrm{sym}}$ if needed; it is not the
canonical \commandname{QInteger}.}
\end{catalogueentry}

\begin{catalogueentry}{generalized-q-number}{Generalized \(q\)-numbers and their inverse}
\entryfield{Branch data}{Fix \(q\in\ComplexNumbers\setminus\{0,1\}\) and a named
logarithm value \(L_q\) satisfying \(\exp(L_q)=q\).  For \(z\in\ComplexNumbers\), define
\(q^z=\exp(zL_q)\).  If \(q\in(0,1)\), the default is the real logarithm
\(L_q=\log q<0\).}
\entryfield{Definition}{The monograph-local command
\commandname{GeneralizedQNumber}\texttt{\{z\}\{q\}} prints
\[
 [z]_q\DefinitionEquals\frac{1-q^z}{1-q}.
\]
At \(q=1\), the separately assigned continuous value is \([z]_1=z\).  For
\(z=n\in\NaturalNumbers\), this agrees with \(\QInteger n q\).}
\entryfield{Zeros and branch dependence}{For \(q\ne1\), \([z]_q=0\) exactly when
\(\exp(zL_q)=1\).  In the real regime \(0<q<1\) and \(z\in\RealNumbers\), this means
\(z=0\).  For complex \(z\), changing \(L_q\) can change both the function and its zero
lattice.}
\entryfield{Inverse on a chosen branch}{If \(w=[z]_q\), then
\[
 q^z=1-(1-q)w,\qquad
 z=\frac{\log_{\!*}(1-(1-q)w)}{L_q},
\]
where \(\log_{\!*}\) is a specifically chosen local branch and its domain excludes zero
and the branch cut.  The formula produces one inverse branch, not a global single-valued
inverse on \(\ComplexNumbers\).}
\entryfield{Classical limit}{For fixed \(z\) and a compatible logarithm branch near
\(q=1\), \([z]_q\to z\).  A uniform limit over unbounded \(z\)-sets is not implicit.}
\end{catalogueentry}

\begin{catalogueentry}{q-falling-power}{\(q\)-falling powers and \(q\)-Newton bases}
\entryfield{Definition}{For scalars \(x,a,q\) in a commutative ring and
\(m\in\NaturalNumbers\),
\[
 (x-a)_{m}^{\langle q\rangle}
 \DefinitionEquals\prod_{j=0}^{m-1}(x-q^ja),
 \qquad (x-a)_{0}^{\langle q\rangle}=1.
\]
The monograph-local command is
\commandname{QFallingPower}\texttt{\{x\}\{a\}\{q\}\{m\}}.  All four inputs are
semantic; in particular the base is never inferred from a surrounding chapter.}
\entryfield{Relation to \(q\)-Pochhammer}{When \(x\) is invertible,
\[
 (x-a)_{m}^{\langle q\rangle}
 =x^m(a/x;q)_m.
\]
The product definition remains valid without invertibility and is therefore primary.}
\entryfield{Jackson derivative}{With the derivative \(D_q\) from the
\texttt{q-difference} entry,
\[
 D_q (x-a)_{m}^{\langle q\rangle}
 =[m]_q(x-a)_{m-1}^{\langle q\rangle}
\]
for \(m\ge1\).  Iteration supplies
\(D_q^j(x-a)_{m}^{\langle q\rangle}
 =[m]_q!/[m-j]_q!\,(x-a)_{m-j}^{\langle q\rangle}\) only for
\(0\le j\le m\) and where the quotient notation has a valid interpretation.}
\entryfield{Basis role}{For fixed \(a,q\), these are monic polynomials of distinct
degrees.  Thus the first \(N+1\) form a basis for polynomials of degree at most \(N\).
The phrase \(q\)-Newton basis refers to this declared family, not to an unspecified
ordinary falling-factorial basis.}
\end{catalogueentry}

\begin{catalogueentry}{gaussian-binomial}{Gaussian binomial coefficients}
\entryfield{Definition}{For $n,k\in\NaturalNumbers$,
\[
 \GaussianBinomial n k q\DefinitionEquals
 \begin{cases}
 \dfrac{\QPochhammer q q n}
       {\QPochhammer q q k\,\QPochhammer q q{n-k}},
       &0\le k\le n,\\[0.8em]
 0,&k>n,
 \end{cases}
\]
interpreted as its polynomial in $q$, so removable singularities are filled.}
\entryfield{Boundary values}{$\GaussianBinomial n0q
=\GaussianBinomial nnq=1$, including $n=0$.}
\entryfield{Recurrence}{For $1\le k\le n$,
\[
 \GaussianBinomial n k q
 =\GaussianBinomial{n-1}kq+
 q^{n-k}\GaussianBinomial{n-1}{k-1}q.
\]}
\entryfield{Retired signatures}{The corpus formerly used incompatible argument orders
for its abbreviated Gaussian-binomial macro.  The canonical
\commandname{GaussianBinomial} receives upper index, lower index, then base.}
\entryfield{Limit}{$\GaussianBinomial n k1=\binom nk$, understood by polynomial
evaluation, not substitution into the displayed rational quotient.}
\end{catalogueentry}

\begin{catalogueentry}{generalized-gaussian-binomial}{Generalized Gaussian binomial
coefficients}
\entryfield{Complex upper parameter}{Fix a nonzero complex base \(q\), a logarithm value
\(L_q\), and \(q^\alpha=\exp(\alpha L_q)\).  For
\(\alpha\in\ComplexNumbers\) and \(k\in\NaturalNumbers\), define
\[
 \genfrac{[}{]}{0pt}{}{\alpha}{k}_{q}
 \DefinitionEquals
 \frac{(q^{\alpha-k+1};q)_k}{(q;q)_k},
\]
where the denominator is nonzero.  The monograph-local command is
\par\smallskip\noindent
{\small\commandname{GeneralizedGaussianBinomial}}.
\par\smallskip\noindent
Its arguments are the upper parameter, lower parameter, and base, in that order.}
\entryfield{Integer specialization}{When \(\alpha=n\in\NaturalNumbers\) and
\(0\le k\le n\), the generalized expression equals
\(\GaussianBinomial n k q\).  For \(k>n\), the canonical finite Gaussian coefficient is
the zero polynomial by definition; this boundary convention is not silently extended to
arbitrary complex \(\alpha\).}
\entryfield{Two complex arguments}{A separately declared meromorphic continuation may use
\[
 \genfrac{[}{]}{0pt}{}{\alpha}{\beta}_{q}
 \DefinitionEquals
 \frac{\Gamma_q(\alpha+1)}
 {\Gamma_q(\beta+1)\Gamma_q(\alpha-\beta+1)}
\]
in a regime where the named \(q\)-gamma function and quotient are defined.  This is not a
polynomial in \(q\), depends on the \(q\)-gamma convention, and has no automatic
out-of-range zero rule.}
\entryfield{No implicit regime change}{Polynomial identities, rational identities,
branch-dependent complex formulas, and meromorphic continuations are four distinct
claims.  A source moving between them states the domain and continuation theorem rather
than reusing the bracket alone.}
\end{catalogueentry}

\begin{catalogueentry}{q-multinomial-path-area}{\(q\)-multinomials, inversions, and lattice
path area}
\entryfield{Definition and domain}{For \(r\in\NaturalNumbers\),
\(n,n_1,\ldots,n_r\in\NaturalNumbers\), and
\(n_1+\cdots+n_r=n\),
\[
 \genfrac{[}{]}{0pt}{}{n}{n_1,\ldots,n_r}_{q}
 \DefinitionEquals
 \frac{(q;q)_n}{(q;q)_{n_1}\cdots(q;q)_{n_r}},
\]
interpreted as its polynomial in \(q\).  For \(r=0\), the constraint forces \(n=0\) and
the value is \(1\).  If an input is negative or the sum constraint fails, the expression
is undefined unless that source separately declares a zero extension.}
\entryfield{Factorization}{For a valid composition,
\[
 \genfrac{[}{]}{0pt}{}{n}{n_1,\ldots,n_r}_{q}
 =\genfrac{[}{]}{0pt}{}{n}{n_1}_{q}
  \genfrac{[}{]}{0pt}{}{n-n_1}{n_2,\ldots,n_r}_{q}.
\]
The monograph-local \commandname{QMultinomial} receives \(n\), the complete lower list,
and \(q\); the list length and sum are not macro-inferred data.}
\entryfield{Word inversion statistic}{Let \(W(n_1,\ldots,n_r)\) be the words on the
ordered alphabet \(1<\cdots<r\) in which \(i\) occurs \(n_i\) times.  For
\(w=w_1\cdots w_n\),
\[
 \operatorname{inv}(w)
 =\CardinalityOf{\{(i,j):1\le i<j\le n,\ w_i>w_j\}}.
\]
Then the \(q\)-multinomial is
\(\sum_{w\in W(n_1,\ldots,n_r)}q^{\operatorname{inv}(w)}\).}
\entryfield{Binary lattice-path convention}{For a binary word, \(0\) is an east step and
\(1\) is a north step.  Its area is any of the equal quantities
\[
 \sum_{\substack{i:w_i=0}}\CardinalityOf{\{j<i:w_j=1\}}
 =\sum_{\substack{j:w_j=1}}\CardinalityOf{\{i>j:w_i=0\}}
 =\operatorname{inv}(w),
\]
the number of unit cells below the path in its \(m\)-by-\(k\) rectangle.  Hence
\(\operatorname{area}(0^m1^k)=0\) and
\(\operatorname{area}(1^k0^m)=mk\).  A complementary-area convention must say so and
replaces \(q\) by \(q^{-1}\) with the corresponding normalizing power.}
\end{catalogueentry}

\begin{catalogueentry}{finite-fields-flags}{Finite fields, general linear groups,
subspaces, and flags}
\entryfield{Prime-power parameter}{Here \(Q=p^e\), where \(p\) is prime and
\(e\in\PositiveIntegers\).  The symbol \(\mathbb F_Q\) is a field with exactly \(Q\)
elements, unique up to noncanonical field isomorphism.  Upper-case \(Q\) prevents a
collision with the \(q\)-series base.}
\entryfield{Vector spaces and zero dimension}{For \(n\in\NaturalNumbers\),
\(\mathbb F_Q^n\) is the \(n\)-dimensional coordinate vector space.
\(\mathbb F_Q^0=\{0\}\) is the one-element zero vector space, not the empty set.}
\entryfield{General linear group}{For \(r\in\NaturalNumbers\),
\[
 \operatorname{GL}_r(\mathbb F_Q)
 =\{A\in\operatorname{Mat}_{r\times r}(\mathbb F_Q):\det A\ne0\}.
\]
Its cardinality is
\(\prod_{j=0}^{r-1}(Q^r-Q^j)\); at \(r=0\), the empty product is \(1\) and
\(\operatorname{GL}_0(\mathbb F_Q)\) contains the unique empty identity matrix.}
\entryfield{Subspace count}{For \(0\le k\le n\),
\[
 \CardinalityOf{\{U\le\mathbb F_Q^n:\dim_{\mathbb F_Q}U=k\}}
 =\GaussianBinomial n k Q.
\]
The bracket is evaluated at the prime power \(Q\); it is not a second definition of the
field.}
\entryfield{Flags}{For \(n_1+\cdots+n_r=n\), the \(Q\)-multinomial counts flags
\[
 0=V_0\subseteq V_1\subseteq\cdots\subseteq V_r=\mathbb F_Q^n,
 \qquad \dim(V_i/V_{i-1})=n_i.
\]
All inclusions and quotient dimensions are part of the object being counted.}
\entryfield{Weight character}{If a finite-dimensional representation \(V\) has an
integral-weight decomposition \(V=\bigoplus_{m\in\IntegerNumbers}V_m\), its
monograph-local character is
\[
 \operatorname{ch}_{x}(V)
 \DefinitionEquals\sum_{m\in\IntegerNumbers}(\dim V_m)x^m.
\]
The sum is finite, repeated weights contribute multiplicity, \(x\) is a formal variable,
and \(\operatorname{ch}_{x}(0)=0\).  The local command
\commandname{WeightCharacter} receives the representation and the formal variable.}
\end{catalogueentry}

\begin{catalogueentry}{subspace-mobius}{Möbius function of a finite subspace lattice}
\entryfield{Poset definition}{For a finite-dimensional \(\mathbb F_Q\)-space \(V\), let
\(\mathcal L(V)\) be its poset of linear subspaces ordered by inclusion.  Its Möbius
function is defined on comparable pairs by
\[
 \mu_{\mathcal L(V)}(U,U)=1,\qquad
 \sum_{U\subseteq W\subseteq Z}\mu_{\mathcal L(V)}(U,W)=0
 \quad(U\subsetneq Z).
\]
The ambient space and both interval endpoints are mandatory inputs.}
\entryfield{Closed form}{If \(U\subseteq Z\) and
\(d=\dim_{\mathbb F_Q}(Z/U)\), then
\[
 \mu_{\mathcal L(V)}(U,Z)=(-1)^dQ^{\binom d2}.
\]
For incomparable \(U,Z\), the incidence-algebra extension is \(0\); alternatively the
two-argument poset Möbius function is simply left undefined off comparable pairs.  The
chosen convention must be stated when off-interval inputs can occur.}
\entryfield{Semantic command}{The monograph-local
\commandname{SubspaceMobius}\texttt{\{V\}\{U\}\{Z\}} prints the ambient lattice and both
endpoints.  A bare \(\mu\) is reserved neither for this function nor for the arithmetic
Möbius function because the corpus also uses \(\mu\) for measures and means.}
\end{catalogueentry}

\begin{catalogueentry}{basic-hypergeometric}{Unilateral and bilateral basic
hypergeometric series}
\entryfield{Unilateral definition}{The unilateral basic hypergeometric series is
\[
 {}_r\phi_s\!\left(
 \begin{matrix}a_1,\ldots,a_r\\b_1,\ldots,b_s\end{matrix};
 q,z\right)
 \DefinitionEquals
 \sum_{n=0}^{\infty}
 \frac{\prod_{j=1}^r\QPochhammer{a_j}qn}
      {\QPochhammer q q n\prod_{j=1}^s\QPochhammer{b_j}qn}
 \left((-1)^nq^{\binom n2}\right)^{1+s-r}z^n,
\]
provided no denominator factor vanishes.  A terminating instance is a finite rational
identity.  A nonterminating analytic instance additionally states a convergence domain,
normally with \(0<|q|<1\).}
\entryfield{Normalization warning}{References differ on the
$((-1)^nq^{\binom n2})^{1+s-r}$ factor.  Every occurrence in this corpus either expands
the definition above or cites this entry.}
\entryfield{Termination}{If an upper parameter is $q^{-N}$ with
$N\in\NaturalNumbers$, the series terminates at $n=N$ unless an earlier denominator
singularity makes the expression undefined.}
\entryfield{Bilateral definition}{With integer-index \(q\)-Pochhammers already defined,
\[
 {}_r\psi_s\!\left(
 \begin{matrix}a_1,\ldots,a_r\\b_1,\ldots,b_s\end{matrix};
 q,z\right)
 \DefinitionEquals
 \sum_{n\in\IntegerNumbers}
 \frac{\prod_{j=1}^{r}(a_j;q)_n}
      {\prod_{j=1}^{s}(b_j;q)_n}
 \left((-1)^nq^{\binom n2}\right)^{s-r}z^n.
\]
This is never identified with the unilateral symbol.  Every analytic use supplies an
annulus of convergence and excludes denominator zeros.}
\entryfield{Balanced series}{A terminating
\({}_{r+1}\phi_r(a_1,\ldots,a_{r+1};b_1,\ldots,b_r;q,q)\) is called balanced when
\[
 b_1\cdots b_r=q\,a_1\cdots a_{r+1}.
\]
If \(a_{r+1}=q^{-N}\), the condition is equivalently
\(b_1\cdots b_r=q^{1-N}a_1\cdots a_r\).  Both the argument \(z=q\) and termination index
\(N\) are part of the label.}
\entryfield{Well-poised and very-well-poised}{The series
\({}_r\phi_{r-1}(a_1,\ldots,a_r;b_1,\ldots,b_{r-1};q,z)\) is well-poised when
\[
 qa_1=a_2b_1=a_3b_2=\cdots=a_rb_{r-1}.
\]
It is very-well-poised when, in addition, a chosen square root satisfies
\(a_2=q\sqrt{a_1}\) and \(a_3=-q\sqrt{a_1}\).  The paired signs make the resulting
parameter pair independent of replacing the square root by its negative, but a derivation
that separates the two parameters still names the choice.}
\entryfield{Ramanujan's \({}_1\psi_1\) summation}{For \(0<|q|<1\) and
\(|b/a|<|z|<1\), with all denominator products nonzero,
\[
 {}_1\psi_1\!\left(\begin{matrix}a\\b\end{matrix};q,z\right)
 =
 \frac{(q,b/a,az,q/(az);q)_\infty}
      {(b,q/a,z,b/(az);q)_\infty}.
\]
The annulus is essential: it controls the two tails of the bilateral series and cannot be
replaced by the unilateral condition \(|z|<1\) alone.}
\end{catalogueentry}

\begin{catalogueentry}{q-difference}{\(q\)-difference operators}
\entryfield{Jackson derivative}{For $q\ne1$ and $x\ne0$,
\[
 D_qf(x)\DefinitionEquals\frac{f(x)-f(qx)}{(1-q)x}.
\]
At $x=0$, $D_qf(0)$ is assigned $f'(0)$ only when that derivative exists and the local
definition explicitly makes the continuous extension.}
\entryfield{Shift operator}{$T_qf(x)=f(qx)$.  Thus
$D_q=(1-T_q)/((1-q)x)$ as an identity on functions where division by $x$ is valid.}
\entryfield{Base direction}{The expression
$(f(qx)-f(x))/((q-1)x)$ is algebraically identical.  The genuinely different operator
$(f(qx)-f(x))/((q-1)qx)$ is labelled separately.}
\entryfield{Classical limit}{If $f$ is differentiable at $x$, then
$D_qf(x)\to f'(x)$ as $q\to1$ through real positive values.  Uniformity in $x$ is never
inferred without a theorem.}
\end{catalogueentry}

\begin{catalogueentry}{jackson-integral}{Jackson \(q\)-integrals and differentials}
\entryfield{Integral from zero}{For \(0<q<1\), a scalar endpoint \(a\), and a function for
which the series converges,
\[
 \int_0^a f(t)\,\mathrm d_qt
 \DefinitionEquals
 a(1-q)\sum_{n=0}^{\infty}q^n f(aq^n).
\]
For \(a=0\) the value is \(0\).  The monograph-local
\commandname{JacksonDifferential}\texttt{\{q\}\{t\}} prints
\(\mathrm d_qt\), so neither the base nor integration variable is implicit.}
\entryfield{General oriented interval}{Define
\[
 \int_a^b f(t)\,\mathrm d_qt
 \DefinitionEquals
 \int_0^b f(t)\,\mathrm d_qt-\int_0^a f(t)\,\mathrm d_qt.
\]
This is an oriented difference, not a measure integral over an empty set when \(b<a\).
For complex endpoints it samples two geometric rays and requires convergence on both.}
\entryfield{Fundamental theorem}{If \(G(aq^n)\to G(0)\) and the telescoping series is
legitimate, then
\[
 \int_0^a D_qG(t)\,\mathrm d_qt=G(a)-G(0).
\]
The limit hypothesis is part of the statement; the isolated formula for \(D_q\) does not
create continuity at zero.}
\entryfield{Integration by parts}{Under absolute convergence sufficient to rearrange the
sampled series,
\[
 \int_0^a f(t)D_qg(t)\,\mathrm d_qt
 =f(a)g(a)-f(0)g(0)
  -\int_0^a g(qt)D_qf(t)\,\mathrm d_qt.
\]
The shifted factor \(g(qt)\) is essential and distinguishes Jackson integration from the
ordinary Leibniz rule.}
\entryfield{Power integral}{For \(a>0\), \(s\in\ComplexNumbers\), and a fixed branch for
\(a^{s+1}\), the geometric series gives
\[
 \int_0^a t^s\,\mathrm d_qt
 =\frac{a^{s+1}}{[s+1]_q}
\]
when \(|q^{s+1}|<1\).  In the real regime \(0<q<1\), this includes
\(\operatorname{Re}s>-1\).}
\end{catalogueentry}

\begin{catalogueentry}{q-gamma-beta}{\(q\)-gamma, \(q\)-digamma, and \(q\)-beta
functions}
\entryfield{\(0<q<1\) gamma convention}{For \(0<q<1\), define the meromorphic function
\[
 \Gamma_q(z)
 \DefinitionEquals
 (1-q)^{1-z}\frac{(q;q)_\infty}{(q^z;q)_\infty},
 \qquad q^z=\exp(z\log q).
\]
The power of \(1-q>0\) uses the real logarithm.  For \(z>0\) it is finite and positive,
\(\Gamma_q(1)=1\), and
\(\Gamma_q(z+1)=[z]_q\Gamma_q(z)\).}
\entryfield{Poles and zeros}{In the real-base convention, the denominator vanishes when
\(q^{z+j}=1\) for some \(j\in\NaturalNumbers\).  Thus the meromorphic pole lattice is
\[
 z=-j+\frac{2\pi\ImaginaryUnit k}{\log q},
 \qquad j\in\NaturalNumbers,\ k\in\IntegerNumbers,
\]
subject to cancellation only if proved.  The function has no zeros because its numerator
is nonzero.}
\entryfield{\(q>1\) reciprocal-base convention}{For real \(q>1\), the separately declared
convention is
\[
 \Gamma_q(z)
 =(q-1)^{1-z}q^{z(z-1)/2}
 \frac{(q^{-1};q^{-1})_\infty}{(q^{-z};q^{-1})_\infty},
\]
with real logarithms for positive bases.  It is not obtained by inserting \(q>1\) into
the divergent \(0<q<1\) product.}
\entryfield{\(q\)-digamma}{At points where \(\Gamma_q\) is finite and nonzero,
\[
 \psi_q(z)\DefinitionEquals\frac{\Differential}{\Differential z}
 \log\Gamma_q(z)=\frac{\Gamma_q'(z)}{\Gamma_q(z)}.
\]
The logarithmic derivative is single-valued even though a local logarithm of
\(\Gamma_q\) was used to motivate it.  This \(\psi_q\) is not a characteristic function
or a wave function.}
\entryfield{\(q\)-beta}{For \(x,y>0\) and \(0<q<1\),
\[
 B_q(x,y)
 \DefinitionEquals\frac{\Gamma_q(x)\Gamma_q(y)}{\Gamma_q(x+y)}
 =\int_0^1 t^{x-1}
   \frac{(qt;q)_\infty}{(q^yt;q)_\infty}\,\mathrm d_qt.
\]
The integral formula includes the Jackson differential and real powers determined by
\(t\in(0,1]\); a complex continuation uses the gamma quotient, not an undeclared
pointwise integral branch.}
\entryfield{Classical limits}{For fixed \(z\) away from poles,
\(\Gamma_q(z)\to\GammaFunction(z)\), \(\psi_q(z)\to\psi(z)\), and
\(B_q(x,y)\to B(x,y)\) as \(q\uparrow1\), on compact subsets where the corresponding
uniform convergence theorem applies.}
\end{catalogueentry}

\begin{catalogueentry}{q-exponentials}{Euler \(q\)-exponentials and factorial-scaled
variants}
\entryfield{Pochhammer-scaled pair}{For \(0<|q|<1\),
\[
 e_q(z)\DefinitionEquals\sum_{n=0}^{\infty}\frac{z^n}{(q;q)_n}
 =\frac1{(z;q)_\infty}\quad(|z|<1),
\]
and
\[
 E_q(z)\DefinitionEquals
 \sum_{n=0}^{\infty}\frac{q^{\binom n2}z^n}{(q;q)_n}
 =(-z;q)_\infty\quad(z\in\ComplexNumbers).
\]
The lower-case series has poles after meromorphic continuation; the upper-case product is
entire in \(z\).}
\entryfield{\(q\)-factorial-scaled pair}{Since
\([n]_q!=(q;q)_n/(1-q)^n\), define
\[
 \mathrm e_q^{\,J}(z)
 =\sum_{n=0}^{\infty}\frac{z^n}{[n]_q!}
 =\frac1{((1-q)z;q)_\infty},
\]
\[
 \mathrm E_q^{\,J}(z)
 =\sum_{n=0}^{\infty}\frac{q^{\binom n2}z^n}{[n]_q!}
 =(-(1-q)z;q)_\infty.
\]
The superscript \(J\) labels the Jackson-factorial scaling; it is not a power.}
\entryfield{Reciprocity}{On the common domain,
\[
 e_q(z)E_q(-z)=1,\qquad
 \mathrm e_q^{\,J}(z)\mathrm E_q^{\,J}(-z)=1.
\]
Changing between the two pairs requires the explicit substitution \(z\mapsto(1-q)z\);
their identical spoken name does not make their arguments interchangeable.}
\entryfield{Classical normalization}{For fixed \(z\),
\(\mathrm e_q^{\,J}(z)\to\EulerE^z\) and
\(\mathrm E_q^{\,J}(z)\to\EulerE^z\) as \(q\uparrow1\).  The unscaled \(e_q\) and
\(E_q\) require an argument rescaling to have that limit.}
\end{catalogueentry}

\begin{catalogueentry}{multiplicative-bilateral-theta}{Multiplicative and bilateral
Jacobi theta normalizations}
\entryfield{Multiplicative product}{For \(0<|q|<1\) and
\(z\in\ComplexNumbers^\times\),
\[
 \vartheta_{\mathrm{mult}}(z;q)
 \DefinitionEquals(z;q)_\infty(q/z;q)_\infty(q;q)_\infty.
\]
The monograph-local \commandname{MultiplicativeTheta} receives \(z\) and \(q\); its
semantic subscript distinguishes this function from the additive Jacobi functions
\(\vartheta_2,\vartheta_3,\vartheta_4\).}
\entryfield{Plus-sign bilateral series}{On the same domain,
\[
 \Theta_{\mathrm{bil}}(z;q)
 \DefinitionEquals\sum_{n\in\IntegerNumbers}q^{\binom n2}z^n.
\]
The series converges normally on compact subsets of
\(\ComplexNumbers^\times\).  Its local command \commandname{BilateralTheta} receives
argument and base.}
\entryfield{Exact normalization crosswalk}{Jacobi's triple product gives
\[
 \vartheta_{\mathrm{mult}}(z;q)=\Theta_{\mathrm{bil}}(-z;q).
\]
The argument negation is mandatory: suppressing it changes every odd Laurent
coefficient.}
\entryfield{Quasi-periodicity and zeros}{One has
\[
 \Theta_{\mathrm{bil}}(qz;q)=z^{-1}\Theta_{\mathrm{bil}}(z;q),
 \qquad
 \vartheta_{\mathrm{mult}}(qz;q)=-z^{-1}
 \vartheta_{\mathrm{mult}}(z;q).
\]
The bilateral zeros are the simple points \(z=-q^m\), while the multiplicative-product
zeros are the simple points \(z=q^m\), \(m\in\IntegerNumbers\).}
\entryfield{Inverse near a zero}{At \(z=q^m\),
\(\partial_z\vartheta_{\mathrm{mult}}(q^m;q)\ne0\), so the local inverse at target
\(y=0\) begins
\[
 z=q^m+\frac{y}{\partial_z\vartheta_{\mathrm{mult}}(q^m;q)}
 +\BigO(y^2).
\]
This is one local analytic branch around a specified zero, not a global inverse of a
quasi-periodic function.}
\end{catalogueentry}

\begin{catalogueentry}{jacobi-theta}{Jacobi theta functions, products, heat equation, and
modular inversion}
\entryfield{Parameters and prescribed lifts}{Let
\(\tau\in\mathbb H\), \(z\in\ComplexNumbers\),
\[
 q=\EulerE^{\pi\ImaginaryUnit\tau},\qquad
 \zeta=\EulerE^{2\pi\ImaginaryUnit z}.
\]
Every rational power is the prescribed lift
\(q^\alpha=\EulerE^{\pi\ImaginaryUnit\alpha\tau}\), and
\(\zeta^{1/2}\) means \(\EulerE^{\pi\ImaginaryUnit z}\).  Neither is obtained from an
unspecified logarithm of its displayed complex value.}
\entryfield{Series definitions}{The three normalizations used in the monograph are
\begin{align*}
 \vartheta_3(z\mid\tau)
 &=\sum_{n\in\IntegerNumbers}
   \EulerE^{\pi\ImaginaryUnit n^2\tau+2\pi\ImaginaryUnit nz}
  =\sum_{n\in\IntegerNumbers}q^{n^2}\zeta^n,\\
 \vartheta_4(z\mid\tau)
 &=\sum_{n\in\IntegerNumbers}(-1)^n
   \EulerE^{\pi\ImaginaryUnit n^2\tau+2\pi\ImaginaryUnit nz},\\
 \vartheta_2(z\mid\tau)
 &=\sum_{n\in\IntegerNumbers}
   \EulerE^{\pi\ImaginaryUnit(n+1/2)^2\tau
   +2\pi\ImaginaryUnit(n+1/2)z}.
\end{align*}
The series and all termwise derivatives converge normally on compact subsets of
\(\ComplexNumbers\times\mathbb H\).}
\entryfield{Product forms}{With the prescribed \(q,\zeta\),
\begin{align*}
 \vartheta_3(z\mid\tau)&=(q^2,-q\zeta,-q/\zeta;q^2)_\infty,\\
 \vartheta_4(z\mid\tau)&=(q^2,q\zeta,q/\zeta;q^2)_\infty,\\
 \vartheta_2(z\mid\tau)&=
 q^{1/4}\zeta^{1/2}(q^2,-q^2\zeta,-\zeta^{-1};q^2)_\infty.
\end{align*}
These formulas fix the nome, the argument exponential, and the quarter- and half-power
lifts simultaneously.}
\entryfield{Heat equation}{For \(j\in\{2,3,4\}\),
\[
 \frac{\partial\vartheta_j}{\partial\tau}
 =\frac{1}{4\pi\ImaginaryUnit}
  \frac{\partial^2\vartheta_j}{\partial z^2}.
\]
The factor \(4\pi\ImaginaryUnit\) belongs to the convention
\(\exp(2\pi\ImaginaryUnit nz)\); changing the argument scale changes the equation.}
\entryfield{Square-root branch}{Because \(-\ImaginaryUnit\tau\) lies in the right
half-plane, \((-\ImaginaryUnit\tau)^{1/2}\) is its principal holomorphic square root,
normalized by
\((-\ImaginaryUnit\cdot\ImaginaryUnit t)^{1/2}=t^{1/2}\) for \(t>0\).}
\entryfield{Modular inversion}{With that square root,
\[
 \vartheta_3\!\left(\frac z\tau\middle|-\frac1\tau\right)
 =(-\ImaginaryUnit\tau)^{1/2}
  \EulerE^{\pi\ImaginaryUnit z^2/\tau}\vartheta_3(z\mid\tau).
\]
At \(z=0\), modular inversion sends
\((\vartheta_2,\vartheta_3,\vartheta_4)\) to
\((\vartheta_4,\vartheta_3,\vartheta_2)\), multiplied componentwise by the same
\((-\ImaginaryUnit\tau)^{1/2}\).}
\entryfield{Quadratic-series crosswalk}{With
\(\Theta_{\mathrm{quad}}(x;Q)\) defined below,
\begin{align*}
 \vartheta_3(z\mid\tau)&=\Theta_{\mathrm{quad}}(\zeta;q)
  =\Theta_{\mathrm{bil}}(q\zeta;q^2),\\
 \vartheta_4(z\mid\tau)&=\Theta_{\mathrm{quad}}(-\zeta;q)
  =\vartheta_{\mathrm{mult}}(q\zeta;q^2),\\
 \vartheta_2(z\mid\tau)&=
 q^{1/4}\zeta^{1/2}\Theta_{\mathrm{bil}}(q^2\zeta;q^2).
\end{align*}}
\end{catalogueentry}

\begin{catalogueentry}{dedekind-eta}{Dedekind eta function}
\entryfield{Definition}{For \(\tau\in\mathbb H\) and
\(q=\EulerE^{\pi\ImaginaryUnit\tau}\),
\[
 \eta(\tau)
 \DefinitionEquals
 \EulerE^{\pi\ImaginaryUnit\tau/12}
 \prod_{n=1}^{\infty}(1-\EulerE^{2\pi\ImaginaryUnit n\tau})
 =q^{1/12}(q^2;q^2)_\infty.
\]
The twelfth power is the prescribed exponential
\(\EulerE^{\pi\ImaginaryUnit\tau/12}\), not a free root choice.}
\entryfield{Nonvanishing}{The product converges and none of its factors vanishes on
\(\mathbb H\); hence \(\eta\) is holomorphic and nonzero there.}
\entryfield{Theta product}{At zero argument,
\[
 \vartheta_2(0\mid\tau)\vartheta_3(0\mid\tau)
 \vartheta_4(0\mid\tau)=2\eta(\tau)^3.
\]}
\entryfield{Modular transformations}{With the same principal square root as for the
Jacobi functions,
\[
 \eta(\tau+1)=\EulerE^{\pi\ImaginaryUnit/12}\eta(\tau),
 \qquad
 \eta(-1/\tau)=(-\ImaginaryUnit\tau)^{1/2}\eta(\tau).
\]
The multiplier in the translation formula is part of the normalization.}
\end{catalogueentry}

\begin{catalogueentry}{ramanujan-quadratic-theta}{Ramanujan's general theta function and
bilateral quadratic theta}
\entryfield{Ramanujan theta}{For \(a,b\in\ComplexNumbers\) with \(|ab|<1\),
\[
 \Theta_{\mathrm{Ram}}(a,b)
 \DefinitionEquals
 \sum_{n\in\IntegerNumbers}
 a^{n(n+1)/2}b^{n(n-1)/2}.
\]
The exponents are integers.  Zeroth powers equal \(1\), including at a zero base.  The
paired, division-free form is
\[
 \Theta_{\mathrm{Ram}}(a,b)
 =1+\sum_{n=1}^{\infty}(a^n+b^n)(ab)^{n(n-1)/2}.
\]}
\entryfield{Product and degenerate base}{Ramanujan's product is
\[
 \Theta_{\mathrm{Ram}}(a,b)=(-a,-b,ab;ab)_\infty.
\]
At base zero, \((c;0)_\infty\DefinitionEquals1-c\), so if \(ab=0\), both sides equal
\(1+a+b\).  The monograph-local \commandname{RamanujanTheta} receives \(a,b\); it never
hides the product base \(ab\).}
\entryfield{Bilateral crosswalk}{When \(a,b\ne0\),
\[
 \Theta_{\mathrm{Ram}}(a,b)=\Theta_{\mathrm{bil}}(a;ab).
\]
The paired definition supplies the continuous algebraic interpretation when one of the
two parameters is zero.}
\entryfield{Quadratic theta}{For \(0<|Q|<1\) and
\(x\in\ComplexNumbers^\times\),
\[
 \Theta_{\mathrm{quad}}(x;Q)
 \DefinitionEquals\sum_{n\in\IntegerNumbers}Q^{n^2}x^n
 =\Theta_{\mathrm{bil}}(Qx;Q^2).
\]
The local \commandname{BilateralQuadraticTheta} receives \(x,Q\).  Upper-case \(Q\) here
is a locally typed theta base, not automatically a finite-field cardinality.}
\entryfield{Schröter lattice dissection}{For \(a,b\in\PositiveIntegers\),
\(0<|q|<1\), and \(x,y\in\ComplexNumbers^\times\),
\begin{align*}
 \Theta_{\mathrm{quad}}(x;q^a)\Theta_{\mathrm{quad}}(y;q^b)
 &=\sum_{k=0}^{a+b-1}y^kq^{bk^2}
 \Theta_{\mathrm{quad}}(xyq^{2bk};q^{a+b})\\
 &\qquad\qquad{}\times
 \Theta_{\mathrm{quad}}(y^ax^{-b}q^{2abk};q^{ab(a+b)}).
\end{align*}
The finite range has exactly \(a+b\) residue classes.  The identity comes from the
bijection \(n=r-bs\), \(m=k+r+as\) on the class
\(m-n\equiv k\pmod{a+b}\); absolute convergence justifies rearrangement.}
\end{catalogueentry}

\begin{catalogueentry}{sums-of-squares}{Theta coefficients, sums of squares, and
Gaussian integers}
\entryfield{Representation count}{For \(d,n\in\NaturalNumbers\),
\[
 r_d(n)
 \DefinitionEquals
 \CardinalityOf{\{(x_1,\ldots,x_d)\in\IntegerNumbers^d:
 x_1^2+\cdots+x_d^2=n\}}.
\]
For \(d=0\), the unique empty tuple gives \(r_0(0)=1\) and \(r_0(n)=0\) for \(n>0\).}
\entryfield{Universal generating function}{Coefficientwise in
\(\IntegerNumbers[[q]]\), and analytically for \(|q|<1\),
\[
 \left(\sum_{m\in\IntegerNumbers}q^{m^2}\right)^d
 =\sum_{n=0}^{\infty}r_d(n)q^n.
\]
A divisor-sum formula for a particular \(d\) requires additional arithmetic or modular
input; it does not follow merely from this Cauchy product.}
\entryfield{Character modulo four}{For \(m\in\IntegerNumbers\),
\[
 \chi_4(m)=
 \begin{cases}
 0,&2\mid m,\\
 1,&m\equiv1\pmod4,\\
 -1,&m\equiv3\pmod4.
 \end{cases}
\]
It is completely multiplicative and periodic modulo \(4\).  The subscript distinguishes
it from a generic characteristic function.}
\entryfield{Two- and four-square formulas}{For \(n\in\PositiveIntegers\),
\[
 r_2(n)=4\sum_{d\mid n}\chi_4(d),\qquad
 r_4(n)=8\sum_{\substack{d\mid n\\4\nmid d}}d.
\]
Separately, \(r_2(0)=r_4(0)=1\); the positive-divisor sums above are not evaluated at
zero.}
\entryfield{Gaussian integers}{The ring
\(\IntegerNumbers[\ImaginaryUnit]=\{x+\ImaginaryUnit y:x,y\in\IntegerNumbers\}\) has
norm
\[
 N(x+\ImaginaryUnit y)=x^2+y^2,
\qquad N(\alpha\beta)=N(\alpha)N(\beta).
\]
It is Euclidean for \(N\), hence a unique-factorization domain, and its units are
\(\{\pm1,\pm\ImaginaryUnit\}\).}
\entryfield{Rational-prime splitting}{One has
\(2=-\ImaginaryUnit(1+\ImaginaryUnit)^2\).  A rational prime \(p\equiv3\pmod4\) remains
irreducible, while a prime \(p\equiv1\pmod4\) factors as
\(p=\varpi\overline\varpi\), where the Gaussian prime \(\varpi\) is unique up to a unit
and conjugation.  The letter \(\varpi\), not \(\pi\), prevents collision with the circle
constant in theta and Fourier formulas.}
\end{catalogueentry}

\begin{catalogueentry}{cyclotomic}{Cyclotomic polynomials and roots of unity}
\entryfield{Primitive root}{$\zeta_N=\EulerE^{2\pi\ImaginaryUnit/N}$ for
$N\in\PositiveIntegers$ is the chosen positively oriented primitive $N$th root.}
\entryfield{Cyclotomic polynomial}{$\Phi_N(x)\in\IntegerNumbers[x]$ is the monic
polynomial whose roots are $\zeta_N^k$ for
$1\le k\le N$ and $\gcd(k,N)=1$.}
\entryfield{Factorization}{$x^N-1=\prod_{d\mid N}\Phi_d(x)$, where $d$ ranges over the
positive divisors of $N$.}
\entryfield{Order of a root}{For $\zeta^N=1$,
$\OrderOperator(\zeta)=\min\{d\in\PositiveIntegers:\zeta^d=1\}$.
This group-theoretic order is distinct from the analytic
$\OrderOperator_{z_0}(f)$ of a zero.}
\entryfield{Evaluation policy}{A rational $q$-series expression is not evaluated at a
root of unity until all denominator zeros and possible cancellations have been audited.}
\entryfield{Cyclotomic valuation}{For \(d\in\PositiveIntegers\) and a nonzero polynomial
\(F\in\IntegerNumbers[q]\),
\[
 \nu^{\mathrm{cyc}}_{d,q}(F)
 \DefinitionEquals
 \max\{e\in\NaturalNumbers:\Phi_d(q)^e\text{ divides }F\}.
\]
It is the exponent of the polynomial \(\Phi_d(q)\), not a \(d\)-adic valuation of an
integer.  It is left undefined at \(F=0\), rather than being silently assigned
\(+\infty\).}
\entryfield{Semantic command}{The monograph-local
\commandname{CyclotomicValuation}\texttt{\{d\}\{q\}\{F\}} prints the cyclotomic index,
polynomial variable, and polynomial.  The variable input is necessary when several
polynomial rings are present.}
\entryfield{Shifted-factorial factorization}{For \(n\in\NaturalNumbers\),
\[
 (q;q)_n=(-1)^n\prod_{d=1}^{n}\Phi_d(q)^{\Floor{n/d}}.
\]
Consequently cyclotomic divisibility of ratios of shifted factorials is checked by
subtracting these explicit nonnegative exponents after polynomiality has been
established.}
\end{catalogueentry}

\begin{catalogueentry}{q-adic-formal}{Coefficientwise convergence, \(q\)-adic order, and
formal congruence}
\entryfield{Coefficientwise convergence}{For formal series
\(F_N(q)=\sum_{j\ge j_0}a_{N,j}q^j\) and
\(F(q)=\sum_{j\ge j_0}a_jq^j\), the notation
\(F_N\to F\) coefficientwise means: for every fixed \(j\), there exists \(N_j\) such that
\(a_{N,j}=a_j\) for all \(N\ge N_j\).  No norm, pointwise evaluation, or uniform
analytic convergence follows from this definition.}
\entryfield{\(q\)-adic order}{For a nonzero formal Laurent series, the \(q\)-adic order is
the already defined \(\OrderOperator_{q=0}(F)\), its least exponent with nonzero
coefficient.  A sequence tends \(q\)-adically to zero when these orders tend to
\(+\infty\).}
\entryfield{Formal congruence}{For \(M\in\IntegerNumbers\) and
\(N\in\NaturalNumbers\),
\[
 F(q)\equiv G(q)\pmod{q^N}
\]
means \(F-G\in q^NR[[q]]\), equivalently equality of the coefficients of
\(q^0,\ldots,q^{N-1}\).  A coefficientwise congruence modulo \(M\) instead means
\(a_j\equiv b_j\pmod M\) for every \(j\); the two moduli must never be conflated.}
\entryfield{\(I\)-adic generalization}{If a commutative ring \(A\) is complete and
separated for an ideal \(I\) and \(q\in I\), infinite products may converge
coefficientwise \(I\)-adically in \(A[[z]]\).  This is an algebraic topology.  It does
not assert convergence after embedding \(A\) into \(\ComplexNumbers\).}
\end{catalogueentry}

\begin{catalogueentry}{bailey-pairs}{Bailey pairs, transforms, lowering, and chains}
\entryfield{Bailey pair}{Fix an explicit base \(q\).  Two sequences
\((\alpha_n)_{n\ge0}\) and \((\beta_n)_{n\ge0}\) form a Bailey pair relative to \(a\)
when
\[
 \beta_n=\sum_{r=0}^{n}
 \frac{\alpha_r}{(q;q)_{n-r}(aq;q)_{n+r}}
 \qquad(n\in\NaturalNumbers),
\]
with all denominators nonzero.  The base is part of the governing declaration even
though traditional terminology names only the relative parameter \(a\).}
\entryfield{Full transform}{Let \(\rho_1\rho_2\ne0\), put
\(A=aq/(\rho_1\rho_2)\), and assume
\((aq;q)_m\), \((aq/\rho_1;q)_m\), and \((aq/\rho_2;q)_m\) are nonzero for
\(m\ge1\).  Define
\begin{align*}
 \alpha_n^{\mathrm{BT}}
 &=\frac{(\rho_1,\rho_2;q)_nA^n}
         {(aq/\rho_1,aq/\rho_2;q)_n}\alpha_n,\\
 \beta_n^{\mathrm{BT}}
 &=\sum_{j=0}^{n}
 \frac{(\rho_1,\rho_2;q)_j(A;q)_{n-j}A^j}
 {(q;q)_{n-j}(aq/\rho_1,aq/\rho_2;q)_n}\beta_j.
\end{align*}
Then the transformed pair is again relative to \(a\).  The superscript
\(\mathrm{BT}\) is a semantic label, not exponentiation.}
\entryfield{Parameter lowering}{Assume \((a;q)_m\ne0\) for \(m\ge1\).  Put
\[
 \alpha_0^{\mathrm{BL}}=\alpha_0,\qquad
 \alpha_n^{\mathrm{BL}}
 =(1-a)\left(
 \frac{\alpha_n}{1-aq^{2n}}
 -\frac{aq^{2n-2}\alpha_{n-1}}{1-aq^{2n-2}}\right)
 \quad(n\ge1),
\]
and \(\beta_n^{\mathrm{BL}}=\beta_n\).  With the displayed denominators nonzero, this
is a Bailey pair relative to \(a/q\).  The label \(\mathrm{BL}\) means Bailey lowering.}
\entryfield{Iterated chain}{For \(t\in\NaturalNumbers\), the depth-zero pair is the
original pair.  For \(t\ge1\),
\[
 \alpha_n^{\mathrm{BC}(t)}=a^{tn}q^{tn^2}\alpha_n,
\]
\[
 \beta_n^{\mathrm{BC}(t)}
 =\sum_{n\ge r_1\ge\cdots\ge r_t\ge0}
 \frac{a^{r_1+\cdots+r_t}q^{r_1^2+\cdots+r_t^2}\beta_{r_t}}
 {(q;q)_{n-r_1}\prod_{\ell=1}^{t-1}(q;q)_{r_\ell-r_{\ell+1}}}.
\]
The finite weakly decreasing index simplex, the chain depth \(t\), and the relative
parameter are never inferred from \(\alpha\) or \(\beta\) alone.}
\entryfield{Local semantic commands}{The monograph uses superscripts
\(\mathrm{BT}\), \(\mathrm{BL}\), and \(\mathrm{BC}(t)\) through local commands so a
transformed sequence cannot be confused with the input Bailey pair or with unrelated
sequences called \(\alpha,\beta\).}
\end{catalogueentry}

\begin{catalogueentry}{rogers-ramanujan}{Rogers--Ramanujan series, continuants, and the
Andrews--Gordon family}
\entryfield{Single Rogers family}{For \(m\in\NaturalNumbers\) and \(|q|<1\),
\[
 F_m^{\mathrm{RR}}(q)
 \DefinitionEquals\sum_{n=0}^{\infty}
 \frac{q^{n^2+mn}}{(q;q)_n}.
\]
Thus \(F_0^{\mathrm{RR}}=G\) and \(F_1^{\mathrm{RR}}=H\) in the traditional notation.
The semantic superscript prevents a collision with Fabius functions and generic
generating functions.}
\entryfield{Recurrence and products}{Absolute convergence permits
\[
 F_m^{\mathrm{RR}}(q)
 =F_{m+1}^{\mathrm{RR}}(q)+q^{m+1}F_{m+2}^{\mathrm{RR}}(q).
\]
The two Rogers--Ramanujan identities are
\[
 F_0^{\mathrm{RR}}(q)=\frac1{(q,q^4;q^5)_\infty},
 \qquad
 F_1^{\mathrm{RR}}(q)=\frac1{(q^2,q^3;q^5)_\infty}.
\]}
\entryfield{Finite continuant}{For \(N\in\NaturalNumbers\),
\[
 C_N^{\mathrm{cont}}
 =1+\cfrac{q}{1+\cfrac{q^2}{\ddots+\cfrac{q^N}{1}}}
 =\frac{P_N^{\mathrm{cont}}}{Q_N^{\mathrm{cont}}},
\]
where
\[
 P_{-1}^{\mathrm{cont}}=P_0^{\mathrm{cont}}=1,\qquad
 Q_{-1}^{\mathrm{cont}}=0,\quad Q_0^{\mathrm{cont}}=1,
\]
\[
 P_N^{\mathrm{cont}}=P_{N-1}^{\mathrm{cont}}+q^NP_{N-2}^{\mathrm{cont}},
 \qquad
 Q_N^{\mathrm{cont}}=Q_{N-1}^{\mathrm{cont}}+q^NQ_{N-2}^{\mathrm{cont}}.
\]
The rational continuant defines the truncation even where a nested intermediate
denominator vanishes but the final rational function has a removable singularity.}
\entryfield{Gaussian finitizations}{The exact polynomials are
\[
 P_N^{\mathrm{cont}}
 =\sum_{j=0}^{\Floor{(N+1)/2}}
 q^{j^2}\GaussianBinomial{N+1-j}{j}{q},
\]
\[
 Q_N^{\mathrm{cont}}
 =\sum_{j=0}^{\Floor{N/2}}
 q^{j^2+j}\GaussianBinomial{N-j}{j}{q}.
\]
Locally uniformly for \(|q|<1\), they tend respectively to
\(F_0^{\mathrm{RR}}(q)\) and \(F_1^{\mathrm{RR}}(q)\).}
\entryfield{Continued fraction and fifth-root branch}{The single-valued ratio satisfies
\[
 \frac{F_1^{\mathrm{RR}}(q)}{F_0^{\mathrm{RR}}(q)}
 =\cfrac1{1+\cfrac q{1+\cfrac{q^2}{1+\cfrac{q^3}{1+\ddots}}}}.
\]
On a simply connected subdomain of \(0<|q|<1\), choose one branch of \(q^{1/5}\) and
define
\[
 R_{\mathrm{RR}}(q)
 \DefinitionEquals q^{1/5}
 \frac{F_1^{\mathrm{RR}}(q)}{F_0^{\mathrm{RR}}(q)}
 =q^{1/5}\frac{(q,q^4;q^5)_\infty}{(q^2,q^3;q^5)_\infty}.
\]
Only the prefactor is branch-dependent; the ratio and continued fraction are
single-valued.}
\entryfield{Andrews--Gordon indexing}{Let \(k\ge2\), \(1\le i\le k\), and
\(n_1,\ldots,n_{k-1}\in\NaturalNumbers\).  Put
\(N_j=n_j+\cdots+n_{k-1}\).  Then
\[
 \sum_{n_1,\ldots,n_{k-1}\ge0}
 \frac{q^{N_1^2+\cdots+N_{k-1}^2+N_i+\cdots+N_{k-1}}}
 {(q;q)_{n_1}\cdots(q;q)_{n_{k-1}}}
 =
 \frac{(q^i,q^{2k+1-i},q^{2k+1};q^{2k+1})_\infty}
 {(q;q)_\infty}.
\]
When \(i=k\), the linear tail \(N_i+\cdots+N_{k-1}\) is the empty sum \(0\).
The modulus is \(2k+1\); an even-modulus identity is not obtained by a nonintegral
substitution for \(k\).}
\entryfield{Recorded-equality marker}{The monograph-local relation
\(\mathrel{\stackrel{\mathrm{record}}{=}}\) marks an identity recorded with provenance
but not proved at that location.  It has the same asserted mathematical equality only
when the accompanying status and hypotheses support it; the marker itself is never a
proof.}
\end{catalogueentry}

\begin{catalogueentry}{borwein-dissection}{Borwein product dissections, reciprocity, and
two-ended stabilization}
\entryfield{Product}{For \(n\in\PositiveIntegers\),
\[
 P_n^{\mathrm{Bor}}(q)
 \DefinitionEquals\frac{(q;q)_{3n}}{(q^3;q^3)_n}
 =\prod_{j=1}^{n}(1-q^{3j-2})(1-q^{3j-1}).
\]
Its degree is \(3n^2\).  The superscript \(\mathrm{Bor}\) distinguishes this product
from generic polynomials and continuant numerators.}
\entryfield{Residue-three dissection}{There are unique
\(A_n^{\mathrm{Bor}},B_n^{\mathrm{Bor}},C_n^{\mathrm{Bor}}\in\IntegerNumbers[x]\)
such that
\[
 P_n^{\mathrm{Bor}}(q)
 =A_n^{\mathrm{Bor}}(q^3)
 -qB_n^{\mathrm{Bor}}(q^3)
 -q^2C_n^{\mathrm{Bor}}(q^3).
\]
The two minus signs are part of the normalization and do not assert a sign for any
coefficient.}
\entryfield{Reciprocity}{For \(n\ge1\),
\begin{align*}
 A_n^{\mathrm{Bor}}(x)&=x^{n^2}A_n^{\mathrm{Bor}}(x^{-1}),\\
 B_n^{\mathrm{Bor}}(x)&=x^{n^2-1}C_n^{\mathrm{Bor}}(x^{-1}),\\
 C_n^{\mathrm{Bor}}(x)&=x^{n^2-1}B_n^{\mathrm{Bor}}(x^{-1}).
\end{align*}
These are Laurent-polynomial identities.  In particular,
\(B_n^{\mathrm{Bor}}(1)=C_n^{\mathrm{Bor}}(1)\) and
\(A_n^{\mathrm{Bor}}(1)=2B_n^{\mathrm{Bor}}(1)\).}
\entryfield{Coefficient names}{For \(r\in\NaturalNumbers\),
\[
 a_{n,r}^{\mathrm{Bor}}=[x^r]A_n^{\mathrm{Bor}}(x),\quad
 b_{n,r}^{\mathrm{Bor}}=[x^r]B_n^{\mathrm{Bor}}(x),\quad
 c_{n,r}^{\mathrm{Bor}}=[x^r]C_n^{\mathrm{Bor}}(x).
\]
Outside the polynomial degree, coefficient extraction gives \(0\).}
\entryfield{Exact stabilization}{For fixed \(d\in\NaturalNumbers\) and \(n\ge d+1\),
\[
 a_{n,d}^{\mathrm{Bor}}=a_{d+1,d}^{\mathrm{Bor}},\quad
 b_{n,d}^{\mathrm{Bor}}=b_{d+1,d}^{\mathrm{Bor}},\quad
 c_{n,d}^{\mathrm{Bor}}=c_{d+1,d}^{\mathrm{Bor}},
\]
and reciprocity gives
\[
 a_{n,n^2-d}^{\mathrm{Bor}}=a_{d+1,d}^{\mathrm{Bor}},\quad
 b_{n,n^2-1-d}^{\mathrm{Bor}}=c_{d+1,d}^{\mathrm{Bor}},\quad
 c_{n,n^2-1-d}^{\mathrm{Bor}}=b_{d+1,d}^{\mathrm{Bor}}.
\]
These identities state eventual coefficient constancy at both ends; they make no
coefficientwise positivity assertion.}
\entryfield{Semantic commands}{The local \commandname{BorweinProduct} names the product.
\commandname{BorweinComponentA}, \commandname{BorweinComponentB}, and
\commandname{BorweinComponentC} name the three residue components.  Their corresponding
coefficient commands display the family, level, and coefficient index.}
\end{catalogueentry}

\begin{catalogueentry}{arithmetic-functions}{Divisibility, gcd, lcm, and standard
arithmetic functions}
\entryfield{Divisibility}{$a\mid b$ for $a,b\in\IntegerNumbers$ means
$b=ac$ for some $c\in\IntegerNumbers$.  In a sum $d\mid n$, the divisor $d$ is positive
unless otherwise stated.}
\entryfield{GCD and LCM}{$\gcd(a,b)\in\NaturalNumbers$ is the nonnegative greatest common
divisor.  The zero-argument command \commandname{LeastCommonMultiple} supplies the
operator token in $\LeastCommonMultiple(a_1,\ldots,a_r)$; its parenthesized integer inputs
are source arguments, not TeX macro arguments.  Its value is the unique nonnegative
generator of the ideal
$a_1\IntegerNumbers\cap\cdots\cap a_r\IntegerNumbers$ and is therefore insensitive to
signs.}
\entryfield{Zero and family conventions}{For $r=0$ the empty intersection is
$\IntegerNumbers$.  Explicitly,
$\LeastCommonMultiple(a,b)=\LeastCommonMultiple(\AbsoluteValue a,\AbsoluteValue b)$,
$\gcd(0,0)=0$, $\LeastCommonMultiple(0,b)=0$ for every $b$, a~singleton family has
LCM $\AbsoluteValue a$, and the empty family has LCM $1$.  The empty-family value is the
unit for adjoining further inputs and is required in the dyadic-denominator formulas.}
\entryfield{Möbius function}{$\mu_{\mathrm{M}}(n)=0$ if a prime square divides $n$, and
$(-1)^r$ if $n$ is a product of $r$ distinct primes; its domain is
$\PositiveIntegers$ and $\mu_{\mathrm{M}}(1)=1$.  The subscript prevents a collision
with measures and means.}
\entryfield{Euler totient}{$\varphi_{\mathrm E}(n)
=|\{1\le k\le n:\gcd(k,n)=1\}|$, $n\in\PositiveIntegers$.
The subscript separates it from characteristic functions.}
\entryfield{Divisor sums}{$\sigma_\alpha(n)=\sum_{d\mid n}d^\alpha$ for
$n\in\PositiveIntegers$ and $\alpha\in\ComplexNumbers$, with positive divisors and
$d^\alpha=\exp(\alpha\log d)$.}
\end{catalogueentry}

\section{Asymptotics, logarithms, Lambert W, and saddle analysis}

\begin{catalogueentry}{asymptotic-regime}{Limits and named asymptotic regimes}
\entryfield{Limit annotation}{A limit variable and destination are part of every
asymptotic assertion: $n\to\infty$, $x\to0^+$, $z\to z_0$ in a sector, or
$q\uparrow1$ are different regimes.  A subscript may abbreviate a regime only after it
has been declared.}
\entryfield{One-sided limits}{$x\to a^+$ means $x>a$ and $x\to a$;
$x\to a^-$ means $x<a$.  For $q\uparrow1$, $q$ is real and approaches $1$ from below.}
\entryfield{Uniformity}{An estimate uniform for $y\in Y$ means its threshold and implied
constant may be chosen independently of $y$.  Pointwise asymptotics never imply this.}
\entryfield{Complex sectors}{A sector is printed as
$S(\theta_0,\alpha,r)=\{z:0<|z|<r,\ |\arg z-\theta_0|<\alpha\}$, with a selected
continuous argument.  Statements at a branch point include the slit and sector.}
\entryfield{Parameter order}{For iterated limits, notation such as
$n\to\infty$ followed by $q\uparrow1$ records the order.  A joint limit instead states a
relation such as $(1-q)n\to\tau$.}
\end{catalogueentry}

\begin{catalogueentry}{big-o}{Big \(O\), little \(o\), and comparison relations}
\entryfield{Big \(O\)}{$f=\BigOAt g{x\to a}$ means that there are $C>0$ and a punctured
neighborhood in the declared domain such that $|f(x)|\le C|g(x)|$.  If $g$ vanishes
arbitrarily near $a$, the inequality definition remains authoritative.}
\entryfield{Little \(o\)}{$f=\LittleOAt g{x\to a}$ means
$f(x)/g(x)\to0$ where the quotient is defined, or equivalently the epsilon inequality
when that avoids isolated zeros.}
\entryfield{Asymptotic equivalence}{$f(x)\sim g(x)$ as $x\to a$ means
$f(x)/g(x)\to1$.  It requires $g$ to be nonzero in a punctured neighborhood.}
\entryfield{Two-sided comparability}{$f\asymp g$ in a regime means
$c|g|\le|f|\le C|g|$ for constants $0<c\le C<\infty$.  This is not asymptotic
equivalence.}
\entryfield{Vinogradov notation}{$f\ll g$ means $f=\BigO(g)$ and $f\gg g$ means
$g\ll f$.  A dependency subscript, for example $\ll_\alpha$, lists every parameter on
which the implied constant may depend.}
\entryfield{Shared commands}{Normative prose prefers
\commandname{BigOAt} and \commandname{LittleOAt}; bare
\commandname{BigO} and \commandname{LittleO} are reserved for formulas whose regime is
printed immediately outside the expression.}
\end{catalogueentry}

\begin{catalogueentry}{asymptotic-expansion}{Asymptotic sequences and expansions}
\entryfield{Asymptotic sequence}{Functions $(\phi_n)$ form an asymptotic sequence in a
regime when $\phi_{n+1}=o(\phi_n)$ for every $n$.}
\entryfield{Expansion}{$f\sim\sum_{n\ge0}a_n\phi_n$ means that for every
$N\in\PositiveIntegers$,
\[
 f-\sum_{n=0}^{N-1}a_n\phi_n=o(\phi_{N-1})
\]
in the stated regime.  It does not assert convergence of the infinite series.}
\entryfield{Poincaré scale}{For powers at infinity,
$f(x)\sim\sum_{n\ge0}a_nx^{-\lambda_n}$ requires
$\lambda_0<\lambda_1<\cdots$ and $x$ in a domain where the powers have been defined.}
\entryfield{Exponentially small terms}{A transseries term such as
$\EulerE^{-A/\epsilon}\epsilon^\beta(\log\epsilon)^m$ lists the sign of
$\operatorname{Re}(A/\epsilon)$, the logarithm branch, and the ordering used to compare
exponential scales.}
\entryfield{Truncation}{An optimally truncated expansion names its truncation index
$N(\epsilon)$ and the remainder norm; the infinite formal series is not assigned its
optimal-truncation error.}
\end{catalogueentry}

\begin{catalogueentry}{principal-log}{Real logarithm, principal logarithm, and arguments}
\entryfield{Real logarithm}{$\log x$ for $x>0$ is the real natural logarithm with base
$\EulerE$.}
\entryfield{Principal branch}{The shared command denotes
\[
 \PrincipalLogarithm z=\log|z|+\ImaginaryUnit\operatorname{Arg}z,\qquad
 \operatorname{Arg}z\in(-\pi,\pi],
\]
for $z\in\ComplexNumbers\setminus(-\infty,0]$, with the boundary convention
$\operatorname{Arg}(-r)=\pi$ when a boundary value is explicitly taken from above.
The analytic domain itself excludes the nonpositive real axis.}
\entryfield{Other branches}{On a simply connected zero-free domain $D$, a branch
$\log_Dz$ is an analytic function satisfying $\exp(\log_Dz)=z$.  Its normalization at
one base point is stated.}
\entryfield{Complex powers}{$z^\alpha\DefinitionEquals
\exp(\alpha\PrincipalLogarithm z)$ when the principal branch is intended.  If
$\alpha\in\IntegerNumbers$, the algebraic power has its global meaning and no branch is
needed.}
\entryfield{Boundary notation}{$\PrincipalLogarithm(x\pm\ImaginaryUnit0)$ denotes a
specified limiting value.  It is not silently substituted into a formula proved only
inside the slit plane.}
\end{catalogueentry}

\begin{catalogueentry}{lambert-w}{Lambert \(W\) and its branches}
\entryfield{Defining equation}{A branch of $\LambertW$ satisfies
\[
 \LambertW_k(z)\EulerE^{\LambertW_k(z)}=z,
\]
where $k\in\IntegerNumbers$ labels the branch.  The inverse is multivalued before a
branch is selected.}
\entryfield{Principal branch}{$\LambertW_0$ is real on $[-\EulerE^{-1},\infty)$,
analytic at $0$, and satisfies $\LambertW_0(0)=0$.  Its standard complex cut is
$(-\infty,-\EulerE^{-1}]$.}
\entryfield{Lower real branch}{$\LambertW_{-1}$ is real on
$[-\EulerE^{-1},0)$, takes values in $(-\infty,-1]$, and approaches $-\infty$ as
$z\uparrow0$.}
\entryfield{Branch point}{At $z=-\EulerE^{-1}$ the two real branches meet at $-1$ and
have a square-root branch point.  A Taylor series in $z+\EulerE^{-1}$ is therefore not
asserted there.}
\entryfield{Asymptotic branch}{For large $z$ away from cuts,
\[
 \LambertW_k(z)=L_1-L_2+\frac{L_2}{L_1}+\cdots,\quad
 L_1=\log z+2\pi\ImaginaryUnit k,\quad L_2=\log L_1,
\]
with both logarithm branches and the sector specified.}
\entryfield{Bare symbol policy}{A branchless $\LambertW(z)$ is permitted only when the
domain makes the real principal branch unambiguous; all complex and negative-real uses
carry a subscript.}
\end{catalogueentry}

\begin{catalogueentry}{polylog-gamma-zeta}{Polylogarithm, gamma, and zeta functions}
\entryfield{Polylogarithm}{For $|z|<1$,
\[
 \Polylogarithm_s(z)\DefinitionEquals\sum_{n=1}^{\infty}\frac{z^n}{n^s},
\qquad s\in\ComplexNumbers.
\]
Elsewhere the symbol denotes its analytically continued branch, with the $z$-plane cut
and logarithm convention stated.}
\entryfield{Gamma function}{For $\operatorname{Re}s>0$,
$\GammaFunction(s)=\int_0^\infty x^{s-1}\EulerE^{-x}\,\Differential x$.
Its meromorphic continuation has simple poles at
$0,-1,-2,\ldots$ and no zeros.  For $n\in\NaturalNumbers$,
$\GammaFunction(n+1)=n!$.}
\entryfield{Riemann zeta}{For $\operatorname{Re}s>1$,
$\RiemannZeta(s)=\sum_{n=1}^\infty n^{-s}$, extended meromorphically to
$\ComplexNumbers$ with a simple pole at $s=1$.}
\entryfield{Hurwitz zeta}{$\HurwitzZeta{s}{a}
=\sum_{n=0}^{\infty}(n+a)^{-s}$ initially for
$\operatorname{Re}s>1$ and $a\notin\{0,-1,-2,\ldots\}$, with the branch of each power
specified if $a$ is complex.  The two arguments of
\commandname{HurwitzZeta} are the order and shift, respectively.}
\entryfield{Collision policy}{The same Greek glyph may label a root of unity or a zeta
function only with a semantic subscript or arguments that make its type immediate.}
\end{catalogueentry}

\begin{catalogueentry}{coefficient-extraction}{Coefficient extraction and formal series}
\entryfield{Coefficient operator}{For
$A(z)=\sum_{n\in\IntegerNumbers}a_nz^n$,
\[
 [z^n]A(z)\DefinitionEquals a_n.
\]
The bracket is an operator, not multiplication by $z^n$.}
\entryfield{Cauchy formula}{If $A$ is analytic on and inside a positively oriented circle
$|z|=r$,
\[
 [z^n]A(z)=\frac{1}{2\pi\ImaginaryUnit}
 \int_{|z|=r}\frac{A(z)}{z^{n+1}}\,\Differential z,
 \qquad n\in\NaturalNumbers.
\]}
\entryfield{Laurent coefficients}{For an annular Laurent series the same operator permits
$n\in\IntegerNumbers$; the annulus and contour winding number are stated.}
\entryfield{Formal interpretation}{In $R[[z]]$, coefficient extraction is algebraic and
does not assert analytic convergence.  Cauchy's formula belongs only to the analytic
interpretation.}
\end{catalogueentry}

\begin{catalogueentry}{residue}{Residues and contour orientation}
\entryfield{Residue}{For meromorphic $f$ near $z_0$,
$\ResidueOperator_{z=z_0}f(z)$ is the coefficient of $(z-z_0)^{-1}$ in the Laurent
expansion.}
\entryfield{Simple pole}{If $f=g/h$, $g,h$ analytic,
$h(z_0)=0$, and $h'(z_0)\ne0$, then
$\ResidueOperator_{z=z_0}f=g(z_0)/h'(z_0)$.}
\entryfield{Contour}{Unless an arrow says otherwise, a simple closed contour is positively
oriented.  Then
$\int_\gamma f(z)\,\Differential z
=2\pi\ImaginaryUnit\sum\operatorname{Ind}(\gamma,z_j)
\ResidueOperator_{z=z_j}f$.}
\entryfield{Indentations and cuts}{A keyhole contour states the upper and lower boundary
values of every multivalued factor and the direction of each segment.  An omitted arc is
discarded only with a displayed bound.}
\end{catalogueentry}

\begin{catalogueentry}{saddle-point}{Saddles, phase functions, and steepest descent}
\entryfield{Phase-amplitude form}{A saddle integral is written
\[
 I(\lambda)=\int_{\Gamma}A(z,\lambda)
 \EulerE^{\lambda\Phi(z)}\,\Differential z,
\]
with $\lambda$'s limiting regime, the oriented contour $\Gamma$, and analytic domains of
$A$ and $\Phi$ stated.}
\entryfield{Saddle}{A simple saddle $z_\ast$ satisfies
$\Phi'(z_\ast)=0$ and $\Phi''(z_\ast)\ne0$.  Degenerate saddles state the first
nonzero derivative order.}
\entryfield{Descent direction}{A steepest-descent path has
$\operatorname{Im}\Phi$ constant and $\operatorname{Re}\Phi$ decreasing away from the
saddle.  The square root in the Gaussian factor is fixed by the contour orientation.}
\entryfield{Leading term}{For a single nondegenerate contributing saddle, the schematic
factor
\[
 A(z_\ast,\lambda)\EulerE^{\lambda\Phi(z_\ast)}
 \sqrt{\frac{2\pi}{-\lambda\Phi''(z_\ast)}}
\]
is meaningful only after selecting the square-root branch and verifying the contour
deformation.  It is not a universal sign-free formula.}
\entryfield{Multiple saddles}{A set $\mathcal S(\lambda)$ lists all saddles; a contributing
subset additionally depends on the original contour and Stokes region.  Dominance is
decided by $\operatorname{Re}(\lambda\Phi(z_\ast))$, not by modulus of $z_\ast$.}
\end{catalogueentry}

\begin{catalogueentry}{lambert-saddles}{Lambert-\(W\) parametrization of saddle families}
\entryfield{Typical reduction}{When a saddle equation is rearranged to
$u\EulerE^u=z(\lambda)$, its solutions are
$u_k=\LambertW_k(z(\lambda))$, $k\in\IntegerNumbers$.}
\entryfield{Back-substitution}{If $u=\alpha z+\beta$, then
$z_k=(\LambertW_k(z(\lambda))-\beta)/\alpha$ with $\alpha\ne0$.
Every scaling and sign is retained; a Lambert solution is not the original saddle until
this substitution is made.}
\entryfield{Branch enumeration}{The integer $k$ labels Lambert branches, not necessarily
the order in which saddle contributions dominate.  A separate ordering permutation is
used if needed.}
\entryfield{Coalescence}{At a branch point, two branch-labelled saddles may coalesce.
Uniform asymptotics then use an Airy or higher canonical model rather than summing two
singular Gaussian approximations.}
\entryfield{Status}{Lambert-\(W\) saddle catalogues in frontier documents are
\textbf{frontier/conjectural} unless an exact contour theorem and error bound are cited.}
\end{catalogueentry}

\begin{catalogueentry}{singularity-analysis}{Singularity analysis and transfer notation}
\entryfield{Dominant singularity}{For a power series with radius $\rho>0$, a dominant
singularity lies on $|z|=\rho$.  There may be more than one; the symbol $\rho$ alone does
not choose one.}
\entryfield{Delta domain}{A dented domain at positive $\rho$ is
\[
 \Delta(\rho,\eta,\theta)=
 \{z:|z|<\rho+\eta,\ z\ne\rho,\
 |\arg(z-\rho)|>\theta\},
\]
with $\eta>0$ and $0<\theta<\pi/2$.  Alternate angular parametrizations must be expanded.}
\entryfield{Basic transfer}{For
$\alpha\notin\{0,-1,-2,\ldots\}$,
\[
 [z^n](1-z/\rho)^{-\alpha}
 \sim\frac{n^{\alpha-1}}{\GammaFunction(\alpha)}\rho^{-n}.
\]
The power uses the branch analytic at $z=0$ in the dented domain.}
\entryfield{Periodic singularities}{Contributions from every dominant singularity are
summed with their phases.  Replacing them by the positive-real contribution can erase
oscillation or exact zeros.}
\entryfield{Remainder}{A transfer theorem states the regularity and domain hypotheses that
turn a local expansion into a coefficient error estimate; a formal local expansion alone
does not supply it.}
\end{catalogueentry}

\section{Fractional calculus, Volterra operators, and regularity classes}

\begin{catalogueentry}{fractional-integral}{Riemann--Liouville fractional integrals}
\entryfield{Left-sided operator}{For $\alpha>0$, lower terminal
$a\in\RealNumbers$, and suitable $f$,
\[
 (I_{a+}^{\alpha}f)(x)\DefinitionEquals
 \frac1{\GammaFunction(\alpha)}
 \int_a^x(x-t)^{\alpha-1}f(t)\,\Differential t,
 \qquad x>a.
\]}
\entryfield{Right-sided operator}{For upper terminal $b$,
\[
 (I_{b-}^{\alpha}f)(x)\DefinitionEquals
 \frac1{\GammaFunction(\alpha)}
 \int_x^b(t-x)^{\alpha-1}f(t)\,\Differential t,
 \qquad x<b.
\]}
\entryfield{Order zero}{$I_{a+}^{0}=I_{b-}^{0}=\IdentityOperator$ by extension, not by
substitution into the singular gamma-integral formula.}
\entryfield{Semigroup}{Under integrability hypotheses,
$I_{a+}^{\alpha}I_{a+}^{\beta}=I_{a+}^{\alpha+\beta}$ for
$\alpha,\beta>0$.  Terminals must match.}
\entryfield{Complex order}{For complex $\alpha$ with $\operatorname{Re}\alpha>0$,
the kernel power uses the real logarithm because $x-t>0$.}
\end{catalogueentry}

\begin{catalogueentry}{fractional-derivative}{Riemann--Liouville and Caputo derivatives}
\entryfield{Riemann--Liouville}{Let $\alpha>0$ and
$m=\Ceiling\alpha$.  The left derivative is
\[
 D_{a+}^{\alpha}f
 \DefinitionEquals\frac{\Differential^m}{\Differential x^m}
 I_{a+}^{m-\alpha}f.
\]}
\entryfield{Caputo}{Under sufficient regularity,
\[
 {}^{\mathrm C}D_{a+}^{\alpha}f
 \DefinitionEquals I_{a+}^{m-\alpha}f^{(m)}.
\]
These operators differ by explicit endpoint terms and are not interchangeable.}
\entryfield{Integer order}{If $\alpha=m\in\NaturalNumbers$, both reduce to the ordinary
$m$th derivative only under the appropriate endpoint/regularity interpretation.}
\entryfield{Endpoint data}{A fractional differential equation declares whether initial
conditions are ordinary derivatives $f^{(j)}(a)$, fractional traces, or integral
conditions.  The operator symbol alone does not determine them.}
\entryfield{Fourier derivative}{On the whole line, a multiplier definition such as
$\FourierAngular{D^\alpha f}(\omega)
=(\ImaginaryUnit\omega)^\alpha\FourierAngular f(\omega)$ requires a branch for the
complex power and is a third, separately named construction.}
\end{catalogueentry}

\begin{catalogueentry}{volterra}{Volterra integral operators and resolvents}
\entryfield{Operator}{For a kernel $K$ on $a\le t\le x\le b$,
\[
 (V_Kf)(x)\DefinitionEquals\int_a^xK(x,t)f(t)\,\Differential t.
\]}
\entryfield{Convolution kernel}{If $K(x,t)=k(x-t)$ and $a=0$, this is the one-sided
convolution $(k\Convolution_+f)(x)=\int_0^xk(x-t)f(t)\,\Differential t$.
It is not the whole-line \commandname{Convolution}.}
\entryfield{Iterates}{$V_K^0=\IdentityOperator$ and
$V_K^{n+1}=V_K\circ V_K^n$.  A superscript here is operator composition, not a pointwise
power.}
\entryfield{Resolvent}{When the Neumann series converges,
$(\IdentityOperator-\lambda V_K)^{-1}
=\sum_{n=0}^{\infty}\lambda^nV_K^n$.  The topology and allowed $\lambda$ are stated.}
\entryfield{Kernel resolvent}{A resolvent kernel $R_\lambda(x,t)$ represents the operator
resolvent after its identity term is separated; conventions that include a Dirac identity
kernel are labelled.}
\end{catalogueentry}

\begin{catalogueentry}{function-spaces}{Lebesgue, Sobolev, and Hölder spaces}
\entryfield{Lebesgue space}{$L^p(E)$, $1\le p<\infty$, consists of almost-everywhere
classes with norm
$\|f\|_{L^p(E)}=(\int_E|f|^p)^{1/p}$; $L^\infty$ uses essential supremum.  The measure
is Lebesgue measure unless printed otherwise.}
\entryfield{Smoothness}{$C^k(E)$ has continuous derivatives through order
$k\in\NaturalNumbers$; $C^\infty(E)=\bigcap_kC^k(E)$.  On a closed interval,
derivatives at endpoints are one-sided restrictions of derivatives on a neighborhood
unless another convention is stated.}
\entryfield{Hölder space}{For $0<\alpha\le1$,
\[
 [f]_{C^{0,\alpha}(E)}
 =\sup_{x\ne y\in E}\frac{|f(x)-f(y)|}{|x-y|^\alpha}.
\]
$C^{k,\alpha}$ applies this seminorm to the $k$th derivative and includes lower sup
norms in its norm.}
\entryfield{Sobolev space}{$W^{k,p}(E)$ uses weak derivatives through order $k$ in
$L^p$.  $H^s$ means the $L^2$-based Sobolev space, with Fourier normalization specified
when a Bessel-potential norm is displayed.}
\entryfield{Local spaces}{The suffix \(\mathrm{loc}\) means membership after restriction
to every compact subset of the open domain; it imposes no boundary control.}
\end{catalogueentry}

\begin{catalogueentry}{gevrey-flatness}{Gevrey classes, flatness, and nonanalyticity}
\entryfield{Gevrey class}{A smooth $f$ is locally Gevrey of order $s\ge1$ when every
compact $K$ admits $C,A>0$ such that
\[
 \sup_{x\in K}|f^{(n)}(x)|\le CA^n(n!)^s
 \quad(n\in\NaturalNumbers).
\]
For several variables, $n!$ is replaced by the multi-index factorial.}
\entryfield{Analytic class}{Order $s=1$ corresponds locally to real analyticity under the
standard derivative criterion.  Membership for $s>1$ does not imply analyticity.}
\entryfield{Flatness}{$f$ is flat at $a$ if $f^{(n)}(a)=0$ for every
$n\in\NaturalNumbers$.  A nonzero smooth function may be flat, so a zero Taylor series
does not identify the function near $a$.}
\entryfield{Fabius endpoints}{$\RvachevUp$ is flat at $\pm1$ and
$\FabiusBounded$ is flat at the transition points $0$ and $1$ relative to its constant
tails.  This is compatible with smoothness and with failure of analyticity there.}
\entryfield{Quasianalyticity}{A Denjoy--Carleman class is declared with its derivative
weight sequence $M_n$; whether it is quasianalytic is a theorem about that sequence, not
synonymous with being Gevrey.}
\end{catalogueentry}

\section{Exact Fabius moments and arithmetic normalizations}

\begin{catalogueentry}{full-moments}{Even full moments of the up-function}
\entryfield{Definition}{For $n\in\NaturalNumbers$,
\[
 c_n\DefinitionEquals
 \int_{\RealNumbers}t^{2n}\RvachevUp(t)\,\Differential t
 =\int_{-1}^{1}t^{2n}\RvachevUp(t)\,\Differential t.
\]
The equality of the two integrals uses the support of $\RvachevUp$.}
\entryfield{Initial value}{$c_0=1$ because $\RvachevUp$ is normalized to integral one.
The full odd moments vanish by evenness.}
\entryfield{Fourier expansion}{With the $2\pi$ convention,
\[
 \FourierTwoPi{\RvachevUp}(z)
 =\sum_{n=0}^{\infty}
 \frac{(-1)^nc_n}{(2n)!}(2\pi z)^{2n}.
\]
The series is entire; $z$ may be complex.}
\entryfield{Canonical integer numerator}{Define
\[
 N_n^{(c)}\DefinitionEquals
 c_n(2n+1)!!\prod_{j=1}^{n}(2^{2j}-1).
\]
The empty product gives $N_0^{(c)}=1$.  Historical prose used a capital-letter
sequence that collided with Fabius functions; new prose uses the semantic superscript.}
\entryfield{Lean declarations}{\lean{moment}~\(n\) and \lean{momentIntegral}~\(n\).}
\entryfield{Lean modules and status}{Exact arithmetic results are in
\modulepath{Arithmetic.lean}; analytic results are in
\modulepath{AnalyticMoments.lean}.  Status: \textbf{Lean-proved}.}
\end{catalogueentry}

\begin{catalogueentry}{half-moments}{Normalized half moments}
\entryfield{Definition}{Set $d_0=1$ and, for $n\in\PositiveIntegers$,
\[
 d_n\DefinitionEquals
 n\int_0^1t^{n-1}\RvachevUp(t)\,\Differential t.
\]
The separate definition at zero avoids the meaningless factor
$0\int_0^1t^{-1}\RvachevUp(t)\,\Differential t$.}
\entryfield{Binomial bridge}{For $n\in\NaturalNumbers$,
\[
 d_n=2^{-n}\sum_{0\le k\le n/2}\binom n{2k}c_k,
\]
where the bound means $k=0,\ldots,\Floor{n/2}$.}
\entryfield{Dyadic-value bridge}{For $n\in\NaturalNumbers$,
\[
 \FabiusBounded(2^{-n})
 =\frac{d_n}{n!\,2^{\binom n2}}.
\]
An inverse-dyadic value means evaluation at an inverse power of two, not a value of
\notelink{fabius-quantile}{the inverse function}.}
\entryfield{Canonical integer numerator}{Define
\[
 N_n^{(d)}\DefinitionEquals
 d_n(n+1)!\prod_{j=1}^{n}(2^j-1).
\]
Again the empty product gives $N_0^{(d)}=1$.  A historical capital-letter numerator
sequence is recorded as a retired alias rather than reused.}
\entryfield{Two-adic value}{For $n\ge1$, the reduced numerator and denominator of
$2d_n$ are odd, equivalently $\TwoAdicValuation(d_n)=-1$ for the rational valuation
extended additively to quotients.}
\entryfield{Lean crosswalk}{\lean{halfMoment}~\(n\) and \lean{halfMomentIntegral}~\(n\).
Status: \textbf{Lean-proved}.}
\end{catalogueentry}

\begin{catalogueentry}{reshetnikov-integers}{Reshetnikov's integer normalization}
\entryfield{Definition}{For $n\in\PositiveIntegers$,
\[
 R_n^{\mathrm{Res}}\DefinitionEquals
 2^{\binom{n-1}{2}}(2n)!\,
 \FabiusBounded(2^{-n})
 \prod_{j=1}^{\Floor{n/2}}(2^{2j}-1).
\]}
\entryfield{Type and theorem}{$R_n^{\mathrm{Res}}\in\PositiveIntegers$ and is odd.
Its superscript is mandatory in cross-topic documents because $R_n$ also denotes
polynomial remainders, resolvents, and recurrence polynomials.}
\entryfield{Factors}{The power of two is triangular, the factorial is $(2n)!$, and the
Mersenne-type product stops at $\Floor{n/2}$.  Removing or shifting any of these produces
a different normalization.}
\entryfield{Lean crosswalk}{\lean{reshetnikov}~\(n\) and the corresponding integrality,
positivity, and oddness theorems in \modulepath{Arithmetic.lean} and
\modulepath{PaperStatements.lean}.  Status: \textbf{Lean-proved}.}
\end{catalogueentry}

\begin{catalogueentry}{dyadic-denominator}{Dyadic common denominators}
\entryfield{Grid}{For $n\in\PositiveIntegers$, let
$A_n=\{a\in\PositiveIntegers:a<2^n,\ a\text{ odd}\}$.}
\entryfield{Reduced denominator}{For $r\in\RationalNumbers$, write
$r=p/q$ with $q\in\PositiveIntegers$ and $\gcd(p,q)=1$; then
$\operatorname{den}(r)=q$, including $\operatorname{den}(0)=1$.}
\entryfield{Definition}{The semantic dyadic denominator is
\[
 D_n^{\mathrm{dyad}}\DefinitionEquals
 \LeastCommonMultiple\{
 \operatorname{den}(\RvachevUp(a/2^n)):a\in A_n\}.
\]
At $n=0$ the index set is empty; the catalogue assigns
$D_0^{\mathrm{dyad}}=1$, the LCM identity, when a zero-level value is needed.}
\entryfield{Reflection}{The same LCM is obtained from the corresponding bounded-Fabius
values after applying the exact reflection/translation relation.}
\entryfield{Collision policy}{A bare $D_n$ is never globally canonical: it also denotes
determinants and asymptotic main terms.  The semantic superscript is required outside a
single-purpose local scope.}
\entryfield{Lean declaration}{\lean{dyadicDenominator}~\(n\).}
\entryfield{Lean modules and status}{The definition is in \modulepath{Arithmetic.lean}.
\par\noindent Bounds are in \modulepath{DenominatorBound.lean}.  Status:
\textbf{Lean-proved}.}
\end{catalogueentry}

\section{Approximation, orthogonal polynomials, splines, and quadrature}

\begin{catalogueentry}{weighted-inner-product}{Weights and polynomial inner products}
\entryfield{Measure form}{For a positive Borel measure $\mu$ with finite needed moments,
\[
 \langle f,g\rangle_\mu\DefinitionEquals
 \int f(x)\overline{g(x)}\,\Differential\mu(x).
\]
The inner product is linear in the first argument and conjugate-linear in the second in
this catalogue.}
\entryfield{Weight form}{If $\Differential\mu(x)=w(x)\Differential x$ on an interval
$I$, write $\langle f,g\rangle_w=\int_I f(x)\overline{g(x)}w(x)\,\Differential x$,
with $w\ge0$ almost everywhere.}
\entryfield{Norm}{$\|f\|_\mu=\sqrt{\langle f,f\rangle_\mu}$.  It is a norm only after
quotienting functions equal $\mu$-almost everywhere or restricting to a nondegenerate
function space.}
\entryfield{Fabius weights}{When the measure has density $\RvachevUp$ on $[-1,1]$, write
$\InnerProduct{f,g}_{\RvachevUp}$ after declaring that measure.  An unlabeled upright-word
subscript is retired because it bypasses the semantic command.  A bounded-Fabius
Stieltjes measure is a different but related object.}
\end{catalogueentry}

\begin{catalogueentry}{orthogonal-polynomials}{Orthogonal, monic, and orthonormal
polynomials}
\entryfield{Orthogonal family}{Polynomials $(p_n)_{n\in\NaturalNumbers}$ satisfy
$\deg p_n=n$ and $\langle p_n,p_m\rangle_\mu=0$ for $n\ne m$.}
\entryfield{Monic normalization}{$P_n$ is monic when its coefficient of $x^n$ is $1$.
This fixes an orthogonal polynomial uniquely when the moment functional is positive
definite.}
\entryfield{Orthonormal normalization}{$\pi_n$ is orthonormal when
$\langle\pi_n,\pi_m\rangle_\mu=\delta_{nm}$ and its leading coefficient is positive.
Thus $\pi_n=P_n/\|P_n\|_\mu$.}
\entryfield{Three-term recurrence}{For monic polynomials,
\[
 P_{n+1}(x)=(x-\alpha_n)P_n(x)-\beta_nP_{n-1}(x),
\]
with $P_{-1}=0$, $P_0=1$, $\alpha_n\in\RealNumbers$, and $\beta_n>0$ for
$n\ge1$.  Orthonormal recurrence coefficients use a different normalization.}
\entryfield{Classical names}{Legendre, Chebyshev, Jacobi, and Gegenbauer polynomials
carry their standard family parameters only after their normalization is displayed;
several incompatible leading-value conventions exist.}
\end{catalogueentry}

\begin{catalogueentry}{legendre-chebyshev}{Legendre and Chebyshev normalizations}
\entryfield{Legendre}{$P_n^{\mathrm{Leg}}$ is defined by
$P_n^{\mathrm{Leg}}(1)=1$ and orthogonality for weight $1$ on $[-1,1]$.  Its leading
coefficient is $2^{-n}\binom{2n}{n}$, so it is not monic for $n>1$.}
\entryfield{First Chebyshev family}{$T_n(\cos\theta)=\cos(n\theta)$,
$T_n(1)=1$, orthogonal for $(1-x^2)^{-1/2}$ on $[-1,1]$.}
\entryfield{Second Chebyshev family}{$U_n(\cos\theta)
=\sin((n+1)\theta)/\sin\theta$, with the continuous endpoint values, orthogonal for
$(1-x^2)^{1/2}$.}
\entryfield{Coefficient convention}{A Legendre coefficient may be the raw inner product
$\langle f,P_n^{\mathrm{Leg}}\rangle$ or the expansion coefficient
$(2n+1)\langle f,P_n^{\mathrm{Leg}}\rangle/2$.  The corpus prints which one it means.}
\entryfield{Collision}{The letter $P_n$ is local: it also denotes probability,
prefix polynomials, and Padé denominators.  A family superscript is required in a
cross-topic statement.}
\end{catalogueentry}

\begin{catalogueentry}{legendre-gaunt-wigner-zero-row}
{Legendre Gaunt coefficients and the squared zero-row Wigner \(3j\) datum}
\entryfield{Indices and normalization}{Throughout this entry
\(i,j,k\in\NaturalNumbers\), and \(P_n^{\mathrm{Leg}}(1)=1\).  The Gaunt coefficient is
\[
 G_{ijk}\DefinitionEquals
 \int_{-1}^{1}P_i^{\mathrm{Leg}}(x)P_j^{\mathrm{Leg}}(x)
 P_k^{\mathrm{Leg}}(x)\,\Differential x.
\]
It is not a coefficient for monic or orthonormal Legendre polynomials.}
\entryfield{Admissible support}{Write \(\mathsf{Adm}(i,j,k)\) for the simultaneous
conditions
\[
 i+j+k\equiv0\pmod2,
 \qquad i\le j+k,\quad j\le i+k,\quad k\le i+j.
\]
Thus admissibility means even total degree and all three weak triangle inequalities;
none is implicit.  On this support, \(s\DefinitionEquals(i+j+k)/2\) is a natural number
and \(s-i,s-j,s-k\in\NaturalNumbers\).}
\entryfield{Total rational square datum}{Define
\[
 W_{ijk}^{2,0}\DefinitionEquals
 \begin{cases}
 \displaystyle
 \frac{
  \binom{2(s-i)}{s-i}
  \binom{2(s-j)}{s-j}
  \binom{2(s-k)}{s-k}}
 {(2s+1)\binom{2s}{s}},
 &\mathsf{Adm}(i,j,k),\\[1.1em]
 0,&\text{otherwise}.
 \end{cases}
\]
This is a total element of \(\RationalNumbers\) and equals the square of the integer-index
zero-row Wigner \(3j\) datum.  The superscript \(2,0\) records ``square'' and ``zero
magnetic row''; it does not select a signed Wigner symbol.}
\entryfield{Gaunt and product-linearization identities}{The exact relations are
\[
 G_{ijk}=2W_{ijk}^{2,0},
 \qquad
 [P_k^{\mathrm{Leg}}]\bigl(P_i^{\mathrm{Leg}}P_j^{\mathrm{Leg}}\bigr)
 =(2k+1)W_{ijk}^{2,0},
\]
where the bracket denotes the coefficient of \(P_k^{\mathrm{Leg}}\) in the finite
Legendre-basis expansion.  Off the admissible support both sides of the first identity
vanish.}
\entryfield{Triangle coordinates and factorial form}{Equivalently, there are
\(a,b,c\in\NaturalNumbers\) with
\(i=b+c\), \(j=a+c\), and \(k=a+b\).  Then
\[
 W_{ijk}^{2,0}
 =\frac{\binom{2a}{a}\binom{2b}{b}\binom{2c}{c}}
 {(2(a+b+c)+1)\binom{2(a+b+c)}{a+b+c}},
\]
and expansion of each central binomial coefficient gives the source's equivalent
factorial quotient.}
\entryfield{Rvachev Gram use family}{Substitution into the finite Rvachev--Legendre
Gaunt expansion gives rational Gram entries, rational Gram-matrix entries, and---under a
bounded Fabius object plus its Fabius-specification proof---real up-law Gram entries as
finite sums of canonical Rvachev--Legendre coefficients times \(W_{ij,2r}^{2,0}\).  This
finite use does not authorize an infinite Legendre-series interchange.}
\entryfield{Formal boundary}{Only the total rational square of the integer-index,
zero-magnetic-row datum is present.  There is no sign or Condon--Shortley phase choice,
no half-integer spin, no nonzero magnetic index, no general \(3j\), \(6j\), or \(9j\)
symbol, and no Wigner orthogonality or recoupling theory.}
\entryfield{Lean crosswalk}{\lean{legendreGauntAdmissible} and
\lean{legendreWignerThreeJZeroSqRat} in \modulepath{LegendreGauntClosedForm.lean}
implement the support predicate and total square datum.  The central identities are
\leanline{legendreGauntRat_eq_two_mul_wignerThreeJZeroSqRat}
\leanline{legendreGaunt_eq_two_mul_wignerThreeJZeroSqRat}
\leanline{legendreProductLinearizationCoeffRat_eq_mul_wignerThreeJZeroSqRat}
\modulepath{FabiusLegendreGauntClosedForm.lean} supplies the finite rational and real
Rvachev Gram-entry sums, with \lean{IsFabius} required for the analytic real formula.
Status: \textbf{Lean-proved with the displayed hypotheses}.}
\end{catalogueentry}

\begin{catalogueentry}{moments-hankel}{Moment matrices and Hankel determinants}
\entryfield{Moment sequence}{For a finite measure $\mu$,
$m_n=\int x^n\,\Differential\mu(x)$ for $n\in\NaturalNumbers$ whenever finite.}
\entryfield{Hankel matrix}{The order-$N$ moment matrix is
\[
 H_N(\mu)\DefinitionEquals(m_{i+j})_{0\le i,j<N}
 \in\ComplexNumbers^{N\times N},
 \qquad N\in\NaturalNumbers.
\]
Thus $H_0$ is the unique $0\times0$ matrix.}
\entryfield{Determinant}{$\Delta_N(\mu)=\det H_N(\mu)$ and
$\Delta_0(\mu)=1$.  A positive measure with infinite support has $\Delta_N>0$ for every
$N\ge1$.}
\entryfield{Shifted matrix}{$H_N^{(r)}=(m_{i+j+r})_{0\le i,j<N}$ for
$r\in\NaturalNumbers$.  The shift belongs in the superscript, not in an unstated change
of the moment sequence.}
\entryfield{Fabius parity}{For the centered up measure, odd moments vanish and even
moments are $c_n$.  Reindexing the even block must show whether its entry is
$c_{i+j}$ or $c_{i+j+r}$.}
\end{catalogueentry}

\begin{catalogueentry}{jacobi-matrix}{Jacobi matrices and spectral measures}
\entryfield{Matrix}{For orthonormal polynomials with
$x\pi_n=a_{n+1}\pi_{n+1}+b_n\pi_n+a_n\pi_{n-1}$,
$a_n>0$, the Jacobi operator is the tridiagonal matrix
\[
 J=\begin{pmatrix}
 b_0&a_1&0&\cdots\\
 a_1&b_1&a_2&\ddots\\
 0&a_2&b_2&\ddots\\
 \vdots&\ddots&\ddots&\ddots
 \end{pmatrix}.
\]}
\entryfield{Indexing}{Rows and columns begin at $0$; $a_0$ does not appear.
A finite truncation $J_N$ uses indices $0,\ldots,N-1$.}
\entryfield{Spectral measure}{The measure $\mu$ satisfies
$\langle e_0,p(J)e_0\rangle=\int p(x)\,\Differential\mu(x)$ for polynomials $p$,
under the spectral theorem.}
\entryfield{Resolvent}{$G_\mu(z)=\langle e_0,(z\IdentityOperator-J)^{-1}e_0\rangle$
matches the Cauchy-transform denominator $z-x$.}
\entryfield{Frontier status}{Any claimed closed recurrence coefficients for the Fabius
measure remain \textbf{frontier/conjectural} unless linked to a proved Lean declaration.}
\end{catalogueentry}

\begin{catalogueentry}{splines}{Knots, B-splines, degree, and order}
\entryfield{Knot vector}{A knot vector is a nondecreasing finite sequence
$\tau=(\tau_0,\ldots,\tau_{m})$.  Repeated knots are allowed and their multiplicity
affects smoothness.}
\entryfield{Degree versus order}{A B-spline of degree $p\in\NaturalNumbers$ has order
$p+1$.  The corpus always says degree or order; a bare index is not assumed to be one
or the other.}
\entryfield{Cox--de Boor base case}{$B_{i,0}(x)=
\IndicatorOf{[\tau_i,\tau_{i+1})}(x)$, except that a designated final endpoint may be
closed to make a partition of unity on the full knot span.}
\entryfield{Recursion}{For $p\ge1$,
\[
 B_{i,p}(x)=
 \frac{x-\tau_i}{\tau_{i+p}-\tau_i}B_{i,p-1}(x)
 +\frac{\tau_{i+p+1}-x}{\tau_{i+p+1}-\tau_{i+1}}B_{i+1,p-1}(x),
\]
where a term with zero denominator is defined to be zero.}
\entryfield{Support}{$\SupportOf{B_{i,p}}\subseteq[\tau_i,\tau_{i+p+1}]$.
Equality needs nondegenerate knot spans.}
\entryfield{Cardinal splines}{Uniform integer knots and a translated normalization are
declared before using the shorter cardinal B-spline symbol.}
\end{catalogueentry}

\begin{catalogueentry}{divided-differences}{Divided differences}
\entryfield{Distinct nodes}{For pairwise distinct $x_0,\ldots,x_n$,
\[
 [x_0,\ldots,x_n]f
 =\sum_{j=0}^n\frac{f(x_j)}
 {\prod_{0\le k\le n,\ k\ne j}(x_j-x_k)}.
\]}
\entryfield{Recursive form}{$[x_j]f=f(x_j)$ and
$[x_0,\ldots,x_n]f=
([x_1,\ldots,x_n]f-[x_0,\ldots,x_{n-1}]f)/(x_n-x_0)$.}
\entryfield{Repeated nodes}{If nodes coalesce and $f$ is sufficiently differentiable,
$[\underbrace{x,\ldots,x}_{n+1}]f=f^{(n)}(x)/n!$.  This is a continuous extension, not
substitution into a zero denominator.}
\entryfield{Spline relation}{B-splines expressed as divided differences state the
truncated-power normalization and knot multiplicities.}
\end{catalogueentry}

\begin{catalogueentry}{quadrature}{Quadrature nodes, weights, and exactness}
\entryfield{Rule}{An $N$-node quadrature rule for $\mu$ is
$Q_N(f)=\sum_{j=1}^{N}w_{j,N}f(x_{j,N})$, with nodes and weights indexed from $1$.}
\entryfield{Degree of exactness}{It has degree of exactness $d$ when
$Q_N(p)=\int p\,\Differential\mu$ for all polynomials of degree at most $d$, and this
fails for at least one polynomial of degree $d+1$.}
\entryfield{Gaussian rule}{For a positive measure, Gaussian nodes are the zeros of the
degree-$N$ orthogonal polynomial; the rule has positive weights and degree $2N-1$.}
\entryfield{Endpoint variants}{Gauss--Radau fixes one endpoint and Gauss--Lobatto fixes
both.  Their exactness degrees and endpoint weights differ from the Gaussian rule and are
printed separately.}
\entryfield{Error}{An expression $E_N(f)=\int f\,\Differential\mu-Q_N(f)$ fixes the sign.
The derivative form of an error requires the exact smoothness interval and normalization
constant.}
\entryfield{Shifted dyadic Rvachev self-sampling}{Let \(F\) be a bounded Fabius object
satisfying the Fabius specification, \(N\in\NaturalNumbers\),
\(P\in\RealNumbers[X]\) with \(\deg P\le N\), and \(\theta\in\RealNumbers\).  Then
\[
 \int_{\RealNumbers}P(x)\RvachevUp_F(x)\,\Differential x
 =2^{-N}\sum_{k\in\IntegerNumbers}
 P\!\left(2^{-N}(\theta+k)\right)
 \RvachevUp_F\!\left(2^{-N}(\theta+k)\right).
\]
The right side is formally a bilateral infinite sum, but compact support makes it finite.
The theorem proves exactness for every polynomial of degree at most \(N\); it does not
assert failure in degree \(N+1\), so it does not by itself establish that the degree of
exactness is exactly \(N\).}
\entryfield{Lean crosswalk}{In \modulepath{PolynomialCombExactness.lean}, the displayed
rule is the declaration
\leanline{integral_polynomial_mul_rvachevUp_eq_dyadic_tsum}
Its normalized source identity is
\leanline{tsum_shifted_polynomial_eq_integral}
Both retain the bounded-Fabius and
Fabius-specification hypotheses.  Status: \textbf{Lean-proved}.}
\end{catalogueentry}

\begin{catalogueentry}{pade}{Padé approximation}
\entryfield{Definition}{For a formal series $f(z)=\sum_{n\ge0}a_nz^n$, a type
$[L/M]$ Padé approximant is a rational function
$P_L(z)/Q_M(z)$ with $\deg P_L\le L$, $\deg Q_M\le M$,
$Q_M(0)=1$, and
\[
 Q_M(z)f(z)-P_L(z)=\BigOAt{z^{L+M+1}}{z\to0}.
\]}
\entryfield{Existence and uniqueness}{A normalized pair may fail to exist with exact
degrees or may have a common factor.  The rational function is unique under the standard
normality hypothesis, which is stated when used.}
\entryfield{Poles}{Zeros of $Q_M$ after cancellation are poles of the approximant.
Uncancelled numerator-denominator zero pairs are not counted as poles.}
\entryfield{Collision policy}{$P,Q$ are local polynomial names; they do not denote
probability or the Fabius quantile inside this entry.  Outside it, the type-$[L/M]$
subscript or semantic label is retained.}
\end{catalogueentry}

\section{Linear algebra, operators, kernels, and spectral notation}

\begin{catalogueentry}{vectors-matrices}{Vectors, matrices, and indices}
\entryfield{Vectors}{A column vector is $x=(x_1,\ldots,x_n)^{\mathsf T}\in
\mathbb K^n$, where $\mathbb K$ is explicitly $\RealNumbers$ or $\ComplexNumbers$.
Algorithmic zero-based coordinates are marked separately.}
\entryfield{Matrices}{$A=(a_{ij})_{1\le i\le m,\,1\le j\le n}
\in\mathbb K^{m\times n}$.  The first dimension is the number of rows.}
\entryfield{Transpose and adjoint}{$A^{\mathsf T}$ is transpose,
$\overline A$ entrywise conjugation, and
$A^\ast=\overline A^{\mathsf T}$ conjugate transpose.  For real matrices,
$A^\ast=A^{\mathsf T}$.}
\entryfield{Identity and zeros}{$I_n$ is the $n\times n$ identity and $0_{m\times n}$
the rectangular zero matrix.  A bare $0$ inherits type only inside a displayed typed
equation.}
\entryfield{Empty size}{The unique $0\times0$ matrix has determinant $1$, trace $0$, and
rank $0$.  Matrix products with a zero inner dimension are zero matrices by empty sums.}
\end{catalogueentry}

\begin{catalogueentry}{matrix-operators}{Trace, rank, determinant, and diagonal}
\entryfield{Trace}{For $A\in\mathbb K^{n\times n}$,
$\TraceOperator A=\sum_{j=1}^{n}a_{jj}$, with $\TraceOperator A=0$ at $n=0$.}
\entryfield{Trace-class extension}{For a trace-class operator $T$ on a real or complex
Hilbert space and any orthonormal basis $(e_j)$,
$\TraceOperator T\DefinitionEquals\sum_j\InnerProduct{Te_j,e_j}$.  Trace class makes
this series absolutely convergent and its value independent of the orthonormal basis; in
finite dimension it agrees with the matrix trace above.  No trace is inferred for an
arbitrary bounded infinite-dimensional operator.}
\entryfield{Rank}{$\RankOperator A$ is the dimension of the image of the associated linear
map.  Row rank and column rank agree over a field.}
\entryfield{Determinant}{$\det A$ is defined only for a square matrix or an endomorphism
of a finite free module with a chosen basis-free determinant construction.}
\entryfield{Diagonal}{$\DiagonalOperator(x_1,\ldots,x_n)$ is the square diagonal matrix
with those diagonal entries.  Applying \commandname{DiagonalOperator} to a matrix in
order to extract a vector is forbidden; extraction is written
$\operatorname{diagvec}(A)$ to prevent an input-type overload.}
\entryfield{Span}{For vectors $v_j$,
$\SpanOperator_{\mathbb K}\{v_j:j\in J\}$ is their linear span over $\mathbb K$.
The scalar field is printed when not obvious.}
\end{catalogueentry}

\begin{catalogueentry}{matrix-products}{Tensor, Kronecker, and Hadamard products}
\entryfield{Tensor product}{$V\otimes_{\mathbb K}W$ is the algebraic tensor product of
vector spaces; completed topological tensor products carry a norm or topology subscript.}
\entryfield{Kronecker product}{For matrices $A=(a_{ij})$ and $B$,
$A\otimes B$ is the block matrix $(a_{ij}B)$.  Its dimensions are the pairwise products
of the input dimensions.}
\entryfield{Hadamard product}{$A\circ B=(a_{ij}b_{ij})$ for matrices of identical shape.
It is not composition.}
\entryfield{Composition}{$S\circ T$ means first $T$, then $S$.  When Hadamard products
occur in the same section, they use $\odot$ instead of $\circ$.}
\entryfield{Powers}{$A^0=I_n$ for square $A\in\mathbb K^{n\times n}$, including the zero
matrix.  Entrywise powers are written $A^{\circ k}$.}
\end{catalogueentry}

\begin{catalogueentry}{norms}{Vector, matrix, and operator norms}
\entryfield{Euclidean norm}{For $x\in\ComplexNumbers^n$,
$\|x\|_2=(\sum_j|x_j|^2)^{1/2}$.}
\entryfield{Induced norm}{For $A:\mathbb K^n\to\mathbb K^m$,
$\|A\|_{p\to q}=\sup_{x\ne0}\|Ax\|_q/\|x\|_p$.  The spectral norm is
$\|A\|_{2\to2}$.}
\entryfield{Frobenius norm}{$\|A\|_{\mathrm F}
=(\sum_{i,j}|a_{ij}|^2)^{1/2}=(\TraceOperator A^\ast A)^{1/2}$.}
\entryfield{Unsubscripted norm}{The shared \commandname{Norm} supplies delimiters only.
Each local scope declares the norm behind them; it is not globally the Euclidean or
supremum norm.}
\entryfield{Condition number}{For invertible $A$,
$\kappa_p(A)=\|A\|_{p\to p}\|A^{-1}\|_{p\to p}$; for singular $A$ it is $+\infty$.}
\end{catalogueentry}

\begin{catalogueentry}{spectrum}{Spectrum, eigenvalues, and resolvents}
\entryfield{Finite spectrum}{For $A\in\ComplexNumbers^{n\times n}$,
$\sigma(A)=\{\lambda\in\ComplexNumbers:\det(\lambda I_n-A)=0\}$.}
\entryfield{Multiplicity}{Algebraic multiplicity is root multiplicity in the
characteristic polynomial; geometric multiplicity is
$\dim\ker(A-\lambda I_n)$.  They are not conflated.}
\entryfield{Spectral radius}{$\rho(A)=\max_{\lambda\in\sigma(A)}|\lambda|$ for
$n>0$; for the $0\times0$ matrix the spectrum is empty and the catalogue sets
$\rho(A)=0$ when a total convention is required.}
\entryfield{Operator spectrum}{For a bounded operator $T$ on a complex Banach space,
$\sigma(T)=\{\lambda:\lambda\IdentityOperator-T\text{ is not boundedly invertible}\}$.}
\entryfield{Resolvent}{$R_T(z)=(z\IdentityOperator-T)^{-1}$ for
$z\notin\sigma(T)$.  This operator-valued $R_T$ is distinct from
$R_n^{\mathrm{Res}}$.}
\end{catalogueentry}

\begin{catalogueentry}{hankel-toeplitz}{Hankel and Toeplitz matrices}
\entryfield{Hankel}{A finite Hankel matrix has entries $H_{ij}=a_{i+j}$, with declared
index origins.  In this catalogue $0\le i,j<N$ is the default.}
\entryfield{Toeplitz}{A Toeplitz matrix has entries $T_{ij}=a_{i-j}$ and therefore
requires a two-sided sequence for zero-based finite indices.}
\entryfield{Moment interpretation}{A Hankel matrix $(m_{i+j})$ is positive semidefinite
when $(m_n)$ is a moment sequence of a positive real measure.}
\entryfield{Operator interpretation}{Infinite Hankel or Toeplitz matrices define
operators only after a sequence space and boundedness/domain conditions are supplied.}
\entryfield{Collision}{A transform, matrix, and determinant may all be called $H$ in
source literature.  Cross-topic prose uses $H_N(\mu)$, $T_N(a)$, or a semantic operator
subscript.}
\end{catalogueentry}

\begin{catalogueentry}{total-positivity}{Minors and total positivity}
\entryfield{Minor}{For row index set $I$ and column index set $J$ of equal finite size,
$A[I,J]$ is the corresponding submatrix and
$\Delta_{I,J}(A)=\det A[I,J]$.  Both index sets are listed in increasing order.}
\entryfield{Totally nonnegative}{A real matrix is totally nonnegative if every minor is
nonnegative.}
\entryfield{Totally positive}{It is totally positive if every minor that is not forced to
vanish by the matrix shape is positive.  For a rectangular matrix, all square orders up
to $\min(m,n)$ are included.}
\entryfield{Strict versus order \(r\)}{The suffix \(\mathrm{TP}_r\) restricts to minors of
order at most $r$; \(\mathrm{TN}_r\) is the nonnegative analogue.  Neither implies
unqualified total positivity when $r$ is finite below matrix size.}
\entryfield{Kernel version}{A kernel $K(x,y)$ is totally positive when determinants
$\det(K(x_i,y_j))$ have the specified sign for strictly ordered point tuples.  Weak
orders and repeated nodes require a confluent formulation.}
\entryfield{Status}{Total-positivity claims for Fabius moment or collocation matrices are
\textbf{frontier/conjectural} unless their exact minor family is linked to a proof.}
\end{catalogueentry}

\begin{catalogueentry}{positive-definite-kernels}{Positive-definite kernels and RKHS
notation}
\entryfield{Positive definite}{A Hermitian kernel $K:X\times X\to\ComplexNumbers$ is
positive semidefinite if
$\sum_{i,j}\overline{c_i}c_jK(x_i,x_j)\ge0$ for all finite choices.}
\entryfield{Strictness}{Strict positive definiteness requires positivity for nonzero
coefficient vectors when the points are pairwise distinct.}
\entryfield{RKHS}{The reproducing-kernel Hilbert space $\mathcal H_K$ satisfies
$f(x)=\langle f,K(\,\cdot\,,x)\rangle_{\mathcal H_K}$ under the catalogue's
first-slot-linear convention.  References using second-slot linearity conjugate this
identity.}
\entryfield{Gram matrix}{$G_X=(K(x_i,x_j))_{i,j=1}^N$.  Its rank may be smaller than
$N$ even for a positive-semidefinite kernel.}
\entryfield{Collision}{Kernel $K$, compact set $K$, and field $\mathbb K$ may coexist only
with type-bearing subscripts or a local declaration.}
\end{catalogueentry}

\section{Research-frontier namespaces}

The entries in this section define the symbols needed to read the semi-formalized frontier
documents.  They do not assert that the proposed Fabius specializations are theorems.
Every such specialization carries one of the proof-status labels in
\cref{sec:status-register}.

\begin{catalogueentry}{appell-sheffer}{Appell and Sheffer polynomial families}
\entryfield{Appell family}{A polynomial sequence $(A_n)_{n\in\NaturalNumbers}$ is Appell
when $A_0$ is a nonzero constant and
$A_n'(x)=nA_{n-1}(x)$ for $n\ge1$.}
\entryfield{Generating function}{With the normalization $A_0=1$,
\[
 \sum_{n=0}^{\infty}A_n(x)\frac{z^n}{n!}
 =g(z)\EulerE^{xz},\qquad g(0)=1,
\]
as a formal identity or analytic identity on a printed disk.}
\entryfield{Associated law}{If $g(z)=1/\MomentGeneratingFunction X(z)$, the family is
associated with a law only where the reciprocal series exists; this does not imply that
all $A_n$ are orthogonal.}
\entryfield{Sheffer family}{A Sheffer sequence has EGF
$g(z)\exp(xf(z))$ with $g(0)\ne0$, $f(0)=0$, and $f'(0)\ne0$.
Appell is the case $f(z)=z$.}
\entryfield{Collision}{The local \(A_n\) is neither a matrix nor an event.  A generic
Appell family retains a declared generator \(g\); the Rvachev reciprocal-moment family is
\(A_n^{\mathrm{Rv}}\) outside its defining section and is specified in the next entry.}
\end{catalogueentry}

\begin{catalogueentry}{rvachev-appell-deconvolution}{Rvachev moments, Appell polynomials,
and polynomial deconvolution}
\entryfield{Probability input}{Let \(F\) be a bounded Fabius object satisfying the
project's Fabius specification, and put \(u_F=\RvachevUp_F\).  The associated probability
measure has density \(u_F(y)\,\Differential y\), total mass \(1\), even symmetry, and
support contained in \([-1,1]\).  Every analytic identity below retains this hypothesis
on \(F\); the notation alone does not manufacture a canonical solution.}
\entryfield{Raw moments}{For \(n\in\NaturalNumbers\),
\[
 m_n^{\mathrm{Rv}}
 \DefinitionEquals\int_{\RealNumbers}y^n u_F(y)\,\Differential y.
\]
Even symmetry gives \(m_{2j+1}^{\mathrm{Rv}}=0\).  The project stores exact rational
values in \lean{rvachevRawMomentRat}; casting that sequence to \(\RealNumbers\) is a
separate typed operation.}
\entryfield{Binomial-convolution reciprocal}{The rational sequence
\((\gamma_n^{\mathrm{Rv}})_{n\ge0}\) is uniquely defined by
\[
 \sum_{k=0}^{n}\binom nk
 m_k^{\mathrm{Rv}}\gamma_{n-k}^{\mathrm{Rv}}
 =\delta_{n,0}\qquad(n\in\NaturalNumbers).
\]
This is reciprocal under exponential-generating-function convolution.  It is not
pointwise reciprocal \(1/m_n^{\mathrm{Rv}}\), and it remains meaningful algebraically
without assuming convergence of an analytic reciprocal moment-generating function.}
\entryfield{Rvachev--Appell family}{Define
\[
 A_n^{\mathrm{Rv}}(x)
 \DefinitionEquals
 \sum_{k=0}^{n}\binom nk
 \gamma_{n-k}^{\mathrm{Rv}}x^k.
\]
Then \(A_n^{\mathrm{Rv}}\) is monic of exact degree \(n\), and
\((A_n^{\mathrm{Rv}})'=nA_{n-1}^{\mathrm{Rv}}\) for \(n\ge1\).  The Lean names are
\lean{rvachevAppellPolynomialRat} over \(\RationalNumbers\) and
\lean{rvachevAppellPolynomial} over \(\RealNumbers\).}
\entryfield{Polynomial deconvolution operator}{For
\(P(x)=\sum_{n=0}^{d}p_nx^n\in\RealNumbers[x]\), define
\[
 (\mathscr D_{\RvachevUp}P)(x)
 \DefinitionEquals
 \sum_{n=0}^{d}p_n A_n^{\mathrm{Rv}}(x).
\]
Thus \(\mathscr D_{\RvachevUp}\) is a real-linear map on polynomials and sends
\(x^n\) to \(A_n^{\mathrm{Rv}}\).  It preserves the natural degree and leading
coefficient of every nonzero polynomial and is injective.  The semantic subscript names
the smoothing law being inverted; \(\mathscr D\), \(\mathcal D\), and
\(\mathscr D_u\) without that subscript are not corpus-wide aliases.}
\entryfield{Smoothing identities}{For \(x\in\RealNumbers\),
\[
 \int_{\RealNumbers}A_n^{\mathrm{Rv}}(x+y)u_F(y)\,\Differential y=x^n,
 \qquad
 \int_{\RealNumbers}(\mathscr D_{\RvachevUp}P)(x+y)u_F(y)\,\Differential y=P(x).
\]
Evenness of \(u_F\) also permits \(x-y\) in place of \(x+y\).  For
\(n\in\PositiveIntegers\),
\[
 \int_{\RealNumbers}A_n^{\mathrm{Rv}}(y)u_F(y)\,\Differential y=0.
\]
These statements are pointwise polynomial identities, not merely almost-everywhere
relations.}
\entryfield{Mesh synthesis}{Let \(M\in\PositiveIntegers\) and
\(\deg P\le\TwoAdicValuation(M)\).  Then
\[
 \sum_{k\in\IntegerNumbers}
 (\mathscr D_{\RvachevUp}P)(k/M)u_F(x-k/M)=MP(x).
\]
Compact support makes the bilateral sum pointwise finite.  For \(x\in[-1,1]\), the exact
uniform truncation is
\[
 P(x)=\frac1M\sum_{-2M<k<2M}
 (\mathscr D_{\RvachevUp}P)(k/M)u_F(x-k/M).
\]
The strict inequalities omit both endpoints; their atoms vanish by support.}
\entryfield{Arbitrary phase}{For every \(\theta,x\in\RealNumbers\), under the same
\(M\) and degree hypotheses,
\[
 P(x)=\frac1M\sum_{k\in\IntegerNumbers}
 (\mathscr D_{\RvachevUp}P)
 \!\left(x-\frac{\theta+k}{M}\right)
 u_F\!\left(\frac{\theta+k}{M}\right).
\]
A summand can be nonzero only if \(|\theta+k|\le M\), so the displayed bilateral sum is
finite.  The shifts contain explicit division by \(M\); juxtaposition does not mean
\(x-\theta-k/M\).}
\entryfield{Dyadic specialization and frontier boundary}{At \(M=2^N\), every polynomial
of degree at most \(N\) is covered.  The proposed extension to degree \(N+1\) at phases
\(\theta\equiv0,1/2\pmod1\) for even \(N\), or
\(\theta\equiv1/4,3/4\pmod1\) for odd \(N\), remains a frontier assertion and is not
included in the Lean-proved theorem above.}
\entryfield{Lean crosswalk}{The exact names are listed one per semantic role:
\par\smallskip\noindent
Raw moment: \lean{rvachevRawMomentRat}.
\par\noindent
Reciprocal moment: \lean{rvachevReciprocalMomentRat}.
\par\noindent
Rational Appell family: \lean{rvachevAppellPolynomialRat}.
\par\noindent
Real Appell family: \lean{rvachevAppellPolynomial}.
\par\noindent
Polynomial deconvolution: \lean{rvachevDeconvolvedPolynomial}.
\par\noindent
Bundled linear map: \lean{rvachevDeconvolutionLinearMap}.
\par\smallskip\noindent
These definitions are in \modulepath{RvachevMomentAppell.lean}.
\par\noindent
Mesh and phase reproduction are in
\modulepath{RvachevPolynomialSynthesis.lean}.  Status:
\textbf{Lean-proved with the displayed hypotheses}; no analytic reciprocal-MGF
convergence or infinite differential-operator expansion is claimed.}
\end{catalogueentry}

\begin{catalogueentry}{polynomial-logarithmic-ambient}
  {Polynomial--logarithmic variables, coefficient classes, and completions}
\entryfield{Commands}{The shared semantic commands are
\commandname{PolynomialLogPowerSeriesRing},
\commandname{PolynomialLogLaurentRing},
\commandname{RationalLogLaurentField},
\commandname{IteratedLaurentLogField}, and
\commandname{NearIdentityPolynomialLogGroup}.  Their arguments name the
coefficient field.  A document-local calligraphic letter is not a synonym for
any of these commands unless that document explicitly defines it.}
\entryfield{Base data}{Fix a field \(K\) of characteristic zero, a large
variable \(X\), and the generators
\[
 t\DefinitionEquals X^{-1},
 \qquad
 L\DefinitionEquals\log X.
\]
For real asymptotics the default regime is \(X\to+\infty\).  For complex
asymptotics a sector and a branch of \(\log\) must be supplied.  In the formal
algebra, \(t\) and \(L\) are independent generators; their analytic relation is
recorded by the distinguished derivation \(D_Xt=-t^2\) and \(D_XL=t\).
The letters \(x,\ell,\lambda\) used in individual reports are local
renamings of \(X,L\), not additional global generators.}
\entryfield{Polynomial-log power-series ring}{Define
\[
 \PolynomialLogPowerSeriesRing{K}
 \DefinitionEquals K[L][[t]].
\]
An element is a unique series
\[
 F=\sum_{n\ge0}p_n(L)t^n,
 \qquad p_n\in K[L].
\]
Each coefficient block is a finite polynomial in \(L\); no uniform bound on
\(\deg_Lp_n\) is implicit.  The ideal of strictly positive outer order is
\(tK[L][[t]]\).}
\entryfield{Polynomial-log Laurent ring}{Define
\[
 \PolynomialLogLaurentRing{K}
 \DefinitionEquals K[L]((t))
 =\bigcup_{m\ge0}t^{-m}K[L][[t]].
\]
Every element has a unique representation
\[
 F=\sum_{n\ge n_0}p_n(L)t^n,
 \qquad n_0\in\IntegerNumbers,\quad p_n\in K[L].
\]
Only finitely many negative powers of \(t\) may occur.  Equivalently, only
finitely many positive powers of \(X\) occur.  This ring is an integral domain
but, in general, not a field.}
\entryfield{Unit criterion}{For nonzero
\[
 F=t^{n_0}\bigl(p_{n_0}(L)+p_{n_0+1}(L)t+\cdots\bigr),
\]
\(F\) is a unit of \(\PolynomialLogLaurentRing{K}\) exactly when
\(p_{n_0}(L)\in K^\times\), that is, when its leading coefficient block is a
nonzero scalar.  A nonzero polynomial such as \(L+1\) is not a unit in
\(K[L]\).  Special divisibility may keep a particular quotient in the
polynomial-log ring, but there is no unconditional division theorem by an
arbitrary nonzero leading polynomial.}
\entryfield{Rational-log Laurent field}{Unrestricted exact division of
nonzero coefficient blocks is performed in
\[
 \RationalLogLaurentField{K}
 \DefinitionEquals K(L)((t)).
\]
This field retains rational dependence on \(L\) exactly.  A scalar such as a
chosen value of \(\log c\) must belong to \(K\); otherwise \(K\) must first be
extended.}
\entryfield{Iterated Laurent-log field}{Expanding rational functions at
\(L=\infty\) gives
\[
 \IteratedLaurentLogField{K}
 \DefinitionEquals K((L^{-1}))((t)).
\]
An element has the form
\[
 \sum_{n\ge n_0}t^n
   \sum_{j\ge j_n}a_{n,j}L^{-j}.
\]
For each fixed outer level \(n\), the inner Laurent series contains only
finitely many positive powers of \(L\), but it may contain infinitely many
negative powers.  The iteration order is part of the type:
\[
 K((L^{-1}))((t))\ne K((t))((L^{-1}))
\]
in general, because the two fields impose different support conditions.
The embedding
\(K(L)((t))\hookrightarrow K((L^{-1}))((t))\) means Laurent expansion of
each rational coefficient at \(L=\infty\); it changes an exact coefficient
such as \(1/(L+1)\) into
\(L^{-1}-L^{-2}+L^{-3}-\cdots\).}
\entryfield{Completion order and finite data}{A finite representation in the
iterated field needs an outer \(t\)-cutoff and an inner \(L^{-1}\)-cutoff, or
one explicitly defined staircase in the lexicographic monomial order.  An
outer cutoff alone does not make an infinite Laurent-log coefficient block
finite.}
\entryfield{Puiseux-log extension}{For a fixed
\(s\in\PositiveIntegers\),
\[
 K[L]\bigl((t^{1/s})\bigr)
\]
allows exponents in \(s^{-1}\IntegerNumbers\).  Formally,
\[
 D_X\!\left(t^\alpha p(L)\right)
 =t^{\alpha+1}\bigl(p'(L)-\alpha p(L)\bigr),
 \qquad \alpha\in s^{-1}\IntegerNumbers.
\]
An analytic realization additionally requires a branch of \(X^{1/s}\).
The union over all \(s\) allows rational exponents but has no single
denominator bound.}
\entryfield{Hahn power-log extension}{Let \(\Gamma\) be an ordered additive
subgroup of \(\RealNumbers\).  A Hahn power-log series has schematic form
\[
 \sum_{\alpha\in S}p_\alpha(L)t^\alpha,
 \qquad p_\alpha\in K[L],
\]
where \(S\subseteq\Gamma\) is well ordered in the direction of increasing
smallness.  The support condition is what makes multiplication locally finite.
If complex exponents are admitted, ordering merely by real part is
insufficient: a tie-breaking order and a support convention must be stated.}
\entryfield{Near-identity group}{Define
\[
 \NearIdentityPolynomialLogGroup{K}
 \DefinitionEquals
 \SetOf{X+A:A\in\PolynomialLogPowerSeriesRing{K}}.
\]
Under the admissibility hypotheses in
\notelink{polynomial-logarithmic-composition}{the composition entry}, this is
a group under composition.  It is not a multiplicative group of coefficient
series.}
\end{catalogueentry}

\begin{catalogueentry}{polynomial-logarithmic-leading-data}
  {Coefficient support, dominance, valuation, and leading data}
\entryfield{Coefficient support command}{The shared command is
\commandname{CoefficientSupportOf}.  It is distinct from
\commandname{SupportOf}, which is reserved for pointwise function support.
For
\[
 F=\sum_{n\ge n_0}p_n(L)t^n
  =\sum_{\substack{n\in\IntegerNumbers\\j\in\NaturalNumbers}}
     a_{n,j}t^nL^j,
\]
define
\[
 \CoefficientSupportOf{F}
 \DefinitionEquals
 \SetOf{(n,j)\in\IntegerNumbers\times\NaturalNumbers:a_{n,j}\ne0}.
\]
Every fiber over a fixed \(n\) is finite, and the projection to the
\(n\)-coordinate is bounded below.  The zero series has empty coefficient
support.  This is support in an exponent-index set, not the pointwise support
of a function on a topological space.}
\entryfield{Dominance order}{On the positive ray as \(X\to+\infty\),
\[
 t^nL^j\succ t^mL^k
 \quad\Longleftrightarrow\quad
 n<m
 \quad\text{or}\quad
 (n=m\ \text{and}\ j>k).
\]
Thus outer inverse-power order is compared first; logarithmic degree refines
only a fixed outer block.  Negative powers of \(L\) belong to an enlarged
Laurent-log coefficient field, not to \(K[L]((t))\).}
\entryfield{Outer valuation}{For nonzero
\(F=\sum_{n\ge n_0}p_n(L)t^n\), define
\[
 \nu_t(F)\DefinitionEquals
 \min\SetOf{n\in\IntegerNumbers:p_n\ne0},
 \qquad
 \nu_t(0)\DefinitionEquals+\infty.
\]
For nonzero inputs,
\[
 \nu_t(FG)=\nu_t(F)+\nu_t(G),
 \qquad
 \nu_t(F+G)\ge
 \min\SetOf{\nu_t(F),\nu_t(G)}.
\]
The second inequality may be strict when the leading blocks cancel.}
\entryfield{Leading block, monomial, and scalar}{Let
\[
 p_{\nu_t(F)}(L)=aL^d+\text{lower powers},
 \qquad a\in K^\times.
\]
Then the leading coefficient block is the whole polynomial
\(p_{\nu_t(F)}\); the leading monomial is
\(t^{\nu_t(F)}L^d\); and the scalar leading coefficient is \(a\).
These are three different objects.  For \(F=0\), the degree, leading monomial,
and leading coefficient are undefined unless a separate extended convention
is declared.}
\entryfield{Rank-two valuation}{The refined monomial valuation is
\[
 v(t^nL^j)=(n,-j)
\]
with lexicographic order.  Hence for nonzero \(F\),
\[
 v(F)=\bigl(\nu_t(F),-\deg_Lp_{\nu_t(F)}\bigr).
\]
The negative sign makes larger logarithmic degree more dominant within a fixed
outer block.}
\entryfield{Log-degree profile}{Extend absent coefficient blocks by zero and
define
\[
 d_F(n)=
 \begin{cases}
   \deg_Lp_n,&p_n\ne0,\\
   -\infty,&p_n=0.
 \end{cases}
\]
The profile \(d_F\) is additional data; it is not recoverable from the coarse
number \(\nu_t(F)\).  Linear bounds such as \(d_F(n)\le an+b\) are hypotheses
or proved invariants, never part of the ring definition.}
\entryfield{Hahn support}{For a Hahn extension the support is a subset of the
chosen exponent group and must be well ordered.  Multiplication uses
Minkowski addition of supports.  Well ordering or the relevant local-finiteness
condition ensures that every output coefficient receives only finitely many
contributions.}
\end{catalogueentry}

\begin{catalogueentry}{polynomial-logarithmic-truncation}
  {Exclusive truncation, formal tails, analytic remainders, and summability}
\entryfield{Exclusive outer truncation}{For \(N\in\IntegerNumbers\), define
\[
 \operatorname{Trunc}_{<N}F
 \DefinitionEquals
 \sum_{n<N}p_n(L)t^n.
\]
It keeps every logarithmic monomial in every retained coefficient block.  If
\(N\le n_0\), the sum is empty and equals \(0\).  The zero series also truncates
to \(0\).  “Through outer order \(t^{N-1}\)” is the exclusive cutoff
\(<N\).}
\entryfield{Inclusive outer truncation}{The report-6 local notation
\[
 [F]_{\le N}\DefinitionEquals\sum_{n\le N}p_n(L)t^n
\]
is inclusive.  Therefore
\[
 [F]_{\le N}=\operatorname{Trunc}_{<N+1}F.
\]
Any transfer between reports must account for this one-block difference.}
\entryfield{Formal tail command}{The command
\commandname{FormalBigOAt} is defined by
\[
 F=G+\FormalBigOAt{t^N}{t}
 \quad\Longleftrightarrow\quad
 \nu_t(F-G)\ge N
 \quad\Longleftrightarrow\quad
 F-G\in t^NK[L][[t]].
\]
It asserts coefficient agreement below outer order \(N\).  It asserts no
numerical bound, convergence, uniformity in \(L\), or analytic realization.
The token \commandname{FormalBigO} is the same operator without supplied
arguments.}
\entryfield{Why the Laurent tail is not an ideal}{In
\(\PolynomialLogLaurentRing{K}=K[L]((t))\), \(t\) is a unit.  Therefore the
principal ideal generated by \(t^N\) is the whole Laurent ring.  The subset
\(t^NK[L][[t]]\) is an additive valuation-filtration tail, not an ideal under
arbitrary Laurent multiplication.  Consequently:
\begin{itemize}
\item in \(K[L][[t]]\), reduction modulo the ideal \((t^N)\) is a genuine ring
quotient;
\item in \(K[L]((t))\), use \(\nu_t(F-G)\ge N\),
\(\FormalBigOAt{t^N}{t}\), or equality of exclusive truncations;
\item quotient-algebra language becomes legitimate after removing the finite
Laurent principal part and restricting to a power-series tail ring.
\end{itemize}}
\entryfield{Coefficient extraction}{For the double expansion above,
\[
 [t^n]F=p_n(L),
 \qquad
 [t^nL^j]F=a_{n,j}.
\]
An absent coefficient is \(0\).  Bracket argument order is not uniform in
the source reports: one local macro prints \([M]F\), another prints
\([F]_M\).  Shared prose must state the series and monomial roles explicitly.}
\entryfield{Separate logarithmic cutoff}{Bounding \(\deg_Lp_n\) is an
additional approximation, not part of the outer \(t\)-adic topology.
Discarding a high power of \(L\) in a retained \(t\)-block requires an explicit
degree bound.  An iterated Laurent-log coefficient requires a second inner
cutoff or a defined two-dimensional staircase.}
\entryfield{Analytic big-O command}{For actual functions on a domain \(S\),
\[
 R(X)=\BigOAt{A(X)}{X\to\infty,\ X\in S}
\]
means that there exist \(C>0\) and a sufficiently late part of the stated
domain such that
\(\AbsoluteValue{R(X)}\le C\AbsoluteValue{A(X)}\) there.  In complex work the
sector, closed subsector if uniformity is intended, and every logarithmic or
fractional-power branch are part of the statement.  The limit must not be
left to context.}
\entryfield{Three different O-roles}{The corpus distinguishes:
\[
 \FormalBigOAt{t^N}{t}
 \quad\text{(formal valuation tail)},\qquad
 \BigOAt{A(X)}{X\to\infty}
 \quad\text{(analytic estimate)},\qquad
 \BigO(N^2)
 \quad\text{(algorithmic complexity)}.
\]
No one of these meanings may be inferred from the other two.}
\entryfield{Completeness and summability}{A sequence is \(t\)-adically Cauchy
when every coefficient block below every fixed cutoff eventually stabilizes.
A family \((S_r)_{r\ge0}\) is formally summable when
\(\nu_t(S_r)\to+\infty\); below any cutoff only finitely many members then
contribute.  This coefficientwise property is the formal substitute for
numerical convergence.  On a Laurent ring, a Cauchy sequence must also have a
common eventual lower support bound; otherwise it need not define a Laurent
series with a finite principal part.}
\end{catalogueentry}

\begin{catalogueentry}{polynomial-logarithmic-differential}
  {Large-variable differentiation and shifted logarithmic operators}
\entryfield{Distinguished derivation}{With \(t=X^{-1}\) and \(L=\log X\),
\[
 D_X=t\,\partial_L-t^2\,\partial_t.
\]
Here \(\partial_t\) and \(\partial_L\) are formal partial derivatives holding
the other generator fixed.  The formula is the unique derivation satisfying
\(D_Xt=-t^2\) and \(D_XL=t\).}
\entryfield{One coefficient block}{For
\(n\in\IntegerNumbers\) and \(p\in K[L]\),
\[
 D_X\bigl(t^np(L)\bigr)
 =t^{n+1}\bigl(\partial_L-n\bigr)p(L).
\]
The Euler derivation preserves outer order:
\[
 XD_X\bigl(t^np(L)\bigr)
 =t^n\bigl(\partial_L-n\bigr)p(L).
\]}
\entryfield{Repeated derivative operator}{For
\(k\in\NaturalNumbers\), define
\[
 \Delta_{n,k}
 \DefinitionEquals
 \prod_{j=0}^{k-1}(\partial_L-n-j),
 \qquad
 \Delta_{n,0}\DefinitionEquals1.
\]
The \(k=0\) product is empty and acts as the identity.  The factors commute
because each is a polynomial in \(\partial_L\).  Then
\[
 D_X^k\bigl(t^np(L)\bigr)
 =t^{n+k}\Delta_{n,k}p(L).
\]}
\entryfield{Valuation behavior}{For nonzero \(F\),
\[
 \nu_t(D_XF)\ge\nu_t(F)+1.
\]
Equality can fail through annihilation or cancellation.  If the leading term
is \(at^nL^d\), then:
\begin{itemize}
\item if \(n\ne0\), the derivative begins with
\(-na\,t^{n+1}L^d\);
\item if \(n=0\) and \(d>0\), it begins with
\(ad\,tL^{d-1}\);
\item if \(n=d=0\), the leading constant is annihilated and the next
nonconstant term decides the order.
\end{itemize}
In particular, \(D_XK=0\) and
\(D_X(K[L][[t]])\subseteq tK[L][[t]]\).}
\entryfield{Formal partial versus analytic derivative}{The formal
\(\partial_t\) used for coefficient extraction holds \(L\) fixed.  It is not
the analytic derivative with respect to \(t\) after substituting
\(L=-\log t\).  Likewise, a bare \(D\) may denote \(D_X\), \(\partial_L\), or
differentiation in a Lambert-polynomial dummy variable; the differentiated
variable must be printed.}
\entryfield{Antiderivative resonance}{Solving
\[
 (\partial_L-n)p(L)=q(L)
\]
for a polynomial \(p\) is triangular and unique when \(n\ne0\).  When \(n=0\)
it reduces to polynomial integration and introduces an arbitrary scalar
constant.  This is the formal source of the constant of integration and must
be handled separately at the resonant block.}
\end{catalogueentry}

\begin{catalogueentry}{polynomial-logarithmic-composition}
  {Admissible composition and near-identity substitution}
\entryfield{Monomial-leading inner argument}{Let
\[
 F(Y)=\sum_{n\ge a}p_n(\log Y)Y^{-n}
\]
and
\[
 G(X)=cX^\rho(1+U(X)),
 \qquad
 c\in K^\times,\quad
 \rho\in\PositiveIntegers,\quad
 U\in tK[L][[t]].
\]
Choose a value of \(\log c\) in the coefficient field and put
\(V=\log(1+U)\).  Then
\[
 (F\circ G)(X)
 =
 \sum_{n\ge a}
 p_n\!\left(\rho L+\log c+V\right)
 c^{-n}t^{\rho n}(1+U)^{-n}.
\]
Every coefficient below a fixed outer cutoff is a finite sum.  If
\(\rho=r/s>0\) is rational, the result belongs to a Puiseux-log extension
after adjoining \(t^{1/s}\); a general ordered real exponent requires a Hahn
exponent group.}
\entryfield{Generator substitution}{For
\[
 G=X+H=X(1+tH),
 \qquad H\in K[L][[t]],
\]
composition changes both generators:
\[
 t\circ G=\frac{t}{1+tH},
 \qquad
 L\circ G=L+\log(1+tH).
\]
The substitution \(t\mapsto t+H\) is incorrect because \(t\) denotes
\(X^{-1}\), not \(X\).  Updating \(L\) while leaving \(t\) fixed is also
incorrect.}
\entryfield{Exact formal Taylor formula}{For
\(F\in K[L]((t))\) and \(H\in K[L][[t]]\),
\[
 F\circ(X+H)
 =
 \sum_{k\ge0}\frac{H^k}{k!}D_X^kF.
\]
This is an exact identity of formal series, not an analytic approximation.
If \(a=\nu_t(F)\), then only finitely many \(k\), in particular
\(k\le N-a\), can affect the coefficient of \(t^N\).}
\entryfield{Finite coefficient formula}{Write
\[
 F=\sum_{n\ge a}p_n(L)t^n,
 \qquad
 H^k=\sum_{m\ge0}h_{k,m}(L)t^m,
\]
with \(h_{0,0}=1\) and \(h_{0,m}=0\) for \(m>0\).  Then
\[
 [t^r]\bigl(F\circ(X+H)\bigr)
 =
 \sum_{\substack{n\ge a,\ k,m\ge0\\n+k+m=r}}
 \frac{h_{k,m}(L)}{k!}\Delta_{n,k}p_n(L).
\]
The displayed affine constraint makes the sum finite.}
\entryfield{Convolution-polynomial boundary cases}{If
\(H=\sum_{j\ge0}h_jt^j\), define
\[
 C_{m,r}
 \DefinitionEquals
 \sum_{\substack{j_1,\ldots,j_r\ge0\\j_1+\cdots+j_r=m}}
 h_{j_1}\cdots h_{j_r}.
\]
The empty-product conventions are
\[
 C_{0,0}=1,\qquad C_{m,0}=0\quad(m>0).
\]
Equivalently, \(H^r=\sum_{m\ge0}C_{m,r}t^m\), and
\[
 C_{m,r+1}=\sum_{j=0}^{m}h_jC_{m-j,r}.
\]
If \(h_0=0\), then \(C_{m,r}=0\) whenever \(m<r\).}
\entryfield{Associativity domain}{The identity
\((F\circ G)\circ H=F\circ(G\circ H)\) is asserted only when every displayed
composition exists in the selected support class.  For substitution maps,
right substitution reverses the apparent order of operator composition.
Formal closure is a theorem about support and local finiteness, not a license
to compose arbitrary transseries.}
\entryfield{Analytic admissibility}{For analytic realizations, the image of
the inner map must eventually lie in the outer expansion's ray or sector,
must avoid every relevant logarithmic branch cut, and must carry the omitted
outer remainder below the requested output order.  These conditions are
additional to the formal coefficient calculation.}
\end{catalogueentry}

\begin{catalogueentry}{polynomial-logarithmic-reversion}
  {Compositional inverse, multiplicative reciprocal, and residual certificates}
\entryfield{Commands}{The shared commands are
\commandname{MultiplicativeInverseOf} and
\commandname{CompositionalInverseOf}.  They render
\[
 \MultiplicativeInverseOf{F}=F^{-1},
 \qquad
 \CompositionalInverseOf{F}
 =F^{\langle-1\rangle_\circ}.
\]
The first is inverse under multiplication and satisfies
\(F\MultiplicativeInverseOf{F}=1\).  The second is inverse under composition
and satisfies
\[
 F\circ\CompositionalInverseOf{F}
 =
 \CompositionalInverseOf{F}\circ F
 =X.
\]
Negative monomial powers, the branch label \(W_{-1}\), inverse transforms, and
time-\((-1)\) iterates are not instances of either replacement merely because
they contain the glyph \(-1\).}
\entryfield{Near-identity inverse}{For
\(F=X+A\in\NearIdentityPolynomialLogGroup{K}\), there is a unique
\(\CompositionalInverseOf{F}=X+B\) in the same group.  The correction \(B\)
is the unique fixed point
\[
 B=-A\circ(X+B).
\]
Starting with \(B_0=0\), the error gains at least one outer block per iteration
because \(D_XA\in tK[L][[t]]\).}
\entryfield{Lagrange--Bürmann formula}{The formal inverse is
\[
 \CompositionalInverseOf{(X+A)}(X)
 =
 X+\sum_{r\ge1}\frac{(-1)^r}{r!}
 D_X^{\,r-1}\!\left(A(X)^r\right).
\]
The sum starts at \(r=1\), where the derivative order is zero.  The \(r\)-th
summand has outer order at least \(r-1\), which makes the family formally
summable.}
\entryfield{Coefficient formula}{If
\[
 A^r=\sum_{m\ge0}A_{r,m}(L)t^m,
 \qquad
 \CompositionalInverseOf{F}
 =X+\sum_{N\ge0}b_N(L)t^N,
\]
then
\[
 b_N(L)=
 \sum_{r=1}^{N+1}\frac{(-1)^r}{r!}
 \left[
   \prod_{s=N-r+1}^{N-1}(\partial_L-s)
 \right]
 A_{r,N-r+1}(L).
\]
When the lower product bound exceeds the upper bound, in particular at
\(r=1\), the product is empty and acts as the identity.}
\entryfield{Residual with all arguments explicit}{For an ambient map \(F\)
and a candidate inverse \(H\), define
\[
 \operatorname{Err}_{\mathrm{rev}}(F,H)
 \DefinitionEquals F\circ H-X.
\]
A report-local one-argument macro may hide \(F\) in a subscript; shared prose
must supply both \(F\) and \(H\) because the residual is not determined by
\(H\) alone.}
\entryfield{Finite certificate}{For
\[
 G_N=X+\sum_{j=0}^{N}b_j(L)t^j,
\]
\(G_N\) is a right inverse through outer block \(N\) exactly when
\[
 F\circ G_N-X=\FormalBigOAt{t^{N+1}}{t}.
\]
The left certificate is
\[
 G_N\circ F-X=\FormalBigOAt{t^{N+1}}{t}.
\]
An exact one-sided inverse in the group is two-sided by uniqueness.  At finite
precision both residuals are useful because they detect different
substitution-order mistakes.}
\entryfield{Triangular correction}{If
\[
 F\circ G_N-X
 =r_N(L)t^N+\FormalBigOAt{t^{N+1}}{t},
\]
then
\[
 G_{N+1}=G_N-r_N(L)t^N
\]
removes the residual at level \(N\) when the leading slope is \(1\).  For a
leading slope \(a\in K^\times\), divide the correction by \(a\).}
\entryfield{Newton reversion}{Newton's update is
\[
 G_{\mathrm{new}}
 =
 G-\frac{F\circ G-X}{(D_XF)\circ G}.
\]
For \(F=X+A\), the denominator is
\(1+\FormalBigOAt{t}{t}\) and is therefore a unit.  The residual order
approximately doubles when all intermediate arithmetic carries sufficient
guard precision.}
\entryfield{Analytic residual transfer}{If \(f\) is continuously
differentiable and one-to-one on an interval, \(x_*=f^{\langle-1\rangle_\circ}(y)\),
and \(\AbsoluteValue{f'(x)}\ge m>0\) between \(x_N\) and \(x_*\), then
\[
 \AbsoluteValue{x_N-x_*}
 \le\frac{\AbsoluteValue{f(x_N)-y}}{m}.
\]
This analytic theorem requires injectivity and a derivative lower bound; a
formal residual alone does not prove existence of an analytic inverse branch.}
\end{catalogueentry}

\begin{catalogueentry}{wright-omega-lambert-polynomials}
  {Wright omega, the canonical Lambert polynomials, and Stirling coefficients}
\entryfield{Commands}{The shared commands are
\commandname{WrightOmega},
\commandname{LambertPolynomial}, and
\commandname{ScaledLambertPolynomial}.  The polynomial command takes only an
index; its evaluation argument is written after the command.}
\entryfield{Real Wright omega}{The map
\(y\mapsto y+\log y\) is a bijection from \((0,\infty)\) to
\(\RealNumbers\).  Its real inverse is
\[
 \WrightOmega(x)+\log\WrightOmega(x)=x,
 \qquad
 \WrightOmega(x)=\LambertW_0(\EulerE^x).
\]
The real function is defined for every \(x\in\RealNumbers\); the
polynomial-log expansion below concerns \(x\to+\infty\).  Complex Wright
omega requires an unwinding-number or analytic-continuation convention.}
\entryfield{Canonical zero-based family}{For
\(n\in\NaturalNumbers\), define
\[
 \LambertPolynomial{n}(u)
 \DefinitionEquals
 \frac{(-1)^{n+1}}{(n+1)!}
 \left[
   \prod_{j=0}^{n-1}(\partial_u-j)
 \right]u^{n+1}.
\]
At \(n=0\) the product is empty, so
\[
 \LambertPolynomial{0}(u)=-u.
\]
For \(n\ge1\), the first values are
\[
 \LambertPolynomial{1}(u)=u,\qquad
 \LambertPolynomial{2}(u)=\frac{u^2}{2}-u,\qquad
 \LambertPolynomial{3}(u)=\frac{u^3}{3}-\frac32u^2+u.
\]}
\entryfield{Omega expansion}{With \(L=\log x\),
\[
 \WrightOmega(x)
 \sim
 x+\sum_{n\ge0}\frac{\LambertPolynomial{n}(L)}{x^n}
 =
 x-L+\sum_{n\ge1}\frac{\LambertPolynomial{n}(L)}{x^n}.
\]
The two displays are identical.  The former is the canonical zero-based
normalization; the latter merely separates the \(n=0\) logarithmic block.}
\entryfield{Recurrence and boundary values}{For \(n\in\NaturalNumbers\),
\[
 \LambertPolynomial{n+1}(u)
 =
 -\LambertPolynomial{n}(u)
 +n\int_0^u\LambertPolynomial{n}(v)\,\Differential v.
\]
Equivalently,
\[
 \frac{\Differential}{\Differential u}\LambertPolynomial{n+1}(u)
 =n\LambertPolynomial{n}(u)
  -\frac{\Differential}{\Differential u}\LambertPolynomial{n}(u),
 \qquad
 \LambertPolynomial{n}(0)=0.
\]
For \(n\ge1\),
\[
 \deg\LambertPolynomial{n}=n,\quad
 [u^n]\LambertPolynomial{n}=\frac1n,\quad
 \left.\frac{\Differential}{\Differential u}\LambertPolynomial{n}(u)
 \right|_{u=0}=(-1)^{n-1}.
\]
The statement \(\deg\LambertPolynomial{0}=0\) would be false because
\(\LambertPolynomial{0}=-u\) has degree \(1\); therefore the degree law is
intentionally restricted to \(n\ge1\).}
\entryfield{Signed and unsigned Stirling numbers}{Define signed Stirling
numbers of the first kind by
\[
 z^{\underline n}
 =z(z-1)\cdots(z-n+1)
 =\sum_{k=0}^{n}s(n,k)z^k.
\]
The \(n=0\) product is empty, so \(s(0,0)=1\).  Set
\(s(n,k)=0\) for \(k<0\) or \(k>n\).  Then
\[
 s(n+1,k)=s(n,k-1)-ns(n,k),
 \qquad
 \UnsignedStirlingFirstKind{n}{k}
 =(-1)^{n-k}s(n,k).
\]
The bracket glyph here denotes an unsigned cycle number only when its two
arguments are integer indices; it is neither coefficient extraction nor a
Gaussian binomial coefficient.}
\entryfield{Explicit coefficients}{For \(n\ge1\),
\[
 \LambertPolynomial{n}(u)
 =
 \sum_{m=1}^{n}
 (-1)^{n-m}
 \UnsignedStirlingFirstKind{n}{n-m+1}
 \frac{u^m}{m!}.
\]
If
\(\LambertPolynomial{n}(u)=\sum_{m=1}^{n}\lambda_{n,m}u^m\), then
\[
 \lambda_{n,m}
 =
 (-1)^{n-m}\frac1{m!}
 \UnsignedStirlingFirstKind{n}{n-m+1},
\]
with \(\lambda_{n,0}=\lambda_{n,n+1}=0\) and
\(\lambda_{1,1}=1\).}
\entryfield{Scaled family}{For \(n\ge1\),
\[
 \ScaledLambertPolynomial{n}(u)
 \DefinitionEquals n!\LambertPolynomial{n}(u)
 \in\IntegerNumbers[u].
\]
Its recurrence is
\[
 \ScaledLambertPolynomial{n+1}(u)
 =-(n+1)\ScaledLambertPolynomial{n}(u)
 +n(n+1)\int_0^u
   \ScaledLambertPolynomial{n}(v)\,\Differential v.
\]
The scaled family is a storage normalization; it does not define a different
asymptotic coefficient sequence.}
\entryfield{Generating equation}{If
\[
 Q(s,u)=\sum_{n\ge1}\LambertPolynomial{n}(u)s^n,
\]
then \(Q(0,u)=0\) and
\[
 Q=-\log(1-su+sQ).
\]
This equation determines the positive-index coefficient family uniquely in
the formal power-series ring \(K[u][[s]]\).}
\end{catalogueentry}

\begin{catalogueentry}{lambert-branch-expansions}
  {Lambert \(W\) branches, lower-branch scaling, and analytic status}
\entryfield{Large complex branch coordinates}{Fix an integer branch label
\(k\), a branch of the logarithm, and a sector on which
\[
 \xi_k\DefinitionEquals\PrincipalLogarithm z
      +2\pi\ImaginaryUnit k,
 \qquad
 \ell_k\DefinitionEquals
 \log_{\mathcal B_k}(\xi_k)
\]
are defined and \(\AbsoluteValue{\xi_k}\to\infty\).  Then the formal
large-branch expansion is
\[
 \LambertW_k(z)
 \sim
 \xi_k-\ell_k+
 \sum_{n\ge1}\frac{\LambertPolynomial{n}(\ell_k)}{\xi_k^n}
 =
 \xi_k+
 \sum_{n\ge0}\frac{\LambertPolynomial{n}(\ell_k)}{\xi_k^n}.
\]
The polynomial family is branch-independent; the coordinates
\(\xi_k,\ell_k\), their branches, and the analytic domain are not.}
\entryfield{Formal versus analytic assertion}{A finite polynomial sum in
\(\xi_k^{-1}\) is a formal truncation until an analytic theorem supplies an
estimate such as
\[
 \LambertW_k(z)-\left(
 \xi_k-\ell_k+
 \sum_{n=1}^{N}
 \frac{\LambertPolynomial{n}(\ell_k)}{\xi_k^n}
 \right)
 =
 \BigOAt{\ell_k^{N+1}/\xi_k^{N+1}}
        {\AbsoluteValue{\xi_k}\to\infty,\ \xi_k\in S_k}.
\]
The sector \(S_k\), distance from branch cuts, and compatibility of the two
logarithms are hypotheses.  A formal residual does not by itself establish a
uniform analytic constant.}
\entryfield{Residual correction}{For
\[
 w_N=\xi_k-\ell_k+
 \sum_{n=1}^{N}\frac{\LambertPolynomial{n}(\ell_k)}{\xi_k^n},
 \qquad
 R_N=w_N+\log_{\mathcal B_k}(w_N)-\xi_k,
\]
one Newton correction is
\[
 \delta_N=-\frac{R_N}{1+1/w_N}.
\]
Taylor's theorem transfers a sufficiently small residual to an inverse error
only in a region avoiding \(w=0\) and the chosen logarithmic cut.}
\entryfield{Principal real branch}{For \(z\to+\infty\), take
\(\xi_0=\log z>0\) and \(\ell_0=\log\log z\) with the real logarithm.  This
specialization yields the familiar principal-branch expansion of
\(\LambertW_0(z)\).}
\entryfield{Lower real branch near zero}{The report-6 local symbol
\[
 \varepsilon_{\mathrm{br}}>0,\qquad
 \varepsilon_{\mathrm{br}}\to0
\]
is a branch-point-independent small parameter for
\(\LambertW_{-1}(-\varepsilon_{\mathrm{br}})\).  Put
\[
 S=\log(1/\varepsilon_{\mathrm{br}}),\qquad L=\log S.
\]
Then the large variable is \(S\), and the expansion of the negative branch is
obtained from the inverse of \(u\mapsto u-\log u\).  Its alternating signs do
not define a second Lambert-polynomial family.}
\entryfield{Square-root branch point}{At \(z=-\EulerE^{-1}\), the derivative
of \(w\mapsto w\EulerE^w\) vanishes at \(w=-1\); the inverse is therefore not
an ordinary power-log series in \(z+\EulerE^{-1}\).  A local coordinate such
as
\[
 p=\sqrt{2(1+\EulerE z)}
\]
produces a Puiseux series with the sign distinguishing the two local branches.
The square-root branch and continuation path must be stated.}
\entryfield{Convergence boundary}{The polynomial recurrence and every
finite residual identity are formal algebra.  Convergence or asymptotic
validity of the infinite series is a separate analytic claim.  A statement of
analytic validity must name the branch, limiting sector or ray, and remainder
estimate; the notation catalogue does not promote an unverified source claim
to a theorem.}
\end{catalogueentry}

\begin{catalogueentry}{polynomial-logarithmic-local-roles}
  {Precision records, epsilon roles, and collision policy in the six reports}
\entryfield{Approximation marker}{The report-6 local command
\(\commandname{ApproximationOf}\{F\}\) prints
\[
 F_{\mathrm{approx}}.
\]
It denotes a declared approximate or truncated representative of \(F\); it
does not specify the cutoff, error norm, or formal valuation.  Those data must
be supplied separately.  It is not a hat accent and is not synonymous with
an estimator unless a statistical model is declared.}
\entryfield{Inclusive truncation marker}{The same report's local
\(\commandname{trunc}\{F\}\{N\}\) prints \([F]_{\le N}\) and is inclusive.
Other reports use exclusive \(\operatorname{Trunc}_{<N}F\).  The relation is
\([F]_{\le N}=\operatorname{Trunc}_{<N+1}F\).}
\entryfield{Nested inverse-log scale}{For
\[
 L_0=X,\qquad L_{j+1}=\log L_j,
\]
the report-6 local symbol
\[
 \varepsilon_{\mathrm{log},j}
 \DefinitionEquals L_j^{-1}
\]
is the reciprocal of the \(j\)-th iterated logarithm.  Its subscript
\(\mathrm{log},j\) is semantic and must not be dropped when other epsilon
roles coexist.}
\entryfield{Lower-branch epsilon}{The symbol
\(\varepsilon_{\mathrm{br}}\) is the positive small argument in
\(\LambertW_{-1}(-\varepsilon_{\mathrm{br}})\), with
\(\varepsilon_{\mathrm{br}}\to0^+\).  It is not an error tolerance.}
\entryfield{Newton residual scale}{The symbol
\(\varepsilon_{\mathrm{Newt}}\) denotes a local residual scale in an analytic
Newton estimate.  Under a nonvanishing derivative, one normalized Newton step
changes a first-order residual scale to a quadratic scale
\(\varepsilon_{\mathrm{Newt}}^2\).  This is neither the branch parameter nor
an inverse-log generator.}
\entryfield{Coefficient bracket collision}{The glyph \([\,\cdot\,]\) has
four unrelated roles in the broader corpus:
\[
 [t^n]F,\qquad [F]_{\le N},\qquad
 \UnsignedStirlingFirstKind{n}{k},\qquad
 \GaussianBinomial{n}{k}{q}.
\]
They mean coefficient extraction, inclusive truncation, an unsigned Stirling
cycle number, and a Gaussian binomial coefficient, respectively.  Argument
types and semantic macros disambiguate them.}
\entryfield{Letter \(L\) collision}{In this namespace \(L=\log X\) is a
logarithmic coordinate.  The Lambert family is always written
\(\LambertPolynomial{n}\), never as a bare \(L_n\).  A local pair
\(L_1=\log z\), \(L_2=\log L_1\) denotes iterated logarithms, not the first two
Lambert polynomials.}
\entryfield{Generic versus Lambert coefficients}{A generic coefficient block
is \(p_n(L)\).  The distinguished Lambert coefficient is
\(\LambertPolynomial{n}(L)\).  The bare family \(P_n\) is not globally
reserved because several source reports used it for both roles.}
\entryfield{Outer \(t\) versus q-series base}{The polynomial-logarithmic outer
coordinate is \(t=X^{-1}\).  A free \(q\) is reserved for a genuine q-series
base when q-series notation is present.  A bound summation index, rational
exponent, or locally declared scalar may still be named \(q\), but it acquires
no q-series meaning from the letter alone.}
\entryfield{Inverse collision}{A reciprocal is
\(\MultiplicativeInverseOf{F}\); a compositional inverse is
\(\CompositionalInverseOf{F}\).  The glyph \(F^{[-1]}\) may denote a
time-\((-1)\) iterate only in an explicitly declared iteration context and
must not be silently identified with reversion.}
\entryfield{Derivative collision}{Write \(D_X\), \(\partial_L\),
\(\partial_t\), or \(\partial_u\) according to the differentiated variable.
The formal partial \(\partial_t\) holds \(L\) fixed; it is not the analytic
derivative after \(L=-\log t\).}
\entryfield{Class-name collision}{The six reports formerly assigned the same
bare calligraphic letters to different rings.  Shared text uses the exact
semantic commands:
\[
\begin{array}{c|c}
\text{command}&\text{defined class}\\
\hline
\PolynomialLogPowerSeriesRing{K}&K[L][[t]]\\
\PolynomialLogLaurentRing{K}&K[L]((t))\\
\RationalLogLaurentField{K}&K(L)((t))\\
\IteratedLaurentLogField{K}&K((L^{-1}))((t))\\
\NearIdentityPolynomialLogGroup{K}&
 \SetOf{X+A:A\in K[L][[t]]}.
\end{array}
\]
A document-local shorthand may expand to one of these commands, but the bare
glyph has no corpus-wide type.}
\entryfield{Sparse implementation record}{A finite one-log representative is
a sparse map \(n\mapsto p_n(L)\) together with:
\begin{itemize}
\item the smallest and largest retained outer exponents;
\item the coefficient field and logarithmic-degree policy;
\item any inner cutoff for Laurent-log coefficients;
\item branch and exponent-group metadata;
\item the declared formal tail or analytic error certificate.
\end{itemize}
No finite record stores the infinite object itself.}
\end{catalogueentry}

\begin{catalogueentry}{lagrange-cardinal}{Lagrange and cardinal interpolation}
\entryfield{Nodes}{Let $x_0,\ldots,x_N$ be pairwise distinct scalars.  The cardinal
polynomial at node $j$ is
\[
 \ell_j(x)=\prod_{\substack{0\le k\le N\\k\ne j}}
 \frac{x-x_k}{x_j-x_k},
 \qquad \ell_j(x_k)=\delta_{jk}.
\]}
\entryfield{Interpolant}{$I_Nf(x)=\sum_{j=0}^{N}f(x_j)\ell_j(x)$ is the unique polynomial
of degree at most $N$ matching $f$ at all nodes.}
\entryfield{Empty case}{For $N=0$, the product is empty and $\ell_0=1$.
There is no $N=-1$ interpolant unless a separate empty-data convention is introduced.}
\entryfield{Barycentric weights}{$w_j=(\prod_{k\ne j}(x_j-x_k))^{-1}$ and
$I_Nf(x)=
(\sum_jw_jf(x_j)/(x-x_j))/(\sum_jw_j/(x-x_j))$ away from nodes, with the value at a
node assigned by continuity.  Scaling every $w_j$ by the same nonzero constant changes
neither interpolant.}
\entryfield{Cardinal infinite grids}{An infinite cardinal function $L$ satisfying
$L(k)=\delta_{0k}$ for $k\in\IntegerNumbers$ is a different object and requires a
function space and convergence mode for $\sum_kf(k)L(x-k)$.}
\entryfield{Finite Rvachev decoder}{Let \(s\) be a finite index set,
\(v:s\to\RealNumbers\) a node map, and \(\ell_i\) the Lagrange basis polynomial for
\(i\in s\).  With \(\mathscr D_{\RvachevUp}\) the Rvachev deconvolution operator,
\(M\in\PositiveIntegers\), and \(k\in\IntegerNumbers\), define
\[
 D_{k i}^{s,v,M}
 \DefinitionEquals
 (\mathscr D_{\RvachevUp}\ell_i)(k/M).
\]
For nodal data \(y:s\to\RealNumbers\), the corresponding atom coefficient is
\[
 a_k^{s,v,y,M}
 \DefinitionEquals
 \frac1M\sum_{i\in s}D_{k i}^{s,v,M}y_i.
\]
The superscripts display the finite node family, node locations, data, and mesh; a source
that suppresses them must declare those objects in the immediately governing scope.}
\entryfield{Exact finite synthesis}{For a bounded Fabius object \(F\) satisfying the
Fabius specification, \(i\in s\), \(M\ne0\),
\(|s|-1\le\TwoAdicValuation(M)\), and \(x\in[-1,1]\),
\[
 \frac1M\sum_{-2M<k<2M}
 D_{k i}^{s,v,M}\,\RvachevUp_F(x-k/M)=\ell_i(x).
\]
The index set is the open integer interval \(-2M<k<2M\), including neither endpoint.
Summing against \(y_i\) reconstructs the complete Lagrange interpolant from the
coefficients \(a_k^{s,v,y,M}\).  Node injectivity is unnecessary for the polynomial
reconstruction identity but is required for the Kronecker-delta cardinal interpretation;
the unit decoder-row-sum theorem additionally assumes \(s\ne\varnothing\).}
\entryfield{Lean crosswalk}{\lean{lagrangeRvachevDecoder} and
\lean{lagrangeRvachevAtomCoefficient} in \modulepath{LagrangeRvachevSynthesis.lean}
implement the two displayed definitions.  The exact reconstruction, node-evaluation,
coefficient/deconvolution, full-interpolation, and row-sum results are the named theorems
\leanline{normalized_sum_Ioo_lagrangeRvachevDecoder_mul_shifted_rvachevUp}
\leanline{normalized_sum_Ioo_lagrangeRvachevDecoder_eval_node}
\leanline{lagrangeRvachevAtomCoefficient_eq_deconvolved_interpolate}
\leanline{sum_Ioo_lagrangeRvachevAtomCoefficient_mul_shifted_rvachevUp}
\leanline{sum_lagrangeRvachevDecoder_eq_one}
Lean's definitions remain total at
\(M=0\) through field division; the reconstruction theorem requires \(M\ne0\).  No
optimized decoder, bundled matrix inverse, Gaussian \(q\)-binomial closed form, or
minimum-variation solution is supplied.  Status: \textbf{Lean-proved with the displayed
hypotheses}.}
\end{catalogueentry}

\begin{catalogueentry}{sampling}{Sampling, reconstruction, and band limitation}
\entryfield{Band limitation}{Under the $2\pi$ Fourier convention, $f$ is bandlimited to
$[-B,B]$ when $\SupportOf{\FourierTwoPi f}\subseteq[-B,B]$.}
\entryfield{Critical lattice}{The Shannon lattice spacing is $(2B)^{-1}$ for $B>0$.
With suitable $L^2$ or regularity hypotheses,
\[
 f(x)=\sum_{k\in\IntegerNumbers}
 f\!\left(\frac{k}{2B}\right)
 \SincPi(2Bx-k).
\]
The theorem states whether convergence is in $L^2$, uniform, or pointwise.}
\entryfield{Oversampling}{A denser lattice adds frame redundancy; it does not change the
definition of \commandname{SincPi}.}
\entryfield{Fabius caveat}{Compact support in physical space makes the Fourier transform
entire, not bandlimited.  Sampling formulas in the inverse-Fabius frontier therefore
require an independently proved interpolation or frame theorem.}
\entryfield{Alias policy}{A sample index $k$, a frequency index, and a Lambert branch
index may all be integers; each sum declares which role its index has.}
\end{catalogueentry}

\begin{catalogueentry}{refinement-dilation}{Refinement equations and dilation operators}
\entryfield{Dilation}{For nonzero real $a$, $(D_af)(x)=f(ax)$.  A unitary
$L^2$ dilation is instead $(U_af)(x)=|a|^{1/2}f(ax)$; the two are not conflated.}
\entryfield{Translation}{$T_bf(x)=f(x-b)$.  Thus $D_aT_b=T_{b/a}D_a$ for
$a\ne0$ under this convention.}
\entryfield{Refinement equation}{A refinable function satisfies
\[
 \phi(x)=\sum_{k\in\IntegerNumbers}h_k\phi(ax-k)
\]
with dilation $|a|>1$, a declared mask $(h_k)$, and a stated convergence/pointwise
interpretation.  A factor such as $\sqrt{|a|}$ or $|a|$ is part of the mask
normalization and is never implicit.}
\entryfield{Mask symbol}{$H(z)=\sum_kh_kz^k$ is a Laurent polynomial or series according
to the mask support.  Its Fourier-domain refinement law depends on the selected Fourier
normalization.}
\entryfield{Fabius scaling}{The equation
$\FabiusGlobal'(x)=2\FabiusGlobal(2x)$ is a differential scaling law, not by itself a
standard refinability equation for $\FabiusGlobal$.}
\end{catalogueentry}

\begin{catalogueentry}{iterates-dynamics}{Function iterates, orbits, and conjugacy}
\entryfield{Iterate}{For a self-map $f:X\to X$,
$f^{\circ0}=\IdentityOperator_X$ and
$f^{\circ(n+1)}=f\circ f^{\circ n}$.  The circle distinguishes iteration from pointwise
power $f(x)^n$.}
\entryfield{Orbit}{The forward orbit of $x$ is
$\mathcal O_f^+(x)=\{f^{\circ n}(x):n\in\NaturalNumbers\}$.}
\entryfield{Periodic point}{$x$ has least period $p\in\PositiveIntegers$ when
$f^{\circ p}(x)=x$ and no smaller positive exponent does.  A fixed point has period one.}
\entryfield{Conjugacy}{Maps $f:X\to X$ and $g:Y\to Y$ are conjugate through a bijection
$h:X\to Y$ when $h\circ f=g\circ h$.  If $h$ is only a surjection, the relation is a
semiconjugacy.}
\entryfield{Fabius inverse}{Iterations of \commandname{FabiusQuantile} stay in $[0,1]$.
Iterations of the clamped inverse on all reals are a different dynamical system and use
\commandname{FabiusClampedQuantile} explicitly.}
\end{catalogueentry}

\begin{catalogueentry}{koopman-transfer}{Koopman and transfer operators}
\entryfield{Koopman}{For a measurable map $T:X\to X$,
$(\mathcal K_Tf)(x)=f(Tx)$.  Its function space and measure are part of the operator
type.}
\entryfield{Perron--Frobenius}{A transfer operator $\mathcal P_T$ is defined by the
duality
\[
 \int(\mathcal P_Tf)g\,\Differential\mu
 =\int f(g\circ T)\,\Differential\mu
\]
when such a density operator exists.  It is not automatically pointwise inverse
composition.}
\entryfield{Invariant measure}{$\mu$ is $T$-invariant when
$T_\#\mu=\mu$, equivalently $\int f\circ T\,\Differential\mu
=\int f\,\Differential\mu$ for bounded measurable $f$.}
\entryfield{Spectrum}{A Koopman eigenfunction satisfies
$\mathcal K_Tf=\lambda f$.  Spectral claims state the Banach or Hilbert space because
the spectrum changes with it.}
\entryfield{Status}{Fabius Koopman models in the frontier corpus are
\textbf{frontier/conjectural} unless an operator, invariant measure, and proved spectral
statement are all linked.}
\end{catalogueentry}

\begin{catalogueentry}{stein-zero-bias}{Stein operators and zero-bias transforms}
\entryfield{Stein operator}{An operator $\mathcal A$ characterizes a law $\mu$ on a test
class $\mathcal T$ when
$\int\mathcal Af\,\Differential\mu=0$ for every $f\in\mathcal T$, together with a
uniqueness theorem.  The operator alone is not a characterization without the class.}
\entryfield{Stein equation}{For a test $h$,
$\mathcal Af_h=h-\int h\,\Differential\mu$.  A bound on $f_h$ specifies its norm and
the class of $h$.}
\entryfield{Zero bias}{For mean-zero $X$ with variance $\sigma^2>0$, a zero-biased
variable $X^\ast$ satisfies
\[
 \Expectation[Xf(X)]=\sigma^2\Expectation[f'(X^\ast)]
\]
for an appropriate class of absolutely continuous $f$.}
\entryfield{Iteration}{A zero-bias tower declares how centering and variance
renormalization are applied at every stage; the star is a distributional transform, not
complex conjugation.}
\entryfield{Status}{Proposed Fabius Stein characterizations and zero-bias towers retain
\textbf{frontier/conjectural} status until uniqueness and domain conditions are proved.}
\end{catalogueentry}

\begin{catalogueentry}{carleman}{Carleman linearization and Carleman matrices}
\entryfield{Analytic map}{Let $f(z)=\sum_{n\ge0}a_nz^n$ near a fixed expansion point.
The Carleman matrix is
\[
 C(f)_{mn}\DefinitionEquals[z^n]f(z)^m,\qquad
 m,n\in\NaturalNumbers.
\]}
\entryfield{Zeroth row}{Since $f(z)^0=1$, $C(f)_{0n}=\delta_{0n}$, including the
convention $0^0=1$ inside the polynomial identity.}
\entryfield{Composition}{Formally, $C(f\circ g)=C(f)C(g)$ when both sides are defined
with compatible expansion points and the infinite matrix product is coefficientwise
finite or convergent.}
\entryfield{Truncation}{$C_N(f)=(C(f)_{mn})_{0\le m,n<N}$ is a finite truncation, not
the exact Carleman operator.  Spectral convergence of truncations requires a theorem.}
\entryfield{Collision}{This Carleman matrix is distinct from the Carleman condition
$\sum_nm_{2n}^{-1/(2n)}=\infty$ for moment determinacy.}
\end{catalogueentry}

\begin{catalogueentry}{holonomic}{Holonomic and \(D\)-finite functions}
\entryfield{\(D\)-finite}{A formal power series $f(z)$ over a characteristic-zero field is
$D$-finite when it satisfies
\[
 p_r(z)f^{(r)}(z)+\cdots+p_0(z)f(z)=0
\]
for polynomials $p_j$, not all zero.}
\entryfield{\(P\)-recursive}{A sequence $(a_n)$ is $P$-recursive when
$\sum_{j=0}^rp_j(n)a_{n+j}=0$ for all sufficiently large $n$, with polynomial
coefficients $p_j(n)$.}
\entryfield{Equivalence}{Over characteristic zero, an ordinary generating function is
$D$-finite exactly when its coefficient sequence is $P$-recursive, subject to the
standard formal-series interpretation.}
\entryfield{Nonholonomicity}{Failure of one guessed differential equation does not prove
nonholonomicity.  A nonholonomicity claim identifies a rigorous obstruction such as
infinitely many singularities or incompatible coefficient asymptotics.}
\entryfield{Status}{Fabius holonomicity questions in the active frontier reports are
\textbf{open} unless a theorem is explicitly cited.}
\end{catalogueentry}

\begin{catalogueentry}{chaos-cumulants}{Chaos expansions, cumulants, and Edgeworth terms}
\entryfield{Orthogonal chaos}{A chaos decomposition
$L^2(\Omega)=\bigoplus_{n\ge0}\mathcal H_n$ is relative to a specified probability
structure.  $J_n$ denotes the orthogonal projection onto $\mathcal H_n$ only within that
scope.}
\entryfield{Cumulant generating function}{For an MGF nonzero near zero,
\[
 K_X(t)=\log\MomentGeneratingFunction X(t)
 =\sum_{r\ge1}\kappa_r(X)\frac{t^r}{r!},
 \qquad K_X(0)=0,
\]
where the analytic logarithm is the branch normalized by the displayed value at zero.}
\entryfield{Standardization}{For variance $\sigma^2>0$,
$Z=(X-\Expectation X)/\sigma$.  Standardized cumulants are
$\lambda_r=\kappa_r(X)/\sigma^r$ for $r\ge3$.}
\entryfield{Edgeworth expansion}{An order-$s$ Edgeworth approximation identifies the
normal density or CDF, Hermite-polynomial normalization, scaling parameter, and a uniform
remainder regime.  The formal polynomial correction alone is not a probability density
theorem.}
\entryfield{Hermite choice}{Probabilists' Hermite polynomials satisfy
$\operatorname{He}_n(x)=(-1)^n\EulerE^{x^2/2}
\frac{\Differential^n}{\Differential x^n}\EulerE^{-x^2/2}$.
Physicists' polynomials are separately labelled $H_n^{\mathrm{phys}}$.}
\end{catalogueentry}

\begin{catalogueentry}{large-deviations}{Large deviations and rate functions}
\entryfield{Log-MGF}{$\Lambda(t)=\log\Expectation[\EulerE^{tX}]$ on its effective domain
$\mathcal D_\Lambda=\{t:\MomentGeneratingFunction X(t)<\infty\}$.}
\entryfield{Legendre--Fenchel transform}{$\Lambda^\ast(x)=
\sup_{t\in\RealNumbers}\{tx-\Lambda(t)\}\in[0,\infty]$.}
\entryfield{LDP}{A sequence $X_n$ obeys a large-deviation principle with speed
$a_n\to\infty$ and rate $I$ when closed-set upper and open-set lower logarithmic bounds
are stated explicitly.}
\entryfield{Good rate function}{$I$ is good when every sublevel set
$\{x:I(x)\le c\}$ is compact.  Convexity is not part of the definition of a general
rate function.}
\entryfield{Endpoint scaling}{An endpoint or small-ball asymptotic for
$\FabiusBounded(x)$ as $x\downarrow0$ is not called an LDP until a family, speed, and
normalization have been identified.}
\end{catalogueentry}

\begin{catalogueentry}{information}{Entropy, relative entropy, and Fisher information}
\entryfield{Differential entropy}{For a density $f$,
$h(f)=-\int f(x)\log f(x)\,\Differential x$ when the positive and negative parts are
well defined, with $0\log0=0$ by continuity.  It can be negative.}
\entryfield{Relative entropy}{For probability measures $\mu\ll\nu$,
\[
 D_{\mathrm{KL}}(\mu\|\nu)
 =\int\log\!\left(\frac{\Differential\mu}{\Differential\nu}\right)\,
 \Differential\mu,
\]
and it is $+\infty$ when absolute continuity fails or the integral diverges.  The order
of the arguments matters.}
\entryfield{Fisher information}{For a differentiable positive density,
$J(f)=\int|(\log f)'|^2f\,\Differential x
=\int|f'|^2/f\,\Differential x$ when finite.  Zeros require a weak or extended-value
interpretation.}
\entryfield{Mutual information}{$I(X;Y)=D_{\mathrm{KL}}(
\LawOf{(X,Y)}\|\LawOf X\otimes\LawOf Y)$.}
\entryfield{Status}{Numerical entropy or Fisher-information tables are
\textbf{computational evidence} unless accompanied by verified error bounds.}
\end{catalogueentry}

\begin{catalogueentry}{zonoids}{Minkowski sums and zonoids}
\entryfield{Minkowski sum}{For $A,B\subseteq\RealNumbers^d$,
$A+B=\{a+b:a\in A,\ b\in B\}$ and
$\lambda A=\{\lambda a:a\in A\}$.}
\entryfield{Zonotope}{A finite zonotope generated by vectors $v_1,\ldots,v_N$ is
\[
 Z_N=\sum_{j=1}^{N}[0,v_j]
 =\left\{\sum_{j=1}^{N}t_jv_j:0\le t_j\le1\right\}.
\]}
\entryfield{Zonoid}{A zonoid is a Hausdorff limit of zonotopes.  A random-vector
representation such as $\Expectation[0,X]$ requires integrability and a declared
selection/Aumann-expectation convention.}
\entryfield{Support function}{$h_K(u)=\sup_{x\in K}\langle u,x\rangle$ for a nonempty
compact convex set $K$.  It satisfies $h_{A+B}=h_A+h_B$.}
\entryfield{Common-digit construction}{When coordinates share binary or uniform digits,
the joint random series and ambient dimension are declared.  Marginal Fabius laws alone
do not determine the zonoid or dependence structure.}
\end{catalogueentry}

\begin{catalogueentry}{radon}{Radon transforms and projections}
\entryfield{Hyperplane Radon transform}{For suitable
$f:\RealNumbers^d\to\ComplexNumbers$,
\[
 (\mathcal Rf)(\theta,s)
 =\int_{\{x:\langle x,\theta\rangle=s\}}f(x)\,
 \Differential\mathcal H^{d-1}(x),
\quad \theta\in\mathbb S^{d-1},\ s\in\RealNumbers.
\]}
\entryfield{Redundancy}{$\mathcal Rf(-\theta,-s)=\mathcal Rf(\theta,s)$.  A parameter
domain may quotient this redundancy or retain it explicitly.}
\entryfield{Measure projection}{For a measure $\mu$, the one-dimensional projection in
direction $\theta$ is $(x\mapsto\langle x,\theta\rangle)_\#\mu$.  Its density, when it
exists, is related to the Radon transform but is not assumed.}
\entryfield{Fourier slice}{With compatible $2\pi$ conventions,
$\FourierTwoPi{(\mathcal Rf)(\theta,\cdot)}(\xi)
=\widehat f_{\,2\pi}(\xi\theta)$, where the transform on the right is
$d$-dimensional and uses exponent $-2\pi\ImaginaryUnit\langle x,\eta\rangle$.}
\entryfield{Status}{Dyadic Radon-profile identifications are
\textbf{frontier/conjectural} unless their measure and slice normalizations are proved.}
\end{catalogueentry}

\begin{catalogueentry}{spectral-heat}{Heat traces, spectral zeta functions, and
determinants}
\entryfield{Eigenvalue list}{For a nonnegative self-adjoint operator $L$ with discrete
spectrum, eigenvalues are
$0\le\lambda_0\le\lambda_1\le\cdots$, repeated by multiplicity.}
\entryfield{Heat trace}{$\Theta_L(t)=\TraceOperator(\EulerE^{-tL})
=\sum_j\EulerE^{-t\lambda_j}$ for $t>0$ when trace class.}
\entryfield{Spectral zeta}{$\zeta_L(s)=\sum_{\lambda_j>0}\lambda_j^{-s}$ initially in
its convergence half-plane.  Zero eigenvalues are excluded and their multiplicity is
recorded separately.}
\entryfield{Zeta determinant}{If $\zeta_L$ is regular at zero,
$\det{}'L=\exp(-\zeta_L'(0))$.  The prime means zero modes are omitted; it is not an
ordinary finite determinant.}
\entryfield{Heat coefficients}{An expansion
$\Theta_L(t)\sim\sum_ja_jt^{\alpha_j}$ as $t\downarrow0$ specifies the exponent order,
possible logarithmic terms, and whether the expansion is proved or numerically fitted.}
\entryfield{Digital spectrum}{A digital or dyadic spectrum defined from a sequence or
matrix is not automatically the spectrum of a self-adjoint operator; the construction
must supply the operator and Hilbert space.}
\end{catalogueentry}

\begin{catalogueentry}{noncommutative}{Noncommutative probability and free cumulants}
\entryfield{Noncommutative probability space}{A pair $(\mathcal A,\tau)$ consists of a
unital algebra and a unital linear functional $\tau$; positivity and traciality are
additional hypotheses, not automatic.}
\entryfield{Moments}{For $X\in\mathcal A$, $m_n=\tau(X^n)$.  Products preserve order;
$XY$ need not equal $YX$.}
\entryfield{Free cumulants}{Free cumulants $(r_n)$ are defined by
\[
 m_n=\sum_{\pi\in NC(n)}\prod_{B\in\pi}r_{|B|},
\]
where $NC(n)$ is the lattice of noncrossing partitions of
$\{1,\ldots,n\}$ and the scalar formula assumes a single variable.}
\entryfield{Freeness}{Subalgebras are free when alternating centered products from
different subalgebras have zero $\tau$.  It is not classical independence.}
\entryfield{Status}{Noncommutative Fabius analogues in the frontier corpus are
\textbf{proposed definitions} or \textbf{conjectural}; they are never crosswalked to the
commutative Lean object without an explicit construction.}
\end{catalogueentry}

\begin{catalogueentry}{matrix-dilation}{Matrix dilations and multivariate scaling}
\entryfield{Expansive matrix}{A real $d\times d$ matrix $A$ is expansive when every
eigenvalue has modulus greater than one.}
\entryfield{Dilation operator}{For $f:\RealNumbers^d\to\ComplexNumbers$,
$(D_Af)(x)=f(Ax)$.  The unitary $L^2$ version is
$(U_Af)(x)=|\det A|^{1/2}f(Ax)$.}
\entryfield{Lattice mask}{A matrix refinement equation has the form
$\phi(x)=\sum_{k\in\IntegerNumbers^d}h_k\phi(Ax-k)$ with a declared mask and convergence
mode.}
\entryfield{Fourier scaling}{Under the $d$-dimensional $2\pi$ convention,
$\widehat{f(A\,\cdot)}_{\,2\pi}(\xi)
=|\det A|^{-1}\widehat f_{\,2\pi}(A^{-\mathsf T}\xi)$.}
\entryfield{Scalar specialization}{For $d=1$ and $A=(2)$ this reduces to the familiar
dyadic scaling.  A matrix theorem does not follow from the scalar theorem without
controlling anisotropy and lattice geometry.}
\end{catalogueentry}

\section{Namespace governance and collision register}
\label{sec:namespace-register}

The shared commands above are global semantic names.  Every other symbol is local and is
governed by the rules in this section.  A local declaration is not optional merely
because a letter is conventional: the active corpus combines probability, analysis,
arithmetic, special functions, approximation theory, spectral theory, and proposed
frontier constructions, so nearly every conventional letter has several legitimate
meanings.

\subsection{Scope hierarchy and declaration schema}

The narrowest enclosing declaration wins among local names.  Scopes, from narrowest to
widest, are: a displayed definition, a theorem/lemma/conjecture environment, an algorithm
or table, a subsection explicitly headed ``Notation'', a section, and finally one
standalone document.  A fragment included with \commandname{input} does not create an
independent preamble or global namespace; it inherits the standalone root's shared input
and may introduce only locally scoped names.  A local meaning ends at the closing
environment or section boundary unless a later sentence explicitly renews it.

Every local notation declaration supplies the following data.  ``Not applicable'' is a
permitted answer; silence is not.

\begingroup\small
\begin{longtable}{@{}p{0.20\textwidth}p{0.74\textwidth}@{}}
\toprule
Field & Required content\\
\midrule
\endhead
Kind & Scalar, set, proposition, sequence, function, measure, matrix, operator, formal
series, equivalence class, or other fully named mathematical kind.\\
Type & Domain and codomain, including scalar field, measurable/topological structure,
and whether equality is pointwise or modulo an equivalence relation.\\
Indices & Complete index set, origin, inclusivity of endpoints, and empty-index behavior.\\
Parameters & Which quantities are fixed and varying; admissible domains, inequalities,
integrality, primality, nonvanishing, and branch indices.\\
Formula & Exact definition with every sum, product, integral, supremum, and infimum bound.\\
Normalization & Fourier frequency/sign, measure mass, polynomial leading coefficient,
orthogonality scale, logarithm/power branch, and any factorial or transform prefactor.\\
Exceptional cases & Values at zero, endpoints, removable singularities, empty sets,
empty matrices, divergent parameters, and nonexistence cases.\\
Status & Definition, Lean-proved, classical, frontier-derived, conjectural, refuted,
historical alias, or numerical evidence.\\
Scope end & The environment, section, or document boundary after which the letter is
available for reassignment.\\
\bottomrule
\end{longtable}
\endgroup

\subsection{High-frequency letter collisions}

This register does not ban ordinary local letters.  It records the qualification needed
when two meanings can meet in one document or in cross-document synthesis.

\begingroup\small
\begin{longtable}{@{}p{0.08\textwidth}p{0.39\textwidth}p{0.47\textwidth}@{}}
\toprule
Glyph & Meanings found in the active corpus & Normative resolution\\
\midrule
\endfirsthead
\toprule
Glyph & Meanings found in the active corpus & Normative resolution\\
\midrule
\endhead
$F$ & Bounded Fabius function, signed global extension, CDF, Fourier transform,
polynomial family, matrix, and optimization functional. & Use $\FabiusBounded$ and
$\FabiusGlobal$ for the two Fabius functions, normalization-labelled Fourier commands
for transforms, and $\DecayOptimizationObjective{n}$ for the decay objective.  A local
$F_a$, $F_m$, or matrix $F_N$ retains its typed subscript and declaration.\\
$Q$ & Rational numbers, Fabius quantile, quadrature rule, polynomial, matrix, and a
$q$-dependent quantity. & Use $\RationalNumbers$ for the number field and
$\FabiusQuantile$ or $\FabiusClampedQuantile$ for inverses.  Quadrature is $Q_N$ only
inside its declared numerical-analysis scope; the base of a $q$-series is lower-case
$q$ and always explicit.\\
$N$ & Natural-number type, positive-integer type, truncation level, matrix size, node
count, and asymptotic parameter. & Use $\NaturalNumbers$ or $\PositiveIntegers$ for
sets.  A local $N$ carries a declaration such as $N\in\PositiveIntegers$ and one role;
zero-based and one-based counts are never inferred.\\
$A$ & Event, matrix, Appell family, algebra, feasible set, and dyadic filtration. & Use
$\DyadicSigmaField_n$ for the canonical dyadic filtration.  An event is typed
$A\in\mathscr A$, a matrix has dimensions, and an Appell polynomial has a degree
subscript and preferably a semantic superscript.\\
$D$ & Derivative, differential operator, dilation, denominator, determinant-derived
sequence, and domain. & The upright differential is $\Differential$; derivatives are
$f'$ or $D^nf$ after declaring $D$.  Dilations use $D_a$ or $D_A$, and dyadic
denominators retain the superscript $D_n^{\mathrm{dyad}}$.\\
$E$ & Expectation, event, energy, error, and exponential generating function. & Use
$\Expectation$ for expectation.  Events are roman-labelled or typed in the event
$\sigma$-field; energies and errors retain descriptive subscripts; an EGF is labelled
``egf''.\\
$G$ & Cauchy transform, generating function, Gram matrix, graph, group, and inverse
alias in historical drafts. & A transform is $G_\mu$, a Gram matrix is $G_N$ or
$G_X$, and a group is typed.  The Fabius inverse is never bare $G$; use the quantile
commands.\\
$H$ & Hilbert space, Hankel matrix, Hermite polynomial, Hessian, entropy, and transfer
function. & Use $\mathcal H$ for a Hilbert/RKHS space, $H_N(\mu)$ for a finite Hankel
matrix, $\operatorname{He}_n$ or $H_n^{\mathrm{phys}}$ for the two Hermite
normalizations, and $\nabla^2f$ for a Hessian.\\
$I$ & Identity, indicator, interpolant, interval, information, and inverse. & Use
$\IdentityOperator$ for an identity map and $\IndicatorOf{A}$ for an indicator.
Interpolants carry a degree such as $I_N$, intervals print endpoints, and neither Fabius
inverse is denoted by bare $I$.\\
$K$ & Kernel, compact set, number field, Koopman operator, and cumulant generating
function. & A kernel is $K(x,y)$, a compact set remains $K\subseteq X$, the scalar field
is $\mathbb K$, the Koopman operator is $\mathcal K_T$, and the cumulant generator is
$K_X(t)$ only inside its probability scope.\\
$L$ & Laplace transform, Lagrange/cardinal function, linear operator, likelihood, and
slowly varying factor. & Use $\LaplaceTransformSymbol$ for Laplace.  Cardinal functions
carry an index, operators declare domain/codomain, and slowly varying functions are
labelled $L_{\mathrm{sv}}$ when transforms are nearby.\\
$M$ & Mellin transform, moment-generating function, matrix, moment cutoff, mask, and
manifold. & Use $\MellinTransformSymbol$ and $\MomentGeneratingFunction{X}$ for the two
transforms.  Matrices carry dimensions; cutoffs and masks are locally typed.\\
$P$ & Probability, polynomial, projection, transition kernel, Padé numerator, and
partition. & Use $\Probability$ for probability.  Polynomials carry degree indices,
projections are $P_V$ with their range, Markov kernels are $P(x,\Differential y)$, and
Padé numerators are $P_L$ only in a declared $[L/M]$ scope.\\
$R$ & Real-number type, radius, residue, resolvent, recurrence polynomial, remainder,
and Reshetnikov integer. & Use $\RealNumbers$ and $\ResidueOperator$ for the first two
semantic roles.  Resolvents are $R_T(z)$, remainders carry an error superscript, and the
arithmetic integer is $R_n^{\mathrm{Res}}$.\\
$S$ & Sum, support, Stieltjes transform, saddle action, spline space, and random series. &
Use $\SupportOf{f}$ for support and $S_\mu$ for Stieltjes transforms.  Saddle actions,
spline spaces, and random series carry semantic subscripts and declared domains.\\
$T$ & Translation, transformation, Toeplitz matrix, random time, Thue--Morse object, and
truncation. & Translations are $T_b$, Toeplitz matrices are $T_N(a)$, and the sign
sequence uses $\ThueMorseSign{n}$.  A dynamical map $T:X\to X$ is typed at first use.\\
$W$ & Lambert function, Wasserstein distance, Walsh function, and weight. & Use
$\LambertW_k$ with branch $k$ and $W_p(\mu,\nu)$ with transport exponent $p$.  Walsh
functions and weights carry their own indices and domains.\\
$\varepsilon$ & Thue--Morse sign, small positive tolerance, perturbation parameter, and
error sequence. & Use $\ThueMorseSign{n}$ for the sign.  A tolerance is quantified
$\varepsilon>0$; perturbations and errors receive semantic subscripts when signs occur
in the same scope.\\
\bottomrule
\end{longtable}
\endgroup

\begin{catalogueentry}{inverse-synthesis-local-namespace}
{Finite-spline, Fourier-product, and endpoint-kernel notation in the inverse-synthesis
volume}
\entryfield{Scope}{The commands \commandname{FiniteSpline},
\commandname{FiniteSplineInverse}, \commandname{Fn}, \commandname{Gn},
\commandname{PhiAng}, \commandname{PhiCyc}, and \commandname{EndpointKernel} are local to
the standalone inverse-analyticity, asymptotics, and computability volume and its included
chapters.  They are not shared-API commands, and their meanings end with that standalone
document.}
\entryfield{Finite-prefix probability data}{For \(k\in\PositiveIntegers\), let
\(U_1,U_2,\ldots\) be independent random variables uniformly distributed on
\([-1,1]\), set
\[
 S_k\DefinitionEquals\sum_{j=1}^{k}2^{-j}U_j,
\]
and let \(p_k:\RealNumbers\to[0,\infty)\) be the Lebesgue density of \(S_k\).}
\entryfield{Finite spline family}{\commandname{FiniteSpline}\texttt{\{k\}} has TeX
arity one and prints \(F_k^{\mathrm{spl}}\).  Its mathematical argument is
\(k\in\PositiveIntegers\), and
\[
 F_k^{\mathrm{spl}}:[0,1]\longrightarrow[0,\infty),\qquad
 F_k^{\mathrm{spl}}(x)\DefinitionEquals p_k(x-1).
\]
The superscript \(\mathrm{spl}\) marks the finite convolution prefix and is never an
iteration exponent.}
\entryfield{Exact box-spline form}{For \(x\in[0,1]\), the same finite spline is the
explicit Thue--Morse box spline
\[
 F_k^{\mathrm{spl}}(x)
 =\frac{2^{\binom{k}{2}}}{(k-1)!}
 \sum_{j=0}^{2^k-1}(-1)^{s_2(j)}
 \PositivePart{x-2^{-k}-j2^{1-k}}^{k-1},
 \qquad s_2(j)\DefinitionEquals\BinaryDigitSum(j).
\]
Here \(\BinaryDigitSum(j)\) is the number of one-bits in the nonnegative integer \(j\),
and \(\PositivePart{t}\DefinitionEquals\max\{t,0\}\).  Thus neither the Thue--Morse
sign nor the truncation at zero is implicit.}
\entryfield{Local inverse family}{\commandname{FiniteSplineInverse}\texttt{\{k\}} has
TeX arity one and prints \(G_k^{\mathrm{spl}}\).  On a stated interval
\(J\subset(0,1)\) on which \(F_k^{\mathrm{spl}}\) is strictly increasing, it denotes the
branch
\[
 G_k^{\mathrm{spl}}:
 F_k^{\mathrm{spl}}(J)\longrightarrow J,\qquad
 G_k^{\mathrm{spl}}\DefinitionEquals
 \left(F_k^{\mathrm{spl}}\big|_J\right)^{-1}.
\]
No inverse value is implied unless the branch interval and monotonicity hypothesis have
been supplied.  In the asymptotic uses, \(J\) is a fixed compact subinterval of
\((0,1)\), and the branch exists for all sufficiently large \(k\).}
\entryfield{Fixed-index aliases}{\commandname{Fn} and \commandname{Gn} each have TeX
arity zero.  With an ambient \(n\in\PositiveIntegers\), they expand respectively to
\(F_n^{\mathrm{spl}}\) and \(G_n^{\mathrm{spl}}\).  They do not denote the forward
composition iterate \(F^{\circ n}\), the inverse composition iterate
\((F^{-1})^{\circ n}\), or either globally catalogued Fabius inverse.}
\entryfield{Finite-prefix use family}{The spline and local-inverse symbols are used for
finite-prefix approximation, uniform inverse-transfer asymptotics, and Richardson
extrapolation.  The prefix depth, evaluation variable, and---for the inverse---the
selected monotone branch interval are ambient mathematical data, not additional TeX
arguments.}
\entryfield{Angular-frequency product}{\commandname{PhiAng} has TeX arity zero; its
evaluation variable is external.  It prints \(\Phi_{\mathrm{ang}}\) and denotes the entire
extension of
\[
 \Phi_{\mathrm{ang}}(\xi)
 \DefinitionEquals\FourierAngular{\RvachevUp}(\xi)
 =\prod_{j=1}^{\infty}\SincRad(2^{-j}\xi),
 \qquad \xi\in\ComplexNumbers,
\]
where \(\SincRad z=\sin z/z\) for \(z\ne0\) and \(\SincRad0=1\).}
\entryfield{Cycles-per-unit product}{\commandname{PhiCyc} has TeX arity zero and an
external variable.  It prints \(\Phi_{2\pi}\) and denotes the same entire transform in
the cycles-per-unit coordinate:
\[
 \Phi_{2\pi}(z)
 \DefinitionEquals\FourierTwoPi{\RvachevUp}(z)
 =\Phi_{\mathrm{ang}}(2\pi z)
 =\prod_{j=0}^{\infty}\SincPi(z/2^j),
 \qquad z\in\ComplexNumbers.
\]
Thus neither the factor \(2\pi\) nor the transform normalization is suppressed.}
\entryfield{Fourier-product use families}{\(\Phi_{\mathrm{ang}}\) is used for
reciprocal-sinc coefficients, moment-generating/Fourier rotation, and polynomial
deconvolution.  \(\Phi_{2\pi}\) is used for integer zero multiplicities, differentiated
transform identities, moments, and dyadic self-sampling.  Both symbols suppress the
fixed normalized Rvachev up-function, but their printed subscripts make the coordinate
normalization explicit; they are never silently interchanged.}
\entryfield{Endpoint variables}{For the endpoint regime \(y\downarrow0\), set
\(L=\log2\), \(T=\log(1/y)\), \(\rho=\sqrt{2T/L}\), \(z=\rho^{-1}\),
\(\ell=\log\rho\), \(\beta=L^{-1}+1/2\), and
\(\theta=\rho+\beta\pmod 1\).  The phase functions \(\Psi,A_1,A_2,\ldots\) are the fixed
one-periodic coefficient functions in the forward expansion
\[
 \log F(x)=-\frac L2\lambda^2+L\beta\lambda-\frac12\log\lambda
 +C_{\sharp}+\Psi(\lambda)
 +\sum_{j=1}^{N-1}\frac{A_j(\lambda)}{\lambda^j}
 +\BigO_N(\lambda^{-N}),
\]
where \(\lambda=-\LambertW_{-1}(-Lx)/L\) and the remainder is uniform in the real phase.}
\entryfield{Endpoint kernel}{\commandname{EndpointKernel} has TeX arity zero and prints
\(\mathcal K\); both displayed arguments are mathematical arguments external to the
macro.  In the formal two-scale coefficient algebra generated by \(\ell\), the
one-periodic functions above, and their derivatives, define
\[
 \begin{aligned}
 \mathcal R(z,w)={}&-\frac L2\beta^2+\frac L2w^2+\frac12\ell
 +\frac12\log(1+\beta z+zw)-C_{\sharp}-\Psi(\theta+w)\\
 &-\sum_{j\ge1}
 \frac{A_j(\theta+w)z^j}{(1+\beta z+zw)^j},\\
 \mathcal K(z,w)\DefinitionEquals{}&-\frac1L\mathcal R(z,w).
 \end{aligned}
\]
The inverse endpoint correction \(D\in z\mathscr A[\ell][[z]]\) is then characterized by
\(D=z\mathcal K(z,D)\).  Here \(w\) is the second formal argument; it is not a hidden
function of \(z\) until the fixed-point substitution \(w=D\) is made.
The same source later writes \(\mathcal K_W=-\mathcal R_W/L\) for a distinct
Wright--omega-carrier kernel.  The form
\commandname{EndpointKernel}\texttt{\_W} prints that local symbol; it does not denote
\(\mathcal K(z,W)\), a derivative, or an automatically shared subscripted family.}
\entryfield{Endpoint-kernel use family}{The kernel is used in the Lagrange diagonal
formulas for the fixed point, the all-orders inverse-endpoint recursion, and comparison
with the distinct Wright--omega carrier.  Its suppressed ambient data are
\(L,\beta,\ell,\theta,C_{\sharp},\Psi\), the full coefficient sequence
\((A_j)_{j\ge1}\), and the chosen two-scale coefficient algebra.}
\entryfield{Script composition}{\commandname{FiniteSpline},
\commandname{FiniteSplineInverse}, \commandname{Fn}, and \commandname{Gn} already own a
subscript and the superscript \(\mathrm{spl}\); a derivative or further order must group
the whole object, as in \((F_n^{\mathrm{spl}})^{(m)}\).  \commandname{PhiAng} and
\commandname{PhiCyc} already own their normalization subscripts; powers and derivatives
may follow the grouped object, but no second raw subscript may be appended.
\commandname{EndpointKernel} owns no script, yet any local suffix such as \(_W\) must be
declared as a new semantic object rather than inferred as evaluation or differentiation.}
\entryfield{Status}{The finite-prefix objects and the two Fourier products are exact
definitions.  The endpoint kernel is an exact formal definition inside the volume's
human-proved all-orders coefficient construction; convergence or analytic equality is
asserted only where a separate theorem supplies it.}
\end{catalogueentry}

\subsection{Accents, scripts, and overloaded punctuation}

\begin{catalogueentry}{accent-collisions}{Hats, tildes, bars, stars, primes, and scripts}
\entryfield{Hats}{A normalization-labelled hat produced by
\commandname{FourierTwoPi} or \commandname{FourierAngular} is a Fourier transform.  A
unit-period coefficient uses \commandname{FourierSeriesCoefficient}; any other period uses
\commandname{FourierSeriesCoefficientAtPeriod} and prints that period.  A hat on a polynomial,
matrix, estimator, accelerated approximant, or lifted object is local: all use sites call
a descriptive document-local macro, and the raw accent occurs only in that macro's
one-line definition.  The immediately preceding source comment has the form
\texttt{LOCAL-NON-FOURIER-HAT: \textbackslash MacroName meaning}, names the same macro,
and explains its mathematical role.  A generic local alias for \commandname{widehat} or
\commandname{hat} is forbidden.  Thus no unlabelled hat is silently interpreted as a
Fourier transform in a document that also uses monic or normalized polynomials.}
\entryfield{Tildes}{A tilde may mark centering, normalization, continuation, lifting, or
an alternative sign convention.  The transformed object's type and the defining map are
given at first use.  It is never a substitute for the signed global Fabius function.}
\entryfield{Bars}{For a scalar, bars mean $\AbsoluteValue{x}$; for a finite set they mean
$\CardinalityOf{A}$; a restriction is $f|_S$; divisibility is $a\mid b$; conditional
probability is $\Probability(A\mid B)$ with $\Probability(B)>0$.  These roles are
distinguished by type and spacing, not guessed from visual similarity.}
\entryfield{Stars}{For a scalar, $\overline z$ is conjugation; $A^\ast$ is an adjoint;
$X^\ast$ in the zero-bias entry is a distributional transform; $X^\ast$ may also denote
a dual space only after a type declaration.  None of these implies optimality; an
optimizer is labelled $x_{\mathrm{opt}}$.}
\entryfield{Primes}{$f'$ is a derivative, $\zeta_L'$ a derivative of the spectral zeta
function, and $\det{}'L$ omits zero modes.  A primed variable such as $x'$ is a new local
variable only after declaration.}
\entryfield{Superscripts}{Parenthesized $f^{(n)}$ means the $n$th derivative;
$f^{\circ n}$ means iteration; $f^{\Convolution n}$ means convolution power; $A^{\mathsf
T}$ means transpose; $A^\ast$ means adjoint.  Ordinary $f^n$ remains pointwise or
algebraic power.}
\entryfield{Subscripts}{A shared symbol that already owns a semantic subscript is not
followed by a second raw subscript.  Use the command's arguments or introduce a wrapper
whose expansion has one syntactically valid script.  This rule covers number systems,
quantiles, transform hats, and any future decorated semantic command.}
\end{catalogueentry}

\subsection{Exhaustive register of document-local non-Fourier hats}
\label{sec:local-non-fourier-hats}

This register inventories every distinct command named by an active
\texttt{LOCAL-NON-FOURIER-HAT:} declaration.  There are 35 marked declarations
and 35 distinct command names; each is declared in exactly one active document scope.
The phrase \emph{prints a wide hat over \(X\)} describes literal output
without reintroducing a generic accent command.  TeX arguments, mathematical arguments,
and scripts appended by a caller are distinguished explicitly.  These commands are
document-local: their definitions and ambient data do not survive the document or
included-fragment scope stated below.  Each heading contains one machine-readable
\commandname{localaccentname} token, so the strict audit rejects a missing, extra, or
duplicate register entry.

\subsubsection{Extensions, certified data, and formal series}

\Needspace{20\baselineskip}
\paragraph{\localaccentname{ClampedExtensionOf}}
TeX arity is one; argument \(1\) is a continuous function \(f\) on
\(I=[0,1]\).  The command prints a wide hat over \(f\).  Its mathematical value is
the continuous map
\[
 E_I f:\RealNumbers\longrightarrow\RealNumbers,\qquad
 (E_I f)(x)
 \DefinitionEquals
 f\!\left(\operatorname{proj}_{I}x\right)
 =
 \begin{cases}
  f(0),&x\le0,\\
  f(x),&0\le x\le1,\\
  f(1),&1\le x,
 \end{cases}
\]
where \(f(x)\) in the middle branch uses the canonical membership
\(x\in I\).  In the Lean construction this is precomposition with
\(\operatorname{Set.projIcc}(0,1)\); admissibility of \(f\) supplies the bounds
and reflection identity used later but is not required to define \(E_I f\).
The only use family is the fixed-point primitive
\(A_f(y)=\int_0^y(E_If)(t)\,\Differential t\).
The accent means endpoint-clamped extension, not Fourier transformation.
Status is an exact definition; the later admissibility and fixed-point theorems are
separate assertions.  The clamping interval, endpoint values, and projection map are suppressed by the
printed symbol.  The function argument must include any script belonging to
\(f\) inside the braces; a derivative or other decoration of the resulting map
must group the complete expansion before adding the decoration.

\paragraph{\localaccentname{FullRvachevLegendreCoefficientRat}}
TeX arity is one; argument \(1\) is the full Legendre degree
\(k\in\NaturalNumbers\).  The command prints a wide hat over \(c\), with
subscript \(k\) and superscript \(\RationalNumbers\).  It denotes the function
\(c^{\mathrm{full},\RationalNumbers}:\NaturalNumbers\to\RationalNumbers\)
defined by
\[
 c_k^{\mathrm{full},\RationalNumbers}
 \DefinitionEquals
 \begin{cases}
  c_{k/2}^{\RationalNumbers},&2\mid k,\\
  0,&2\nmid k,
 \end{cases}
 \qquad
 c_{2r}^{\mathrm{full},\RationalNumbers}=c_r^{\RationalNumbers},
 \qquad
 c_{2r+1}^{\mathrm{full},\RationalNumbers}=0.
\]
Here \(c_r^{\RationalNumbers}\) is the canonical even Rvachev--Legendre
coefficient.  Equivalently,
\[
 c_k^{\mathrm{full},\RationalNumbers}
 =
 \frac{2k+1}{2}\,
 \mathcal M_{\mathrm{Rv}}\!\left(P_k^{\RationalNumbers}\right),
\]
where \(P_k^{\RationalNumbers}\) is the rational Legendre polynomial and
\(\mathcal M_{\mathrm{Rv}}\) is the moment functional with moment sequence
\texttt{rvachevRawMomentRat}.  The use family consists of the rational
Legendre--Gaunt formulas and their real-cast comparison.  The accent means
parity completion of an even-indexed sequence.  The canonical moment sequence,
Legendre normalization, and coefficient family are suppressed ambient data.
Status is an exact rational sequence definition; comparison with analytic coefficients
retains the source's explicit Fabius-function hypotheses.
The command already owns both its subscript and its
\(\RationalNumbers\)-superscript; appending any further raw script would be a
script-composition error.

\paragraph{\localaccentname{CertifiedEnclosureCenter}}
TeX arity is one; argument \(1\) is the exact quantity whose validated numerical
center is being named.  The command prints a wide hat over that argument.
Its sole active use has \(Q,Q_{\mathrm{num}}\in\ComplexNumbers\) and means only
the center component of the certificate
\[
 \varepsilon_Q\ge0,\qquad
 \AbsoluteValue{Q-Q_{\mathrm{num}}}\le\varepsilon_Q.
\]
In the validated argument-principle scope, let \(\Omega\) be a bounded Jordan domain,
let \(f\) be holomorphic on a neighborhood of \(\overline\Omega\), assume
\(0\notin f(\partial\Omega)\), and decompose the positively oriented boundary into
finitely many \(C^3\) panels \(\gamma_k:[0,1]\to\partial\Omega\).  Define
\[
 \begin{aligned}
 g_k(s)&=
 \frac{f'(\gamma_k(s))\gamma_k'(s)}{f(\gamma_k(s))},
 &Q&=\sum_k\frac{g_k(0)+g_k(1)}2,\\
 M_k&\ge\sup_{0\le s\le1}\AbsoluteValue{g_k''(s)},
 &E&=\frac1{12}\sum_kM_k.
 \end{aligned}
\]
The panelwise trapezoid estimate gives
\[
 \AbsoluteValue{\oint_{\partial\Omega}\frac{f'(z)}{f(z)}\,\Differential z-Q}
 \le E.
\]
\par\Needspace{10\baselineskip}\noindent
Outward-rounded arithmetic supplies \(Q_{\mathrm{num}}\) and
\(\varepsilon_Q\).  The resulting certified zero-count disk has center
\(Q_{\mathrm{num}}/(2\pi\ImaginaryUnit)\) and radius
\((E+\varepsilon_Q)/(2\pi)\).  Status is therefore a certificate schema:
the accented center alone is not a certificate.  The ambient normed field,
domain \(\Omega\), function \(f\), boundary orientation and panelization
\((\gamma_k)\), derivative bounds \((M_k)\), trapezoid rule, rounding method,
error radius, and exact target are all suppressed data.
The accent means validated numerical center, not Fourier transformation.
The complete exact-quantity name belongs inside the single argument; callers
must group the expansion before attaching any further script.

\paragraph{\localaccentname{FormalGevreySeriesSymbol}}
TeX arity is one; argument \(1\) is the name \(R\) of a formal series, not its
formal variable.  The command prints a wide hat over \(R\).  In its cyclotomic
resurgence scope it denotes
\[
 R_{\mathrm{for}}(t)
 \DefinitionEquals
 \sum_{n\ge0}c_nt^n
 \in\ComplexNumbers[[t]],
 \qquad
 \AbsoluteValue{c_n}\le CA^n n!
 \quad(n\in\NaturalNumbers)
\]
for some constants \(C,A>0\).  The document's exact formal Borel convention is
\[
 \mathscr B R_{\mathrm{for}}(\xi)
 \DefinitionEquals
 \sum_{n\ge0}\frac{c_{n+1}}{n!}\xi^n;
\]
the constant \(c_0\) is carried separately by Borel--Laplace summation.
Status is a formal-series definition used immediately before a conjectural
resurgence statement; the Gevrey estimate itself does not assert Borel
summability or analytic continuation.  The coefficient field, formal variable,
coefficients, Gevrey order, sector, and Borel convention are suppressed by the
printed symbol.  The accent means formal asymptotic series prior to Borel
transformation.  The series name, including any semantic script, must be placed
inside the one argument; any subsequent decoration must group the expansion.

\Needspace{16\baselineskip}
\paragraph{\localaccentname{FormalAsymptoticExpansionOf}}
TeX arity is one; argument \(1\) is the name \(f\) of an actual function.  The command
prints a wide hat over \(f\); its evaluation variable remains outside the macro.  In the
incoming transseries volume, fix \(X>0\), \(t=X^{-1}\), \(\lambda=\log X\), a
characteristic-zero coefficient field \(K\), and polynomial blocks
\(A_n\in K[\lambda]\).  The formal object is
\[
 \widehat f(X)
 \DefinitionEquals
 \sum_{n\ge n_0}A_n(\log X)X^{-n}
 \in K[\lambda]((t)).
\]
It is a coefficient array in the \(t\)-adic completion.  The separate analytic assertion
that it is a Poincaré expansion of \(f\) means that, for every truncation index \(N\), the
difference between \(f(X)\) and the sum through the block \(N-1\) is little-oh of the
last retained scale along a declared ray or sector.  The field, limit direction, branch
of \(\log X\), block degrees, and precise remainder scale are suppressed data.  The
accent means formal asymptotic expansion, not Fourier transformation and not equality of
the formal sum with \(f\).

\Needspace{18\baselineskip}
\paragraph{\localaccentname{LambertAsymptoticApproximation}}
TeX arity is one; argument \(1\) is the truncation order
\(N\in\PositiveIntegers\).  It prints \(\widehat W_N\), already owning the subscript
\(N\).  On the positive principal-branch ray, let \(z=\EulerE^L\), \(L>1\), put
\(\ell=\log L\), and let \(P_n\) be the explicitly defined Lambert polynomial family of
that document.  Then
\[
 \widehat W_N
 \DefinitionEquals
 L-\ell+\sum_{n=1}^{N}\frac{P_n(\ell)}{L^n}.
\]
This is a finite real scalar approximation to \(W_0(z)\), not the Lambert function
itself and not a Fourier transform.  Its logarithmic residual is
\[
 \rho_N\DefinitionEquals\widehat W_N+\log\widehat W_N-L
\]
whenever \(\widehat W_N>0\).  If \(W=W_0(\EulerE^L)\), the mean-value theorem gives
\(\widehat W_N-W=\rho_N/(1+1/\theta_N)\) for some \(\theta_N\) between them, provided
that interval avoids zero.  The branch, positive ray, variables \(L,\ell\), polynomial
normalization, and residual domain are suppressed.  A semantic superscript such as
\(\mathrm{new}\) for one Newton correction may follow the complete command expansion;
no second raw subscript may be appended.

\Needspace{13\baselineskip}
\paragraph{\localaccentname{wideeq}}
TeX arity is zero; the command prints a hatted equality sign as a binary relation.  For
two formal polynomial--logarithmic series
\[
 A=\sum_{\gamma}a_\gamma\mathfrak m_\gamma,\qquad
 B=\sum_{\gamma}b_\gamma\mathfrak m_\gamma
\]
written over the same coefficient field, ordered monomial family, and locally finite
support convention,
\[
 A\mathrel{\widehat=}B
 \quad\Longleftrightarrow\quad
 a_\gamma=b_\gamma\ \text{for every index }\gamma.
\]
It is coefficientwise formal identity, not equality of convergent numerical sums,
asymptotic equivalence, equality modulo a truncation ideal, or Fourier equality.  The
ambient coefficient field, monomial group, ordering, support condition, and formal
variables must already be declared.

\Needspace{16\baselineskip}
\paragraph{\localaccentname{DulacFormalSeriesOf}}
TeX arity is one; argument \(1\) is the name \(f\) of a function or germ.  The command
prints a wide hat over \(f\), with the small evaluation variable supplied after the
expansion.  For \(x\to0^+\), put
\[
 X=x^{-1},\qquad \ell=\log(1/x)=\log X,
\]
choose a strictly increasing exponent sequence
\(\lambda_0<\lambda_1<\cdots\to+\infty\), and polynomial blocks
\(P_j\in K[\ell]\).  Then
\[
 \widehat f(x)
 \DefinitionEquals\sum_{j\ge0}x^{\lambda_j}P_j(\ell)
\]
is the formal Dulac or power--log series, subject to the declared local-finiteness
condition that makes each formal coefficient operation finite.  The dominance order
compares the exponent \(\lambda_j\) first and logarithmic degree second.  The coefficient
field, exponent group, branch and sign of the logarithm, support condition, and whether
the formal series is an actual asymptotic expansion of a germ are suppressed.  The
accent means formal Dulac expansion, not Fourier transformation.

\subsubsection{Inverse, sampling, and prediction constructions}

\paragraph{\localaccentname{FactorialScaledInverseTransferCoefficient}}
TeX arity is one; argument \(1\) is the coefficient index
\(m\in\{1,\ldots,J-1\}\).  The command prints a wide hat over \(c\) with
subscript \(m\).  The evaluation variable \(y\) is external to the macro and
is suppressed at the bare-symbol use sites.

Fix \(0<\eta<1/2\), \(J\in\PositiveIntegers\), \(q=1/4\), and
\(\varepsilon_n=q^n\).  In the uniform inverse-transfer expansion
\[
 G_n^{\mathrm{spl}}(y)
 =
 \FabiusQuantile(y)
 +\sum_{j=1}^{J-1}c_j(y)\varepsilon_n^j
 +\BigO_{\eta,J}(\varepsilon_n^J),
 \qquad y\in[\eta,1-\eta],
\]
the theorem constructs
\[
 c_m\in C^\infty([\eta,1-\eta];\RealNumbers)
 \qquad(1\le m\le J-1).
\]
With \(x=\FabiusQuantile(y)\), the accented coefficient is exactly
\[
 \widehat c_m(y)\DefinitionEquals m!\,c_m(y).
\]
The source normally suppresses the displayed argument \(y\) and writes only
\(\widehat c_m\) inside partial exponential Bell polynomials.

More explicitly, put
\[
 a_j(x)\DefinitionEquals(-1)^j b_jF^{(2j)}(x),
\]
and let \(\mathcal B_{r,k}\) be the partial exponential Bell polynomial with
\(\mathcal B_{0,0}=1\).  Then
\[
\begin{aligned}
 c_m(y)=-\frac1{F'(x)}\Bigg[&
 \frac1{m!}\sum_{k=2}^{m}F^{(k)}(x)
 \mathcal B_{m,k}
   (\widehat c_1,\ldots,\widehat c_{m-k+1})\\
 &+\sum_{j=1}^{m}\frac1{(m-j)!}
   \sum_{k=0}^{m-j}a_j^{(k)}(x)
 \mathcal B_{m-j,k}
   (\widehat c_1,\ldots,\widehat c_{m-j-k+1})
 \Bigg],
\end{aligned}
\]
where the term with \(m=j\) is interpreted as \(a_m(x)\).

Status is an exact factorial rescaling inside the source's proved uniform
inverse-transfer theorem.  The source supplies a proof but makes no Lean
formalization claim.  The printed symbol suppresses \(y\), \(\eta\), \(J\),
\(q\), the spline inverse \(G_n^{\mathrm{spl}}\), the base function \(F\),
the perturbation coefficients \((a_j)\), and the Bell-polynomial convention.
The accent means conversion from ordinary inverse-expansion coefficients to
the factorial scaling used by exponential Bell polynomials.  It is not a
Fourier transform and not the tomographic Fourier-coefficient estimator
elsewhere printed as a hatted \(c\).

The macro already owns the subscript \(m\).  An explicit evaluation may be
written after the complete expansion as \(\widehat c_m(y)\); callers must not
encode \(y\) as a second raw subscript.  Any derivative, superscript, or
other decoration must group the complete expanded symbol.

\paragraph{\localaccentname{FormalTaylorGerm}}
TeX arity is one; argument \(1\) is a function or germ label, not the center or
local variable.  The command prints a wide hat over that label.  For a smooth
local map \(h:\RealNumbers\to\RealNumbers\) at \(x_0\), its full formal Taylor
series in the increment \(z\) is
\[
 \mathcal T_{x_0}h(z)
 \DefinitionEquals
 \sum_{m\ge0}\frac{h^{(m)}(x_0)}{m!}z^m
 \in\RealNumbers[[z]].
\]
When both source and target centers are translated, the document instead uses
the zero-constant germ
\[
 \mathcal T^{0}_{x_0}h(z)
 \DefinitionEquals
 \mathcal T_{x_0}h(z)-h(x_0)
 =
 \sum_{m\ge1}\frac{h^{(m)}(x_0)}{m!}z^m.
\]
The active use families are formal compositional inversion of smooth local
inverses, the affine center jets of \(G_n\) and \(K_n\), numerical
coefficient-reversion residuals, and the exact quarter-point germ of \(F\).
Status is formal algebra: positive radius and equality with the original
smooth map are separate assertions.  The center, increment variable, whether
the output center was subtracted, coefficient field, and convergence status
are all suppressed and must be stated at each use.  The accent means Taylor
germ, not Fourier transformation.  A scripted function label such as
\(F_{1/4}\) belongs wholly inside the one argument; any further decoration
must group the complete expansion.

\paragraph{\localaccentname{GeometricPrefixReproduction}}
TeX arity is one; argument \(1\) is the prefix depth
\(m\in\PositiveIntegers\).  The command prints a wide hat over \(X\) with
subscript \(m\).  For
\[
 X_q=(1-q)\sum_{j\ge0}q^jU_j,\qquad
 U_j\stackrel{\mathrm{iid}}{\sim}\operatorname{Unif}[-1,1],
 \qquad0<q<1,
\]
it denotes the real square-integrable random variable
\[
 S_{q,m}
 \DefinitionEquals
 (1-q)\sum_{j=0}^{m-1}q^jU_j
 \in L^2(\Omega;\RealNumbers).
\]
This is the deterministic first-\(m\)-digit reproduction used in the
squared-error test channel, with
\[
 I(X_q;S_{q,m})=m\log(1/q),
 \qquad
 \Expectation[(X_q-S_{q,m})^2]=q^{2m}\Variance(X_q).
\]
Status is an exact definition and exact rate/distortion calculation.
The parameter \(q\), probability space, digit sequence, digit law, and
distortion measure are suppressed.  The accent means reproduction or
estimation, not Fourier transformation.  The argument already becomes the
subscript of \(X\); no additional raw subscript or superscript may follow the
expanded command.

\paragraph{\localaccentname{LeastSquaresProjectionEstimate}}
TeX arity is one; argument \(1\) is the complete target random-variable
expression.  The command prints a wide hat over that target.  In its sole
scope,
\[
 X_q=(1-q)\sum_{n\ge0}q^nU_n,\qquad
 U_n\stackrel{\mathrm{iid}}{\sim}\operatorname{Unif}[0,1],
 \qquad
 Y_q=
 \frac{X_q-\tfrac12}
 {\sqrt{(1-q)/(12(1+q))}}.
\]
For \(0<q,q_1,\ldots,q_m<1\), with the comparison nodes pairwise distinct and
\(q\notin\{q_1,\ldots,q_m\}\), the accented target is
the unique real \(L^2\)-orthogonal projection
\[
 \Pi_{\boldsymbol q}Y_q
 \DefinitionEquals
 \operatorname*{argmin}_{
 Z\in\SpanOperator(Y_{q_1},\ldots,Y_{q_m})}
 \Expectation[(Y_q-Z)^2],
\]
and
\[
 \Expectation[(Y_q-\Pi_{\boldsymbol q}Y_q)^2]
 =
 \prod_{j=1}^{m}
 \left(\frac{q-q_j}{1-qq_j}\right)^2.
\]
Status is an exact Hilbert-space projection and prediction-error identity.
The common probability space and digits, standardization, comparison nodes,
coefficient field, and projection subspace are suppressed.  The accent means
least-squares predictor.  The full target, including its own subscript, must
be enclosed in the one argument; any subsequent decoration must group the
complete expansion.

\paragraph{\localaccentname{NestedLambertLogApproximation}}
TeX arity is one; argument \(1\) is the iterate index
\(r\in\{0,\ldots,n\}\).  The command prints a wide hat over \(T\) with
subscript \(r\).  The endpoint input, recursion depth, truncation order, and
evaluation phases are not TeX arguments.

In its sole active scope, \(I=\FabiusQuantile\) on \((0,1)\).  Put
\[
 y_0=y,\qquad y_{r+1}=I(y_r),\qquad
 T_r\DefinitionEquals\log(1/y_r),
\]
in the endpoint regime where every attained \(y_r\) is positive and belongs
to the domain of \(I\).  The source does not assert a larger maximal domain
for this construction.  For a fixed depth \(n\ge1\), put
\[
 q_n\DefinitionEquals T_0^{-2^{-n}},\qquad
 \ell\DefinitionEquals\log T_0,
\]
and set \(L=\log2\) and \(\beta=L^{-1}+1/2\).  For a chosen truncation order
\(N\in\PositiveIntegers\) and the one-step coefficient functions \(h_j\)
supplied by the preceding all-orders theorem, define
\[
 \mathcal I_N(T)
 \DefinitionEquals
 L\rho+L\beta-\log\rho
 -\sum_{j=1}^{N}
   \frac{h_j(\log\rho,\rho+\beta)}{\rho^j},
 \qquad
 \rho\DefinitionEquals\sqrt{2T/L}.
\]
The accented sequence is the procedural recursion
\[
 \widehat T_0\DefinitionEquals T_0,
\]
followed, for each \(r\), by expansion of
\(\mathcal I_N(\widehat T_r)\) in powers of \(q_n\) and deletion of every
term whose \(q_n\)-exponent is at least \(N\); the retained expression is
called \(\widehat T_{r+1}\).

The source presents this construction inside the open
\emph{Nested Lambert phase closure} problem.  It asks for, but does not prove,
the uniform estimate
\[
 T_r-\widehat T_r
 =
 \BigO_{n,N}\!\left((1+\ell)^{D'_{n,N}}q_n^N\right)
 \qquad(0\le r\le n),
\]
the phase-substitution bounds, nonconstant transport of each periodic phase,
and closure or minimal enlargement of the exponent lattice
\(q_n^{\IntegerNumbers}\).  Because that closure is itself part of the
problem, the source does not yet supply a proved ambient formal-series ring
or a proved global function domain and codomain for every \(\widehat T_r\).
It intends real asymptotic expressions along the attained endpoint regime.
No Lean formalization is claimed.

The printed symbol suppresses \(y\), \(I\), \(n\), \(N\), \(T_0\), \(q_n\),
\(\ell\), \(L\), \(\beta\), the coefficient family \((h_j)\), the phase
vector, and the expansion/truncation convention.  The accent means recursive
asymptotic truncation, not Fourier transformation and not the
lacunary-product tail extrapolation elsewhere printed with a hatted \(T\).
The macro already owns the subscript \(r\); callers must not append a second
raw subscript.  Any further decoration must group the complete expanded
symbol and receive a separate semantic declaration.

\subsubsection{Omission marks, endpoint matrices, and lifted operators}

\paragraph{\localaccentname{OmittedSymmetricArgument}}
TeX arity is one; argument \(1\) is the complete variable to delete from a
symmetric-polynomial argument list.  The command prints a wide hat over that
variable.  It is omission metanotation, not a map on the scalar \(q_j\).
Its exact active meaning is
\[
 T_k
 \DefinitionEquals
 \sum_{j\ge1}q_j^2
 e_{k-1}(q_1,\ldots,q_{j-1},q_{j+1},\ldots),
 \qquad
 q_j=Q^j,\qquad Q=\frac14,
\]
so the \(j\)-th entry is absent from the arguments of \(e_{k-1}\).  The
identity
\[
 e_1e_k=T_k+(k+1)e_{k+1}
\]
supplies the first correction in the small-field \(q\)-binomial chaos
asymptotic.  Status is exact combinatorial deletion notation.  The ambient
infinite sequence, deleted position, convergence of the symmetric sums, and
value \(Q=1/4\) are suppressed.  The accent means omission, not Fourier
transformation.  A scripted item such as \(q_j\) belongs inside the one
argument, and the resulting omission marker receives no further script.

\paragraph{\localaccentname{GeneratedEndpointCoefficientMatrix}}
TeX arity is zero.  The command prints a wide hat over \(B\); the degree
subscript \(2\) in the generated fragment is appended externally rather than
passed as an argument.  The fixed evidence object is
\[
 B_2^{\mathrm{end}}
 =
 \begin{bmatrix}
  \frac{392}{9}&-\frac{496}{9}&\frac{176}{9}\\
  -\frac{752}{9}&\frac{1072}{9}&-\frac{392}{9}\\
  \frac{680}{9}&-\frac{1072}{9}&\frac{464}{9}\\
  -\frac{464}{9}&\frac{496}{9}&-\frac{104}{9}
 \end{bmatrix}
 \in\RationalNumbers^{4\times3}.
\]
It is the endpoint-basis coefficient matrix for degree \(d=2\), scale
\(M=2^d=4\), nodes \(x=(-1,0,1)\), transported nodes
\(t=(0,\tfrac12,1)\), and retained integer centers \(k=(1,2,3,4)\).
Its columns are the exact endpoint-atom coefficient vectors of the three
Lagrange cardinal polynomials.  The generator constructs those cardinals,
applies the reciprocal-moment differential operator at scale \(4\), samples
the active shifts, removes the annihilator-recurrence null component, and
retains the final \(d+2\) coefficients.  Status is an exact generated
companion-evidence fragment, not an independent globally parameterized
matrix family.  All degree, node, scale, basis, and row/column data are
suppressed.  The accent means endpoint augmentation/compression.  The external
degree subscript is presently valid because the zero-argument expansion owns
no script; regeneration must preserve that invariant or replace the command
with an indexed wrapper.

\paragraph{\localaccentname{AugmentedEvaluationMatrix}}
TeX arity is zero.  The command prints a wide hat over \(U\); entry indices
such as \(i,r\) are appended externally.  Fix \(d\in\NaturalNumbers\),
distinct nodes \(x_0,\ldots,x_d\in[-1,1]\), and the \(d+2\) normalized endpoint
atoms
\[
 c_{d,r}=2^{d+1}-2d-3+2r,\qquad
 \Phi_{d,r}(x)
 \DefinitionEquals
 2^{-d}\RvachevUp\!\left(\frac{x-c_{d,r}}{2^{d+1}}\right),
 \qquad0\le r\le d+1.
\]
The matrix is
\[
 U_X^{\mathrm{end}}
 \DefinitionEquals
 [\,\Phi_{d,r}(x_i)\,]_{
 0\le i\le d,\ 0\le r\le d+1}
 \in\RealNumbers^{(d+1)\times(d+2)}.
\]
It maps an endpoint-basis coefficient vector to its \(d+1\) nodal values.
Status is an exact finite evaluation matrix.  The degree, ordered node set,
endpoint basis, index ranges, and normalization of the atoms are suppressed.
The accent means augmentation of the polynomial evaluation system by one
endpoint atom.  A single external entry subscript is valid because the
zero-argument expansion has no internal script; callers must not create a
second subscript, and a future scripted body would require an entry wrapper.

\paragraph{\localaccentname{AugmentedCoefficientMatrix}}
TeX arity is zero.  The command prints a wide hat over \(B\); transpose or
entry scripts are external.  With the degree, nodes, atoms, and evaluation
matrix fixed as above, let \(\ell_j\in\Pi_d\) be the Lagrange cardinal
satisfying \(\ell_j(x_i)=\delta_{ij}\).  The coefficient matrix
\[
 B_X^{\mathrm{end}}
 =
 [\,b_{rj}\,]_{
 0\le r\le d+1,\ 0\le j\le d}
 \in\RealNumbers^{(d+2)\times(d+1)}
\]
is defined by the unique endpoint-basis expansions
\[
 \ell_j=\sum_{r=0}^{d+1}b_{rj}\Phi_{d,r},
 \qquad
 U_X^{\mathrm{end}}B_X^{\mathrm{end}}=I_{d+1}.
\]
Define the canonical cofactor vectors by
\[
 v_r=(-1)^r\det U_X^{(r)},\qquad
 \zeta_r=(-1)^r\det B_X^{(r)},
\]
where \(U_X^{(r)}\) deletes column \(r\) from \(U_X^{\mathrm{end}}\) and
\(B_X^{(r)}\) deletes row \(r\) from \(B_X^{\mathrm{end}}\).  Then
\(U_X^{\mathrm{end}}v=0\), \(\zeta^TB_X^{\mathrm{end}}=0\), and Cauchy--Binet
gives the normalization \(\zeta^Tv=1\).
With
\[
 \mathcal M_X=
 \begin{pmatrix}U_X^{\mathrm{end}}\\ \zeta^T\end{pmatrix},
\]
then the defect-completed inverse formula is
\[
 B_X^{\mathrm{end}}
 =
 \mathcal M_X^{-1}
 \begin{pmatrix}I_{d+1}\\0\end{pmatrix},
 \qquad
 \zeta^TB_X^{\mathrm{end}}=0.
\]
Status is an exact finite right inverse selected by the endpoint-basis
polynomiality constraint.  The degree, nodes, cardinals, endpoint atoms,
polynomiality row, and matrix dimensions are suppressed.  The accent means
endpoint-atom augmentation.  An external transpose is valid only because the
zero-argument expansion has no script; do not combine it with an internal or
second external superscript.

\paragraph{\localaccentname{OmittedMatrixIndex}}
TeX arity is one; argument \(1\) is the index
\(r\in\{0,\ldots,d+1\}\) to omit.  The command prints a wide hat over \(r\).
It is deletion metanotation, not a map on indices.  In the evaluation-matrix
context,
\[
 U_X^{(r)}
 \in\RealNumbers^{(d+1)\times(d+1)}
\]
means \(U_X^{\mathrm{end}}\) with column \(r\) deleted; in the coefficient-matrix
context,
\[
 B_X^{(r)}
 \in\RealNumbers^{(d+1)\times(d+1)}
\]
means \(B_X^{\mathrm{end}}\) with row \(r\) deleted.  The cofactor vectors are
\[
 v_r=(-1)^r\det U_X^{(r)},\qquad
 \zeta_r=(-1)^r\det B_X^{(r)},
\]
and Cauchy--Binet yields \(\zeta^Tv=1\).
Status is exact maximal-minor notation.  Which matrix is present, whether a
row or column is deleted, the matrix dimensions, and the permitted index set
are suppressed.  The accent means omission.  The argument is used as a nested
marker inside the matrix's one subscript; it must not itself acquire any
further script.

\paragraph{\localaccentname{LiftedCoefficientProjector}}
TeX arity is two.  Argument \(1\) is the interpolation level \(q\), and
argument \(2\) is the terminal coefficient-space level \(s\), in that order.
The command prints a wide hat over \(\mathcal P\), with subscript \(q\) and
parenthesized superscript \(s\).  Fix nested node sets
\(X^{(0)}\subset\cdots\subset X^{(s)}\), interpolation projectors
\(I_q:C[-1,1]\to\Pi_{d_q}\), a terminal coefficient space \(C_s\), a synthesis
map \(T_s:C_s\to C[-1,1]\), and a canonical coefficient map
\(E_s:\Pi_{d_s}\to C_s\) satisfying \(T_sE_s=I\) on \(\Pi_{d_s}\).  Then
\[
 \mathcal P_q^{(s)}
 \DefinitionEquals
 E_sI_qT_s:C_s\longrightarrow C_s,
 \qquad0\le q\le s.
\]
These operators satisfy
\[
 \mathcal P_q^{(s)}\mathcal P_r^{(s)}
 =
 \mathcal P_{\min(q,r)}^{(s)},
 \qquad
 \RankOperator\mathcal P_q^{(s)}=d_q+1.
\]
Status is an exact coefficient-space lift of an interpolation projector.
The nested grids, degrees, coefficient dictionary, synthesis and analysis
maps, and terminal level are suppressed.  The accent means coefficient lift.
Both available scripts are already built from the two arguments; no raw
subscript, superscript, or prime may be appended.

\paragraph{\localaccentname{LiftedCoefficientDetail}}
TeX arity is two.  Argument \(1\) is the detail level \(q\), and argument
\(2\) is the terminal coefficient-space level \(s\), in that order.  The
command prints a wide hat over \(\Delta\), with subscript \(q\) and
parenthesized superscript \(s\).  Using the lifted projectors just defined,
\[
 \Delta_0^{(s)}
 \DefinitionEquals
 \mathcal P_0^{(s)},
 \qquad
 \Delta_q^{(s)}
 \DefinitionEquals
 \mathcal P_q^{(s)}-\mathcal P_{q-1}^{(s)}
 \quad(1\le q\le s).
\]
Each \(\Delta_q^{(s)}:C_s\to C_s\) is idempotent, distinct details annihilate
one another, and
\[
 \RankOperator\Delta_0^{(s)}=d_0+1,\qquad
 \RankOperator\Delta_q^{(s)}=d_q-d_{q-1}\quad(q\ge1).
\]
Status is an exact finite multiresolution-detail definition.  The nested node
sets, interpolation projectors, terminal dictionary, degree sequence, and
coefficient-space maps are suppressed.  The accent means coefficient-space
lift, not frequency transformation.  Both scripts are owned by the two
arguments; appending any further raw script is invalid.

\subsubsection{Monic and hypothetical polynomial families}

\paragraph{\localaccentname{MonicDeconvolvedLegendre}}
TeX arity is one; argument \(1\) is the degree
\(n\in\NaturalNumbers\).  The command prints a wide hat over \(Q\), with
 subscript \(n\) and superscript minus.  Let
\[
 Z\coloneqq\sum_{j\ge1}2^{-j}U_j,
 \qquad U_j\stackrel{\mathrm{iid}}{\sim}\operatorname{Unif}[-1,1],
\]
so \(Z\) has the normalized Rvachev up-density on \([-1,1]\).  Let \(P_n\)
be the classical Legendre polynomial and set
\[
 M_u(z)
 \DefinitionEquals
 \Expectation[\EulerE^{zZ}]
 =
 \prod_{j\ge1}
 \frac{\sinh(z/2^j)}{z/2^j},
 \qquad
 \frac1{M_u(z)}
 =
 \sum_{r\ge0}\gamma_r\frac{z^r}{r!},
 \qquad
 \mathscr D_u
 \DefinitionEquals
 \sum_{r\ge0}\gamma_r\frac{D^r}{r!}.
\]
The deconvolved polynomial is \(Q_n^-=\mathscr D_uP_n\).  Since
\([x^n]P_n=2^{-n}\binom{2n}{n}\) and \(\mathscr D_u\) preserves leading
coefficients, its exact monic normalization is
\[
 Q_n^{-,\mathrm{mon}}
 \DefinitionEquals
 \frac{2^n}{\binom{2n}{n}}\,
 \mathscr D_uP_n
 \in\RationalNumbers[x].
\]
For example,
\[
 Q_2^{-,\mathrm{mon}}=x^2-\frac49,\qquad
 Q_3^{-,\mathrm{mon}}=x^3-\frac{14}{15}x,
\]
\[
 Q_4^{-,\mathrm{mon}}
 =x^4-\frac{32}{21}x^2+\frac{1072}{4725}
 =xQ_3^{-,\mathrm{mon}}
  -\frac{62}{105}Q_2^{-,\mathrm{mon}}
  -\frac8{225}Q_0^{-,\mathrm{mon}}.
\]
Status is an exact rational polynomial definition; the final equality is the
explicit degree-four Favard obstruction.  The up-law, moment-generating
function, reciprocal moments, differential operator, Legendre normalization,
and coefficient field are suppressed.  The accent means monic normalization.
The command already owns its degree subscript and minus superscript; no
additional raw script may follow it.

\paragraph{\localaccentname{GenericMonicOrthogonalPolynomial}}
TeX arity is one; argument \(1\) is the degree
\(n\in\NaturalNumbers\).  The command prints a wide hat over \(Q\) with
subscript \(n\).  It is a schema for a hypothetical monic,
parity-alternating scalar orthogonal-polynomial family
\[
 Q_n^{\mathrm{mon}}\in\RealNumbers[x],\qquad
 \deg Q_n^{\mathrm{mon}}=n,\qquad
 [x^n]Q_n^{\mathrm{mon}}=1.
\]
For a scalar quasi-definite moment functional, Favard's theorem would force
\[
 Q_0^{\mathrm{mon}}=1,\qquad Q_1^{\mathrm{mon}}=x,
 \qquad \lambda_n\ne0\quad(n\ge1),
\]
and the parity-reduced three-term recurrence
\[
 Q_{n+1}^{\mathrm{mon}}
 =
 xQ_n^{\mathrm{mon}}-\lambda_nQ_{n-1}^{\mathrm{mon}}
 \qquad(n\ge1).
\]
The symbol does not name a separately constructed global family; it is used
only to contrast this three-term schema with the nonzero degree-zero defect
in the actual monic deconvolved Legendre family.  Status is therefore a
generic proof schema, not a definition of new canonical polynomials.  The
scalar functional, orthogonality normalization, recurrence coefficients,
coefficient field, and parity assumptions are suppressed.  The accent means
monic normalization, not Fourier transformation.  The degree is already
placed by the one argument, so no second subscript or conflicting superscript
may be appended.

\subsubsection{Lacunary-product extrapolants}

\paragraph{Common lacunary-product data.}
For the product extrapolants below, let \(b>1\), \(N\in\NaturalNumbers\),
\(q=b^{-2}\in(0,1)\), and
\[
 P_b(t)\coloneqq\prod_{k=1}^{\infty}\SincRad(t/b^k),\qquad
 P_{b,N}(t)\coloneqq\prod_{k=1}^{N}\SincRad(t/b^k),
\]
where \(P_{b,0}=1\).  Put
\[
 T_N(t)\coloneqq\prod_{k=N+1}^{\infty}\SincRad(t/b^k),\qquad
 T_{N,j}(t)\coloneqq\prod_{k=N+1}^{N+j}\SincRad(t/b^k),
 \qquad T_{N,0}=1,
\]
and
\[
 D_N\coloneqq
 \{t\in\ComplexNumbers:\AbsoluteValue{t}<\pi b^{N+1}\}.
\]
Every \(T_{N,j}\), including \(T_N\), is zero-free on \(D_N\).  Write
\(L_{N,j}\) for the unique holomorphic logarithm of \(T_{N,j}\) on \(D_N\)
satisfying \(L_{N,j}(0)=0\).  For \(r\in\NaturalNumbers\) and
\(0\le j\le r\), define
\[
 (q;q)_m\coloneqq\prod_{\ell=1}^{m}(1-q^\ell),\qquad (q;q)_0\coloneqq1,
 \qquad
 \lambda_{r,j}(q)\coloneqq
 \frac{(-1)^{r-j}q^{(r-j)(r-j+1)/2}}
      {(q;q)_j(q;q)_{r-j}}.
\]

\begin{catalogueentry}{local-hat-factorized-tail}
{\localaccentname{FactorizedTailExtrapolant}: factorized tail extrapolant}
\entryfield{TeX interface}{Arity \(2\), with arguments \(N,r\) in that order;
it prints \(\widehat T_{N,r}\).  The evaluation variable \(t\) is written
after the macro and is not a TeX argument.}
\entryfield{Definition and type}{For \(N,r\in\NaturalNumbers\),
\[
 \widehat T_{N,r}(t)\coloneqq
 \exp\!\left(\sum_{j=0}^{r}\lambda_{r,j}(q)L_{N,j}(t)\right),
 \qquad
 \widehat T_{N,r}:D_N\longrightarrow\ComplexNumbers^\times.
\]
It extrapolates the exact omitted tail \(T_N\), is holomorphic and zero-free
on \(D_N\), and is positive on \(D_N\cap\RealNumbers\).}
\entryfield{Scope and status}{The printed notation suppresses the fixed base
\(b\), hence also \(q=b^{-2}\), and leaves \(t\) external.  The construction,
its exact error, and its zero-free properties are proved in the factorized
frontier report; that report claims no Lean counterpart at its snapshot.}
\entryfield{Accent and scripts}{The hat means logarithmic \(q\)-Richardson
extrapolation, not Fourier transformation.  The symbol owns only the
subscript \(N,r\), so a later grouped or direct superscript is syntactically
available.}
\end{catalogueentry}

\begin{catalogueentry}{local-hat-factorized-product}
{\localaccentname{FactorizedProductExtrapolant}: zero-preserving product
extrapolant}
\entryfield{TeX interface}{Arity \(3\), with arguments \(b,N,r\) in that
order; it prints \(\widehat P^{\mathrm{fac}}_{b,N,r}\).  The evaluation
variable \(t\) is external.}
\entryfield{Definition and type}{
\[
 \widehat P^{\mathrm{fac}}_{b,N,r}(t)
 \coloneqq
 P_{b,N}(t)
 \exp\!\left(\sum_{j=0}^{r}\lambda_{r,j}(b^{-2})L_{N,j}(t)\right),
 \qquad
 \widehat P^{\mathrm{fac}}_{b,N,r}:D_N\longrightarrow\ComplexNumbers.
\]
It has exactly the zero divisor of \(P_b\) in \(D_N\), including
multiplicities, and agrees with the unfactorized logarithmic product
extrapolant wherever \(P_{b,N}\ne0\).}
\entryfield{Scope and status}{The derived parameter \(q=b^{-2}\) and the
argument \(t\) are not printed as indices.  Holomorphic continuation,
zero-divisor preservation, exact zero-normalized quotients, and local jet
identities are proved in the frontier report but are not claimed there as
Lean theorems.}
\entryfield{Accent and scripts}{The hat means a factorized approximation to
\(P_b\), not a Fourier transform.  The macro already owns the superscript
\(\mathrm{fac}\); any derivative, power, or further semantic superscript must
be attached to the grouped expression
\((\widehat P^{\mathrm{fac}}_{b,N,r})\), never appended directly.}
\end{catalogueentry}

\begin{catalogueentry}{local-hat-relative-tail}
{\localaccentname{RelativeTailExtrapolant}: dyadic relative-tail
extrapolant}
\entryfield{TeX interface}{Arity \(2\), with arguments \(N,r\) in that order;
it prints \(\widehat{\mathcal T}_{N,r}\).  The evaluation variable \(x\) is
external, and the dyadic base is fixed rather than printed.}
\entryfield{Definition and type}{Put
\[
 \mathcal T_{N,j}(x)\coloneqq
 \prod_{h=N+1}^{N+j}\SincRad(\pi x/2^h),\qquad
 \mathcal T_{N,0}=1,
\]
and let \(L^{\mathrm{dy}}_{N,j}\) be its holomorphic logarithm on
\(\AbsoluteValue{x}<2^{N+1}\), normalized by
\(L^{\mathrm{dy}}_{N,j}(0)=0\).  Then
\[
 \widehat{\mathcal T}_{N,r}(x)\coloneqq
 \exp\!\left(
   \sum_{j=0}^{r}\lambda_{r,j}(1/4)L^{\mathrm{dy}}_{N,j}(x)
 \right),
 \qquad
 \widehat{\mathcal T}_{N,r}:
 \{x\in\ComplexNumbers:\AbsoluteValue{x}<2^{N+1}\}
 \longrightarrow\ComplexNumbers^\times.
\]
Its target is
\(\mathcal T_N(x)=\prod_{h=N+1}^{\infty}\SincRad(\pi x/2^h)\).}
\entryfield{Scope and status}{The fixed data \(b=2\), \(q=1/4\), and the
argument \(x\) are suppressed.  The source introduces this object by
cross-reference rather than by a standalone equality; the formula above is
the explicit specialization used in its proved dyadic bracket theorem.}
\entryfield{Accent and scripts}{The hat means extrapolation of a product tail
that is zero-free on the complex disk and positive on its real slice, not Fourier
transformation.  The symbol owns only the subscript \(N,r\).}
\end{catalogueentry}

\begin{catalogueentry}{local-hat-dyadic-factorized-product}
{\localaccentname{DyadicFactorizedProductExtrapolant}: zero-preserving dyadic
Fourier-product approximant}
\entryfield{TeX interface}{Arity \(2\), with arguments \(N,r\) in that order;
it prints \(\widehat\Phi_{N,r}\).  The evaluation variable \(x\) is external.}
\entryfield{Definition and type}{Let
\[
 \Phi(x)\coloneqq\FourierTwoPi{\RvachevUp}(x)
 =\prod_{h=0}^{\infty}\SincRad(\pi x/2^h),\qquad
 \Phi_N(x)\coloneqq\prod_{h=0}^{N}\SincRad(\pi x/2^h).
\]
Then
\[
 \widehat\Phi_{N,r}(x)\coloneqq
 \Phi_N(x)\widehat{\mathcal T}_{N,r}(x),
 \qquad
 \widehat\Phi_{N,r}:
 \{x\in\ComplexNumbers:\AbsoluteValue{x}<2^{N+1}\}
 \longrightarrow\ComplexNumbers.
\]
It preserves every included integer zero of \(\Phi\), its multiplicity, and
the real Thue--Morse lobe sign.}
\entryfield{Scope and status}{The fixed dyadic data \(b=2\), \(q=1/4\), the
relative-tail logarithms, and \(x\) are suppressed.  The source defines the
object by reference to factorization; its zero, jet, bracket, and sign
properties are proved in the manuscript but are not claimed as Lean results.}
\entryfield{Accent and scripts}{Although \(\Phi\) is itself the
\(2\pi\)-normalized Fourier transform of \(\RvachevUp\), the hat on
\(\widehat\Phi_{N,r}\) means numerical/analytic extrapolation of \(\Phi\), not
a second Fourier transform.  The macro owns only a subscript, so the current
derivative notation \(\widehat\Phi_{N,r}^{(\mu)}\) is syntactically safe.}
\end{catalogueentry}

\subsubsection{Radial--angular and unfactorized product estimators}

\begin{catalogueentry}{local-hat-radial-jet}
{\localaccentname{RadialJetEstimator}: radial endpoint-coefficient estimator}
\entryfield{TeX interface}{Arity \(1\), whose sole argument is the coefficient
index \(j\); it prints \(\widehat A_j\).  The phase \(\theta\) is an external
function argument.}
\entryfield{Definition and type}{Let \(a,d\in\IntegerNumbers\) satisfy
\(d\ge1\) and \(0\le a<d\), put \(\theta=a/d\), and let the integer \(\rho>1\)
satisfy \(\rho\equiv1\pmod d\).  Choose \(n\in\PositiveIntegers\) with
\(n+\theta\ge1/\log2\).  For fixed \(p\in\NaturalNumbers\), put
\[
 q=\rho^{-1},\qquad H=(n+\theta)^{-1},\qquad
 y_s=\mathscr E(\rho^s(n+\theta))
     =\sum_{m=0}^{p}A_m(\theta)H^m q^{sm}+r_s,
\]
and
\[
 \ell_s(z)\coloneqq
 \prod_{\substack{0\le u\le p\\u\ne s}}
 \frac{z-q^u}{q^s-q^u}.
\]
For \(0\le j\le p\),
\[
 \widehat A_j(\theta)\coloneqq
 H^{-j}[z^j]\sum_{s=0}^{p}y_s\ell_s(z).
\]
For the real Fabius endpoint expansion this is real-valued; the same finite
interpolation formula is valid over \(\ComplexNumbers\).}
\entryfield{Scope and status}{The printed symbol suppresses
\(p,n,\rho,q,H\), the sample vector \((y_s)_{s=0}^{p}\), and the normalized
endpoint residual \(\mathscr E\).  Exact recovery in the polynomial model and
the remainder bound are conditional on this externally supplied normalization
and sample vector: the source intentionally leaves the precise nonperiodic
subtraction inside \(\mathscr E\) schematic.  The resulting conditional theorem
is proved in the tomography manuscript, not formalized there in Lean.}
\entryfield{Accent and scripts}{The hat means recovery of \(A_j(\theta)\) by
geometric radial interpolation.  It is not a Fourier transform.  The symbol
owns only the subscript \(j\).}
\end{catalogueentry}

\begin{catalogueentry}{local-hat-tomographic-coefficient}
{\localaccentname{TomographicCoefficientEstimator}: radial--angular
Fourier-coefficient estimator}
\entryfield{TeX interface}{Arity \(1\), whose sole macro argument is the
Fourier mode \(k\); the base form is \(\widehat c_k\).  The source appends the
external configuration superscript \((d,p,n)\), printing
\(\widehat c_k^{(d,p,n)}\).}
\entryfield{Definition and type}{If
\[
 A_0(u)=\sum_{m\in\IntegerNumbers}c_m\EulerE^{2\pi\ImaginaryUnit m u},
\qquad
 \mathcal R_{0;p,n,\rho}(\theta)
 \coloneqq [z^0]\sum_{s=0}^{p}y_s\ell_s(z),
\]
where \(y_s,\ell_s\) are the radial data in the preceding entry, then
\[
 \widehat c_k^{(d,p,n;\rho)}
 \coloneqq
 \frac1d\sum_{a=0}^{d-1}
 \mathcal R_{0;p,n,\rho}(a/d)
 \EulerE^{-2\pi\ImaginaryUnit k a/d}
 \in\ComplexNumbers.
\]
The semicolon-\(\rho\) augmentation is catalogue-only disambiguation: the active
source prints \(\widehat c_k^{(d,p,n)}\) and suppresses \(\rho\).
The theorem uses \(d\ge2\), \(p\in\NaturalNumbers\),
\(n\in\PositiveIntegers\), \(k\in\IntegerNumbers\),
\(\AbsoluteValue{k}\le(d-1)/2\), and an integer
\(\rho>1\) satisfying \(\rho\equiv1\pmod d\).}
\entryfield{Scope and status}{The printed superscript records \(d,p,n\) but
suppresses \(\rho\), the radial samples, and the endpoint normalization.  Under
the stated uniform endpoint expansion and Wiener-strip hypothesis, the report proves
\[
 \widehat c_k^{(d,p,n)}-c_k
 =\BigO(n^{-p-1})
  +\BigO\!\left(\exp[-2\pi\sigma(d-\AbsoluteValue{k})]\right)
\]
as \(n\to\infty\), uniformly over the grid phases \(a/d\), with
\(d,p,\rho\), and the strip width \(\sigma\) fixed.  It does not claim a Lean
formalization.}
\entryfield{Accent and scripts}{The hat means an estimate of the already
defined Fourier coefficient \(c_k\).  The finite Fourier transform occurs
inside the construction, but the hat is not the Fourier transform of \(c_k\).
Because the macro owns only \(k\) as a subscript, the externally appended
configuration superscript is currently safe; changing the macro to own a
superscript would require a dedicated wrapper.}
\end{catalogueentry}

\begin{catalogueentry}{local-hat-geometric-product}
{\localaccentname{GeometricProductExtrapolant}: multiplicative
\(q\)-Richardson product extrapolant}
\entryfield{TeX interface}{Arity \(3\), with arguments \(b,N,r\) in that
order; it prints \(\widehat P_{b;N,r}\).  The semicolon separates the base
from the two discrete indices.  The evaluation variable \(t\) is external.}
\entryfield{Definition and type}{For real \(t\) satisfying
\[
 \frac{\AbsoluteValue{t}}{b^{N+1}}<\pi,\qquad P_{b,N}(t)\ne0,
\]
define
\[
 \widehat P_{b;N,r}(t)\coloneqq
 \Sign(P_{b,N}(t))
 \exp\!\left(
   \sum_{j=0}^{r}\lambda_{r,j}(b^{-2})
   \log\AbsoluteValue{P_{b,N+j}(t)}
 \right)\in\RealNumbers^\times.
\]
Equivalently,
\[
 \AbsoluteValue{\widehat P_{b;N,r}(t)}
 =
 \prod_{j=0}^{r}
 \AbsoluteValue{P_{b,N+j}(t)}^{\lambda_{r,j}(b^{-2})}.
\]}
\entryfield{Scope and status}{The derived value \(q=b^{-2}\) and \(t\) are
not printed as indices.  This is the real logarithmic construction away from
prefix zeros, unlike the holomorphic factorized continuation.  Its exact
error series, sign, brackets, and fixed-parameter convergence are proved in
the report; the associated analytic claims remain separate Lean obligations.}
\entryfield{Accent and scripts}{The hat means multiplicative extrapolation of
\(P_b(t)\), not Fourier transformation.  The macro owns only a subscript, so
the hybrid-correction superscript \([s]\) used in the source is syntactically
safe and denotes correction depth, not differentiation.}
\end{catalogueentry}

\begin{catalogueentry}{local-hat-dyadic-product}
{\localaccentname{DyadicProductExtrapolant}: dyadic multiplicative product
extrapolant}
\entryfield{TeX interface}{Arity \(2\), with arguments \(N,r\) in that order;
it prints \(\widehat P_{N,r}\).  The base \(b=2\) is fixed, and \(t\) is an
external function argument.}
\entryfield{Definition and type}{With
\[
 P_N(t)\coloneqq\prod_{k=1}^{N}\SincRad(t/2^k),
\]
define, whenever
\(\AbsoluteValue{t}/2^{N+1}<\pi\) and \(P_N(t)\ne0\),
\[
 \widehat P_{N,r}(t)\coloneqq
 \Sign(P_N(t))
 \exp\!\left(
   \sum_{j=0}^{r}\lambda_{r,j}(1/4)
   \log\AbsoluteValue{P_{N+j}(t)}
 \right)\in\RealNumbers^\times.
\]
Thus it is exactly the \(b=2\), \(q=1/4\) specialization of the preceding
geometric estimator.}
\entryfield{Scope and status}{The base, \(q=1/4\), and \(t\) are suppressed.
The target is the tail product \(P_2(t)\), not the full Fourier product
\(\Phi(x)=\SincRad(\pi x)P_2(\pi x)\).  Explicit low-order formulas, error
bounds, and the finite Thue--Morse realization are proved in the report.  The raw
finite complex-exponential identity and the sine/sinc Thue--Morse normalization
are Lean-formalized in the complex-product bridge modules; the analytic Richardson
estimator, error, bracket, and convergence claims are not.}
\entryfield{Accent and scripts}{The hat means dyadic \(q\)-Richardson
extrapolation, not Fourier transformation.  The symbol owns only the
subscript \(N,r\).}
\end{catalogueentry}

\subsubsection{Phase-locked lower-Lambert estimators}

\paragraph{Common phase-locked data.}
Let \(L=\log2\).  For \(\lambda>1/L\), set
\[
 x(\lambda)\coloneqq\lambda2^{-\lambda},
 \qquad
 \lambda(x)\coloneqq-\frac1L\LambertW_{-1}(-Lx)
 \quad\left(0<x<\frac1{\EulerE L}\right),
\]
so that \(\lambda(x(\lambda))=\lambda\).  The phase-locked nodes and weights are
\[
 x_j(\lambda)\coloneqq(\lambda+j)2^{-(\lambda+j)},\qquad
 \omega_{r,j}(\lambda)\coloneqq
 \frac{(-1)^{r-j}}{r!}\binom rj(\lambda+j)^r
 \quad(0\le j\le r).
\]
They satisfy
\[
 \sum_{j=0}^{r}\omega_{r,j}(\lambda)=1,\qquad
 \sum_{j=0}^{r}\omega_{r,j}(\lambda)(\lambda+j)^{-m}=0
 \quad(1\le m\le r).
\]
With the endpoint theorem's fixed constant \(C_F\), define
\[
 \mathcal W_{\mathrm{MSE}}(x)
 \coloneqq
 -\frac L2\lambda(x)^2+
 \left(1+\frac L2\right)\lambda(x)
 -\frac12\log\lambda(x)+C_F,
 \qquad
 Y(x)\coloneqq
 \log\FabiusBounded(x)-\mathcal W_{\mathrm{MSE}}(x).
\]
For every fixed \(M\in\NaturalNumbers\), as \(x\downarrow0\) (equivalently
\(\lambda(x)\to\infty\)), uniformly along a fixed compact phase set, the
all-orders input has
\[
 Y(x)=\Psi(\lambda(x))
 +\sum_{m=1}^{M}A_m(\lambda(x))\lambda(x)^{-m}
 +\BigO(\lambda(x)^{-M-1}),
\]
where \(\Psi\) and every \(A_m\) are one-periodic in \(\lambda\).

\begin{catalogueentry}{local-hat-phase-locked-log}
{\localaccentname{PhaseLockedLogEstimator}: phase-locked logarithmic endpoint
estimator}
\entryfield{TeX interface}{Arity \(1\), whose argument is the extrapolation
order \(r\); it prints \(\widehat\Psi_r\).  The lower-Lambert coordinate
\(\lambda\) is an external function argument.}
\entryfield{Definition and type}{For \(r\in\NaturalNumbers\) and
\(\lambda>1/\log2\),
\[
 \widehat\Psi_r(\lambda)\coloneqq
 \sum_{j=0}^{r}\omega_{r,j}(\lambda)Y(x_j(\lambda))
 \in\RealNumbers.
\]
The weights preserve the periodic constant term and cancel the first \(r\)
inverse-\(\lambda\) corrections exactly, so the target is
\(\Psi(\lambda)\).}
\entryfield{Scope and status}{The printed form suppresses \(\lambda\), the
nodes \(x_j\), the reduced observable \(Y\), the main-saddle normalization
\(C_F\), and the weights.  At each fixed \(r\) and finite residual depth, the
logarithmic extraction and its first-omitted-term estimate are Lean-backed
under the repository's all-orders endpoint input.  No convergence of the
infinite residual series and no uniformity in growing \(r\) are asserted.}
\entryfield{Accent and scripts}{The hat means a phase-locked reciprocal-scale
estimate of the one-periodic correction \(\Psi(\lambda)\), not a Fourier
transform.  The symbol owns only the subscript \(r\).}
\end{catalogueentry}

\begin{catalogueentry}{local-hat-phase-locked-exponential}
{\localaccentname{PhaseLockedExponentialEstimator}: exponentiated
phase-locked endpoint estimator}
\entryfield{TeX interface}{Arity \(1\), whose argument is the extrapolation
order \(r\); it prints \(\widehat{\mathcal A}_r\).  The coordinate
\(\lambda\) is external.}
\entryfield{Definition and type}{For \(r\in\NaturalNumbers\),
\(\lambda>1/\log2\), and the common nodes, weights, and reduced observable,
\[
 \widehat{\mathcal A}_r(\lambda)\coloneqq
 \exp\!\left(
   \sum_{j=0}^{r}\omega_{r,j}(\lambda)Y(x_j(\lambda))
 \right)
 =
 \prod_{j=0}^{r}
 \left(
   \frac{\FabiusBounded(x_j(\lambda))}
        {\exp(\mathcal W_{\mathrm{MSE}}(x_j(\lambda)))}
 \right)^{\omega_{r,j}(\lambda)}
 \in\RealNumbers_{>0}.
\]
Its exact target is the periodic endpoint amplitude
\(\mathcal A(\lambda)\coloneqq\exp(\Psi(\lambda))\).}
\entryfield{Scope and status}{The printed form suppresses \(\lambda\),
\(x_j\), \(Y\), \(\mathcal W_{\mathrm{MSE}}\), and \(\omega_{r,j}\).  The
displayed equality is exact exponentiation of the logarithmic estimator; the
multiplicative estimator's Bell relative-error hierarchy is not part of the
closed Lean phase-extraction tranche.}
\entryfield{Accent and scripts}{The hat means an estimated positive periodic
amplitude, not Fourier transformation.  The symbol owns only the subscript
\(r\).}
\end{catalogueentry}

\begin{catalogueentry}{local-hat-phase-locked-coefficient}
{\localaccentname{PhaseLockedCoefficientEstimator}: phase-locked endpoint
coefficient extractor}
\entryfield{TeX interface}{Arity \(2\), with arguments \(s,r\) in that exact
order: \(s\) is the coefficient index and \(r\) is the interpolation order.
It prints \(\widehat A_s^{(r)}\).  The coordinate \(\lambda\) is external.}
\entryfield{Definition and type}{For \(\lambda>1/\log2\), set
\[
 h_j\coloneqq(\lambda+j)^{-1},\qquad
 \ell_j(h)\coloneqq
 \prod_{\substack{0\le k\le r\\k\ne j}}
 \frac{h-h_k}{h_j-h_k}.
\]
For \(r\in\NaturalNumbers\) and \(0\le s\le r\),
\[
 \widehat A_s^{(r)}(\lambda)\coloneqq
 \frac1{s!}\sum_{j=0}^{r}\ell_j^{(s)}(0)Y(x_j(\lambda))
 \in\RealNumbers.
\]
The factor \(1/s!\) makes this the coefficient of \(h^s\) in the interpolating
polynomial, rather than its unnormalized \(s\)th derivative.  Locally,
\(\widehat A_0^{(r)}=\widehat\Psi_r\); for \(s\ge1\), the target is the
periodic endpoint coefficient \(A_s(\lambda)\).}
\entryfield{Scope and status}{The printed symbol suppresses \(\lambda\), the
nodes, the observable \(Y\), and the Lagrange basis.  The extraction formula
is exact finite interpolation, but the higher coefficient extractors remain
a formalization target; only the logarithmic \(s=0\) tranche has the compiled
fixed-order status described above.}
\entryfield{Accent and scripts}{The hat means coefficient reconstruction, not
Fourier transformation.  The built-in superscript \((r)\) records
interpolation order and is not a derivative.  Any additional derivative,
power, or semantic superscript must be attached to the grouped expression
\((\widehat A_s^{(r)})\).}
\end{catalogueentry}

\subsubsection{Tail and inverse-root approximations}

\begin{catalogueentry}{local-hat-simple-tail}
{\localaccentname{SimpleTailEstimate}: first-order negative-power tail
estimate}
\entryfield{TeX interface}{Arity \(1\), whose argument is the tail-start
index \(N\); it prints \(\widehat R_N\).  The exponent parameter \(a\) is an
external function argument.}
\entryfield{Definition and type}{Let \(X\) be an \([0,1]\)-valued random variable
with distribution function
\(\FabiusBounded\), put \(d_n\coloneqq\Expectation X^n\), and, for \(a>0\),
define
\[
 (a)^{\overline k}\coloneqq
 \begin{cases}1,&k=0,\\a(a+1)\cdots(a+k-1),&k\ge1,\end{cases}
 \qquad
 \binom{a+k-1}{k}=\frac{(a)^{\overline k}}{k!},
\]
\[
 t_k(a)\coloneqq
 \binom{a+k-1}{k}\frac{d_{k+1}}{k+1},\qquad
 R_N(a)\coloneqq\sum_{k\ge N}t_k(a).
\]
With \(L=\log2\), the simple hazard estimate is
\[
 \widehat R_N(a)\coloneqq
 \frac{N\,t_N(a)}
      {\log N/L-(a-\tfrac12)}.
\]
It is a real number for \(N\in\PositiveIntegers\) whenever the denominator is
nonzero, and has the intended positive-tail interpretation for \(N\)
sufficiently large that the denominator is positive.}
\entryfield{Scope and status}{The printed form suppresses \(a\), the moment
sequence \(d_n\), and the fact that the more precise conjectural hazard also
contains \(-\Psi'(\log_2N)/L\).  This is a heuristic first-order estimator:
as \(N\to\infty\) with fixed \(a>0\), the displayed numerical ratios are
structural evidence, not certification,
and the corresponding uniform tail asymptotic remains conjectural.}
\entryfield{Accent and scripts}{The hat means approximation of the positive
series remainder \(R_N(a)\), not Fourier transformation.  The symbol owns
only the subscript \(N\).}
\end{catalogueentry}

\begin{catalogueentry}{local-hat-lambert-root}
{\localaccentname{LambertRootApproximation}: numerical Lambert-root
approximation}
\entryfield{TeX interface}{Arity \(0\); it prints the ambient scalar
\(\widehat w\).  The input \(x\), branch \(k\), numerical method, and iteration
count are not TeX arguments.}
\entryfield{Definition and domain}{The exact target is
\[
 r=\LambertW_k(x),\qquad r\EulerE^r=x,
\]
where
\[
 \LambertW_0:[-1/\EulerE,\infty)\longrightarrow[-1,\infty),
 \qquad
 \LambertW_{-1}:[-1/\EulerE,0)\longrightarrow(-\infty,-1].
\]
The value \(\widehat w\in\RealNumbers\) is any numerical approximation chosen
on the same monotone real inverse interval as \(r\); the symbol itself fixes
no iteration.  If \(I\) is the closed interval joining \(r\) and
\(\widehat w\), \(I\) avoids \(-1\), and
\[
 m_I\coloneqq
 \min_{t\in I}\AbsoluteValue{\EulerE^t(1+t)}>0,
\]
then the ordinary residual gives
\[
 \AbsoluteValue{\widehat w-r}
 \le
 \frac{\AbsoluteValue{\widehat w\EulerE^{\widehat w}-x}}{m_I}.
\]
For \(x\ne0\), if \(I\) also avoids \(0\), define
\[
 G_x(w)\coloneqq w+\log\AbsoluteValue{w}-\log\AbsoluteValue{x},
 \qquad
 \mu_I\coloneqq\min_{t\in I}\AbsoluteValue{1+1/t}>0.
\]
The logarithmic residual then gives the alternative certificate
\[
 \AbsoluteValue{\widehat w-r}
 \le\frac{\AbsoluteValue{G_x(\widehat w)}}{\mu_I}.
\]}
\entryfield{Scope and status}{The branch and input are mandatory ambient data.
Both residual inequalities are proved a posteriori theorems.  Near the branch
point their denominators correctly record the conditioning loss.  At
\(x=-1/\EulerE\), where \(r=-1\), the interval necessarily meets \(-1\), so
neither certificate applies; the guide returns \(-1\) exactly or changes to the
signed square-root branch coordinate.  No particular
Newton, Halley, series, or rational approximation is privileged.}
\entryfield{Accent and scripts}{The hat means an approximate inverse root, not
Fourier transformation.  The symbol owns no subscript or superscript, but any
later decoration must state whether it records branch, iteration, precision,
or another numerical parameter.}
\end{catalogueentry}

\subsection{Local alphabet register}

The following roles are intentionally local because assigning them global commands would
create more collisions than it removes.  Each occurrence must nevertheless instantiate
the declaration schema above.

\begin{longtable}{@{}p{0.22\textwidth}p{0.72\textwidth}@{}}
\toprule
Family & Mandatory local data\\
\midrule
\endhead
Polynomials $p_n,P_n,A_n$ & Coefficient field, variable, degree, leading-coefficient
normalization, orthogonality measure if any, and index origin.\\
Matrices $A_N,H_N,T_N,G_N$ & Row and column index types, dimensions, scalar field,
entry formula, and zero-size convention.\\
Kernels $K(x,y)$ & Two domain spaces, scalar codomain, symmetry/Hermitian convention,
regularity, and positive-definiteness status.\\
Measures $\mu,\nu$ & Underlying measurable space, positivity/signedness, total mass,
and density only if one exists.\\
Random variables $X,Y,U_j$ & Probability space or a statement that only their joint law
is relevant, codomain, independence/dependence, and distribution.\\
Generating functions $A(z),G(z),M(z)$ & Ordinary, exponential, Dirichlet, or other
series type; coefficient indexing; formal versus analytic interpretation; convergence
domain.\\
Recurrences & Initial values, index range on which the recurrence applies, coefficient
ring, and treatment of singular leading coefficients.\\
Optimization variables & Feasible set, topology/norm, objective codomain, existence and
uniqueness status, and whether infima are attained.\\
Asymptotic phases and amplitudes & Limit variable and direction, uniform parameters,
branch choices, order of truncation, and quantified remainder.\\
Frontier constructions & Complete proposed definition plus one of the status labels in
\cref{sec:status-register}; a suggestive name does not confer a theorem.\\
\bottomrule
\end{longtable}

\section{Retired-alias index}
\label{sec:retired-aliases}

This index is a reading aid for historical sources and a migration contract, not a second
namespace.  Names in the first column are printed through \commandname{commandname} so
that the retired control sequences are not executed.  Migration is semantic: determine
the old occurrence's type and normalization first, then choose the target.  A textual
replacement performed before that classification is invalid whenever the target column
contains ``or''.

\begingroup\small
\begin{longtable}{@{}p{0.29\textwidth}p{0.30\textwidth}p{0.35\textwidth}@{}}
\toprule
Retired source name(s) & Canonical target & Required disambiguation\\
\midrule
\endfirsthead
\toprule
Retired source name(s) & Canonical target & Required disambiguation\\
\midrule
\endhead
\commandname{R} & \commandname{RealNumbers} & Number field only; a radius remains a
typed local variable.\\
\commandname{C} & \commandname{ComplexNumbers} & Number field only; constants and
contours remain local.\\
\commandname{Q} & \commandname{RationalNumbers} & Number field only; do not rewrite a
quantile, quadrature rule, or polynomial.\\
\commandname{N} & \commandname{NaturalNumbers} or \commandname{PositiveIntegers} & Decide
from whether zero belongs to the set; never append another raw subscript to the chosen
command.\\
\commandname{Z} & \commandname{IntegerNumbers} & Number field only.\\
\commandname{Up}, \commandname{up}, \commandname{upf} & \commandname{RvachevUp} & The
even compactly supported function, not the bounded or global Fabius function.\\
\commandname{sinc} & \commandname{SincRad} or \commandname{SincPi} & Inspect the defining
formula or zero set.  Radian zeros are at nonzero multiples of $\pi$; normalized zeros
are at nonzero integers.\\
\commandname{sincpi}, \commandname{sincp}, \commandname{sincpiPartc} &
\commandname{SincPi} & These source-local names were explicitly normalized by $\pi$.\\
\commandname{sincPartx}, \commandname{sincPartl}, \commandname{sincPartii},
\commandname{sincPartxx}, \commandname{sincPartll}, \commandname{sincPartcc},
\commandname{sincPartvvv} & \commandname{SincRad} & These generated-fragment aliases
were classified from their defining formulas as radian sinc.\\
\commandname{wt} & \commandname{BinaryDigitSum} & Only for binary Hamming weight; a
general graph or code weight remains local.\\
\commandname{tm}, \commandname{TM} & \commandname{ThueMorseSignSymbol} or
\commandname{ThueMorseSign} & Use the symbol-only command in declarations and the
argument-taking command in formulas.\\
\commandname{e}, \commandname{ee} & \commandname{EulerE} & Only the exponential base;
an error variable remains local.\\
\commandname{ii} & \commandname{ImaginaryUnit} & The complex unit, not an index.\\
\commandname{dd} & \commandname{Differential} & Upright integral differential; a
derivative operator is declared separately.\\
\commandname{Li} & \commandname{Polylogarithm} & The special function, with its order
printed as a subscript.\\
\commandname{Var} & \commandname{Variance} & Probability variance only.\\
\commandname{Cov} & \commandname{Covariance} & Probability covariance only.\\
\commandname{E} & \commandname{Expectation} & Expectation only; events and energies are
local and semantically labelled.\\
\commandname{Prob}, \commandname{Pp} & \commandname{Probability} & Probability measure,
not polynomial or projection.\\
\commandname{bigO}, \commandname{Oh}, \commandname{OO},
\commandname{bigOPartl} &
\commandname{BigO} or \commandname{BigOAt} & Prefer
the explicit-regime command; retain the token-only form only when the same sentence
quantifies the limit.\\
\commandname{smallo}, \commandname{smalloPartl}, \commandname{littleo},
\commandname{littleoh} &
\commandname{LittleO} or \commandname{LittleOAt} & Prefer the explicit-regime command
and state uniformity parameters.\\
\commandname{supp} & \commandname{SupportOperator} or \commandname{SupportOf} & Use the
argument-delimited form in prose; distinguish pointwise, topological, measure, and
essential support.\\
\commandname{sgn} & \commandname{Sign} & Real sign with value zero at zero.\\
\commandname{TV} & \commandname{TotalVariationOf} or
\commandname{TotalVariationDistance} & Count the mathematical arguments: one typed
object means total variation; two probability measures mean the fixed distance.\\
\commandname{dist} & \commandname{MetricDistance} or \commandname{EqualInLaw} & Inspect
the relation: a numeric metric distance is not equality in distribution.\\
\commandname{ord} & \commandname{OrderOperator} & Multiplicative order, vanishing order,
or another order operator only after the local arguments define which notion is meant.\\
\commandname{lcm} & \commandname{LeastCommonMultiple} & State the finite family and
empty-family convention.\\
\commandname{vtwo} & \commandname{TwoAdicValuation} & State the value at zero or restrict
to nonzero arguments.\\
\commandname{qbinom} & \commandname{GaussianBinomial} & Supply all three arguments:
upper index, lower index, and base.\\
\commandname{qpoch}, \commandname{poch} & \commandname{QPochhammer} & Supply element,
base, and finite length or $\infty$ explicitly.\\
\commandname{Law} & \commandname{LawOperator} or \commandname{LawOf} & Use the
argument-delimited form for a random variable.\\
\commandname{Lop} & \commandname{LaplaceTransformSymbol} or
\commandname{LaplaceTransformOf} & Decide between an operator token and application; mark
bilateral transforms separately.\\
\commandname{abs} & \commandname{AbsoluteValue} & Scalar absolute value or complex
modulus; do not use for set cardinality.\\
\commandname{norm} & \commandname{Norm} & Declare the normed space and norm species.\\
\commandname{floor}, \commandname{ceil} & \commandname{Floor},
\commandname{Ceiling} & Preserve exact endpoint inequalities.\\
\commandname{pospart} & \commandname{PositivePart} & Ordered-scalar positive part with
value zero on nonpositive inputs.\\
\commandname{set} & \commandname{SetOf} & Supply the ambient type and predicate/list
grammar.\\
\commandname{angles} & \commandname{InnerProduct} & Declare inner product versus dual
pairing and the linear slot.\\
\commandname{diag} & \commandname{DiagonalOperator} & Matrix constructor only; diagonal
extraction uses a separately named local operator.\\
\commandname{Span}, \commandname{spanop} & \commandname{SpanOperator} & State the scalar
field.\\
\commandname{Tr} & \commandname{TraceOperator} & Matrix or trace-class operator trace;
state the space.\\
\commandname{rank} & \commandname{RankOperator} & Rank over the stated field.\\
\commandname{Res} & \commandname{ResidueOperator} & Complex residue; a named arithmetic
sequence retains its semantic superscript.\\
\commandname{Id} & \commandname{IdentityOperator} & Identity map on a stated domain.\\
\commandname{Log} & \commandname{PrincipalLogarithm} & Only the principal complex
branch; a formal logarithm or another branch is separately labelled.\\
\commandname{defeq} & \commandname{DefinitionEquals} & Introduction of a definition,
not a proved equality.\\
\commandname{Fglobal}, \commandname{Fs}, \commandname{extF}, \commandname{Fext} &
\commandname{FabiusGlobal} & Signed global Fabius extension only.\\
\commandname{InvF} & \commandname{FabiusQuantile} or
\commandname{FabiusClampedQuantile} & Decide from the input domain: $[0,1]$ for the
ordinary inverse, all reals for the clamped totalization.\\
\bottomrule
\end{longtable}
\endgroup

\begin{boxedremark}
The alias index deliberately does not assign global replacements to bare local letters,
generic hats, primes, stars, or calligraphic alphabets.  Those collisions require a type-
and-scope decision, recorded in \cref{sec:namespace-register}, rather than a lexical
substitution.
\end{boxedremark}

\section{Symbol and shared-command indexes}
\label{sec:symbol-index}

\subsection{Rendered-symbol index}

The rendered column is intentionally redundant with the semantic command name.  The
redundancy makes a PDF search by glyph useful while preserving a stable source-level
identifier.

\begingroup\small
\begin{longtable}{@{}L{0.22\textwidth}S{0.38\textwidth}L{0.34\textwidth}@{}}
\toprule
Rendered symbol family & Shared command(s) & Distinguishing data\\
\midrule
\endfirsthead
\toprule
Rendered symbol family & Shared command(s) & Distinguishing data\\
\midrule
\endhead
$\FabiusNotationVersion$ & \commandname{FabiusNotationVersion} & Date-valued
snapshot stamp for the shared notation contract; it is metadata, not a mathematical
variable.\\
$\NaturalNumbers,\PositiveIntegers$ & \commandname{NaturalNumbers},
\commandname{PositiveIntegers} & Zero included versus excluded.\\
$\IntegerNumbers,\RationalNumbers,\RealNumbers,\ComplexNumbers,
\ExtendedRealNumbers$ & \commandname{IntegerNumbers}, \commandname{RationalNumbers},
\commandname{RealNumbers}, \commandname{ComplexNumbers},
\commandname{ExtendedRealNumbers} & Algebraic number system versus ordered extension by
two infinities.\\
$\FabiusBounded,\FabiusGlobal,\RvachevUp$ & \commandname{FabiusBounded},
\commandname{FabiusGlobal}, \commandname{RvachevUp} & Bounded CDF, signed global
extension, and even compactly supported density-shaped function.\\
$\FabiusQuantile,\FabiusClampedQuantile$ & \commandname{FabiusQuantile},
\commandname{FabiusClampedQuantile} & Inverse on $[0,1]$ versus clamped totalization on
$\RealNumbers$.\\
$\EulerE,\ImaginaryUnit,\Differential x$ & \commandname{EulerE},
\commandname{ImaginaryUnit}, \commandname{Differential} & Exponential base, complex
unit, and integration differential.\\
$\AbsoluteValue{x},\Norm{x},\CardinalityOf{A}$ & \commandname{AbsoluteValue},
\commandname{Norm}, \commandname{CardinalityOf} & Scalar modulus, norm, and set
cardinality; the types disambiguate identical bars.\\
$\Floor{x},\Ceiling{x},\FractionalPart{x},\PositivePart{x}$ &
\commandname{Floor}, \commandname{Ceiling}, \commandname{FractionalPart},
\commandname{PositivePart} & Integer rounding, fractional residue, and ordered positive
part.\\
$\IndicatorOf{A},\SetOf{x:P(x)},\InnerProduct{x,y}$ &
\commandname{IndicatorOf}, \commandname{SetOf}, \commandname{InnerProduct} & Function,
set-builder delimiters, and inner-product/declared duality delimiters.\\
$\DefinitionEquals,\Sign,\IdentityOperator$ & \commandname{DefinitionEquals},
\commandname{Sign}, \commandname{IdentityOperator} & Definition relation, real sign,
and identity map.\\
\(\MultiplicativeInverseOf{F},\CompositionalInverseOf{F}\) &
\commandname{MultiplicativeInverseOf}, \commandname{CompositionalInverseOf} &
Inverse under multiplication versus inverse under composition.\\
$\SpanOperator,\TraceOperator,\RankOperator,\DiagonalOperator$ &
\commandname{SpanOperator}, \commandname{TraceOperator}, \commandname{RankOperator},
\commandname{DiagonalOperator} & Linear-algebra operators with an explicit field and
space.\\
$\ResidueOperator,\LeastCommonMultiple,\OrderOperator,\PrincipalLogarithm$ &
\commandname{ResidueOperator}, \commandname{LeastCommonMultiple},
\commandname{OrderOperator}, \commandname{PrincipalLogarithm} & Complex residue,
arithmetic LCM, qualified order species, and principal complex branch.\\
$\FallingFactorial{x}{n},\RisingFactorial{x}{n}$ &
\commandname{FallingFactorial}, \commandname{RisingFactorial} & Unit-step falling versus
rising products; both equal one at length zero.\\
\(\SignedStirlingFirstKind{n}{k},
  \UnsignedStirlingFirstKind{n}{k},
  \StirlingSecondKind{n}{k}\) &
\commandname{SignedStirlingFirstKind},
\commandname{UnsignedStirlingFirstKind},
\commandname{StirlingSecondKind} & Signed first-kind, unsigned cycle, and
second-kind partition triangles.\\
\(\LahNumber{n}{k}\) & \commandname{LahNumber} & Unsigned Lah number with
ordered blocks and an unordered block collection.\\
\(\TypeAEulerianNumber{n}{k},\TypeBEulerianNumber{n}{k},
  \SecondOrderEulerianNumber{n}{k}\) &
\commandname{TypeAEulerianNumber}, \commandname{TypeBEulerianNumber},
\commandname{SecondOrderEulerianNumber} & Type-A, type-B, and second-order
Eulerian numbers.\\
\(\TypeAEulerianPolynomial{n},\TypeBEulerianPolynomial{n},
  \SecondOrderEulerianPolynomial{n}\) &
\commandname{TypeAEulerianPolynomial}, \commandname{TypeBEulerianPolynomial},
\commandname{SecondOrderEulerianPolynomial} & Corresponding row polynomials,
with the evaluation variable following.\\
\(\BellNumber{n},\ComplementaryBellNumber{n},\TouchardPolynomial{n}\) &
\commandname{BellNumber}, \commandname{ComplementaryBellNumber},
\commandname{TouchardPolynomial} & Set-partition count, signed block-parity
sum, and Touchard row polynomial.\\
\(\ExponentialPartialBellPolynomial{n}{k},
  \ExponentialCompleteBellPolynomial{n},
  \OrdinaryPartialBellPolynomial{n}{k}\) &
\commandname{ExponentialPartialBellPolynomial},
\commandname{ExponentialCompleteBellPolynomial},
\commandname{OrdinaryPartialBellPolynomial} & Exponential partial/complete and
ordinary partial Bell-polynomial normalizations.\\
\(\CoefficientExtraction{z}{n}\) & \commandname{CoefficientExtraction} &
Coefficient of \(z^n\) in the following formal or analytic series operand.\\
\(\CycleCountOperator,\SetPartitionOperator,a\BitwiseAnd b\) &
\commandname{CycleCountOperator}, \commandname{SetPartitionOperator},
\commandname{BitwiseAnd} & Cycle count, set-partition constructor, and
nonnegative-integer bitwise AND.\\
\(\ExponentialGeneratingFunctionOf{A},\OrdinaryGeneratingFunctionOf{A}\) &
\commandname{ExponentialGeneratingFunctionOf},
\commandname{OrdinaryGeneratingFunctionOf} & Explicit EGF versus OGF labels.\\
\(\NormalizedReversionCoefficient{f}{j},
  \OrdinaryGeneratingCoefficient{x}{j}\) &
\commandname{NormalizedReversionCoefficient},
\commandname{OrdinaryGeneratingCoefficient} & Normalized exponential-series
reversion input versus ordinary-series coefficient ledger.\\
\(\PartitionLatticeMinimum,\PartitionLatticeMaximum,\TraceBellArgument{k}\) &
\commandname{PartitionLatticeMinimum}, \commandname{PartitionLatticeMaximum},
\commandname{TraceBellArgument} & Partition-lattice endpoints and scaled trace
input for Bell identities.\\
\(\FiniteField_q\) & \commandname{FiniteField} & Finite field with declared
prime-power order \(q\).\\
\(\CatalanNumber{n},\GregoryCoefficient{n},
  \GenocchiNumber{n},\GenocchiPolynomial{n}\) &
\commandname{CatalanNumber}, \commandname{GregoryCoefficient},
\commandname{GenocchiNumber}, \commandname{GenocchiPolynomial} & Distinct
Catalan, Gregory, and Genocchi families.\\
\(\BarnesGFunction\) & \commandname{BarnesGFunction} & Barnes \(G\)-function,
not a Gregory or Genocchi coefficient.\\
\(\BellHessenbergMatrixH{n},\BellHessenbergMatrixK{n}\) &
\commandname{BellHessenbergMatrixH}, \commandname{BellHessenbergMatrixK} &
Two Hessenberg determinant models for complete Bell polynomials.\\
\(\AssociahedronByDimension{d},\LagrangeShiftCoordinate\) &
\commandname{AssociahedronByDimension}, \commandname{LagrangeShiftCoordinate}
& Dimension-indexed associahedron and scope-local Lagrange shift coordinate.\\
$\SupportOperator,\SupportOf{f},\allowbreak\TopologicalSupportOf{f}$ &
\commandname{SupportOperator}, \commandname{SupportOf},
\commandname{TopologicalSupportOf} & Operator token, pointwise support, and its
topological closure.\\
\(\CoefficientSupportOf{F}\) & \commandname{CoefficientSupportOf} &
Coefficient-exponent support in an explicitly stated index set.\\
$\Convolution,\MetricDistance,\TotalVariationOf{a},\TotalVariationDistance$ &
\commandname{Convolution}, \commandname{MetricDistance},
\commandname{TotalVariationOf}, \commandname{TotalVariationDistance} &
Algebraic/measure convolution, generic metric, a one-object variation functional, and
the fixed two-measure probability TV convention.\\
$\Probability,\Expectation,\Variance,\Covariance$ & \commandname{Probability},
\commandname{Expectation}, \commandname{Variance}, \commandname{Covariance} & Measure,
integral operator, and second-order statistics.\\
$\LawOperator,\LawOf{X},\EqualInLaw,\ConvergesInLaw$ & \commandname{LawOperator},
\commandname{LawOf}, \commandname{EqualInLaw}, \commandname{ConvergesInLaw} & Law
operator/application, law equality, and weak convergence of laws.\\
$\DyadicSigmaField_n$ & \commandname{DyadicSigmaField} & Level-$n$ filtration
$\sigma$-field with explicitly stated generators.\\
$\DecayOptimizationObjective{n},\LinearizedDecayObjective{n}$ &
\commandname{DecayOptimizationObjective}, \commandname{LinearizedDecayObjective} &
Exact versus linearized Fourier-decay objectives.\\
$\FourierTwoPiOperator,\InverseFourierTwoPiOperator$ &
\commandname{FourierTwoPiOperator}, \commandname{InverseFourierTwoPiOperator} & Forward
and inverse cycles-per-unit operators.\\
$\FourierAngularOperator,\InverseFourierAngularOperator$ &
\commandname{FourierAngularOperator}, \commandname{InverseFourierAngularOperator} &
Forward and inverse angular-frequency operators.\\
$\FourierTwoPiTransformOf{f},\FourierTwoPi{f}$ &
\commandname{FourierTwoPiTransformOf}, \commandname{FourierTwoPi} & Bracketed operator
application and hat rendering of the same $2\pi$ transform.\\
$\FourierAngularTransformOf{f},\FourierAngular{f}$ &
\commandname{FourierAngularTransformOf}, \commandname{FourierAngular} & Bracketed
operator application and hat rendering of the same angular transform.\\
$\FourierSeriesCoefficient{h}{k},\FourierSeriesCoefficientAtPeriod{h}{k}{T}$ &
\commandname{FourierSeriesCoefficient},
\commandname{FourierSeriesCoefficientAtPeriod} & Unit-period and explicitly
$T$-periodic normalized coefficients; neither is a continuous line transform.\\
$\LaplaceTransformSymbol,\LaplaceTransformOf{f}$ &
\commandname{LaplaceTransformSymbol}, \commandname{LaplaceTransformOf} & One-sided
operator token and application; bilateral use is subscripted.\\
$\MellinTransformSymbol,\MellinTransformOf{f}$ &
\commandname{MellinTransformSymbol}, \commandname{MellinTransformOf} & Operator token and
application on a stated convergence strip.\\
$\MomentGeneratingFunction{X},\CharacteristicFunction{X}$ &
\commandname{MomentGeneratingFunction}, \commandname{CharacteristicFunction} &
Exponential moment transform and positive-sign probabilistic Fourier transform.\\
$\SincRad,\SincPi$ & \commandname{SincRad}, \commandname{SincPi} & Radian versus
$\pi$-normalized entire sinc.\\
$\BinaryDigitSum,\ThueMorseSignSymbol,\ThueMorseSign{n}$ &
\commandname{BinaryDigitSum}, \commandname{ThueMorseSignSymbol},
\commandname{ThueMorseSign} & Bit count, reserved sequence glyph, and indexed sign.\\
$\PadicValuation{p},\TwoAdicValuation$ & \commandname{PadicValuation},
\commandname{TwoAdicValuation} & Prime-parametric and specialized two-adic valuations;
zero convention local.\\
$\QPochhammer{a}{q}{n},\GaussianBinomial{n}{k}{q},\QInteger{n}{q}$ &
\commandname{QPochhammer}, \commandname{GaussianBinomial}, \commandname{QInteger} & Every
base and finite/infinite length explicit.\\
\(\Polylogarithm,\LambertW,\GammaFunction,\RiemannZeta,
  \HurwitzZeta{s}{a}\) &
\commandname{Polylogarithm}, \commandname{LambertW}, \commandname{GammaFunction},
\commandname{RiemannZeta}, \commandname{HurwitzZeta} & Named special functions;
order, branch, shift, and continuation domains printed as applicable.\\
\(\WrightOmega,\LambertPolynomial{n},\ScaledLambertPolynomial{n}\) &
\commandname{WrightOmega}, \commandname{LambertPolynomial},
\commandname{ScaledLambertPolynomial} & Wright omega, canonical zero-based Lambert
polynomial, and its factorial scaling.\\
\(\PolynomialLogPowerSeriesRing{K},\PolynomialLogLaurentRing{K},
\RationalLogLaurentField{K},\IteratedLaurentLogField{K}\) &
\commandname{PolynomialLogPowerSeriesRing},
\commandname{PolynomialLogLaurentRing},
\commandname{RationalLogLaurentField},
\commandname{IteratedLaurentLogField} & Power-series, polynomial-coefficient Laurent,
rational-coefficient Laurent, and iterated Laurent-log classes.\\
\(\NearIdentityPolynomialLogGroup{K}\) &
\commandname{NearIdentityPolynomialLogGroup} & Near-identity composition group over
the named coefficient field.\\
$\BigO,\LittleO,\BigOAt{g}{a},\LittleOAt{g}{a}$ & \commandname{BigO},
\commandname{LittleO}, \commandname{BigOAt}, \commandname{LittleOAt} & Token-only and
explicit-limit asymptotic operators.\\
\(\FormalBigO,\FormalBigOAt{t^N}{t}\) &
\commandname{FormalBigO}, \commandname{FormalBigOAt} & Formal valuation-tail token and
argument-bearing form; neither is an analytic estimate.\\
\bottomrule
\end{longtable}
\endgroup

\subsection{Alphabetical shared-command index}

The following is the complete source-level index for version
\FabiusNotationVersion.  Argument counts are part of the API; braces displayed in a
command name are mathematical arguments, not optional typography.

\begingroup\small
\begin{longtable}{@{}S{0.35\textwidth}L{0.08\textwidth}L{0.51\textwidth}@{}}
\toprule
Command & Arity & Normative entry\\
\midrule
\endfirsthead
\toprule
Command & Arity & Normative entry\\
\midrule
\endhead
\commandname{AbsoluteValue} & 1 & \notelink{scalar-size}{Absolute value, norm, distance}.\\
\commandname{AssociahedronByDimension} & 1 &
\notelink{combinatorial-letter-collisions}{Associahedron indexing}.\\
\commandname{BarnesGFunction} & 0 &
\notelink{combinatorial-letter-collisions}{Barnes G-function}.\\
\commandname{BellHessenbergMatrixH} & 1 &
\notelink{combinatorial-letter-collisions}{Bell Hessenberg matrices}.\\
\commandname{BellHessenbergMatrixK} & 1 &
\notelink{combinatorial-letter-collisions}{Bell Hessenberg matrices}.\\
\commandname{BellNumber} & 1 &
\notelink{bell-and-touchard-families}{Bell numbers}.\\
\commandname{BigO} & 0 & \notelink{big-o}{Asymptotic operators}.\\
\commandname{BigOAt} & 2 & \notelink{big-o}{Asymptotic operators}.\\
\commandname{BinaryDigitSum} & 0 & \notelink{binary-weight}{Binary digit sum}.\\
\commandname{BitwiseAnd} & 0 &
\notelink{combinatorial-series-normalizations}{Bitwise AND}.\\
\commandname{CardinalityOf} & 1 & \notelink{semantic-delimiters}{Semantic delimiters}.\\
\commandname{CatalanNumber} & 1 &
\notelink{combinatorial-letter-collisions}{Catalan numbers}.\\
\commandname{Ceiling} & 1 & \notelink{rounding-positive-part}{Rounding and positive part}.\\
\commandname{CharacteristicFunction} & 1 & \notelink{characteristic-function}{Characteristic functions}.\\
\commandname{CoefficientExtraction} & 2 &
\notelink{combinatorial-series-normalizations}{Coefficient extraction}.\\
\commandname{CoefficientSupportOf} & 1 &
\notelink{polynomial-logarithmic-leading-data}{Coefficient support}.\\
\commandname{ComplementaryBellNumber} & 1 &
\notelink{bell-and-touchard-families}{Complementary Bell numbers}.\\
\commandname{ComplexNumbers} & 0 & \notelink{number-systems}{Number systems}.\\
\commandname{CompositionalInverseOf} & 1 &
\notelink{polynomial-logarithmic-reversion}{Compositional inverse}.\\
\commandname{ConvergesInLaw} & 0 & \notelink{law}{Law and weak convergence}.\\
\commandname{Convolution} & 0 & \notelink{convolution}{Convolution}.\\
\commandname{Covariance} & 0 & \notelink{expectation}{Expectation and moments}.\\
\commandname{CycleCountOperator} & 0 &
\notelink{combinatorial-series-normalizations}{Permutation cycle count}.\\
\commandname{DecayOptimizationObjective} & 1 & \notelink{fourier-decay-objectives}{Decay objectives}.\\
\commandname{DefinitionEquals} & 0 & \notelink{semantic-delimiters}{Semantic delimiters}.\\
\commandname{DiagonalOperator} & 0 & \notelink{matrix-operators}{Matrix operators}.\\
\commandname{Differential} & 0 & \notelink{calculus}{Calculus}.\\
\commandname{DyadicSigmaField} & 0 & \notelink{dyadic-filtration}{Dyadic filtration}.\\
\commandname{EqualInLaw} & 0 & \notelink{law}{Law and weak convergence}.\\
\commandname{EulerE} & 0 & \notelink{elementary-constants}{Elementary constants}.\\
\commandname{Expectation} & 0 & \notelink{expectation}{Expectation and moments}.\\
\commandname{ExponentialCompleteBellPolynomial} & 1 &
\notelink{bell-and-touchard-families}{Complete exponential Bell polynomials}.\\
\commandname{ExponentialGeneratingFunctionOf} & 1 &
\notelink{combinatorial-series-normalizations}{Exponential generating functions}.\\
\commandname{ExponentialPartialBellPolynomial} & 2 &
\notelink{bell-and-touchard-families}{Partial exponential Bell polynomials}.\\
\commandname{ExtendedRealNumbers} & 0 & \notelink{number-systems}{Number systems}.\\
\commandname{FabiusBounded} & 0 & \notelink{fabius-bounded}{Bounded Fabius function}.\\
\commandname{FabiusClampedQuantile} & 0 & \notelink{fabius-quantile}{Fabius inverses}.\\
\commandname{FabiusGlobal} & 0 & \notelink{fabius-global}{Global extension}.\\
\commandname{FabiusNotationVersion} & 0 & \notelink{contract-version}{Contract version}.\\
\commandname{FabiusQuantile} & 0 & \notelink{fabius-quantile}{Fabius inverses}.\\
\commandname{FallingFactorial} & 2 & \notelink{powers-factorials}{Powers and factorials}.\\
\commandname{FiniteField} & 0 &
\notelink{combinatorial-series-normalizations}{Finite fields}.\\
\commandname{Floor} & 1 & \notelink{rounding-positive-part}{Rounding and positive part}.\\
\commandname{FormalBigO} & 0 &
\notelink{polynomial-logarithmic-truncation}{Formal valuation tails}.\\
\commandname{FormalBigOAt} & 2 &
\notelink{polynomial-logarithmic-truncation}{Formal valuation tails}.\\
\commandname{FourierAngular} & 1 & \notelink{fourier-angular}{Angular Fourier transform}.\\
\commandname{FourierAngularOperator} & 0 & \notelink{fourier-angular}{Angular Fourier transform}.\\
\commandname{FourierAngularTransformOf} & 1 & \notelink{fourier-angular}{Angular Fourier transform}.\\
\commandname{FourierSeriesCoefficient} & 2 & \notelink{periodic-fourier}{Fourier series}.\\
\commandname{FourierSeriesCoefficientAtPeriod} & 3 & \notelink{periodic-fourier}{Fourier series}.\\
\commandname{FourierTwoPi} & 1 & \notelink{fourier-two-pi}{$2\pi$ Fourier transform}.\\
\commandname{FourierTwoPiOperator} & 0 & \notelink{fourier-two-pi}{$2\pi$ Fourier transform}.\\
\commandname{FourierTwoPiTransformOf} & 1 & \notelink{fourier-two-pi}{$2\pi$ Fourier transform}.\\
\commandname{FractionalPart} & 1 & \notelink{rounding-positive-part}{Rounding and positive part}.\\
\commandname{GammaFunction} & 0 & \notelink{polylog-gamma-zeta}{Gamma function}.\\
\commandname{GaussianBinomial} & 3 & \notelink{gaussian-binomial}{Gaussian binomial}.\\
\commandname{GenocchiNumber} & 1 &
\notelink{combinatorial-letter-collisions}{Genocchi numbers}.\\
\commandname{GenocchiPolynomial} & 1 &
\notelink{combinatorial-letter-collisions}{Genocchi polynomials}.\\
\commandname{GregoryCoefficient} & 1 &
\notelink{combinatorial-letter-collisions}{Gregory coefficients}.\\
\commandname{HurwitzZeta} & 2 &
\notelink{polylog-gamma-zeta}{Hurwitz zeta function}.\\
\commandname{IdentityOperator} & 0 & \notelink{maps}{Maps and the identity map}.\\
\commandname{ImaginaryUnit} & 0 & \notelink{elementary-constants}{Elementary constants}.\\
\commandname{IndicatorOf} & 1 & \notelink{indicator}{Indicators}.\\
\commandname{InnerProduct} & 1 & \notelink{semantic-delimiters}{Semantic delimiters}.\\
\commandname{IntegerNumbers} & 0 & \notelink{number-systems}{Number systems}.\\
\commandname{InverseFourierAngularOperator} & 0 & \notelink{fourier-angular}{Angular Fourier transform}.\\
\commandname{InverseFourierTwoPiOperator} & 0 & \notelink{fourier-two-pi}{$2\pi$ Fourier transform}.\\
\commandname{IteratedLaurentLogField} & 1 &
\notelink{polynomial-logarithmic-ambient}{Iterated Laurent-log field}.\\
\commandname{LagrangeShiftCoordinate} & 0 &
\notelink{combinatorial-letter-collisions}{Lagrange shift coordinate}.\\
\commandname{LahNumber} & 2 &
\notelink{stirling-and-lah-families}{Lah numbers}.\\
\commandname{LambertPolynomial} & 1 &
\notelink{wright-omega-lambert-polynomials}{Canonical Lambert polynomials}.\\
\commandname{LambertW} & 0 & \notelink{lambert-w}{Lambert (W)}.\\
\commandname{LaplaceTransformOf} & 1 & \notelink{laplace-transform}{Laplace transforms}.\\
\commandname{LaplaceTransformSymbol} & 0 & \notelink{laplace-transform}{Laplace transforms}.\\
\commandname{LawOf} & 1 & \notelink{law}{Law and weak convergence}.\\
\commandname{LawOperator} & 0 & \notelink{law}{Law and weak convergence}.\\
\commandname{LeastCommonMultiple} & 0 & \notelink{arithmetic-functions}{Arithmetic functions}.\\
\commandname{LinearizedDecayObjective} & 1 & \notelink{fourier-decay-objectives}{Decay objectives}.\\
\commandname{LittleO} & 0 & \notelink{big-o}{Asymptotic operators}.\\
\commandname{LittleOAt} & 2 & \notelink{big-o}{Asymptotic operators}.\\
\commandname{MellinTransformOf} & 1 & \notelink{mellin-transform}{Mellin transform}.\\
\commandname{MellinTransformSymbol} & 0 & \notelink{mellin-transform}{Mellin transform}.\\
\commandname{MetricDistance} & 0 & \notelink{scalar-size}{Absolute value, norm, distance}.\\
\commandname{MomentGeneratingFunction} & 1 & \notelink{mgf}{Moment-generating functions}.\\
\commandname{MultiplicativeInverseOf} & 1 &
\notelink{polynomial-logarithmic-reversion}{Multiplicative reciprocal}.\\
\commandname{NaturalNumbers} & 0 & \notelink{number-systems}{Number systems}.\\
\commandname{NearIdentityPolynomialLogGroup} & 1 &
\notelink{polynomial-logarithmic-ambient}{Near-identity polynomial-log group}.\\
\commandname{Norm} & 1 & \notelink{scalar-size}{Absolute value, norm, distance}.\\
\commandname{NormalizedReversionCoefficient} & 2 &
\notelink{combinatorial-series-normalizations}{Normalized reversion coefficients}.\\
\commandname{OrderOperator} & 0 & \notelink{logarithm-order}{Logarithm and order}.\\
\commandname{OrdinaryGeneratingCoefficient} & 2 &
\notelink{combinatorial-series-normalizations}{Ordinary generating coefficients}.\\
\commandname{OrdinaryGeneratingFunctionOf} & 1 &
\notelink{combinatorial-series-normalizations}{Ordinary generating functions}.\\
\commandname{OrdinaryPartialBellPolynomial} & 2 &
\notelink{bell-and-touchard-families}{Partial ordinary Bell polynomials}.\\
\commandname{PadicValuation} & 1 & \notelink{padic-valuation}{Valuations}.\\
\commandname{PartitionLatticeMaximum} & 0 &
\notelink{partition-lattice-bounds}{Partition-lattice maximum}.\\
\commandname{PartitionLatticeMinimum} & 0 &
\notelink{partition-lattice-bounds}{Partition-lattice minimum}.\\
\commandname{Polylogarithm} & 0 & \notelink{polylog-gamma-zeta}{Polylogarithm}.\\
\commandname{PolynomialLogLaurentRing} & 1 &
\notelink{polynomial-logarithmic-ambient}{Polynomial-log Laurent ring}.\\
\commandname{PolynomialLogPowerSeriesRing} & 1 &
\notelink{polynomial-logarithmic-ambient}{Polynomial-log power-series ring}.\\
\commandname{PositiveIntegers} & 0 & \notelink{number-systems}{Number systems}.\\
\commandname{PositivePart} & 1 & \notelink{rounding-positive-part}{Rounding and positive part}.\\
\commandname{PrincipalLogarithm} & 0 & \notelink{logarithm-order}{Logarithm and order}.\\
\commandname{Probability} & 0 & \notelink{probability-space}{Probability spaces}.\\
\commandname{QInteger} & 2 & \notelink{q-integer}{(q)-integers}.\\
\commandname{QPochhammer} & 3 & \notelink{q-pochhammer}{(q)-Pochhammer symbols}.\\
\commandname{RankOperator} & 0 & \notelink{matrix-operators}{Matrix operators}.\\
\commandname{RationalLogLaurentField} & 1 &
\notelink{polynomial-logarithmic-ambient}{Rational-log Laurent field}.\\
\commandname{RationalNumbers} & 0 & \notelink{number-systems}{Number systems}.\\
\commandname{RealNumbers} & 0 & \notelink{number-systems}{Number systems}.\\
\commandname{ResidueOperator} & 0 & \notelink{residue}{Residues}.\\
\commandname{RiemannZeta} & 0 & \notelink{polylog-gamma-zeta}{Zeta functions}.\\
\commandname{RisingFactorial} & 2 & \notelink{powers-factorials}{Powers and factorials}.\\
\commandname{RvachevUp} & 0 & \notelink{rvachev-up}{Rvachev up-function}.\\
\commandname{ScaledLambertPolynomial} & 1 &
\notelink{wright-omega-lambert-polynomials}{Scaled Lambert polynomials}.\\
\commandname{SecondOrderEulerianNumber} & 2 &
\notelink{eulerian-families}{Second-order Eulerian numbers}.\\
\commandname{SecondOrderEulerianPolynomial} & 1 &
\notelink{eulerian-families}{Second-order Eulerian row polynomials}.\\
\commandname{SetOf} & 1 & \notelink{semantic-delimiters}{Semantic delimiters}.\\
\commandname{SetPartitionOperator} & 0 &
\notelink{combinatorial-series-normalizations}{Set partitions}.\\
\commandname{Sign} & 0 & \notelink{rounding-positive-part}{Rounding and positive part}.\\
\commandname{SignedStirlingFirstKind} & 2 &
\notelink{stirling-and-lah-families}{Signed first-kind Stirling numbers}.\\
\commandname{SincPi} & 0 & \notelink{sinc}{Sinc normalizations}.\\
\commandname{SincRad} & 0 & \notelink{sinc}{Sinc normalizations}.\\
\commandname{SpanOperator} & 0 & \notelink{matrix-operators}{Matrix operators}.\\
\commandname{StirlingSecondKind} & 2 &
\notelink{stirling-and-lah-families}{Second-kind Stirling numbers}.\\
\commandname{SupportOf} & 1 & \notelink{support}{Support}.\\
\commandname{SupportOperator} & 0 & \notelink{support}{Support}.\\
\commandname{ThueMorseSign} & 1 & \notelink{thue-morse}{Thue--Morse sign}.\\
\commandname{ThueMorseSignSymbol} & 0 & \notelink{thue-morse}{Thue--Morse sign}.\\
\commandname{TopologicalSupportOf} & 1 & \notelink{support}{Support}.\\
\commandname{TotalVariationDistance} & 0 & \notelink{probability-metrics}{Probability metrics}.\\
\commandname{TotalVariationOf} & 1 & \notelink{probability-metrics}{Probability metrics}.\\
\commandname{TouchardPolynomial} & 1 &
\notelink{bell-and-touchard-families}{Touchard polynomials}.\\
\commandname{TraceBellArgument} & 1 &
\notelink{combinatorial-series-normalizations}{Scaled Bell trace inputs}.\\
\commandname{TraceOperator} & 0 & \notelink{matrix-operators}{Matrix operators}.\\
\commandname{TwoAdicValuation} & 0 & \notelink{padic-valuation}{Valuations}.\\
\commandname{TypeAEulerianNumber} & 2 &
\notelink{eulerian-families}{Type-A Eulerian numbers}.\\
\commandname{TypeAEulerianPolynomial} & 1 &
\notelink{eulerian-families}{Type-A Eulerian row polynomials}.\\
\commandname{TypeBEulerianNumber} & 2 &
\notelink{eulerian-families}{Type-B Eulerian numbers}.\\
\commandname{TypeBEulerianPolynomial} & 1 &
\notelink{eulerian-families}{Type-B Eulerian row polynomials}.\\
\commandname{UnsignedStirlingFirstKind} & 2 &
\notelink{stirling-and-lah-families}{Unsigned first-kind Stirling numbers}.\\
\commandname{Variance} & 0 & \notelink{expectation}{Expectation and moments}.\\
\commandname{WrightOmega} & 0 &
\notelink{wright-omega-lambert-polynomials}{Wright omega}.\\
\bottomrule
\end{longtable}
\endgroup

\section{Lean crosswalk and formalization boundary}
\label{sec:lean-crosswalk}

The crosswalk maps semantic mathematical objects to actual current Lean declarations.
It does not imply that every identity printed in a frontier document has been formalized.
A row marked ``notation only'' records precisely that no single global Lean object is
asserted by the shared command.  Lean names are case-sensitive and module paths are
relative to \modulepath{Analysis/FabiusFunction/Lean/FabiusFunction/}.

\begingroup\small
\begin{longtable}{@{}L{0.22\textwidth}S{0.40\textwidth}L{0.32\textwidth}@{}}
\toprule
Catalogue object & Lean declaration and module & Exact boundary\\
\midrule
\endfirsthead
\toprule
Catalogue object & Lean declaration and module & Exact boundary\\
\midrule
\endhead
$\FabiusBounded$ & \lean{fabiusReal} in \modulepath{Basic.lean}; canonical witness
\lean{fabius} and proof \lean{fabius_spec} in \modulepath{PaperStatements.lean} &
\lean{fabiusReal} is the real-valued view of a \lean{BoundedFabius}; the witness and
\lean{IsFabius} proof select the canonical object.\\
Abstract Fabius specification & \lean{IsFabius} in \modulepath{Basic.lean} & The declaration's
actual type is authoritative for hypotheses; the catalogue summary does not manufacture
a weaker premise-free theorem.\\
$\RvachevUp$ & \lean{rvachevUp} in \modulepath{Basic.lean} & Takes a
\lean{BoundedFabius} argument.  Evenness, bounds, and other properties are separate
theorems, not fields silently attached to an arbitrary function.\\
$\FabiusGlobal$ & \lean{extendedFabius} in \modulepath{Basic.lean} & Takes a
\lean{BoundedFabius}; translate-cell and calculus results occur in later dedicated
modules.\\
$\FabiusQuantile$ and $\FabiusClampedQuantile$ & \lean{fabiusOrderIso} and
\lean{fabiusInv} in \modulepath{FabiusInverse.lean} & The formal definitions take both a
bounded object and its \lean{IsFabius} proof.  Domain behavior must be checked against
the exact declaration before claiming ordinary versus clamped inverse laws.\\
\(F_n^{\mathrm{spl}}\) and its local inverse \(G_n^{\mathrm{spl}}\) &
\lean{reportFiniteFabiusApproximant} in \modulepath{QuarterSplineLocalPolynomial.lean};
\lean{fabiusUniformSpline} in \modulepath{FabiusUniformSpline.lean} &
The exact index bridge identifies the report approximant at \(n\) with the uniform spline
at \(n-1\).  Strict monotonicity is proved only for
\(n\ge3\) on \([1/4,1/4+2^{-n}]\).  No global \(G_n^{\mathrm{spl}}\) is defined in Lean;
the quantile theorem assumes a supplied local left inverse and is not \lean{fabiusInv}.\\
\(\Phi_{\mathrm{ang}}\) and \(\Phi_{2\pi}\) & \lean{centeredSincProduct} in
\modulepath{FourierLaplaceRotation.lean}; \lean{rvachevFourierProduct} in
\modulepath{Basic.lean}; bridge
\lean{centeredSincProduct_eq_}\allowbreak\lean{rvachevFourierProduct} & The bridge is
\(\operatorname{centeredSincProduct}(z)=
\operatorname{rvachevFourierProduct}(z/(2\pi))\).  Identifying the product with the
integral transform uses \lean{rvachevFourier_eq_product} and retains
\(F:\lean{BoundedFabius}\) and \(h_F:\lean{IsFabius}\,F\).\\
Endpoint kernel \(\mathcal K(z,w)\) & Notation only & Formal frontier construction in the
two-scale endpoint coefficient algebra.  No exact current Lean definition of
\(\mathcal K\), \(\mathcal R\), or its fixed-point algebra is claimed.\\
$\ThueMorseSign{n}$ & \lean{thueMorseSign} in \modulepath{Arithmetic.lean} & Integer-valued
sign sequence on \lean{Nat}; digit, recurrence, and product identities live in their
named Thue--Morse modules.\\
Prime-power Pascal-row valuations & \lean{primePowerChoose_padicValNat_add},
\lean{primePowerChoose_padicValNat}, and \lean{twoPowChoose_padicValNat} in
\modulepath{PrimePowerBinomialValuation.lean} & Arbitrary-prime formulas include
\(j=p^m\) and \(m=0\); the dyadic theorem uses \(0<j<2^m\).  No theorem for the distinct
endpoint-flat weights \(\binom{2^m-2}{j-1}\) is inferred.\\
Exact full moments & \lean{moment} in \modulepath{Arithmetic.lean};
\lean{momentIntegral} in \modulepath{Basic.lean} & Rational recurrence and real integral are
distinct objects; \modulepath{AnalyticMoments.lean} supplies bridge theorems with hypotheses.\\
Exact half moments & \lean{halfMoment} in \modulepath{Arithmetic.lean};
\lean{halfMomentIntegral} in \modulepath{Basic.lean} & Rational recurrence and real integral
are not identified by notation alone.\\
$R_n^{\mathrm{Res}}$ & \lean{reshetnikov} in \modulepath{Arithmetic.lean} & The Lean value is
rational at the definition boundary; integrality, positivity, and oddness require their
named theorems.\\
$D_n^{\mathrm{dyad}}$ & \lean{dyadicDenominator} in \modulepath{Arithmetic.lean} & A natural
number definition; bounds and divisibility results occur in \modulepath{DenominatorBound.lean}
and related modules.\\
Number systems and elementary algebra & Mathlib types \lean{Nat}, \lean{Int},
\lean{Rat}, \lean{Real}, \lean{Complex} and their standard APIs & Printed coercions are
mathematical exposition.  Lean coercions and side conditions remain those required by
the elaborated theorem.\\
Intervals & \lean{Set.Icc}, \lean{Set.Ioo}, \lean{Set.Ico}, \lean{Set.Ioc} & Endpoint
order in the Lean constructor matches the four catalogue interval conventions.\\
Probability laws and pushforwards & Mathlib \lean{Measure.map} and probability-measure
APIs & A prose random-series representation is not Lean-proved unless a project theorem
states the corresponding measure equality.\\
Fourier commands & Mathlib Fourier-transform APIs plus project-specific theorems in
dedicated Fourier modules & The shared commands choose print normalization.  They are
not names of one new Lean definition and do not erase Mathlib's hypotheses.\\
Finite-node Lagrange--Rvachev decoder & \lean{lagrangeRvachevDecoder} and
\lean{lagrangeRvachevAtomCoefficient} in \modulepath{LagrangeRvachevSynthesis.lean} &
Exact synthesis uses \(M\ne0\), \(|s|-1\le\TwoAdicValuation(M)\), and \(x\in[-1,1]\).
Node injectivity is required for cardinal evaluation, not polynomial reconstruction;
row-sum normalization also requires a nonempty node family.\\
Shifted dyadic polynomial self-sampling &
\lean{integral_polynomial_mul_}\allowbreak\lean{rvachevUp_eq_dyadic_tsum} in
\modulepath{PolynomialCombExactness.lean} & Exact for every polynomial of degree at most
\(N\), every real shift, and a bounded Fabius object satisfying \lean{IsFabius}.  Compact
support makes the displayed bilateral sum finite; exactness failure at degree \(N+1\) is
not claimed.\\
Legendre Gaunt and zero-row Wigner square & \lean{legendreGauntAdmissible} and
\lean{legendreWignerThreeJZeroSqRat} in \modulepath{LegendreGauntClosedForm.lean}; finite
Rvachev Gram sums in \modulepath{FabiusLegendreGauntClosedForm.lean} & Only the total
rational square for natural-number indices and zero magnetic row is defined.  No signed
symbol, phase, half-integer/nonzero-magnetic theory, or general \(3j/6j/9j\) API is
claimed; the real Gram formula retains \lean{IsFabius}.\\
$\DyadicSigmaField_n$ & Notation only & Each document declares its generators and
probability space; no single project-wide Lean filtration is asserted by the command.\\
$\DecayOptimizationObjective{n}$ and $\LinearizedDecayObjective{n}$ & Notation only &
They name document-local typed functionals in the Fourier-decay frontier.  No general
formalization or equality between exact and linearized objectives is claimed.\\
Generic $q$-series, fractional, matrix, spectral, and frontier entries & Definition or
classical notation unless a row or document cites an exact declaration & Their presence
in this encyclopedic catalogue is not evidence that the general theory or a proposed
Fabius specialization is formalized.\\
\bottomrule
\end{longtable}
\endgroup

\begin{boxedremark}
The safe direction of use is Lean-to-prose: inspect a declaration and then state exactly
that theorem with all hypotheses.  Similar names, compiled object files, an earlier audit,
or a catalogue entry are not substitutes for inspecting the current declaration.
\end{boxedremark}

\section{Active-source manifest and reproducibility}
\label{sec:source-manifest}

\subsection{Snapshot boundary}

The inventory below is the recursive active tree observed on 2026-09-01 at the
source-validation checkpoint.  It excludes every path below \modulepath{docs/archive/} and
counts this catalogue exactly once.  It contains 172 mathematical content sources: 67
standalone documents and 105 included or generated fragments.  The separate shared
contract \modulepath{docs/fabius-notation.tex} raises the total before this catalogue to 173
active TeX files.  Including this standalone catalogue gives 174 active TeX files and
68 standalone documents.  ``Standalone'' means that the uncommented source
contains \commandname{documentclass}; every other content source is a fragment inheriting
its root preamble.

\begingroup\small
\begin{longtable}{@{}L{0.46\textwidth}@{\hspace{1.5em}}r@{\hspace{1.5em}}r@{\hspace{1.5em}}r@{}}
\toprule
Manifest group & Standalone & Fragments & Total\\
\midrule
\endfirsthead
\toprule
Manifest group & Standalone & Fragments & Total\\
\midrule
\endhead
\modulepath{Fabius_Function_and_Rvachev_Up/} & 1 & 0 & 1\\
\modulepath{fabius_lean_walkthrough/} & 1 & 0 & 1\\
\modulepath{FabiusFunction_Mathematical_Glossary/} & 1 & 0 & 1\\
\modulepath{papers/} & 7 & 0 & 7\\
Frontier synthesis root & 1 & 0 & 1\\
\modulepath{drafts/combinatorial-coefficient-calculus/} & 6 & 0 & 6\\
\modulepath{drafts/exponents-and-q-series/} & 7 & 20 & 27\\
\modulepath{drafts/frontier-compilations/} & 3 & 1 & 4\\
\modulepath{drafts/integration-and-transforms/} & 1 & 7 & 8\\
\modulepath{drafts/inverse-and-sampling/} & 8 & 37 & 45\\
\modulepath{drafts/lambert-w/} & 7 & 4 & 11\\
\modulepath{drafts/representations/} & 12 & 26 & 38\\
\modulepath{drafts/rvachev_up_fourier_decay/} & 1 & 0 & 1\\
\modulepath{drafts/spectra-and-arithmetic/} & 10 & 10 & 20\\
\modulepath{drafts/thue-morse/} & 1 & 0 & 1\\
\midrule
Mathematical content subtotal & 67 & 105 & 172\\
Shared notation contract & 0 & 1 & 1\\
This catalogue & 1 & 0 & 1\\
\midrule
Active TeX total including catalogue & 68 & 106 & 174\\
\bottomrule
\end{longtable}
\endgroup

The frontier paths in the table are all below
\modulepath{semi-formalized-research-frontiers/}; the root row is its standalone synthesis.
Counts describe files, not distinct mathematical claims or publication units.  Corrected
and source-faithful historical-paper variants are counted separately because each is an
active TeX source requiring an interpretable notation namespace.

\subsection{Reproducible inventory and strict gate}

The repository audit is invoked from the repository root as follows.

\begin{lstlisting}[style=pseudo]
python Analysis/FabiusFunction/scripts/tex_notation_audit.py --strict --list
python Analysis/FabiusFunction/scripts/tex_notation_audit.py --semantic --strict --list
\end{lstlisting}

For each active source it records the relative path, byte and line counts, SHA-256 digest,
standalone/fragment classification, canonical-input count, input targets, and lexical
command counts.  It excludes the archive by resolved path, not by a fragile filename
pattern.  A focused catalogue check supplies this file's path as the positional argument.

The strict gate detects retired commands, local redefinitions of shared semantic
commands, missing or duplicated canonical inputs, fragment imports of the shared input,
selected literal bypasses of semantic operators, and known script-composition hazards.
The opt-in semantic gate additionally requires classification of raw hats, delimiter
spellings, number systems, indicator symbols, digit sums, asymptotic operators, sinc
normalizations, and probability operators.  Its sole raw-hat exemption is a descriptive
one-line local macro definition whose immediately preceding
\texttt{LOCAL-NON-FOURIER-HAT:} comment names that macro and states its meaning; direct
raw-hat use remains an error.  The audit is deliberately not a TeX parser and cannot infer
the meaning of every one-letter variable.  The local declaration schema and collision
register close that semantic gap.  On any scan whose resolved file set is the complete
active corpus, the audit also compares valid marker/declaration pairs with the reserved
\commandname{localaccentname} metadata in this catalogue and rejects missing, extra, or
duplicate rows.  The dedicated metadata token is not available for examples or literal
listings.  Rendered-PDF inspection later closes the composition and layout gap.

\section{Completeness and maintenance contract}
\label{sec:maintenance}

At shared-contract version \FabiusNotationVersion:

\begin{itemize}
\item \modulepath{docs/fabius-notation.tex} defines exactly 91 shared commands.
\item Every one of those 91 command names occurs in the canonical quick reference and
      in the alphabetical command index, with its exact arity.
\item Every command belongs to a detailed normative entry defining its type-level role,
      normalization, exceptional cases, or the local data the command intentionally
      requires at each use.
\item Both line-Fourier conventions, unit-period and explicit-period Fourier-series
      coefficients, both sinc normalizations, both Fabius inverses, the
      bounded/global/up-function triad, zero-inclusive/positive naturals, exact/linearized
      decay objectives, and pointwise/topological support have distinct printed forms.
\item The retired-alias index is exhaustive with respect to the audit's retired-command
      registry.  Ambiguous aliases map only after semantic classification.
\item The active corpus has 35 validated \texttt{LOCAL-NON-FOURIER-HAT:}
      marker/declaration pairs representing 35 distinct document-local command names.
      Every distinct name occurs exactly once in the exhaustive accent register, and the
      whole-corpus audit rejects missing, extra, or duplicate registry metadata.
\item Symbols not promoted to the shared API remain governed by the local declaration
      schema; they have no implicit corpus-wide type or normalization.
\end{itemize}

A future shared-command change is incomplete unless it updates, in the same logical
change, the version stamp, quick reference, detailed entry, rendered-symbol index,
alphabetical index, alias mapping if applicable, audit registry, and at least one focused
strict-audit result.  A future source that introduces a genuinely new mathematical
convention updates the relevant encyclopedic entry even when no new shared command is
needed.  A convention found only in an explicit historical quotation is entered in the
retired-alias register, not revived as normative project notation.

\begin{boxedremark}
This catalogue is a notation contract and an interpretive reference.  It does not convert
frontier assertions into theorems, does not weaken hypotheses found in Lean, and does not
substitute lexical or rendered checks for mathematical verification.  Its promise is
narrower and exact: every symbol has a declared semantic route, every global command has
one normative meaning, and every unavoidable local ambiguity has an explicit scope rule.
\end{boxedremark}

\end{document}
